\documentclass[11pt]{article}

\usepackage[margin=1.1in]{geometry}

\usepackage{amssymb,mathtools,natbib}

% Anchored formal environments for causalsmith papers.
% First mandatory argument = obj_id (machine anchor joining the paper to the
% Lean crosswalk; invisible in the rendered PDF). Optional argument = name.
% Each environment also defines \label{obj:<obj_id>} so prose can use
% \cref{obj:T-1}; cleveref derives the printed kind and number from the target.
\usepackage{amsmath}
\usepackage{mathrsfs}
\usepackage{amsthm}
\usepackage{booktabs}
% Long displays may break across pages; let TeX stretch a little harder before an
% overfull \hbox so wide formulas/tables are less likely to run off the page.
\allowdisplaybreaks
\setlength{\emergencystretch}{3em}
\newtheorem{theorem}{Theorem}
\newtheorem{lemma}{Lemma}
\newtheorem{assumption}{Assumption}
\newtheorem{definition}{Definition}
\newtheorem{proposition}{Proposition}
\newtheorem{remark}{Remark}
\newtheorem{citedresult}{Cited result}
% The amsthm counter is `algorithmbox` (not `algorithm`) so a later `\usepackage{algorithm}` cannot clash.
\newcounter{algorithmbox}
\renewcommand{\thealgorithmbox}{\arabic{algorithmbox}}
\NewDocumentEnvironment{theoremv}{m o s}{\begin{theorem}\label{obj:#1}{\normalfont[\texttt{\detokenize{#1}}]\IfValueT{#2}{\ (#2)}\IfBooleanT{#3}{\verificationfootnotemark}\ }}{\end{theorem}}
\NewDocumentEnvironment{lemmav}{m o s}{\begin{lemma}\label{obj:#1}{\normalfont[\texttt{\detokenize{#1}}]\IfValueT{#2}{\ (#2)}\IfBooleanT{#3}{\verificationfootnotemark}\ }}{\end{lemma}}
\NewDocumentEnvironment{assumptionv}{m o s}{\begin{assumption}\label{obj:#1}{\normalfont[\texttt{\detokenize{#1}}]\IfValueT{#2}{\ (#2)}\IfBooleanT{#3}{\verificationfootnotemark}\ }}{\end{assumption}}
\NewDocumentEnvironment{definitionv}{m o s}{\begin{definition}\label{obj:#1}{\normalfont[\texttt{\detokenize{#1}}]\IfValueT{#2}{\ (#2)}\IfBooleanT{#3}{\verificationfootnotemark}\ }}{\end{definition}}
% A `propositionv` is a proved result tiered as a Proposition (theorem-grade, machine-verified).
\NewDocumentEnvironment{propositionv}{m o s}{\begin{proposition}\label{obj:#1}{\normalfont[\texttt{\detokenize{#1}}]\IfValueT{#2}{\ (#2)}\IfBooleanT{#3}{\verificationfootnotemark}\ }}{\end{proposition}}
% A `remarkv` holds NON-load-bearing interpretive / open-question content (e.g. a causalsmith `oeq:`
% open question). It is neither proved nor assumed; no theorem may depend on it.
\NewDocumentEnvironment{remarkv}{m o s}{\begin{remark}\label{obj:#1}{\normalfont[\texttt{\detokenize{#1}}]\IfValueT{#2}{\ (#2)}\IfBooleanT{#3}{\verificationfootnotemark}\ }}{\end{remark}}
% An `algorithmv` is a DEFINITION-kind object presented as a boxed, numbered-step procedure (an
% estimator / selection rule built in stages). Same anchor contract as `definitionv`; its body is
% ordinary LaTeX whose steps are `\begin{enumerate}` items, so tex2html needs no extra machinery.
% Presented in the CONVENTIONAL algorithm style readers expect from `algorithm`+`algorithmic`:
% a rule above, an "Algorithm N (Title)" caption line, a rule under the caption, the numbered
% steps, and a closing rule. It is NOT a float — the anchor contract requires the object to stay
% where the outline places it, and floats drift. The obj-id tag is suppressed here (it is
% bookkeeping, and a caption line carrying it reads as debug output in a finished paper).
\NewDocumentEnvironment{algorithmv}{m o s}%
  {\par\addvspace{\medskipamount}\noindent\rule{\linewidth}{0.8pt}\par\nobreak\vspace{2pt}%
   \refstepcounter{algorithmbox}\label{obj:#1}%
   \noindent\textbf{Algorithm \thealgorithmbox}\IfValueT{#2}{ \textnormal{(#2)}}%
   \IfBooleanT{#3}{\verificationfootnotemark}\par\nobreak\vspace{2pt}%
   \noindent\rule{\linewidth}{0.4pt}\par\nobreak\vspace{2pt}\normalfont}%
  {\par\nobreak\vspace{2pt}\noindent\rule{\linewidth}{0.8pt}\par\addvspace{\medskipamount}}
% A `citedv` is an IMPORTED external result the paper relies on but does NOT prove
% (a causalsmith `gate_class:"cited"` node). It is machine-checked against the cited
% source at F2.5, never proved here; the fixed provenance note makes that explicit
% so the environment can never be read as a contribution of this paper. The \citep
% attribution is supplied by the surrounding prose (P2), keyed to references.bib.
\NewDocumentEnvironment{citedv}{m o s}{\begin{citedresult}\label{obj:#1}{\normalfont[\texttt{\detokenize{#1}}]\IfValueT{#2}{\ (#2)}\IfBooleanT{#3}{\verificationfootnotemark}\ \textit{(imported result; machine-checked against the cited source, not proved here.)}\ }}{\end{citedresult}}
% \leanref{obj_id}{display text}: an inline first-mention link to a from-note object's
% verified Lean code. In the PDF it renders the display text plainly (the Lean link is a
% web-only affordance); tex2html turns it into a drawer-opening span.
\newcommand{\leanref}[2]{#2}
% For a theorem/lemma whose Lean declaration accepts a source-matched published proposition as an
% input, the starred anchored environment puts the mark in its heading; the generated text remains
% outside the frozen mathematical environment. Keep the combined macro for older paper bundles.
\newcommand{\verificationfootnotemark}{\footnotemark}
\newcommand{\verificationfootnotetext}[1]{\footnotetext{#1}}
\newcommand{\verificationfootnote}[1]{\verificationfootnotemark\verificationfootnotetext{#1}}
% xcolor before hyperref (hyperref would otherwise pull in plain `color` for colorlinks, and a
% later xcolor load clashes). The page-1 identity stamp at the end of this file needs a grey.
\usepackage{xcolor}
\usepackage{graphicx}
\usepackage{eso-pic}
% hyperref follows natbib and the environment macros; cleveref must follow hyperref.
\usepackage[colorlinks=true,linkcolor=blue,citecolor=blue,urlcolor=blue]{hyperref}
\usepackage[capitalise,nameinlink,noabbrev]{cleveref}
% Pin the names explicitly: labels are placed inside the underlying amsthm environments, and these
% declarations make the PDF contract independent of cleveref's theorem-name inference. House style:
% reference names always print with a capital first letter ("Theorem 1", "Definition 2"), in any
% sentence position — hence `capitalise` above and capitalized \crefname forms below.
\crefname{theorem}{Theorem}{Theorems}
\Crefname{theorem}{Theorem}{Theorems}
\crefname{lemma}{Lemma}{Lemmas}
\Crefname{lemma}{Lemma}{Lemmas}
\crefname{assumption}{Assumption}{Assumptions}
\Crefname{assumption}{Assumption}{Assumptions}
\crefname{definition}{Definition}{Definitions}
\Crefname{definition}{Definition}{Definitions}
\crefname{proposition}{Proposition}{Propositions}
\Crefname{proposition}{Proposition}{Propositions}
\crefname{remark}{Remark}{Remarks}
\Crefname{remark}{Remark}{Remarks}
\crefname{citedresult}{Cited result}{Cited results}
\Crefname{citedresult}{Cited result}{Cited results}
\crefname{algorithmbox}{Algorithm}{Algorithms}
\Crefname{algorithmbox}{Algorithm}{Algorithms}
% ---------------------------------------------------------------------------
% Page-1 identity stamp (arXiv style) and the pinned title-block date.
%
% Split of responsibility: THIS file owns the FORMAT (placement, font, colour, punctuation),
% the generated `paper_stamp.tex` owns only the VALUES (number+version, area, revised date,
% DOI, link target). P4 refreshes this template on every emit, so a layout fix reaches every
% bundle from a recompile; and because `paper_stamp.tex` is not one of the P2 assembly
% digest's authored inputs, a DOI that arrives after publication needs only that one small
% file rewritten plus a recompile — never a redraft, never a reassembly.
%
% Defaults first, so a bundle WITHOUT a stamp file compiles exactly as it always did:
% `\cspaperdate` falls back to `\today` and no stamp is drawn.
\providecommand*{\cspaperdate}{\today}  % the \date{} target
\providecommand*{\csstampid}{}          % e.g. CSWP-2026-008v3 (empty ⇒ no stamp at all)
\providecommand*{\csstamparea}{}        % e.g. Stat
\providecommand*{\csstampdate}{}        % e.g. 27 Aug 2026
\providecommand*{\csstampdoi}{}         % e.g. 10.5281/zenodo.123 (empty ⇒ DOI part omitted)
\providecommand*{\csstamplink}{}        % href target (DOI url, else the paper's page; may be empty)
\IfFileExists{paper_stamp.tex}{% GENERATED by CausalSmith P4 — paper identity stamp values. Do not hand-edit.
%
% Field order on the stamp (owned by paper_macros.tex):
%   <number>v<version> · doi:<doi> · [<Area>] <date>   — the DOI is never the last token, so a
%   text extractor cannot run it into whatever follows. Without a DOI that part is omitted.
% Format lives in paper_macros.tex; only the values are here, and this file is NOT one of
% the P2 assembly digest's authored inputs. A DOI arriving after publication therefore needs
% this file rewritten plus a `latexmk` recompile — never a redraft or a reassembly.
\renewcommand*{\csstampid}{CSWP-2026-038v3}
\renewcommand*{\csstamparea}{Stat}
\renewcommand*{\csstampdate}{3 Oct 2026}
\renewcommand*{\csstampdoi}{}
\renewcommand*{\csstamplink}{https://causalsmith.org/papers/stat_logodds_lowsmooth_frontier_v1}
\renewcommand*{\cspaperdate}{3 October 2026}
}{}
\makeatletter
% Field order is deliberate: the DOI is NEVER the last token on the line. A trailing identifier
% runs into whatever a text extractor emits next, which makes it unsafe to match mechanically
% (the Zenodo stamp matcher reads exactly this line); the date is a safe terminator.
%   <number>v<version> · doi:<doi> · [<Area>] <date>      — with a DOI
%   <number>v<version> · [<Area>] <date>                  — without
\newcommand{\cs@stampsep}{\,\textperiodcentered\,}
\newcommand{\cs@stamptext}{%
  \csstampid
  \ifx\csstampdoi\@empty\else\cs@stampsep doi:\csstampdoi\fi
  \ifx\csstamparea\@empty\else\cs@stampsep[\csstamparea]\fi
  \ifx\csstampdate\@empty\else\quad\csstampdate\fi
}
\ifx\csstampid\@empty\else
  \definecolor{csstampgrey}{gray}{0.45}
  % Background shipout material: it occupies no space, so the text block and the \thanks
  % footnote are byte-identical to an unstamped compile. The starred form applies to the
  % NEXT page shipped only — i.e. page 1. Placed in the left margin (the text block starts
  % at 1.1in), reading bottom-to-top, in a Times-like face at 8pt.
  \AddToShipoutPictureBG*{%
    \AtPageLowerLeft{%
      \put(\LenToUnit{0.42in},\LenToUnit{1.1in}){%
        \rotatebox{90}{%
          \fontfamily{ptm}\fontsize{8}{10}\selectfont
          \ifx\csstamplink\@empty
            \textcolor{csstampgrey}{\cs@stamptext}%
          \else
            \expandafter\href\expandafter{\csstamplink}{\textcolor{csstampgrey}{\cs@stamptext}}%
          \fi
        }%
      }%
    }%
  }
\fi
\makeatother


\title{Honest confidence intervals for homogeneous causal log odds under low smoothness}

\author{CausalSmith\thanks{Machine-generated by the CausalSmith research pipeline.}}

\date{\cspaperdate}

\begin{document}

\maketitle

\begin{abstract}
We characterize the sharp expected length of honest confidence intervals for a homogeneous conditional log-odds coefficient with binary treatment and outcome. The experiment has a known uniform scalar covariate design, an effect bounded in absolute value by \(1/2\), and propensity and control-risk logits with supremum norms at most \(1\) and Hölder seminorms at most \(2\). For public propensity and prognosis smoothness exponents \(\alpha,\beta\) satisfying \(0<\beta<1/4\) and \(\beta<\alpha\le1\), coverage is required uniformly over this ambient model, while expected length is evaluated on laws with effect magnitude at most \(r\), the evaluation radius. For every sample size \(n\ge1\) and \(r\in[0,1/2]\), the minimax expected length is, up to absolute constants,
\[
n^{-\min\{1/2,\,2(\alpha+\beta)/(2\alpha+2\beta+1)\}}
+r\,n^{-4\beta/(4\beta+1)}.
\]
One radius-independent interval attains this profile simultaneously across all radii with finite-sample \(90\%\) coverage. Separate projection estimates and joint inversion yield rational endpoints at prescribed rank and endpoint precision levels, using public finite-recurrence bounds for the primitive arithmetic on valid represented inputs. Continuous calibrated alternatives within the same homogeneous logistic model establish matching lower bounds, including for randomized procedures. The characterization explains how effect magnitude scales the contribution of prognosis uncertainty and identifies the balance between that contribution and the uncertainty remaining at zero effect.
\end{abstract}

\section{Introduction}

This paper establishes the sharp expected-length profile for honest inference on a homogeneous conditional log-odds coefficient in a prescribed binary treatment and outcome experiment. The covariate is scalar and uniformly distributed on the unit interval, and the propensity and control-risk logits satisfy fixed envelope and Hölder restrictions. Their smoothness exponents are public inputs. Within this model, we construct one interval sequence with finite-sample \(90\%\) coverage and characterize its optimal expected length simultaneously across effect-radius subclasses. The characterization includes evaluation at zero effect, arbitrary shrinking-radius sequences, and fixed positive effect radii.

The target is \(\theta(P)\), the common conditional log-odds coefficient under an observed law \(P\). It is the treated-arm log odds minus the control-arm log odds at every covariate value. Under consistency and conditional exchangeability, this coefficient is also the homogeneous conditional causal log-odds contrast; \cref{obj:lem:observable-odds} records the corresponding full-data interpretation. Conditional odds contrasts have a longstanding role in stratified association analysis \citep{Mantel1959}, and potential-outcome identification supplies their causal interpretation \citep{Rubin1974,Rosenbaum1983}. Their conditional meaning is particularly relevant because confounding and odds-ratio collapsibility are distinct issues \citep{Greenland1999}.

Our precision criterion separates the coverage obligation from the laws at which interval length is evaluated. The ambient model in \cref{obj:def:model} imposes a known uniform design, a homogeneous effect with absolute value at most \(1/2\), nuisance-logit supremum norms at most \(1\), and Hölder seminorms at most \(2\). Let \(\alpha\) denote the public propensity-logit exponent and \(\beta\) the public control-risk-logit exponent. We study
\[
0<\beta<1/4,\qquad \beta<\alpha\le1.
\]
Every admissible procedure must cover the coefficient with probability at least \(90\%\) throughout this ambient model. For \(r\in[0,1/2]\), the effect evaluation radius, its worst-case expected length is assessed over the subclass with \(|\theta(P)|\le r\). This ambient-coverage and subclass-length distinction follows the honest-confidence-interval perspective of \citet{Low1997} and \citet[Section 2, Theorem 1]{CaiLow2004ConfidenceAdaptation}.

Write \(J_n(\alpha,\beta;r)\) for the minimax expected interval length under this criterion, with \(n\) the sample size. Our main result, \cref{obj:thm:full-honest-frontier-answer}, establishes
\[
J_n(\alpha,\beta;r)
\asymp
n^{-\mathsf a(\alpha,\beta)}
+r\,n^{-\mathsf b(\beta)},
\]
where \(\mathsf a\) and \(\mathsf b\), the numerator and denominator rate exponents, are
\[
\mathsf a(\alpha,\beta)
=\min\left\{\frac12,\frac{2(\alpha+\beta)}{2\alpha+2\beta+1}\right\},
\qquad
\mathsf b(\beta)=\frac{4\beta}{4\beta+1}.
\]
The positive comparison constants are absolute and apply simultaneously to every admissible exponent pair, every integer \(n\ge1\), and every evaluation radius. A single procedure, tuned by the sample size and public exponents, attains the upper comparison across all radii. The lower comparison applies to the entire ambiently honest class, including procedures with independent randomization.

The two terms have an observable-coordinate interpretation. The odds identity in \cref{obj:lem:observable-odds} expresses the coefficient through a numerator and a uniformly positive denominator. Numerator uncertainty contributes directly to interval length, whereas denominator uncertainty is scaled by effect magnitude. Separate projection resolutions and joint inversion turn this identity into the attaining interval of \cref{obj:def:upper-handle}. At zero radius, the sharp length is \(n^{-\mathsf a(\alpha,\beta)}\), attaining the parametric order precisely when \(\alpha+\beta\ge1/2\). Throughout the stated exponent domain, \(\mathsf a>\mathsf b\), so the denominator contribution eventually dominates at each fixed positive radius. For shrinking radii, the two contributions balance at
\[
r\asymp n^{-\{\mathsf a(\alpha,\beta)-\mathsf b(\beta)\}}.
\]
These comparisons describe effect-dependent precision while retaining the same ambient coverage requirement.

The construction also specifies a finite rational realization. An explicit interval-arithmetic engine and a dyadic naming policy select a concrete interval with ordered rational endpoints in the effect envelope on every sample. Under the quantitative primitive and representation contracts in \cref{obj:def:arithmetic-engine,obj:def:dyadic-name-convention}, every admissible engine and Borel naming policy satisfies the same coverage and expected-length comparisons. Public exponent representations are fixed before sampling. Prescribed rank and endpoint precision budgets establish termination on valid represented inputs; the resulting intervals may vary across admissible implementations.

The result connects estimand-specific logistic inference to projection-based functional estimation and honest interval optimality. \citet[Propositions 1--3]{Tan2019LogisticDR} develops robust logistic coefficient scores and effect-dependent efficiency comparisons near zero. \citet[arXiv version 1, Assumptions REG and ML1, Theorem 2]{LiuZhangZhou2020LogisticDML} establishes regular asymptotic inference under its nuisance-estimation and regularity conditions. Binary DINA likewise targets conditional log-odds contrasts \citep[Section 2.1.1, equation (5)]{GaoHastie2025DINA}. Projection-based conditional covariance estimation provides a closely related numerator construction \citep[version 5, Theorem 1 and Corollary 1]{McGrath2026}. Here, the contribution is the simultaneous sharp expected-length characterization under ambient honesty in the specified homogeneous logistic model. Its converse uses calibrated continuous alternatives whose every draw satisfies that model's homogeneity, envelope, and native-logit smoothness restrictions.

The results supporting \cref{obj:thm:full-honest-frontier-answer} hold under their stated assumptions.

The paper first reviews related work, then specifies the model and honest interval criterion in \cref{obj:def:model,obj:def:length-objective}. It next presents the sharp comparison in \cref{obj:thm:full-honest-frontier-answer}, develops the interval construction in \cref{obj:def:upper-handle}, and discusses interpretation and limitations. The appendices supply identification, projection bounds, rational realization, calibrated alternatives, original-record lower comparisons, the verification note, and proofs.


\section{Related work}

The target of this paper is the homogeneous conditional log-odds coefficient \leanref{sym:\theta}{\ensuremath{\theta(P)}} in \cref{obj:def:logistic-class}. The statistical experiment combines this pointwise homogeneity restriction with a uniform covariate design and prescribed bounds on the nuisance logits and their Hölder seminorms, as specified in \cref{obj:def:model}. Precision is measured by expected interval length on effect-radius subclasses, while coverage is required over the ambient model; \cref{obj:def:radius-model,obj:def:honest-procedures,obj:def:length-objective} formalize these distinct roles. These choices connect the analysis to logistic coefficient estimation, estimation of nonlinear functionals under low smoothness, and honest confidence intervals.

The closest logistic coefficient results develop robustness and efficiency within partially linear logistic models. \citet[Propositions 1--3]{Tan2019LogisticDR} characterizes nuisance-orthogonal estimating functions and obtains unbiasedness when either the baseline-logit model or the conditional exposure-mean model given covariates and a fixed outcome value is correct. The latter nuisance is an outcome-conditioned exposure mean. When both nuisance models are correct, the efficiency comparison identifies an optimal weight within the proposed estimating-function class. The accompanying near-null comparison explains why a simpler weight approaches that optimal choice as the coefficient approaches zero. This supplies a useful precedent for effect-dependent precision: proximity to the null changes the efficiency comparison within a regular coefficient-estimation problem. Here, proximity to the null indexes subclasses on which the length of an ambiently honest interval is evaluated.
% [Tan (2019), Propositions 1--3 and efficiency discussion](https://arxiv.org/pdf/1901.09138)

Direct odds-ratio antecedents make the nuisance comparison especially clear. \citet{TchetgenTchetgen2010} constructs doubly robust estimators in a semiparametric odds-ratio model by pairing alternative baseline conditional-density specifications. For binary exposure, \citet{TchetgenTchetgen2013} gives a closed-form adjusted odds-ratio estimator with the corresponding double-robustness interpretation. Those papers provide estimation and large-sample inference for an odds-ratio parameter under their nuisance specifications. The present result instead fixes a homogeneous scalar conditional log odds, requires coverage over the full prescribed Hölder model, evaluates length over every effect-radius subclass, and proves same-model converses as well as attainment.

Regular asymptotic inference for the logistic coefficient is developed by \citet[arXiv version 1, Assumptions REG and ML1, Theorem 2]{LiuZhangZhou2020LogisticDML}. Their regularity conditions impose compact parameter and exposure domains together with nondegenerate score sensitivity and variance. For sample size \leanref{sym:n}{\ensuremath{n}}, Assumption ML1 requires uniform consistency of the baseline-logit and outcome-conditioned exposure-mean estimators, with both \(L^2\) errors of order \(o_P(n^{-1/4})\). Under these conditions, the coefficient estimator admits an asymptotically linear representation and a normal limit at the parametric rate. The present expected-length criterion evaluates precision over the prescribed Hölder classes, including their low-smoothness regimes, subject to the coverage requirement in \cref{obj:ass:global-coverage}.
% [Liu, Zhang, and Zhou (2020), arXiv version 1](https://arxiv.org/pdf/2009.14461v1)

The causal interpretation of the estimand also aligns with DINA, the difference in natural parameters. For binary responses, \citet[Section 2.1.1, equation (5)]{GaoHastie2025DINA} defines DINA as the conditional log-odds contrast between potential outcomes. A constant treatment-effect basis specializes this formulation to a homogeneous effect. Their finite-dimensional coefficient analysis assumes a correctly represented treatment-effect function, bounded covariates and nuisance functions, and a score derivative bounded away from singularity. Under those conditions, nuisance estimation errors enter the coefficient bound quadratically, alongside a parametric sampling term \citep[Section 3.3, Proposition 1]{GaoHastie2025DINA}. This comparison separates the shared conditional log-odds target from the precision criterion: their coefficient bound controls estimation error, while the present objective optimizes expected interval length under ambient coverage.
% [Gao and Hastie (2025), Sections 2.1.1 and 3.3](https://academic.oup.com/biometrics/article/81/4/ujaf162/8403946)

The observable numerator in \cref{obj:lem:observable-odds} connects the construction to expected conditional covariance estimation. In the known-design setting, \citet[version 5, Theorem 1 and Corollary 1]{McGrath2026} gives projection bias of order \(k^{-s/d}\) and variance bounded by a constant times \(n^{-1}+k/n^2\), with \(s\) the sum of nuisance smoothness exponents and \(d\) the covariate dimension; the resulting resolution choices attain the covariance minimax rate. Related results extend to local averaging methods: \citet[arXiv version 3, Assumption 4 and Theorem 2]{McCleanBalakrishnanKennedyWasserman2024DCDR} obtains expected absolute error \(n^{-2s/(2s+d)}\) when \(s\le d/2\) and root-\(n\) inference when \(s>d/2\), using independent nuisance-training samples, a known smooth covariate density bounded away from zero, and suitable undersmoothing. Below the regular threshold, additional undersmoothing yields a slower normal limit under boundedness and continuity of conditional variances \citep[arXiv version 3, Theorem 3]{McCleanBalakrishnanKennedyWasserman2024DCDR}. These results distinguish rate-optimal estimation from normal inference with negligible bias. The present construction combines those projection ingredients for two observable coordinates with ambient honesty, all-radius length evaluation, and a homogeneous-logistic converse, as specified in \cref{obj:def:coordinate-estimates,obj:def:upper-handle}.
% [McGrath, version 5, Theorem 1 and Corollary 1](https://arxiv.org/pdf/2212.14857v5)
% [McClean et al., arXiv version 3, Assumption 4 and Theorems 2--3](https://arxiv.org/html/2403.15175v3)

Higher-order inference for implicitly defined parameters provides a further comparison. \citet[Section 3.3, Assumptions 1 and 3--9, Theorems 1--4]{ZhangLiuZhang2026GeneralTreatmentHOIF} analyzes solutions to corrected moment equations under identification, a nonsingular moment derivative, boundedness and overlap, complexity and continuity conditions on the moment class, and localized, nondegenerate sieve bases. Their oracle error bound combines a stochastic term with the product of two residual projection errors. Feasible estimation adds a product involving nuisance errors, the initial parameter estimate, and the error in estimating the projection Gram matrix. Normal inference requires the locally uniform projection bias to vanish after normalization, an additional continuity-rate condition, and, for the feasible estimator, the stated Gram-matrix error condition. The normalization accommodates sieve dimensions both below and above sample size, with dimension growing more slowly than \(n^2\). This decomposition clarifies the shared role of projection in reducing nuisance bias. In the prescribed logistic experiment, the coordinate identity permits an explicit interval construction whose bias and variance bounds enter the finite-sample coverage argument in \cref{obj:thm:finite-honest-upper}.
% [Zhang, Liu, and Zhang (2026), Section 3.3](https://arxiv.org/html/2606.01706v1)

Converse results must be compared at the level of both estimand and statistical class. \citet[Appendix B]{JinMackeySyrgkanis2025HardNormal} studies structure-agnostic bounds for binary treatment in a partially linear mean model, making explicit the roles of overlap and feasible nuisance perturbations. For average conditional log-odds differences, \citet[version 2, Theorem 7.7]{JinSyrgkanis2025GeneralLinearLower} obtains an absolute-error lower bound containing a squared regression-error term, a mixed nuisance-error term, and the parametric sampling term. That result uses unrestricted log-odds contrasts in nuisance-error neighborhoods, under its density, overlap, nondegeneracy, and nonzero-curvature conditions. Applying a converse to the present coefficient requires alternatives that satisfy the homogeneous logistic relation throughout the covariate domain, together with the prescribed envelopes and smoothness restrictions. The calibrated families and original-record mixture comparisons in \cref{obj:def:continuous-calibrated-priors,obj:lem:calibrated-support,obj:lem:original-record-mixtures} supply that comparison within the stated model.
% [Jin, Mackey, and Syrgkanis (2025), Appendix B](https://arxiv.org/html/2507.02275v1)
% [Jin and Syrgkanis, version 2, Theorem 7.7](https://arxiv.org/html/2512.17341v2)

Finally, the distinction between the coverage class and the length-evaluation class belongs to the honest-interval literature. \citet{Low1997} relates expected length at a distribution to coverage over a family of distributions. In Gaussian experiments for linear functionals, \citet[Section 2, Theorem 1]{CaiLow2004ConfidenceAdaptation} expresses a corresponding subclass expected-length lower bound through a between-class modulus of continuity. The present criterion adopts this separation for a homogeneous logistic experiment: the public smoothness exponents specify the ambient class, and effect radii specify nested subclasses for evaluating length. The contribution of \cref{obj:thm:full-honest-frontier-answer} is the joint all-radius characterization under its stated conditions, combining an attaining ambiently honest procedure with matching expected-length lower comparisons, including evaluation at the null.
% [Cai and Low (2004), Section 2, including the discussion of Low (1997)](https://arxiv.org/pdf/math/0503662)

The research provenance also includes two internal supporting problems: the known-fair-radius development \texttt{stat\_dina\_logodds\_honest\_radius} and the fixed-exponent \((0.3,0.1)\) development \texttt{stat\_logodds\_anisotropic\_radius\_frontier}. They supplied antecedent calibration and lower-bound architectures; in particular, the parametric null-floor reduction is credited to the known-fair development. The present theorem establishes the complete public-exponent domain and all-radius profile. The neighboring internal problem \texttt{stat\_logodds\_roughdesign\_radius\_frontier} studies a different experiment with fixed exponents and an unknown bounded design density, whereas the current design is known and uniform.


\section{Model and honest interval criterion}

One observed record is \(O=(X,A,Y)\) in
\[
\Omega=[0,1]\times\{0,1\}\times\{0,1\},
\]
equipped with its product Borel structure. The data
\(\mathcal O=(O_1,\ldots,O_n)\) are independent draws from an observed law
\(P\); sample sizes are positive integers, and \(A,Y\in\{0,1\}\). Write
\(p_{ay}(P,x)=P(A=a,Y=y\mid X=x)\). The observed-law carrier used below
consists of normalized conditional four-cell tables whose entries are
continuous and strictly positive on \([0,1]\), with a covariate marginal of
full support. Thus all conditional quantities have pointwise continuous
versions throughout the covariate interval.

Define the arm risks and treatment probability by
\[
\mu_a(P,x)=\frac{p_{a1}(P,x)}{p_{a0}(P,x)+p_{a1}(P,x)},
\qquad
e(P,x)=p_{10}(P,x)+p_{11}(P,x),
\]
and define the native logits
\[
g(P,x)=\operatorname{logit}\{e(P,x)\},
\qquad
\nu(P,x)=\operatorname{logit}\{\mu_0(P,x)\}.
\]
For \(0<\gamma\le1\), our Hölder convention is
\[
[f]_\gamma=\sup_{x\ne z}\frac{|f(x)-f(z)|}{|x-z|^\gamma}.
\]
The public exponent domain is
\(\mathcal D=\{(\alpha,\beta):0<\beta<1/4,\ \beta<\alpha\le1\}\).
Generic positive constants may change from line to line; \(\|\cdot\|_\infty\)
denotes the supremum norm, \(\lesssim\) hides a positive constant independent
of sample size and evaluation radius, and \(\asymp\) denotes two such bounds.

The comparison output space \(\mathcal C\) consists of every nonempty connected
Borel subset of \([-1/2,1/2]\). Encode an element by its lower and upper
endpoints together with two endpoint-inclusion indicators, excluding the empty
encoding; equip this space with the induced product Borel structure. The
attaining construction below returns the closed member \([l,u]\) with rational
endpoints. Randomized procedures also receive an independent seed
\(U\sim\operatorname{Unif}[0,1]\).
This specifies the record, sampling, and output spaces before the model and
honesty conditions.

The design law is known and fixed as follows.

\begin{assumptionv}{ass:uniform-design}[Uniform covariate design]
The covariate \(X\) has the uniform probability law \(\lambda\) on the unit interval under the observed law \(P\):
\[
P\circ X^{-1}=\lambda.
\]
\end{assumptionv}

The known uniform design is a fixed-density specialization of the known-density setting \citep{McCleanBalakrishnanKennedyWasserman2024DCDR}. It fixes the covariate integration measure across observed laws, allowing the model restrictions to focus on treatment assignment and conditional outcome risks.

We express those risks on the log-odds scale. The log-odds map \leanref{sym:\operatorname{logit}}{\ensuremath{\operatorname{logit}}} is defined first, followed by the native control-risk logit.

\begin{definitionv}{synth_2}[Log-odds map definition]
The log-odds map \(\operatorname{logit}\colon\mathbb R\to\mathbb R\) is defined by
\[
\operatorname{logit}(t)=\log\!\left(\frac{t}{1-t}\right),
\qquad t\in\mathbb R.
\]
Here division by zero has value zero, and the real logarithm is extended by \(\log 0=0\) and \(\log t=\log|t|\) for \(t<0\). For \(0<t<1\), this is the usual log-odds transformation.
\end{definitionv}

The conditional risks in \cref{obj:def:model} lie in \((0,1)\), where this transformation has its usual statistical meaning. The native prognosis logit \leanref{sym:\nu}{\ensuremath{\nu(P,x)}} applies it to the control-arm risk.

\begin{definitionv}{synth_4}[Native prognosis logit]
The native prognosis logit \(\nu(P,x)\) is the log odds of the conditional outcome risk under control, for \(x\in[0,1]\) and an observed probability law \(P\) with binary treatment and outcome whose conditional four-cell table satisfies:
\begin{itemize}
\item \textbf{(Conditional probabilities.)} At each covariate value, the four cells sum to one and give a conditional representation of \(P\) given the covariate.
\item \textbf{(Continuity and interiority.)} Each cell probability is continuous in the covariate and lies strictly between zero and one at every covariate value.
\item \textbf{(Full support.)} The covariate marginal assigns positive probability to every nonempty relatively open subset of \([0,1]\).
\end{itemize}
Its defining formula is
\[
\nu(P,x)
:=
\operatorname{logit}\!\left(
\frac{\Pr_P(\text{treatment}=0,\text{outcome}=1\mid X=x)}
{\Pr_P(\text{treatment}=0,\text{outcome}=0\mid X=x)
+\Pr_P(\text{treatment}=0,\text{outcome}=1\mid X=x)}
\right),
\qquad
\operatorname{logit}(t):=\log\!\left(\frac{t}{1-t}\right).
\]
\end{definitionv}

Consequently, \(\nu(P,x)=\operatorname{logit}\{\mu_0(P,x)\}\). The logistic map \leanref{sym:\ell}{\ensuremath{\ell}} converts a real argument \leanref{sym:t}{\ensuremath{t}} back to a probability.

\begin{definitionv}{synth_1}[The logistic map]
The logistic map \(\ell:\mathbb R\to\mathbb R\) is defined by
\[
\ell(t)=\frac{1}{1+\exp(-t)},\qquad t\in\mathbb R.
\]
\end{definitionv}

On interior probabilities, the logistic and log-odds maps are inverses. We can therefore describe treatment assignment through its native logit and fix the scalar outcome contrast before imposing homogeneity.



The public smoothness exponents govern regularity of these native logits. We fix their notation and the seminorm convention before stating the bounds.



The model combines a homogeneous conditional effect with fixed scalar and nuisance-function envelopes. Its structural restriction is the following logistic relation.

\begin{assumptionv}{ass:homogeneous-logit}[Homogeneous conditional log-odds relation]
The treated-arm conditional outcome risk \(\mu_1(P,x)\) satisfies
\[
\mu_1(P,x)=\ell\{\nu(P,x)+\theta(P)\}
\qquad (x\in[0,1]),
\]
where \(\ell\) is the logistic map, \(\nu(P,x)\) is the native prognosis logit, and \(\theta(P)\) is the homogeneous conditional log-odds coefficient.
\end{assumptionv}

The logistic partially linear coefficient with binary exposure is a standard specification \citep{Tan2019LogisticDR}. Here it makes the conditional log-odds contrast constant across covariate values:
\[
\operatorname{logit}\{\mu_1(P,x)\}
-
\operatorname{logit}\{\mu_0(P,x)\}
=\theta(P).
\]
Equivalently, the treated-to-control conditional odds ratio is \(\exp\{\theta(P)\}\). Thus the coefficient describes a common conditional odds contrast while allowing the baseline risk to vary with the covariate.

Under consistency and conditional independence of treatment and potential outcomes given the covariate, this contrast has the corresponding causal interpretation \citep{Rubin1974,Rosenbaum1983}. The \leanref{S-4}{full-data representation} associated with the observed law is specified in \cref{obj:def:causal-extension}; \cref{obj:lem:observable-odds} supplies the identification result and its causal interpretation in the appendix.

We maintain a compact range for the common coefficient.

\begin{assumptionv}{ass:effect-envelope}[Compact scalar effect envelope]
For an observed law \(P\), the log-odds contrast at covariate zero,
\[
\theta(P)=\operatorname{logit}\bigl(\mu_1(P,0)\bigr)-\nu(P,0),
\]
satisfies
\[
|\theta(P)|\le \frac12.
\]
\end{assumptionv}

A compact coefficient region is standard in logistic coefficient analysis \citep{LiuZhangZhou2020LogisticDML}. The numerical bound \(1/2\) is maintained here as part of the model and gives a common effect range for the interval comparison.

The treatment mechanism is controlled by a separate logit envelope.

\begin{assumptionv}{ass:propensity-envelope}[Propensity logit envelope]
The native propensity-logit function \(g(P,\cdot)\) satisfies
\[
\|g(P,\cdot)\|_\infty\le 1.
\]
\end{assumptionv}

The bounded propensity logit is a quantitative bounded-overlap specialization, with the constant \(1\) maintained here \citep{RobinsTchetgenLiVanderVaart2009Minimax}. Since \(e(P,x)=\ell\{g(P,x)\}\), it places treatment probabilities between \(\ell(-1)\) and \(\ell(1)\) uniformly over covariate values.

A corresponding envelope controls the baseline outcome risk.

\begin{assumptionv}{ass:prognosis-envelope}[Prognosis logit envelope]
The native prognosis-logit function \(\nu(P,\cdot)\) satisfies
\[
\|\nu(P,\cdot)\|_\infty\le 1.
\]
\end{assumptionv}

The bounded prognostic logit belongs to the bounded logistic framework, with this numerical envelope maintained for the present analysis \citep{LiuZhangZhou2020LogisticDML}. Together with \cref{obj:ass:effect-envelope}, it keeps both arm-specific outcome risks uniformly in the interior of the probability interval.

The remaining nuisance restrictions describe variation across covariate values, beginning with treatment assignment.

\begin{assumptionv}{ass:propensity-holder}[Propensity Hölder seminorm bound]
The Hölder seminorm of the native propensity-logit function \(g(P,\cdot)\) at the public smoothness exponent \(\alpha\) satisfies
\[
[g(P,\cdot)]_\alpha\le 2.
\]
\end{assumptionv}

A Hölder ball provides the standard smoothness restriction on the treatment mechanism \citep{RobinsTchetgenLiVanderVaart2009Minimax}. Here the ball is imposed on the native propensity logit, with prescribed radius \(2\); its exponent controls how rapidly treatment assignment can vary over nearby covariate values.

Baseline prognosis has its own public exponent.

\begin{assumptionv}{ass:prognosis-holder}[Prognosis Hölder bound]
The native prognosis-logit function \(\nu(P,\cdot)\) satisfies the Hölder seminorm bound
\[
[\nu(P,\cdot)]_\beta\le 2.
\]
\end{assumptionv}

The prognostic Hölder ball is likewise a standard regularity condition, imposed here on the control-risk logit with prescribed radius \(2\) \citep{LiuMukherjeeRobinsTchetgen2021Adaptive}. Distinct exponents permit treatment assignment and baseline prognosis to have different degrees of smoothness.

We collect the structural relation and the quantitative restrictions into two classes. The homogeneous logistic class \leanref{sym:\mathcal L}{\ensuremath{\mathcal L}} records the common conditional effect relation.

\begin{definitionv}{def:logistic-class}[Homogeneous logistic class \(\mathcal L\)]
\[
\mathcal L
=
\left\{
P:\mu_1(P,x)=\ell\{\nu(P,x)+\theta(P)\}
\ \text{for every }x\in[0,1]
\right\}.
\]
Here \(\mu_1(P,x)\) is the treated-arm conditional outcome risk, \(\ell\) is the logistic map, \(\nu(P,x)\) is the native prognosis logit, and \(\theta(P)\) is the homogeneous conditional log-odds coefficient. The defining relation is that of \cref{obj:ass:homogeneous-logit}.
\end{definitionv}

For a fixed public exponent pair, the ambient model \leanref{sym:\mathcal M}{\ensuremath{\mathcal M(\alpha,\beta)}} adds the known design law and the prescribed envelopes and seminorm bounds.

\begin{definitionv}{def:model}[Ambient model \(\mathcal M(\alpha,\beta)\)]
\(\mathcal M(\alpha,\beta)\) is the class of observed laws \(P\) satisfying the following conditions, for each \((\alpha,\beta)\in\mathcal D\), the prescribed domain of public smoothness exponents:
\begin{itemize}
\item \textbf{(Uniform design.)} With \(\lambda\) denoting the uniform probability law on the unit interval,
\[
P\circ X^{-1}=\lambda.
\]
See \cref{obj:ass:uniform-design}.
\item \textbf{(Homogeneous logit.)} The treated-arm conditional outcome risk \(\mu_1(P,x)\), logistic map \(\ell\), native prognosis logit \(\nu(P,x)\), and homogeneous conditional log-odds coefficient \(\theta(P)\) satisfy
\[
\mu_1(P,x)=\ell\{\nu(P,x)+\theta(P)\}
\qquad (x\in[0,1]).
\]
See \cref{obj:ass:homogeneous-logit}.
\item \textbf{(Effect envelope.)}
\[
|\theta(P)|\le \frac12.
\]
See \cref{obj:ass:effect-envelope}.
\item \textbf{(Propensity envelope.)} The native propensity-logit function \(g(P,\cdot)\) satisfies
\[
\|g(P,\cdot)\|_\infty\le 1.
\]
See \cref{obj:ass:propensity-envelope}.
\item \textbf{(Propensity smoothness.)}
\[
[g(P,\cdot)]_\alpha\le 2.
\]
See \cref{obj:ass:propensity-holder}.
\item \textbf{(Prognosis envelope.)}
\[
\|\nu(P,\cdot)\|_\infty\le 1.
\]
See \cref{obj:ass:prognosis-envelope}.
\item \textbf{(Prognosis smoothness.)}
\[
[\nu(P,\cdot)]_\beta\le 2.
\]
See \cref{obj:ass:prognosis-holder}.
\end{itemize}
The bracket notation denotes the Hölder seminorm at the indicated public exponent.
\end{definitionv}

This ambient class determines the coverage requirement. To evaluate precision near different effect magnitudes, we introduce an evaluation radius \leanref{sym:r}{\ensuremath{r}} in \([0,1/2]\).

\begin{assumptionv}{ass:evaluation-radius}[Effect evaluation radius]
The homogeneous conditional log-odds coefficient \(\theta(P)\) satisfies
\[
|\theta(P)|\le r.
\]
\end{assumptionv}

The nested evaluation subregion follows the ambient-coverage and subclass-length criterion of the honest-interval literature \citep{CaiLow2004ConfidenceAdaptation}. The radius indexes the laws over which expected length is assessed, while the ambient class remains the coverage domain.

The effect-radius slice \leanref{sym:\mathcal M_r}{\ensuremath{\mathcal M_r(\alpha,\beta)}} makes this evaluation class explicit.

\begin{definitionv}{def:radius-model}[Effect-radius slice \(\mathcal M_r(\alpha,\beta)\)]
\[
\mathcal M_r(\alpha,\beta)
=
\{P\in\mathcal M(\alpha,\beta):|\theta(P)|\le r\}
\qquad
\bigl((\alpha,\beta)\in\mathcal D,\ r\in[0,1/2]\bigr).
\]
Here \(\mathcal D\) is the prescribed domain of public smoothness exponents and \(\theta(P)\) is the homogeneous conditional log-odds coefficient. This defines the slice of the ambient model in \cref{obj:def:model} satisfying \cref{obj:ass:evaluation-radius}.
\end{definitionv}

These slices are nested in the radius. At \(r=0\), they contain the ambient laws with zero conditional log-odds contrast; at \(r=1/2\), the slice equals the ambient model. The nuisance restrictions remain fixed throughout this comparison.

Before stating honesty, we specify the sample, interval outputs, and independent randomization used by the comparison class.



Including an independent seed permits randomized procedures in the comparison class. A deterministic procedure is represented by a map that is constant in its seed argument. Coverage probabilities and expected lengths average over both the original records and this independent randomization.

We require coverage uniformly over the ambient model.

\begin{assumptionv}{ass:global-coverage}[Global interval coverage]
The interval procedure \(I(\mathcal O,U)\), evaluated on the tuple of observed records \(\mathcal O\) and an independent uniform randomization seed \(U\), satisfies
\[
\inf_{P\in\mathcal M(\alpha,\beta)}
(P^{\otimes n}\otimes\lambda)
\{\theta(P)\in I(\mathcal O,U)\}\ge 0.90,
\]
where \(\mathcal M(\alpha,\beta)\) is the ambient observed-law model, \(\lambda\) is the uniform probability law on the unit interval, and \(\theta(P)\) is the homogeneous conditional log-odds coefficient.
\end{assumptionv}

Global honesty with subclass expected-length evaluation is the standard comparison criterion, specialized here to the binary treatment and outcome experiment \citep{CaiLow2004ConfidenceAdaptation}. It is a procedure-admissibility requirement: for each public exponent pair and sample size, an admissible interval covers the coefficient with probability at least \(90\%\) at every law in the ambient model.

The honest procedure class \leanref{sym:\mathcal A}{\ensuremath{\mathcal A_n(\alpha,\beta)}} collects the maps satisfying this requirement.

\begin{definitionv}{def:honest-procedures}[Honest interval procedures $\mathcal A_n(\alpha,\beta)$]
For every \(n\ge1\) and \((\alpha,\beta)\in\mathcal D\), define
\[
\mathcal A_n(\alpha,\beta)
=
\left\{
I:\Omega^n\times[0,1]\to\mathcal C:
\inf_{P\in\mathcal M(\alpha,\beta)}
(P^{\otimes n}\otimes\lambda)
\{\theta(P)\in I(\mathcal O,U)\}
\ge0.90
\right\},
\]
where the maps \(I\) are Borel and have measurable coverage events and interval length, as required by \cref{obj:ass:global-coverage}. Here \(\Omega\) is the measurable space of one observed record, \(\mathcal O\) is the tuple of original observed records, \(\mathcal C\) is the space of admissible nonempty connected Borel effect intervals, and \(\lambda\) is the uniform probability law on \([0,1]\). The variable \(U\) is an independent randomization seed, \(\mathcal D\) is the prescribed domain of public smoothness exponents, and \(\theta(P)\) is the homogeneous conditional log-odds coefficient.
\end{definitionv}

For a fixed sample size and public exponent pair, this same admissible class is used at every evaluation radius. Precision is measured by the minimax expected length \leanref{sym:J}{\ensuremath{J_n(\alpha,\beta;r)}} over the corresponding slice.

\begin{definitionv}{def:length-objective}[Minimax expected length $J_n(\alpha,\beta;r)$]
The minimax expected interval length on the effect-radius slice, subject to honesty over the ambient model, is
\[
J_n(\alpha,\beta;r)
=
\inf_{I\in\mathcal A_n(\alpha,\beta)}
\sup_{P\in\mathcal M_r(\alpha,\beta)}
E_{P^{\otimes n}\otimes\lambda}
\bigl|I(\mathcal O,U)\bigr|.
\]
Here \(\mathcal A_n(\alpha,\beta)\) is the class of Borel interval procedures with global \(90\%\) coverage, \(\mathcal M_r(\alpha,\beta)\) is the effect-radius slice, \(\mathcal O\) is the tuple of original observed records, and \(U\) is an independent seed with uniform law \(\lambda\) on \([0,1]\). The quantity \(\lvert I(\mathcal O,U)\rvert\) is the length of the returned interval.
\end{definitionv}

The objective combines ambient coverage with local evaluation of expected length. In particular, evaluation at the null retains coverage over all effects in the ambient envelope. Varying the radius therefore measures how much precision an ambiently honest procedure can attain across nested effect regions.


\section{Sharp expected length across effect radii}

The expected-length objective in \cref{obj:def:length-objective} admits a sharp comparison with two contributions: uncertainty in an observable numerator and effect-scaled uncertainty in an observable denominator. Their rates depend on the public propensity smoothness exponent \leanref{sym:\alpha}{\ensuremath{\alpha}} and prognosis smoothness exponent \leanref{sym:\beta}{\ensuremath{\beta}}. A single concrete interval sequence, \leanref{sym:I^\star}{\ensuremath{I_n^\star}}, attains the comparison simultaneously across evaluation radii while maintaining coverage over the ambient model. We first define the diagnostic rate and then state the complete comparison. The statistical construction and numerical realization follow in the next section and appendices.

The diagnostic envelope \leanref{sym:d_n}{\ensuremath{d_n(\alpha,\beta;r)}} separates the numerator contribution from the radius-scaled denominator contribution.

\begin{definitionv}{def:diagnostic-envelope}[Diagnostic envelope \(d_n(\alpha,\beta;r)\)]
The diagnostic envelope is
\[
d_n(\alpha,\beta;r)
=
n^{-\mathsf a(\alpha,\beta)}
+r\,n^{-\mathsf b(\beta)},
\]
where the numerator and denominator rate exponents are, respectively,
\[
\mathsf a(\alpha,\beta)
=
\min\left\{\frac12,\frac{2(\alpha+\beta)}{2\alpha+2\beta+1}\right\},
\qquad
\mathsf b(\beta)=\frac{4\beta}{4\beta+1}.
\]
\end{definitionv}

\paragraph{Main statistical result.}
For every public pair \(0<\beta<1/4\), \(\beta<\alpha\le1\), every
\(n\ge1\), and every \(r\in[0,1/2]\), the minimax expected length over
procedures with \(90\%\) coverage on the ambient model satisfies
\[
J_n(\alpha,\beta;r)\asymp
n^{-\min\{1/2,\,2(\alpha+\beta)/(2\alpha+2\beta+1)\}}
+r n^{-4\beta/(4\beta+1)}.
\]
The comparison constants are absolute. One procedure, using the sample size and
public exponents and independent of \(r\), attains the upper bound for every
radius simultaneously; the converse covers every globally honest Borel
procedure, including randomized procedures. The exact finite-sample statement,
including the rational implementation guarantee, is
\cref{obj:thm:full-honest-frontier-answer} below. The observable identity
\[
C(P)=\{\exp(\theta(P))-1\}S(P),\qquad S(P)\ge1/512,
\]
explains the two terms: numerator uncertainty remains at the null, while
denominator uncertainty is multiplied by the effect radius.

\begin{definitionv}{synth_3}[Cosine projection and its kernel]
Let \(\lambda\) be the uniform probability measure on \([0,1]\). The orthonormal cosine basis of \(L^2(\lambda)\) is
\[
\varphi_0(x)=1,\qquad
\varphi_j(x)=\sqrt{2}\cos(\pi jx)\quad(j\ge1).
\]
For each integer \(k\ge1\), define the rank-\(k\) orthogonal projection onto \(\operatorname{span}\{\varphi_0,\ldots,\varphi_{k-1}\}\) by
\[
\Pi_k f
=\sum_{j=0}^{k-1}\langle f,\varphi_j\rangle_{L^2(\lambda)}\varphi_j,
\qquad f\in L^2(\lambda),
\]
where \(\langle f,h\rangle_{L^2(\lambda)}=\int_0^1 f(z)h(z)\,d\lambda(z)\). Its projection kernel is
\[
H_k(x,z)=\sum_{j=0}^{k-1}\varphi_j(x)\varphi_j(z)
=1+2\sum_{j=1}^{k-1}\cos(\pi jx)\cos(\pi jz),
\qquad x,z\in[0,1].
\]
Thus the continuous representative of \(\Pi_k f\) satisfies
\[
(\Pi_k f)(x)=\int_0^1 H_k(x,z)f(z)\,d\lambda(z),
\qquad x\in[0,1].
\]
\end{definitionv}

\begin{definitionv}{def:projection-statistic}[Projection statistic $\mathbb U_{n,k}(W,V)$]
For \(n\ge2\), the ordered-pair projection statistic for two bounded record marks \(W\) and \(V\) is
\[
\mathbb U_{n,k}(W,V)
=
\frac{1}{n(n-1)}
\sum_{\substack{1\le i,j\le n\\ i\ne j}}
H_k(X_i,X_j)W(O_i)V(O_j).
\]
Here \(O_i\) is the \(i\)th observed record, \(X_i\) is its covariate, and \(H_k\) is the kernel of the projection at rank \(k\).
\end{definitionv}

\begin{definitionv}{def:dyadic-name-convention}[Dyadic naming policies \(\mathsf N,\mathsf N_0\)]
A dyadic name for a real number \(v\) is a sequence of nested closed intervals with dyadic rational endpoints that contain \(v\) and have width at most \(2^{-q}\) at each integer precision \(q\ge0\). The product coding of these sequences uses the discrete measurable structure on rational responses.

An admissible naming policy \(\mathsf N\) is a sequence of policies, one for each sample size \(n\), satisfying the following conditions.
\begin{itemize}
\item \textbf{(Public exponents.)} Valid names for \(\alpha,\beta\) are selected by Borel functions of the public pair \((\alpha,\beta)\) alone, before sampling.
\item \textbf{(Observed covariates.)} Valid names for the \(n\) observed covariates are selected by Borel functions of \((\alpha,\beta,\mathcal O)\), where \(\mathcal O\) is the tuple of original observed records.
\item \textbf{(Separate arithmetic input.)} The arithmetic engine is independently fixed, and name selection supplies the real-input representations used with that engine.
\end{itemize}
The interval procedure receives these selected names and the exact binary labels. Name selection is an input convention.

The dyadic-floor policy \(\mathsf N_0\) assigns every real input \(v\) the intervals
\[
\mathsf N_0(v)_q
=
\left[
2^{-q}\lfloor2^qv\rfloor,\;
2^{-q}\bigl(\lfloor2^qv\rfloor+1\bigr)
\right],
\qquad q\ge0.
\]
This is a valid name also at dyadic inputs and selects one concrete member of the policy-indexed interval family.

At the rank step, the queried exponent intervals are intersected with the known compact boxes
\[
\alpha\in[0,1],\qquad \beta\in[0,1/4].
\]
A syntactically malformed queried prefix or an empty intersection may yield the full effect interval. At the endpoint precision, the supplied exponent and observed-covariate intervals are passed raw to the endpoint expression tree, and the integer ranks selected at the rank step are retained. Outputs may vary across valid names.
\end{definitionv}

\begin{definitionv}{def:arithmetic-engine}[Arithmetic engine $\mathsf E$]
An arithmetic engine \(\mathsf E\) is a deterministic collection of maps on closed rational intervals, with integer precision input \(q\ge 0\) and tolerance
\[
\epsilon_q=2^{-q}.
\]
It supplies rational enclosures of \(\pi\), the exponential and cosine on every bounded rational interval, the logarithm on rational intervals \([a,b]\) with \(0<a\le b\), the square root on rational intervals \([a,b]\) with \(0\le a\le b\), and the power \(b^v\) for a positive integer base \(b\) and a rational exponent interval contained in \([-3,3]\). The symbols \(a,b,v,q\) are bound primitive arguments; in the power primitive, \(b\) denotes the integer base, supplied separately from the interval endpoints.

The engine also supplies exact rational range operations for addition, subtraction, multiplication, division on denominator intervals avoiding zero, minimum, maximum, intersection, and clipping. In particular, the minimum of two intervals has range
\[
\min\bigl([a,b],[c,d]\bigr)
=
[\min(a,c),\min(b,d)],
\]
and multiplication has range
\[
[a,b]\,[c,d]
=
[\min\{ac,ad,bc,bd\},\max\{ac,ad,bc,bd\}].
\]
Division uses the reciprocal range followed by multiplication, and clipping uses the exact minimum and maximum operations.

Admissibility requires the following primitive contracts.
\begin{itemize}
\item \textbf{(Enclosure.)} Every returned interval contains the real image of its input interval.
\item \textbf{(Monotone primitives.)} For exponential, logarithm, square root, and positive-integer-base power, each returned endpoint differs outward from the corresponding exact range endpoint by at most \(\epsilon_q\).
\item \textbf{(Constant enclosure.)} The returned enclosure of \(\pi\) lies in \([3,4]\), with outward excess at each endpoint at most \(\epsilon_q\).
\item \textbf{(Cosine.)} At the rational midpoint of the input interval, the engine encloses the cosine value with outward endpoint excess at most \(\epsilon_q\). It enlarges this enclosure by half the input width on each side and intersects the result with \([-1,1]\). The returned width is at most the input width plus \(2\epsilon_q\).
\item \textbf{(Primitive domains and signs.)} Exponential and power outputs have positive lower endpoints, and square-root outputs have nonnegative lower endpoints. Logarithm is evaluated on intervals with positive rational lower endpoints, and division is evaluated on denominator intervals avoiding zero.
\item \textbf{(Finite rational computation.)} Every map is defined by a finite rational recurrence whose iteration count is bounded by a specified public integer function of the rational range bounds, the integer base when present, and \(q\). The recurrence may use rational comparisons and integer arithmetic; its access to a real input is exclusively through the supplied rational interval.
\item \textbf{(Independent choice.)} The engine is fixed independently of all public exponents, sample sizes, observed records, laws, naming policies, evaluation radii, and unknown nuisance quantities.
\end{itemize}
Distinct admissible engines may return distinct rational enclosures. Admissibility is specified entirely by these primitive contracts, with rank conclusions, final-endpoint conclusions, final precision budgets, and statistical properties outside its defining requirements.
\end{definitionv}

\begin{algorithmv}{def:resolutions}[Public projection rank selection]
The projection ranks are the nonnegative integer ceilings of the upper endpoints produced by the following construction. For \(n\ge2\) and real public exponents \(\alpha,\beta\), define the public precision budget and target resolutions by
\[
q_{\mathrm{rank}}(n)=7+3\lceil\log_2 n\rceil,\qquad
R_C=n^{2/(2\alpha+2\beta+1)},\qquad
R_S=n^{2/(4\beta+1)}.
\]
For a positive integer base \(b\), the logarithmic ceiling is defined by integer comparisons:
\[
\lceil\log_2 b\rceil
=\min\{q\in\mathbb Z_{\ge0}:2^q\ge b\}.
\]

Fix a rational interval-arithmetic engine \(\mathsf E\) and a naming policy \(\mathsf N\). Query the public exponent names at precision \(q_{\mathrm{rank}}(n)\) and intersect them with \([0,1]\) and \([0,1/4]\), respectively. Using these intersections, form the ranges of
\[
\frac{2}{2\alpha+2\beta+1},
\qquad
\frac{2}{4\beta+1}
\]
by exact rational interval operations, then apply the power map of \(\mathsf E\) with integer base \(n\) at the same precision. If these operations succeed, denote the resulting boxes by \([a_C,b_C]\) and \([a_S,b_S]\), and define
\[
\begin{aligned}
k_C(n,\alpha,\beta;\mathsf E,\mathsf N)
&=\max\{\lceil b_C\rceil,0\},\\
k_S(n,\beta;\mathsf E,\mathsf N)
&=\max\{\lceil b_S\rceil,0\}.
\end{aligned}
\]
Here \(b_C,b_S\) denote resolution-box upper endpoints. If either an intersection or an interval division fails, define both ranks to be \(1\). Both ranks use the public inputs \((n,\alpha,\beta)\) through \(\mathsf N\), including in the abbreviated notation for \(k_S\), and are fixed before sampling.

For the admissible construction, take:
\begin{itemize}
\item \textbf{(Engine.)} \(\mathsf E\) satisfying \cref{obj:def:arithmetic-engine}.
\item \textbf{(Names.)} \(\mathsf N\) satisfying \cref{obj:def:dyadic-name-convention}, with public exponent names depending on \(n,\alpha,\beta\).
\item \textbf{(Exponent domain.)} \(0<\beta<1/4\) and \(\beta<\alpha\le1\).
\end{itemize}
For every \(n\ge2\), the construction succeeds and its boxes satisfy
\[
R_C\in[a_C,b_C],\qquad R_S\in[a_S,b_S],
\qquad b_C-a_C\le1,\qquad b_S-a_S\le1.
\]
Consequently, the selected ranks satisfy
\[
\begin{aligned}
R_C&\le k_C(n,\alpha,\beta;\mathsf E,\mathsf N)\le3R_C,\\
R_S&\le k_S(n,\beta;\mathsf E,\mathsf N)\le3R_S.
\end{aligned}
\]
\end{algorithmv}

\begin{definitionv}{def:coordinate-estimates}[Coordinate estimates $\widehat C_n$ and $\widehat S_n$]
The exact real sample coordinates are
\[
\widehat C_n
=\frac1n\sum_{i=1}^n A_iY_i-\mathbb U_{n,k_C}(A,Y),\qquad
\widehat S_n
=\mathbb U_{n,k_S}\bigl(A(1-Y),(1-A)Y\bigr),
\]
where, for \(n\ge2\), the retained ranks are
\[
k_C=k_C(n,\alpha,\beta;\mathsf E,\mathsf N),\qquad
k_S=k_S(n,\alpha,\beta;\mathsf E,\mathsf N).
\]
Both ranks are selected by the fixed-budget resolution construction with the same admissible engine \(\mathsf E\) and the public exponent approximations supplied by \(\mathsf N\). The upper construction uses this same engine and the supplied covariate approximations to enclose the finite sums at these retained ranks. The statistics \(\widehat C_n,\widehat S_n\) estimate the observable numerator and denominator coordinates \(C(P),S(P)\), respectively.
\end{definitionv}

\begin{definitionv}{def:causal-extension}[Full-data law $\widetilde P_{P,\Gamma}$]
The full-data law \(\widetilde P_{P,\Gamma}\) is defined by the successive draws
\[
X\sim\lambda,\qquad
(Y^0,Y^1)\mid X=x\sim\Gamma(P,x),\qquad
A\mid(X,Y^0,Y^1)\sim\operatorname{Bernoulli}\{e(P,X)\},
\]
and the observed outcome
\[
Y=(1-A)Y^0+AY^1.
\]
Here \(\lambda\) is the uniform probability law on \([0,1]\), \(Y^0\) and \(Y^1\) are the control and treated potential outcomes, \(\Gamma(P,x)\) is an admissible conditional coupling of these outcomes, and \(e(P,X)\) is the treatment probability conditional on the covariate under the observed law \(P\).
\end{definitionv}

\begin{definitionv}{synth_5}[Frontier comparison constants and sample-size threshold]
Fix positive finite absolute constants \(c,C\) satisfying the simultaneous comparisons in \cref{obj:thm:sharp-honest-frontier}. Define the selected lower and upper comparison functions and sample-size threshold on \(\mathcal D\) by
\[
c_\star(\alpha,\beta)=c,\qquad
C_\star(\alpha,\beta)=C,\qquad
n_0(\alpha,\beta)=1.
\]
Thus, for every admissible arithmetic engine \(\mathsf E\), admissible Borel naming policy \(\mathsf N\), \((\alpha,\beta)\in\mathcal D\), integer \(n\ge n_0(\alpha,\beta)\), and \(r\in[0,1/2]\),
\[
c_\star(\alpha,\beta)\rho_n(\alpha,\beta;r)
\le J_n(\alpha,\beta;r)
\le
\sup_{P\in\mathcal M_r(\alpha,\beta)}
E_{P^{\otimes n}}\!
\left|I^\star_{n,\mathsf E,\mathsf N}(\alpha,\beta)\right|
\le C_\star(\alpha,\beta)\rho_n(\alpha,\beta;r).
\]
The constants \(c,C\) are selected once, independently of the exponents, sample size, radius, engine, and naming policy.
\end{definitionv}

\begin{definitionv}{def:calibration-maps}[Calibration maps $\mathfrak c,\mathfrak B,\mathfrak q,\mathsf T,\mathfrak H,\mathfrak p$]
The covariance adjustment and its effect-normalized counterpart are
\[
\mathfrak c(t,\xi,\upsilon)
=\frac{2dv_*}{L_*+\sqrt{L_*^2-4d^2v_*}},\qquad
\mathfrak B(t,\xi,\upsilon)
=\frac{2D(t)v_*}{L_*+\sqrt{L_*^2-4d^2v_*}},
\]
where \(\xi,\upsilon\) are scalar margins near \(1/2\), and
\[
d=\exp(t)-1,\qquad
v_*=\xi(1-\xi)\upsilon(1-\upsilon),
\]
\[
L_*=1+d\{\xi(1-\upsilon)+(1-\xi)\upsilon\},\qquad
D(t)=\int_0^1\exp(st)\,ds.
\]
Here \(s\) is an integration variable. These formulas apply on the domain where the displayed square root and denominator are positive.

For two independent uniform signs \(\epsilon_0,\epsilon_1\), define
\[
Z_u=\epsilon_0\cos(\pi u/2)+\epsilon_1\sin(\pi u/2).
\]
Expectations below are over these signs. Let \(\varepsilon_{\mathrm{mix}}>0\) be the radius of the fixed signed open amplitude square on which the mixed calibration branch exists uniquely. On
\[
|\eta|<\varepsilon_{\mathrm{mix}},\qquad
|\zeta|<\varepsilon_{\mathrm{mix}},\qquad
u\in[0,1],
\]
define the mixed calibration root by
\[
\mathfrak q(\eta,\zeta,u)
=
\text{the unique }q\in(3/4,5/4)
\text{ satisfying }
16E\mathfrak B
\bigl(32\eta\zeta,1/2-\eta qZ_u,1/2+\zeta Z_u\bigr)-q=0.
\]
The smaller closed amplitude squares used for support and mixture bounds are contained in this domain.

For the separate fair-prior maps, define the baseline risk and local variance by
\[
p_*=2/5,\qquad v_*=p_*(1-p_*)=6/25.
\]
Here \(v_*\) denotes the fair-branch Bernoulli variance. Define
\[
\mathsf G(t,\xi)
=\frac{\exp(t)\xi}{1+\{\exp(t)-1\}\xi},\qquad
B_*=1+\{\exp(t)-1\}p_*,
\]
\[
\mathsf T(t,\delta)
=t+
\log\left(1-\frac{\{\exp(t)-1\}\delta^2}{p_*B_*}\right)
-\log\left(1+\frac{\{\exp(t)-1\}\delta^2}{(1-p_*)B_*}\right).
\]
On the fair core \(0\le t\le1/4\), for \(t\delta\ne0\), the normalized fair-matching residual is
\[
\mathfrak H(t,\delta,\xi,u)
=
\frac{
E\mathsf G(t,\xi+\delta Z_u)
-\mathsf G(\mathsf T(t,\delta),\xi)
}{t\delta^2}.
\]
The symbol \(\mathfrak H\) also denotes its removable smooth extension to the fixed open calibration neighborhood, including across \(t=0\), \(\delta=0\), and the spatial endpoints. The displayed quotient identity applies on the fair core.

Let \(\varepsilon_{\mathrm{fair}}>0\) be the signed fair-amplitude radius, and let \(b\in(0,1/10)\) be the fixed sufficiently small bracket half-width about \(p_*\). Here \(b\) denotes a bracket half-width. The fair parameter domain contains
\[
0\le t\le1/4,\qquad
|\delta|<\varepsilon_{\mathrm{fair}},\qquad
u\in[0,1],
\]
together with the fixed open neighborhood of these parameters, including signed parameters across \(t=0\) and the spatial endpoints. The denominators and logarithm arguments in the fair formulas are positive on this neighborhood. At \(t=0\) or \(\delta=0\), \(\mathfrak H\) is interpreted through its removable smooth extension.

On this neighborhood, define the fair calibration root by
\[
\mathfrak p(t,\delta,u)
=
\text{the unique }\xi\in(p_*-b,p_*+b)
\text{ satisfying }\mathfrak H(t,\delta,\xi,u)=0.
\]
Selection in the larger bracket \((3/10,1/2)\) gives this same root on the stated neighborhood. The smaller closed amplitude neighborhood used for support and mixture bounds is contained in this fair domain, uniformly over \(t\in[0,1/4]\) and \(u\in[0,1]\). The local branches are defined on the stated domains, where their existence, uniqueness, regularity, and matching properties hold.
\end{definitionv}

\begin{definitionv}{def:continuous-calibrated-priors}[Continuous calibrated priors \(P_{b,\sigma},P_{t,\sigma},Q_b,Q_t\)]
The continuous sign field on an equal partition is
\[
\mathsf Z_{k,\sigma}(x)
=
\sigma_j\cos(\pi u/2)+\sigma_{j+1}\sin(\pi u/2),
\qquad
u=\frac{x-j\Delta}{\Delta},
\qquad
\Delta=\frac1k,
\]
for an integer \(k\ge1\), independent uniform signs
\(\sigma=(\sigma_0,\ldots,\sigma_k)\), and
\(j\Delta\le x\le(j+1)\Delta\), \(0\le j<k\). Either adjoining cell gives the value at a shared endpoint.

For signed amplitudes \(\eta,\zeta\), define
\[
e_0(x)=\frac12+\eta\mathfrak q(\eta,\zeta,u)\mathsf Z_{k,\sigma}(x),
\qquad
e_1(x)=\frac12-\eta\mathfrak q(\eta,\zeta,u)\mathsf Z_{k,\sigma}(x),
\]
\[
m_\sigma(x)=\frac12+\zeta\mathsf Z_{k,\sigma}(x),
\qquad
c_0(x)=0,
\qquad
c_1(x)=\mathfrak c(32\eta\zeta,e_1(x),m_\sigma(x)).
\]
For the family index \(b\in\{0,1\}\), the law \(P_{b,\sigma}\) has uniform covariate law \(\lambda\) on the unit interval and conditional cell probabilities
\[
\begin{aligned}
p_{11}&=e_bm_\sigma+c_b,
&
p_{10}&=e_b(1-m_\sigma)-c_b,\\
p_{01}&=(1-e_b)m_\sigma-c_b,
&
p_{00}&=(1-e_b)(1-m_\sigma)+c_b.
\end{aligned}
\]
Here \(e_b,m_\sigma,c_b\) are evaluated at the current covariate; \(p_{11},p_{10},p_{01},p_{00}\) correspond to treated success, treated failure, control success, and control failure.

For \(0\le t\le1/4\), the separate fair-treatment laws \(P_{t,\sigma}\) and \(P_{*,t,\delta}\) have covariate law \(\lambda\), treatment probability \(1/2\), and control and treated outcome risks
\[
\begin{array}{c|cc}
&\mu_0(x)&\mu_1(x)\\ \hline
P_{t,\sigma}
&
\mathfrak p(t,\delta,u)+\delta\mathsf Z_{k,\sigma}(x)
&
\mathsf G\bigl(t,\mathfrak p(t,\delta,u)+\delta\mathsf Z_{k,\sigma}(x)\bigr)
\\
P_{*,t,\delta}
&
\mathfrak p(t,\delta,u)
&
\mathsf G\bigl(\mathsf T(t,\delta),\mathfrak p(t,\delta,u)\bigr).
\end{array}
\]
These formulas define probability laws on the calibration neighborhood where the conditional cells are positive.

The sample mixtures are
\[
Q_b
=
2^{-(k+1)}
\sum_{\sigma\in\{-1,1\}^{k+1}}
P_{b,\sigma}^{\otimes n},
\qquad
Q_t
=
2^{-(k+1)}
\sum_{\sigma\in\{-1,1\}^{k+1}}
P_{t,\sigma}^{\otimes n}.
\]
The signs are drawn once for the entire sample, and each mixture consists of exactly \(n\) original observed records.
\end{definitionv}


The first term depends on the sum of the two public smoothness exponents and reaches the parametric rate when that sum is at least \(1/2\). The second depends on prognosis smoothness and is multiplied by the evaluation radius. This decomposition will describe both the attainable expected length and the lower bound for every globally honest procedure.

We now collect the rate profile \leanref{sym:\rho}{\ensuremath{\rho_n(\alpha,\beta;r)}} and the attaining procedure. The interval \leanref{sym:I^{\mathrm{up}}}{\ensuremath{I^{\mathrm{up}}_{n,\mathsf E,\mathsf N}}} is the statistical output of the joint inversion in \cref{obj:def:upper-handle}; fixing the explicit engine and dyadic-floor policy selects the concrete sequence.

\begin{definitionv}{def:frontier-handle}[Rate profile $\rho$ and attaining interval family $I^\star$]
The rate profile and its attaining interval family are defined, for every admissible engine \(\mathsf E\) and admissible naming policy \(\mathsf N\), by
\[
\rho_n(\alpha,\beta;r)=d_n(\alpha,\beta;r),
\qquad
I^\star_{n,\mathsf E,\mathsf N}
=
I^{\mathrm{up}}_{n,\mathsf E,\mathsf N}.
\]
Here \(d_n(\alpha,\beta;r)\) is the diagnostic expected-length envelope from \cref{obj:def:diagnostic-envelope}, and \(I^{\mathrm{up}}_{\mathsf E,\mathsf N}\) is the procedure sequence from \cref{obj:def:upper-handle} for the fixed engine and policy. The frontier pair consists of this rate profile and procedure sequence:
\[
\bigl(\rho,I^{\mathrm{up}}_{\mathsf E,\mathsf N}\bigr).
\]
The concrete attaining sequence is
\[
I_n^\star
=
I^{\mathrm{up}}_{n,\mathsf E_0,\mathsf N_0},
\]
where \(\mathsf E_0\) is the explicitly constructed rational interval-arithmetic engine supplied by \cref{obj:lem:concrete-arithmetic-engine}, and \(\mathsf N_0\) is the dyadic-floor policy for selecting numerical representations in \cref{obj:def:dyadic-name-convention}.

Each interval is the literal output of the same fixed-budget expression tree, evaluated by its supplied fixed engine on its supplied valid numerical representations. Each choice of engine and representations determines its own rational output, subject to the endpoint enclosure bounds in \cref{obj:lem:fixed-budget-realization}. The construction uses the observable numerator and denominator coordinates \(C(P),S(P)\) from \cref{obj:lem:observable-odds}, separate projection resolutions from \cref{obj:def:resolutions}, and joint inversion from \cref{obj:def:upper-handle}.

The represented inputs are \((n,\alpha,\beta,\mathcal O)\), with numerical representations selected by \(\mathsf N\) and an independently fixed engine \(\mathsf E\) as in \cref{obj:def:arithmetic-engine}. The radius \(r\) indexes evaluation of the expected-length profile.

The converse uses the continuous calibrated priors in \cref{obj:def:continuous-calibrated-priors}, with admissible native propensity and prognosis perturbations, one common scalar effect per draw, analytic singleton matching, and factorization across disjoint latent-sign components conditional on the original covariates. It applies independently of the engine and naming policy.

The model and honesty conditions are specified in \cref{obj:synth_1,obj:synth_2,obj:synth_3,obj:ass:uniform-design,obj:ass:homogeneous-logit,obj:ass:effect-envelope,obj:ass:propensity-envelope,obj:ass:prognosis-envelope,obj:ass:propensity-holder,obj:ass:prognosis-holder,obj:ass:evaluation-radius,obj:ass:global-coverage,obj:def:logistic-class,obj:def:model,obj:def:radius-model,obj:def:honest-procedures,obj:def:length-objective}. The associated constructions are given in \cref{obj:def:causal-extension,obj:def:projection-statistic,obj:def:coordinate-estimates,obj:def:calibration-maps}. Their bounds, calibrations, and sharp comparisons are stated in \cref{obj:lem:projection-control,obj:thm:finite-honest-upper,obj:lem:parametric-null-floor,obj:thm:sharp-honest-frontier,obj:lem:scalar-implicit-function,obj:lem:exact-calibrations,obj:lem:calibrated-support,obj:lem:original-record-mixtures,obj:thm:full-honest-frontier-solution,obj:thm:full-honest-frontier-answer}.
\end{definitionv}

The evaluation radius indexes the expected-length criterion, while the procedure uses the observed sample and public exponents. Thus the same interval sequence can attain the comparison over all radius slices. The following result combines concrete attainment, the lower bound over the entire honest class, and the supporting guarantee for every admissible engine and naming policy.

\begin{theoremv}{thm:full-honest-frontier-answer}[The full honest frontier]
There exist constants \(c,C\in\mathbb R\), with \(c>0\) and \(C>0\), satisfying the following guarantees simultaneously. The public exponent domain and the sharp expected-length profile are
\[
\mathcal D=\{(\alpha,\beta)\in\mathbb R^2:0<\beta<1/4,\ \beta<\alpha\le1\},
\]
\[
\rho_n(\alpha,\beta;r)=d_n(\alpha,\beta;r)
=n^{-\min\{1/2,\,2(\alpha+\beta)/(2\alpha+2\beta+1)\}}
+r\,n^{-4\beta/(4\beta+1)}.
\]

\begin{itemize}
\item \textbf{(Concrete attainment.)}
For every \((\alpha,\beta)\in\mathcal D\) and every integer \(n\ge1\), the concrete interval
\[
I_n^\star(\alpha,\beta)
=I^{\mathrm{up}}_{n,\mathsf E_0,\mathsf N_0}(\alpha,\beta),
\]
constructed with the explicit rational arithmetic engine \(\mathsf E_0\) and the dyadic-floor naming policy \(\mathsf N_0\), belongs to the globally honest class \(\mathcal A_n(\alpha,\beta)\) of \cref{obj:def:honest-procedures}. For every sample and randomization seed, its output is a closed interval \([l,u]\) with
\[
l,u\in\mathbb Q,\qquad -1/2\le l\le u\le1/2.
\]
For every \(r\in[0,1/2]\), the minimax expected length \(J_n(\alpha,\beta;r)\) of \cref{obj:def:length-objective} satisfies
\[
c\,\rho_n(\alpha,\beta;r)
\le J_n(\alpha,\beta;r)
\le
\sup_{P\in\mathcal M_r(\alpha,\beta)}
E_{P^{\otimes n}\otimes\lambda}
\bigl|I_n^\star(\alpha,\beta)(\mathcal O,U)\bigr|
\le C\,\rho_n(\alpha,\beta;r),
\]
where \(\mathcal M_r(\alpha,\beta)\) is the prescribed radius slice and \(\lambda\) is the uniform law of the independent randomization seed.

\item \textbf{(Universal construction.)}
For every admissible rational interval-arithmetic engine \(\mathsf E\), every admissible Borel naming policy \(\mathsf N\), both satisfying the quantitative contracts of \cref{obj:def:arithmetic-engine,obj:def:dyadic-name-convention}, every \((\alpha,\beta)\in\mathcal D\), and every integer \(n\ge1\), the fixed-budget output
\(I^{\mathrm{up}}_{n,\mathsf E,\mathsf N}(\alpha,\beta)\)
is Borel, globally \(90\%\)-honest, and radius independent. On every valid named input, the construction terminates within its prescribed fixed budgets and returns a nonempty closed interval with ordered rational endpoints in \([-1/2,1/2]\), using the observed data and public inputs. For every \(r\in[0,1/2]\),
\[
c\,\rho_n(\alpha,\beta;r)
\le J_n(\alpha,\beta;r)
\le
\sup_{P\in\mathcal M_r(\alpha,\beta)}
E_{P^{\otimes n}\otimes\lambda}
\bigl|I^{\mathrm{up}}_{n,\mathsf E,\mathsf N}(\alpha,\beta)(\mathcal O,U)\bigr|
\le C\,\rho_n(\alpha,\beta;r).
\]
The same constants \(c,C\) apply to all these engines, policies, exponents, sample sizes, and radii.

\item \textbf{(Compact uniformity.)}
Taking the comparison constants \(c_\star=c\), \(C_\star=C\), and the sample-size threshold \(n_0=1\), every nonempty compact set \(K\subseteq\mathcal D\) satisfies
\[
0<\inf_{(\alpha,\beta)\in K}c_\star,\qquad
\sup_{(\alpha,\beta)\in K}C_\star<\infty,\qquad
\sup_{(\alpha,\beta)\in K}n_0=1.
\]

\item \textbf{(Elbow identities.)}
Define the numerator and denominator exponents by
\[
\mathsf a(\alpha,\beta)
=\min\{1/2,\,2(\alpha+\beta)/(2\alpha+2\beta+1)\},
\qquad
\mathsf b(\beta)=\frac{4\beta}{4\beta+1}.
\]
For every \((\alpha,\beta)\in\mathcal D\),
\[
\mathsf a(\alpha,\beta)-\mathsf b(\beta)>0,
\]
and
\[
\mathsf a(\alpha,\beta)-\mathsf b(\beta)
=
\begin{cases}
\displaystyle
\frac{2(\alpha-\beta)}
{(2\alpha+2\beta+1)(4\beta+1)},
&\alpha+\beta\le1/2,\\[6pt]
\displaystyle
\frac{1-4\beta}{2(4\beta+1)},
&\alpha+\beta\ge1/2.
\end{cases}
\]
At \(\alpha+\beta=1/2\), the two displayed expressions are equal.

\item \textbf{(Universal converse.)}
For every \((\alpha,\beta)\in\mathcal D\), every \(r\in[0,1/2]\), every integer \(n\ge1\), and every procedure \(I\in\mathcal A_n(\alpha,\beta)\), including procedures using independent randomization,
\[
c\,\rho_n(\alpha,\beta;r)
\le
\sup_{P\in\mathcal M_r(\alpha,\beta)}
E_{P^{\otimes n}\otimes\lambda}
\bigl|I(\mathcal O,U)\bigr|.
\]
\end{itemize}
\end{theoremv}

The coordinate identity behind the comparison is
\[
C(P)=\{\exp(\theta(P))-1\}S(P),
\qquad
\theta(P)=\log\{1+C(P)/S(P)\},
\qquad S(P)\ge1/512,
\]
as proved in \cref{obj:lem:observable-odds}. Numerator uncertainty therefore contributes directly to interval length, while denominator uncertainty is scaled by the effect radius. Separate projection resolutions accommodate these distinct sources of uncertainty. The matching lower comparison establishes that their sum describes the best possible worst-case expected length under ambient honesty, including for randomized procedures.

The converse has one experiment for each term. In the mixed family, propensity
and prognosis amplitudes \(\eta,\zeta\) produce an effect separation at their
product scale; singleton matching removes first-order marginal evidence, and
the original-record Hellinger bound determines the admissible projection
resolution. This yields the numerator contribution. In the fair-treatment
family, a prognosis amplitude \(\delta\) changes a comparator effect by order
\(t\delta^2\) at center \(t\), while its Hellinger distance is controlled at
fourth order in \(\delta\); taking \(t\) at the evaluation-radius scale yields
the radius-weighted denominator contribution. The continuous sign field shares
endpoint signs across neighboring cells. After conditioning on the observed
covariates, the sign-incidence graph is therefore factored into connected
components and controlled by a spanning-tree bound. The exact matching,
same-model support, and resulting mixture bounds are stated in
\cref{obj:lem:exact-calibrations,obj:lem:calibrated-support,obj:lem:original-record-mixtures};
\cref{obj:lem:parametric-null-floor,obj:thm:sharp-honest-frontier} complete the
two-term comparison.

At the null radius, \cref{obj:thm:full-honest-frontier-answer} gives
\[
c\,n^{-\mathsf a(\alpha,\beta)}
\le J_n(\alpha,\beta;0)
\le C\,n^{-\mathsf a(\alpha,\beta)}.
\]
Coverage remains imposed over the ambient model through \cref{obj:def:honest-procedures}, even though expected length is evaluated at zero effect. The null comparison therefore measures the precision attainable at zero by a procedure that also covers every effect in the prescribed ambient envelope. Its rate is governed by the numerator contribution.

For any sequence \(r_n\in[0,1/2]\), including an arbitrary shrinking-radius sequence, the simultaneous comparison gives
\[
c\left\{
n^{-\mathsf a(\alpha,\beta)}
+r_n n^{-\mathsf b(\beta)}
\right\}
\le J_n(\alpha,\beta;r_n)
\le
C\left\{
n^{-\mathsf a(\alpha,\beta)}
+r_n n^{-\mathsf b(\beta)}
\right\}.
\]
The same concrete interval sequence attains these bounds across the sequence. Equating the two contributions yields the radius balance
\[
r_n=n^{-\{\mathsf a(\alpha,\beta)-\mathsf b(\beta)\}}.
\]
For fixed public exponents, radii smaller than this scale make the numerator contribution dominant; radii larger than this scale make the denominator contribution dominant. At the balance scale, both contributions have the same order.

The numerator transition occurs at \(\alpha+\beta=1/2\). Below this boundary, its exponent is \(2(\alpha+\beta)/(2\alpha+2\beta+1)\), and the radius-balance exponent depends on the difference \(\alpha-\beta\). At and above the boundary, the numerator exponent equals \(1/2\), and the balance exponent depends on \(\beta\) alone. The two formulas agree at the shared boundary, so the radius balance is continuous across the numerator transition.

The propensity endpoint \(\alpha=1\) belongs to the public exponent domain in \cref{obj:thm:full-honest-frontier-answer}. At this endpoint, every admissible prognosis exponent places the numerator in the parametric regime, with balance exponent
\[
\mathsf a(1,\beta)-\mathsf b(\beta)
=
\frac{1-4\beta}{2(4\beta+1)}.
\]
Along the prognosis boundary approached from \(\beta<1/4\), this expression tends to zero, while the denominator exponent tends to \(1/2\). Within the stated domain, the positive exponent gap ensures that for every fixed \(r\in(0,1/2]\), the radius-scaled denominator contribution eventually dominates. Consequently, the sharp expected length at a fixed positive radius has order \(r n^{-\mathsf b(\beta)}\).

Concrete attainment is a guarantee for a statistical interval procedure with rational endpoints on every sample. The engine and naming-policy clause attaches the same coverage and expected-length comparison to every admissible numerical realization, with common constants. The following construction specifies the interval itself in \cref{obj:def:upper-handle,obj:thm:finite-honest-upper}, and \cref{obj:lem:concrete-arithmetic-engine,obj:lem:fixed-budget-realization} supplies its arithmetic realization.


\section{Constructing an honest interval}

The statistical ingredients used by the interval are defined first. The primitive arithmetic contracts and their finite rational realization appear in \cref{sec:rational-realization}. In reader-facing prose, \(d_n\) is the rate and \(I^{\mathrm{up}}\) is the interval; the formal layer also uses the aliases \(\rho_n=d_n\) and \(I^\star=I^{\mathrm{up}}\). It reuses \(b_C,b_S\) locally for resolution-box upper endpoints and, later, for coordinate error radii; we refer to these quantities by role outside the formal environments.

The attaining procedure combines separate estimates of two observable coordinates with a joint inversion of their confidence bounds. Write an observed record as \leanref{sym:O}{\ensuremath{O}}, with covariate \leanref{sym:X}{\ensuremath{X}}, binary treatment \leanref{sym:A}{\ensuremath{A}}, and binary outcome \leanref{sym:Y}{\ensuremath{Y}}. Under an observed law \leanref{sym:P}{\ensuremath{P}} in \cref{obj:def:model}, the numerator and denominator coordinates of \cref{obj:lem:observable-odds} have the representations
\[
C(P)
=
E_P[AY]
-
\int_0^1
E_P[A\mid X=x]E_P[Y\mid X=x]\,dx,
\]
\[
S(P)
=
\int_0^1
E_P[A(1-Y)\mid X=x]
E_P[(1-A)Y\mid X=x]\,dx.
\]
That result supplies the identity
\[
C(P)=\bigl(\exp\{\theta(P)\}-1\bigr)S(P),
\qquad
\theta(P)=\log\!\left(1+\frac{C(P)}{S(P)}\right),
\]
together with the uniform denominator lower bound \leanref{sym:s_0}{\ensuremath{s_0=1/512}}. These coordinates therefore reduce effect estimation to estimating products of conditional means.

The projection family in \cref{obj:synth_3} approximates those conditional means under the uniform covariate design. At integer rank \leanref{sym:k}{\ensuremath{k\ge1}}, its cosine projection and associated kernel satisfy
\[
\Pi_k f
=
\sum_{j=0}^{k-1}
\langle f,\varphi_j\rangle_{L^2(\lambda)}\varphi_j,
\qquad
H_k(x,z)
=
1+2\sum_{j=1}^{k-1}\cos(\pi jx)\cos(\pi jz).
\]
With this family fixed, the ordered-pair statistic of \cref{obj:def:projection-statistic} estimates a projected product using distinct records:
\[
\mathbb U_{n,k}(W,V)
=
\frac{1}{n(n-1)}
\sum_{\substack{1\le i,j\le n\\i\ne j}}
H_k(X_i,X_j)W(O_i)V(O_j).
\]
The coordinate estimates in \cref{obj:def:coordinate-estimates} apply this statistic to the marks \(A,Y\) for the numerator and \(A(1-Y),(1-A)Y\) for the denominator.

Separate resolutions accommodate the different projection errors in \cref{obj:lem:projection-control}. The numerator bias is controlled at order \(k^{-(\alpha+\beta)}\), whereas the denominator bias is controlled at order \(k^{-2\beta}\). Balancing each against the rank-dependent stochastic term gives the two target resolutions in \cref{obj:def:resolutions}. Their numerical selection uses the admissible arithmetic and naming conventions of \cref{obj:def:arithmetic-engine,obj:def:dyadic-name-convention}; \cref{obj:lem:concrete-arithmetic-engine} supplies an explicit admissible engine.

\paragraph{Rank selection.}
For \(n\ge2\), the numbered implementation of \cref{obj:def:resolutions} is as follows.
\begin{enumerate}
\item Set the public rank precision \leanref{sym:q_{\mathrm{rank}}}{\ensuremath{q_{\mathrm{rank}}(n)}} and the two target resolutions:
\[
q_{\mathrm{rank}}(n)=7+3\lceil\log_2 n\rceil,
\qquad
R_C=n^{2/(2\alpha+2\beta+1)},
\qquad
R_S=n^{2/(4\beta+1)}.
\]
\item Query the public exponent names at this precision, intersect their responses with the prescribed exponent boxes, and use the fixed engine to enclose the two target resolutions.
\item Retain the nonnegative integer ceilings of the enclosure upper endpoints as \(k_C\) and \(k_S\). For valid inputs in the stated exponent domain, \cref{obj:def:resolutions} gives
\[
R_C\le k_C\le3R_C,
\qquad
R_S\le k_S\le3R_S.
\]
These ranks are fixed before sampling and retained throughout endpoint evaluation.
\end{enumerate}
\paragraph{Interval construction.}
The observed sample \leanref{sym:\mathcal O}{\ensuremath{\mathcal O}} determines the two coordinate estimates. Bias allowances and variance bounds give a confidence rectangle for their population targets. Clipping its denominator endpoints to the known positive range makes the logarithmic inversion well defined. For each positive denominator, the inversion is increasing in the numerator; for each fixed numerator, its extrema over a denominator interval occur at the endpoints. Evaluating the appropriate corners therefore produces the ideal effect interval.

The upper interval \leanref{sym:I^{\mathrm{up}}}{\ensuremath{I^{\mathrm{up}}_{n,\mathsf E,\mathsf N}(\alpha,\beta)}} is the rational output of the following construction. Its exact-real expressions specify the quantities to be enclosed, while the endpoint precision \leanref{sym:q_{\mathrm{end}}}{\ensuremath{q_{\mathrm{end}}(n,k_C,k_S)}} governs their finite rational evaluation.

\begin{algorithmv}{def:upper-handle}[Upper interval $I^{\mathrm{up}}_{n,\mathsf E,\mathsf N}(\alpha,\beta)$]
The inputs are the sample size \(n\ge1\), public smoothness exponents \(\alpha,\beta\), and observed sample \(\mathcal O\); the tuning choices are an independently fixed admissible rational interval-arithmetic engine \(\mathsf E\) and a Borel policy \(\mathsf N\) selecting valid numerical representations of the exponents and observed covariates.
\begin{enumerate}
\item For \(n=1\), return
\[
I^{\mathrm{up}}_{1,\mathsf E,\mathsf N}(\alpha,\beta)
=[-1/2,1/2].
\]
For \(n\ge2\), proceed through the remaining steps.

\item Select the numerator and denominator projection ranks using \(\mathsf E\), the public exponent representations selected by \(\mathsf N\), and the fixed public precision budget \(q_{\mathrm{rank}}(n)\):
\[
k_C=k_C(n,\alpha,\beta;\mathsf E,\mathsf N),
\qquad
k_S=k_S(n,\alpha,\beta;\mathsf E,\mathsf N).
\]
Use these retained integers in the coordinate estimates \(\widehat C_n\), the numerator estimate, and \(\widehat S_n\), the denominator estimate, from \cref{obj:def:coordinate-estimates}, with ranks defined in \cref{obj:def:resolutions}.

\item Define the smoothness sum, variance envelope, and coordinate error radii by
\[
s=\alpha+\beta,\qquad
V_n(k)=\frac{10}{n}+\frac{4k}{n(n-1)},
\]
\[
b_C=8k_C^{-s}+\sqrt{20V_n(k_C)},
\qquad
b_S=25k_S^{-2\beta}+\sqrt{20V_n(k_S)}.
\]
Here \(b_C\) and \(b_S\) denote error radii, distinct from the resolution-enclosure endpoints in \cref{obj:def:resolutions}.

\item Define scalar clipping, the fixed denominator lower bound from \cref{obj:lem:observable-odds}, and the confidence-rectangle endpoints by
\[
\operatorname{clip}_{[l,u]}(v)=\max\{l,\min\{u,v\}\},
\qquad s_0=\frac{1}{512},
\]
\[
c_\pm=\widehat C_n\pm b_C,
\qquad
s_\pm=\operatorname{clip}_{[s_0,1]}(\widehat S_n\pm b_S).
\]
Define the clipped logarithmic inversion and its ideal interval endpoints by
\[
F(c,d)=
\log\!\left(
\operatorname{clip}_{[\exp(-1/2),\exp(1/2)]}(1+c/d)
\right),
\]
\[
L=\min\{F(c_-,s_-),F(c_-,s_+)\},
\qquad
R=\max\{F(c_+,s_-),F(c_+,s_+)\}.
\]
Here \(R\) denotes the ideal upper endpoint, distinct from the generic rank target in \cref{obj:def:resolutions}. These exact-real expressions specify the quantities enclosed in the next step.

\item Define the largest retained rank and the endpoint precision budget by
\[
K=\max\{k_C,k_S\},
\qquad
q_{\mathrm{end}}(n,k_C,k_S)
=44+\lceil\log_2(n+1)\rceil
+4\lceil\log_2(K+1)\rceil.
\]
At this single precision, query the supplied numerical representations of both exponents and all observed covariates. Evaluate the displayed finite expressions with the same engine \(\mathsf E\), its primitive maps at precision \(q_{\mathrm{end}}\), and exact rational range operations in the displayed expression tree, retaining \(k_C,k_S\) from the rank step. Obtain rational endpoint enclosures satisfying
\[
L_-\le L\le L_+,
\qquad
R_-\le R\le R_+.
\]
For every valid input with \(0<\beta<1/4\), \(\beta<\alpha\le1\), and \(n\ge2\), \cref{obj:lem:fixed-budget-realization} gives
\[
L_+-L_-\le\frac1n,
\qquad
R_+-R_-\le\frac1n.
\]

\item Output the closed rational interval
\[
I^{\mathrm{up}}_{n,\mathsf E,\mathsf N}(\alpha,\beta)
=
\left[
\max\{-1/2,L_-\},
\min\{1/2,R_+\}
\right].
\]
For every valid input with \(0<\beta<1/4\), \(\beta<\alpha\le1\), and \(n\ge2\), the output is nonempty, contains the ideal interval, and satisfies
\[
[L,R]\subseteq
I^{\mathrm{up}}_{n,\mathsf E,\mathsf N}(\alpha,\beta),
\]
\[
0\le
\min\{1/2,R_+\}-\max\{-1/2,L_-\}-(R-L)
\le\frac2n.
\]
The computation terminates after the fixed rank and endpoint precision budgets.
\end{enumerate}
For fixed \(\mathsf E,\mathsf N\), this construction defines a sample map with inputs \((n,\alpha,\beta,\mathcal O)\). Its rational endpoints may depend on the admissible engine and valid numerical representations supplied.

The model and honesty conditions are specified in \cref{obj:synth_1,obj:synth_2,obj:synth_3,obj:ass:uniform-design,obj:ass:homogeneous-logit,obj:ass:effect-envelope,obj:ass:propensity-envelope,obj:ass:prognosis-envelope,obj:ass:propensity-holder,obj:ass:prognosis-holder,obj:ass:evaluation-radius,obj:ass:global-coverage,obj:def:logistic-class,obj:def:model,obj:def:radius-model,obj:def:honest-procedures,obj:def:length-objective}. The associated constructions are given in \cref{obj:def:causal-extension,obj:def:projection-statistic,obj:def:diagnostic-envelope,obj:def:frontier-handle,obj:def:calibration-maps,obj:def:continuous-calibrated-priors,obj:def:dyadic-name-convention,obj:def:arithmetic-engine}. Their bounds and realization properties are stated in \cref{obj:lem:projection-control,obj:thm:finite-honest-upper,obj:lem:parametric-null-floor,obj:thm:sharp-honest-frontier,obj:lem:scalar-implicit-function,obj:lem:exact-calibrations,obj:lem:calibrated-support,obj:lem:original-record-mixtures,obj:thm:full-honest-frontier-solution,obj:lem:concrete-arithmetic-engine,obj:thm:full-honest-frontier-answer}.
\end{algorithmv}

The exact endpoints \(L,R\) describe the interval used in the statistical analysis. On valid named inputs, the delivered endpoints are rational outward enclosures, clipped to the effect range, with total additional length at most \(2/n\). The realization guarantee in \cref{obj:lem:fixed-budget-realization} applies at the single prescribed endpoint precision. Different admissible engines or names may produce different rational endpoints within these guarantees.

Every tuning choice uses the sample size and public exponents, so the same procedure serves all evaluation radii. Clipping uses the known denominator and effect ranges from \cref{obj:lem:observable-odds,obj:ass:effect-envelope}. At sample size one, the full effect interval supplies the specified output. The following result establishes global honesty and the effect-sensitive expected-length bound for this finite-sample procedure.

\begin{theoremv}{thm:finite-honest-upper}[Uniform honesty and length bound]
There exists an absolute constant \(K_0>0\) such that the following guarantees hold for every choice satisfying:
\begin{itemize}
\item \textbf{(Arithmetic engine.)} \(\mathsf E\) is an admissible, independently fixed deterministic rational interval-arithmetic engine satisfying the quantitative contracts of \cref{obj:def:arithmetic-engine}.
\item \textbf{(Naming policy.)} \(\mathsf N\) is an admissible Borel policy satisfying \cref{obj:def:dyadic-name-convention}, selecting valid names for the public exponents and observed covariates, with public exponent names fixed independently of the sample.
\item \textbf{(Public exponents.)} The real exponents \(\alpha,\beta\) satisfy
\[
0<\beta<\frac14,\qquad \beta<\alpha\le1.
\]
\item \textbf{(Sample size.)} The integer sample size \(n\) satisfies \(n\ge1\).
\end{itemize}

Let \(I^{\mathrm{up}}_{n,\mathsf E,\mathsf N}(\alpha,\beta)\) denote the measurable interval procedure given by the literal rational output of the upper construction, whose inputs are the engine, naming policy, sample size, public exponents, and observed records. This procedure satisfies
\[
I^{\mathrm{up}}_{n,\mathsf E,\mathsf N}(\alpha,\beta)
\in\mathcal A_n(\alpha,\beta),
\]
where \(\mathcal A_n(\alpha,\beta)\) is the class of globally honest procedures in \cref{obj:def:honest-procedures}.

For every \(r\in[0,1/2]\) and every observed law \(P\in\mathcal M_r(\alpha,\beta)\), the effect-radius model of \cref{obj:def:radius-model},
\[
E_{P^{\otimes n}\otimes\lambda}
\left[
\left|I^{\mathrm{up}}_{n,\mathsf E,\mathsf N}(\alpha,\beta)(\mathcal O,U)\right|
\right]
\le K_0\,d_n(\alpha,\beta;r).
\]
Here \(\mathcal O\) is the sample of original observed records, \(U\) is the independent randomization seed with uniform law \(\lambda\) on the unit interval, and \(d_n(\alpha,\beta;r)\) is the diagnostic envelope of \cref{obj:def:diagnostic-envelope}. The same constant \(K_0\) applies to all the engines, naming policies, exponents, sample sizes, radii, and laws quantified above.

For every sample and seed, there exist rational endpoints \(l,u\) such that the output is the closed interval \([l,u]\) and
\[
-\frac12\le l\le u\le\frac12.
\]
When \(n=1\), every sample and seed yields the full interval \([-1/2,1/2]\).
\end{theoremv}

The length bound reflects the structure of the odds identity. Numerator uncertainty contributes even at a zero effect, while denominator uncertainty enters the inversion through a numerator proportional to \(\exp\{\theta(P)\}-1\). This explains the numerator term and the radius-scaled denominator term in \cref{obj:def:diagnostic-envelope}. Separate ranks accommodate their distinct smoothness requirements, and outward numerical enclosure preserves the ideal interval. Consequently, the literal rational output has global \(90\%\) coverage over the ambient model and the stated unconditional expected-length bound on every radius slice, with one common constant across the admissible arithmetic engines and naming policies.


\section{Interpretation and limitations}

The sharp expected-length characterization in \cref{obj:thm:full-honest-frontier-answer} separates two sources of uncertainty in the homogeneous causal log odds. The numerator contribution depends on the combined smoothness of the propensity and prognosis logits, whereas the denominator contribution depends on prognosis smoothness and is multiplied by the evaluation radius. Through \cref{obj:def:diagnostic-envelope}, these contributions give
\[
d_n(\alpha,\beta;r)
=
n^{-\min\{1/2,\,2(\alpha+\beta)/(2\alpha+2\beta+1)\}}
+
r\,n^{-4\beta/(4\beta+1)}.
\]
The sharp comparison gives this decomposition an optimality interpretation under its stated conditions: it describes the smallest worst-case expected length on each radius slice among procedures satisfying ambient honesty. The construction in \cref{obj:thm:finite-honest-upper} attains the corresponding upper bound with a single procedure whose tuning uses the sample size and public smoothness exponents.

The distinction between coverage and length evaluation is essential to this interpretation. In \cref{obj:def:length-objective}, the evaluation radius restricts the laws over which expected length is assessed, while admissibility continues to require global \(90\%\) coverage over the ambient model in \cref{obj:def:honest-procedures}. Thus, as the evaluation radius approaches zero, the procedure remains accountable for coverage throughout that model. At radius zero, the denominator contribution vanishes and the numerator contribution remains. The sharp comparison therefore quantifies the precision available at a zero effect while retaining the same coverage obligation for nearby and larger effects.

This behavior follows the structure of the observable odds identity in \cref{obj:lem:observable-odds}. The numerator is proportional to the denominator times the exponential effect increment. Consequently, denominator uncertainty has less influence on the recovered log odds when the effect is small, while numerator uncertainty continues to enter the inversion. The separate coordinate estimates and ranks in \cref{obj:def:coordinate-estimates,obj:def:resolutions} accommodate these different roles. Joint inversion then converts the coordinate bounds into an interval whose expected length responds to effect magnitude.

The balance can be expressed using the numerator exponent \(\mathsf a(\alpha,\beta)\) and the denominator exponent \(\mathsf b(\beta)\) introduced in \cref{obj:def:diagnostic-envelope}. The two contributions are comparable when
\[
r\asymp n^{-\{\mathsf a(\alpha,\beta)-\mathsf b(\beta)\}}.
\]
When \(\mathsf b(\beta)<\mathsf a(\alpha,\beta)\), this scale decreases with sample size. Above it, the effect-scaled denominator contribution is larger; below it, the numerator contribution is larger. Rough prognosis can therefore govern expected length at appreciable effect magnitudes even when combined propensity and prognosis smoothness gives a faster numerator rate. Concentrating evaluation near zero progressively reduces that contribution.

The known uniform covariate design and public exponents make this experiment precise. The design in \cref{obj:ass:uniform-design} supplies the integration measure for the projection construction, fixing the geometry used to control approximation error. The public exponents specify the smoothness restrictions in \cref{obj:def:model} and determine the two resolutions before sampling. Together with the prescribed envelopes, these features support the finite-sample coverage guarantee and the comparison of expected lengths across effect radii. The rational realization in \cref{obj:lem:fixed-budget-realization} also gives a concrete numerical meaning to the construction: prescribed precision budgets produce outward endpoint enclosures with the stated additional-length guarantee.

\subsection{Limitations and future work}

The precision characterization concerns a known uniform design, public smoothness exponents, a scalar covariate, and a homogeneous conditional log-odds coefficient. Unknown design would require accounting for the covariate distribution in the projection construction. Unknown smoothness raises a separate adaptation problem, including the coverage requirements appropriate to a collection of smoothness classes. Higher-dimensional covariates would change the approximation problem and its balance with sampling uncertainty.

Heterogeneous effects would require specifying a new target, such as an average contrast or an effect function, together with an appropriate coverage criterion. The scalar inversion and radius-based length objective studied here provide a starting point for formulating such questions; their extensions require separate precision characterizations.

Practical calibration and computational optimization also merit further study. The proved construction supplies coverage and expected-length guarantees with prescribed arithmetic precision. Empirical assessment could examine interval lengths at realistic sample sizes, sensitivity to conservative constants, and numerical performance across admissible implementations. Software development could optimize evaluation of the projection statistics and endpoint enclosures while preserving the stated guarantees. These prospective assessments concern practical performance beyond the proved precision characterization.


\appendix

\section*{Appendices}

\section{Identification and projection bounds}

The causal full-data extension used for interpretation is the following.

The statistical construction uses observable product coordinates to recover the homogeneous conditional log-odds coefficient. Its causal interpretation follows from the potential-outcome coupling in \cref{obj:def:causal-extension}. An admissible coupling for an observed law \(P\) is a Markov kernel \(\Gamma\) from \([0,1]\) to \(\{0,1\}^2\) whose first and second margins, at every covariate value, are the observed control-arm and treated-arm Bernoulli outcome laws. This specifies the conditional joint distribution of the two potential outcomes. The full-data construction then draws the covariate from the uniform law \leanref{sym:\lambda}{\ensuremath{\lambda}}, draws the potential-outcome pair from this coupling, and draws treatment using the observed conditional treatment probability. The following statement gives the observable coordinate identity and the properties of this construction.

\begin{lemmav}{lem:observable-odds}[Observable odds and causal interpretation]
Let \(0<\beta<1/4\) and \(\beta<\alpha\leq1\), and let \(P\in\mathcal M(\alpha,\beta)\), the model class in \cref{obj:def:model}. Write \(\lambda\) for the uniform probability law on \([0,1]\), and define the observable coordinates
\[
C(P)
=
E_P[AY]
-
\int_{[0,1]} e(P,x)\,E_P[Y\mid X=x]\,\lambda(dx),
\qquad
S(P)
=
\int_{[0,1]}
P(A=1,Y=0\mid X=x)\,
P(A=0,Y=1\mid X=x)\,\lambda(dx),
\]
where \(e(P,x)\) is the conditional treatment probability. With
\[
\theta(P)=\operatorname{logit}\mu_1(P,0)-\nu(P,0),
\qquad s_0=\frac{1}{512},
\]
where \(\mu_1(P,0)\) is the treated-arm risk at zero and \(\nu(P,0)\) is the control-arm prognosis logit at zero, one has
\[
C(P)=\{\exp(\theta(P))-1\}S(P),
\qquad
S(P)\geq s_0,
\qquad
\theta(P)=\log\{1+C(P)/S(P)\}.
\]

For every Markov kernel \(\Gamma\) from \([0,1]\) to \(\{0,1\}^2\) whose first and second margins at every \(x\) are Bernoulli laws with means
\[
P(Y=1\mid A=0,X=x)
\quad\text{and}\quad
P(Y=1\mid A=1,X=x),
\]
respectively, the full-data law \(\widetilde P_{P,\Gamma}\) of \cref{obj:def:causal-extension} has the following properties:
\begin{itemize}
\item \textbf{(Observed law and consistency.)}
The image of \(\widetilde P_{P,\Gamma}\) under
\[
((x,(y^0,y^1)),a)\longmapsto
\bigl(x,a,(1-a)y^0+ay^1\bigr)
\]
is \(P\). For every full-data record,
\[
Y=(1-A)Y^0+AY^1.
\]

\item \textbf{(Construction and conditional exchangeability.)}
The law \(\widetilde P_{P,\Gamma}\) is the successive kernel composition obtained by drawing \(X\) from \(\lambda\), drawing \((Y^0,Y^1)\) from \(\Gamma(P,X)\), and drawing \(A\) from the observed-law treatment kernel with Bernoulli mean \(e(P,X)\). Under its full conditional law at every \(x\in[0,1]\), the pair \((Y^0,Y^1)\) and \(A\) are independent.

\item \textbf{(Conditional causal log-odds.)}
For every \(x\in[0,1]\),
\[
\operatorname{logit}
\left(\int_{\{0,1\}^2} y^1\,\Gamma(P,x)(d(y^0,y^1))\right)
-
\operatorname{logit}
\left(\int_{\{0,1\}^2} y^0\,\Gamma(P,x)(d(y^0,y^1))\right)
=
\theta(P).
\]

\item \textbf{(Full conditional disintegration.)}
Composing \(\lambda\) with the full conditional laws of \(((Y^0,Y^1),A)\) at \(x\), then taking the image under the canonical reassociation
\[
\bigl(x,((y^0,y^1),a)\bigr)
\longmapsto
\bigl((x,(y^0,y^1)),a\bigr),
\]
gives exactly \(\widetilde P_{P,\Gamma}\).
\end{itemize}
\end{lemmav}

\begin{proof}[Proof of \cref{obj:lem:observable-odds}]
Fix \(P\in\mathcal M(\alpha,\beta)\).

\begin{enumerate}
\item Define, for \(a,y\in\{0,1\}\) and \(x\in[0,1]\),
\[
p_{ay}(P,x)=P(A=a,Y=y\mid X=x),\qquad
\mu_a(P,x)=\frac{p_{a1}(P,x)}{p_{a0}(P,x)+p_{a1}(P,x)},
\]
and
\[
m_{\mathrm{obs}}(P,x)=E_P[Y\mid X=x].
\]
The strictly positive, normalized conditional cells give
\[
0<e(P,x)<1,\qquad 0<\mu_a(P,x)<1,
\]
and
\[
\begin{aligned}
e(P,x)&=p_{10}(P,x)+p_{11}(P,x),\\
m_{\mathrm{obs}}(P,x)&=p_{01}(P,x)+p_{11}(P,x),\\
p_{11}(P,x)&=e(P,x)\mu_1(P,x),&
p_{10}(P,x)&=e(P,x)\{1-\mu_1(P,x)\},\\
p_{01}(P,x)&=\{1-e(P,x)\}\mu_0(P,x),&
p_{00}(P,x)&=\{1-e(P,x)\}\{1-\mu_0(P,x)\}.
\end{aligned}
\]
Indeed, each arm risk is the positive success cell divided by its arm total; normalization identifies the control-arm total with \(1-e(P,x)\). Substitution gives the four factorizations.

Normalization also gives
\[
p_{00}(P,x)=1-p_{11}(P,x)-p_{10}(P,x)-p_{01}(P,x).
\]
Expanding the product of the two margins therefore yields
\[
p_{11}(P,x)-e(P,x)m_{\mathrm{obs}}(P,x)
=p_{11}(P,x)p_{00}(P,x)-p_{10}(P,x)p_{01}(P,x).
\]

For the logistic map, write
\[
\ell(\vartheta)=\frac{1}{1+\exp(-\vartheta)}
\qquad(\vartheta\in\mathbb R).
\]
Its odds and log-odds satisfy
\[
\frac{\ell(\vartheta)}{1-\ell(\vartheta)}=\exp(\vartheta),
\qquad
\operatorname{logit}\ell(\vartheta)=\vartheta.
\]
Since \(\nu(P,x)=\operatorname{logit}\mu_0(P,x)\), the homogeneous relation supplied by \cref{obj:ass:homogeneous-logit} gives
\[
\operatorname{logit}\mu_1(P,x)
-\operatorname{logit}\mu_0(P,x)=\theta(P).
\]
Exponentiating this identity at the interior arm risks gives
\[
\exp\{\theta(P)\}
=
\frac{\mu_1(P,x)/\{1-\mu_1(P,x)\}}
     {\mu_0(P,x)/\{1-\mu_0(P,x)\}}.
\]
Using the four cell factorizations and clearing the positive denominators,
\[
p_{11}(P,x)p_{00}(P,x)
=\exp\{\theta(P)\}\,p_{10}(P,x)p_{01}(P,x).
\]

The cell disintegration and the uniform design in \cref{obj:ass:uniform-design} imply
\[
E_P[AY]=\int_{[0,1]}p_{11}(P,x)\,\lambda(dx).
\]
Here the binary product \(AY\) selects exactly the cell \(A=Y=1\). Continuity on the compact covariate interval ensures integrability of the cell functions and their products. Consequently,
\[
\begin{aligned}
C(P)
&=\int_{[0,1]}
 \{p_{11}(P,x)-e(P,x)m_{\mathrm{obs}}(P,x)\}\,\lambda(dx)\\
&=\int_{[0,1]}
 \{p_{11}(P,x)p_{00}(P,x)-p_{10}(P,x)p_{01}(P,x)\}\,\lambda(dx)\\
&=\{\exp(\theta(P))-1\}
  \int_{[0,1]}p_{10}(P,x)p_{01}(P,x)\,\lambda(dx)\\
&=\{\exp(\theta(P))-1\}S(P).
\end{aligned}
\]
% lean: CausalSmith.Stat.LogoddsLowsmoothFrontier.covarianceCoordinate_eq_odds

\item We next establish the numerical denominator floor. The numerical exponential bound gives
\[
\exp(1)<2.7182818286<\frac{2719}{1000}<\frac{11}{4}.
\]
% lean: Real.exp_one_lt_d9
Moreover,
\[
\exp(3/2)^2=\exp(1)^3
<\left(\frac{2719}{1000}\right)^3
<\frac{81}{4}.
\]
Positivity of the exponential implies \(\exp(3/2)<9/2\). Hence
\[
\ell(-1)=\frac{1}{1+\exp(1)}\ge\frac4{15},
\qquad
\ell(-3/2)=\frac{1}{1+\exp(3/2)}\ge\frac2{11}.
\]

The bounds in
\cref{obj:ass:propensity-envelope,obj:ass:prognosis-envelope,obj:ass:effect-envelope}
supply, at every \(x\),
\[
-1\le g(P,x)\le1,\qquad
-1\le\nu(P,x)\le1,\qquad
-\frac12\le\theta(P)\le\frac12.
\]
Inverting the native logits at the interior probabilities gives
\[
e(P,x)=\ell\{g(P,x)\},\qquad
\mu_0(P,x)=\ell\{\nu(P,x)\}.
\]
The logistic map is increasing and satisfies
\[
1-\ell(\vartheta)=\ell(-\vartheta)
\qquad(\vartheta\in\mathbb R).
\]
Together with the homogeneous relation, these facts yield
\[
\begin{aligned}
e(P,x)&\ge\ell(-1)\ge\frac4{15},\\
1-e(P,x)&\ge\ell(-1)\ge\frac4{15},\\
\mu_0(P,x)&\ge\ell(-1)\ge\frac4{15}\ge\frac2{11},\\
1-\mu_1(P,x)
&=\ell\{-\nu(P,x)-\theta(P)\}
 \ge\ell(-3/2)\ge\frac2{11}.
\end{aligned}
\]
Multiplying these nonnegative lower bounds,
\[
\begin{aligned}
p_{10}(P,x)p_{01}(P,x)
&=e(P,x)\{1-\mu_1(P,x)\}\{1-e(P,x)\}\mu_0(P,x)\\
&\ge\frac4{15}\frac2{11}\frac4{15}\frac2{11}
=\frac{64}{27225}
>\frac1{512}=s_0.
\end{aligned}
\]
Since \(\lambda\) is a probability law, integration gives
\[
S(P)\ge\int_{[0,1]}s_0\,\lambda(dx)=s_0.
\]
% lean: CausalSmith.Stat.LogoddsLowsmoothFrontier.denominatorCoordinate_lower

\item In particular, \(S(P)\ge s_0>0\), so the coordinate identity permits division:
\[
1+\frac{C(P)}{S(P)}
=1+\exp\{\theta(P)\}-1
=\exp\{\theta(P)\}>0.
\]
Applying the logarithm proves
\[
\theta(P)=\log\{1+C(P)/S(P)\}.
\]
% lean: CausalSmith.Stat.LogoddsLowsmoothFrontier.effect_eq_log_of_coordinates

\item Fix an admissible kernel \(\Gamma\). Define its full conditional law by
\[
\mathcal K_{P,\Gamma}(x)
=\Gamma(P,x)\otimes\operatorname{Bernoulli}\{e(P,x)\},
\qquad x\in[0,1],
\]
and define the reassociation map by
\[
\begin{aligned}
\mathsf r_{\mathrm{assoc}}:\;
[0,1]\times(\{0,1\}^2\times\{0,1\})
&\longrightarrow([0,1]\times\{0,1\}^2)\times\{0,1\},\\
\mathsf r_{\mathrm{assoc}}\bigl(x,((y^0,y^1),a)\bigr)
&=\bigl((x,(y^0,y^1)),a\bigr).
\end{aligned}
\]
For fixed \(x\), the treatment kernel assigns
\(\operatorname{Bernoulli}\{e(P,x)\}\) to every potential-outcome pair. Its composition with \(\Gamma(P,x)\) is therefore precisely \(\mathcal K_{P,\Gamma}(x)\). This also makes \(x\mapsto\mathcal K_{P,\Gamma}(x)\) a measurable probability kernel.

Associativity of successive kernel composition, applied to the construction in \cref{obj:def:causal-extension}, now gives
\[
\widetilde P_{P,\Gamma}
=(\mathsf r_{\mathrm{assoc}})_\#
 \bigl\{\lambda(dx)\,
 \mathcal K_{P,\Gamma}(x)(d((y^0,y^1),a))\bigr\}.
\]
Here the measure inside braces is the joint law obtained by drawing \(x\) from \(\lambda\) and then \(((y^0,y^1),a)\) from \(\mathcal K_{P,\Gamma}(x)\), and the subscript \(\#\) denotes pushforward. This proves the asserted full conditional disintegration.
% lean: CausalSmith.Stat.LogoddsLowsmoothFrontier.observable_odds

\item Define the observation map by
\[
\begin{aligned}
\mathsf o_{\mathrm{obs}}:\;
([0,1]\times\{0,1\}^2)\times\{0,1\}
&\longrightarrow[0,1]\times\{0,1\}^2,\\
\mathsf o_{\mathrm{obs}}\bigl((x,(y^0,y^1)),a\bigr)
&=\bigl(x,a,(1-a)y^0+ay^1\bigr).
\end{aligned}
\]
Let
\[
\varphi:[0,1]\times\{0,1\}^2\longrightarrow[0,\infty]
\]
be any nonnegative measurable test function. At each fixed \(x\), expand the treatment Bernoulli law into its two atoms. The first margin of \(\Gamma(P,x)\) handles \(a=0\), and the second handles \(a=1\). Thus
\[
\begin{aligned}
&\int_{\{0,1\}^2}\int_{\{0,1\}}
 \varphi\bigl(x,a,(1-a)y^0+ay^1\bigr)\,
 \operatorname{Bernoulli}\{e(P,x)\}(da)\,
 \Gamma(P,x)(d(y^0,y^1))\\
&\quad=\{1-e(P,x)\}\int_{\{0,1\}^2}
 \varphi(x,0,y^0)\,\Gamma(P,x)(d(y^0,y^1))\\
&\qquad\quad
 +e(P,x)\int_{\{0,1\}^2}
 \varphi(x,1,y^1)\,\Gamma(P,x)(d(y^0,y^1))\\
&\quad=\{1-e(P,x)\}
 \bigl[\{1-\mu_0(P,x)\}\varphi(x,0,0)
       +\mu_0(P,x)\varphi(x,0,1)\bigr]\\
&\qquad\quad
 +e(P,x)
 \bigl[\{1-\mu_1(P,x)\}\varphi(x,1,0)
       +\mu_1(P,x)\varphi(x,1,1)\bigr]\\
&\quad=\sum_{a=0}^1\sum_{y=0}^1p_{ay}(P,x)\varphi(x,a,y).
\end{aligned}
\]
Integrating over \(\lambda\), successive kernel integration and the observed-law cell disintegration give
\[
\int\varphi\circ\mathsf o_{\mathrm{obs}}\,
 d\widetilde P_{P,\Gamma}
=
\int_{[0,1]}\sum_{a=0}^1\sum_{y=0}^1
 p_{ay}(P,x)\varphi(x,a,y)\,\lambda(dx)
=
\int\varphi\,dP.
\]
Equality for all such tests proves
\[
(\mathsf o_{\mathrm{obs}})_\#\widetilde P_{P,\Gamma}=P.
\]
% lean: CausalSmith.Stat.LogoddsLowsmoothFrontier.causalExtension_observed_marginal

\item For a full-data record, the two treatment cases give
\[
Y=
\begin{cases}
Y^0,&A=0,\\
Y^1,&A=1,
\end{cases}
\qquad\text{and hence}\qquad
Y=(1-A)Y^0+AY^1.
\]
The successive-draw construction is exactly that of \cref{obj:def:causal-extension}. At every \(x\), its full conditional law is the product probability measure
\[
\mathcal K_{P,\Gamma}(x)
=\Gamma(P,x)\otimes\operatorname{Bernoulli}\{e(P,x)\}.
\]
The coordinate projections of a product probability measure are independent, so
\[
(Y^0,Y^1)\ \perp\ A\mid X=x
\qquad(x\in[0,1]).
\]
% lean: CausalSmith.Stat.LogoddsLowsmoothFrontier.observedMap_consistency
% lean: CausalSmith.Stat.LogoddsLowsmoothFrontier.fullConditionalLaw_independent

\item Finally, pushing the coordinate integrals through the prescribed Bernoulli margins yields
\[
\begin{aligned}
\int_{\{0,1\}^2}y^0\,\Gamma(P,x)(d(y^0,y^1))
&=0\{1-\mu_0(P,x)\}+1\mu_0(P,x)=\mu_0(P,x),\\
\int_{\{0,1\}^2}y^1\,\Gamma(P,x)(d(y^0,y^1))
&=0\{1-\mu_1(P,x)\}+1\mu_1(P,x)=\mu_1(P,x).
\end{aligned}
\]
Both means are interior. Substitution into the homogeneous log-odds identity established above gives, at every \(x\in[0,1]\),
\[
\operatorname{logit}
 \left(\int_{\{0,1\}^2}y^1\,\Gamma(P,x)(d(y^0,y^1))\right)
-
\operatorname{logit}
 \left(\int_{\{0,1\}^2}y^0\,\Gamma(P,x)(d(y^0,y^1))\right)
=\theta(P).
\]
% lean: CausalSmith.Stat.LogoddsLowsmoothFrontier.admissibleCoupling_means
% lean: CausalSmith.Stat.LogoddsLowsmoothFrontier.homogeneous_logit_difference
\end{enumerate}
\end{proof}

The coordinate identity in \cref{obj:lem:observable-odds} identifies the effect through an observable ratio, with a denominator bounded uniformly away from zero over the model. Its causal clauses establish that every admissible coupling produces the observed law and the same conditional causal log-odds coefficient. Conditional exchangeability is a property of the specified full-data construction, while the marginal restrictions ensure that its potential-outcome risks equal the observed arm risks.

To control the coordinate estimates in \cref{obj:def:coordinate-estimates}, retain the rank-\(k\) cosine projection \(\Pi_k\) and ordered-pair statistic from \cref{obj:def:projection-statistic}. The approximation bound below concerns a continuous function \leanref{sym:f}{\ensuremath{f}} at Hölder exponent \leanref{sym:\gamma}{\ensuremath{\gamma}}, with Hölder seminorm \leanref{sym:[\cdot]_\gamma}{\ensuremath{[f]_\gamma}}. For two bounded record marks, let \leanref{sym:w}{\ensuremath{w}} and \leanref{sym:v}{\ensuremath{v}} denote their conditional means. The next statement connects approximation of these means to the bias and variance of the product statistic.

\begin{lemmav}{lem:projection-control}[Projection approximation and moments]
Let \(P\) be an observed law with continuous, strictly interior conditional four-cell probabilities, satisfying \cref{obj:ass:uniform-design}. Let \(k\ge1\) be an integer, and let \(\Pi_k\) denote the rank-\(k\) cosine projection used in \cref{obj:def:projection-statistic}.

For every \(0<\gamma\le1\) and every continuous function \(f:[0,1]\to\mathbb R\),
\[
\|f-\Pi_k f\|_{L^2(\lambda)}
\le 5[f]_\gamma k^{-\gamma},
\qquad
[f]_\gamma
=
\sup_{\substack{x,z\in[0,1]\\x\ne z}}
\frac{|f(x)-f(z)|}{|x-z|^\gamma}.
\]
Here \(\lambda\) is the uniform probability law on \([0,1]\), and the norm and Hölder seminorm are interpreted in the extended nonnegative reals, allowing \([f]_\gamma=+\infty\).

For every pair of measurable record marks \(W,V\) taking values in \([0,1]\), and every integer \(n\ge2\), define their conditional means using the four-cell conditional law under \(P\):
\[
w(x)=E_P[W\mid X=x],
\qquad
v(x)=E_P[V\mid X=x].
\]
For the ordered-pair statistic \(\mathbb U_{n,k}(W,V)\) of \cref{obj:def:projection-statistic}, evaluated on \(n\) independent original records with law \(P\),
\[
E_P\mathbb U_{n,k}(W,V)
=
\langle\Pi_k w,\Pi_k v\rangle_{L^2(\lambda)},
\]
\[
\langle w,v\rangle_{L^2(\lambda)}
-
E_P\mathbb U_{n,k}(W,V)
=
\langle w-\Pi_k w,v-\Pi_k v\rangle_{L^2(\lambda)},
\]
and
\[
\operatorname{Var}_P\{\mathbb U_{n,k}(W,V)\}
\le
\frac4n+\frac{2k}{n(n-1)}.
\]
\end{lemmav}

\begin{proof}[Proof of \cref{obj:lem:projection-control}]
Write \(\|\cdot\|_2\) and \(\langle\cdot,\cdot\rangle\) for the norm and inner product in \(L^2(\lambda)\). We use the cosine basis and projection kernel of \cref{obj:synth_3}.

\begin{enumerate}
\item Fix \(0<\gamma\le1\), a continuous function \(f:[0,1]\to\mathbb R\), and an integer \(k\ge1\). If \([f]_\gamma=+\infty\), then
\[
5[f]_\gamma k^{-\gamma}=+\infty,
\]
because \(k^{-\gamma}>0\), and the extended-real approximation inequality follows. Otherwise define
\[
L_f=[f]_\gamma\in[0,\infty).
\]
The definition of the seminorm gives
\[
|f(x)-f(z)|\le L_f|x-z|^\gamma
\qquad(x,z\in[0,1]).
\]
For \(x\ne z\), this follows by multiplying the defining quotient bound by the positive denominator; for \(x=z\), both sides vanish.
% lean: Causalean.Mathlib.Analysis.Approximation.Chebyshev.CosineProjection.cosineProjection_error_le_seminorm

\item Define the folding map and the even extension by
\[
\chi(\vartheta)=\frac{\arccos(\cos\vartheta)}{\pi},
\qquad
\bar f(\vartheta)=f(\chi(\vartheta)),
\qquad \vartheta\in\mathbb R.
\]
The map \(\chi\) takes values in \([0,1]\), is continuous, even, and \(2\pi\)-periodic, and satisfies
\[
\chi(\pi x)=x\qquad(x\in[0,1]).
\]
We also have
\[
|\chi(\vartheta)-\chi(\vartheta')|
\le \frac{|\vartheta-\vartheta'|}{\pi}
\qquad(\vartheta,\vartheta'\in\mathbb R).
\]
To see this, choose an integer \(\ell_{\mathrm{per}}\) such that
\[
\vartheta'-2\pi\ell_{\mathrm{per}}\in[-\pi,\pi).
\]
The elementary identities and bounds
\[
\arccos(\cos\vartheta)=|\vartheta|
\quad\text{if }|\vartheta|\le\pi,
\qquad
\arccos(\cos\vartheta)\le|\vartheta|
\quad\text{for every }\vartheta\in\mathbb R
\]
give, by periodicity and the triangle inequality,
\[
\begin{aligned}
\arccos(\cos\vartheta)
&\le|\vartheta-2\pi\ell_{\mathrm{per}}|\\
&\le|\vartheta'-2\pi\ell_{\mathrm{per}}|
   +|\vartheta-\vartheta'|\\
&=\arccos(\cos\vartheta')+|\vartheta-\vartheta'|.
\end{aligned}
\]
Interchanging the two angles proves the asserted Lipschitz bound. Consequently,
\[
|\bar f(\vartheta)-\bar f(\vartheta')|
\le L_f\left(\frac{|\vartheta-\vartheta'|}{\pi}\right)^\gamma.
\]
The extension \(\bar f\) is continuous, even, and \(2\pi\)-periodic.
% lean: Causalean.Mathlib.Analysis.Approximation.Chebyshev.CosineProjection.evenExtension_holder

\item For an integer \(m_{\mathrm J}\ge1\), introduce the continuous sine quotient
\[
\mathcal S_{m_{\mathrm J}}(\vartheta)
=
\begin{cases}
\dfrac{\sin(m_{\mathrm J}\vartheta/2)}{\sin(\vartheta/2)},
   &\sin(\vartheta/2)\ne0,\\[6pt]
m_{\mathrm J}(-1)^{(m_{\mathrm J}-1)\iota_{\mathrm{per}}},
   &\vartheta=2\pi\iota_{\mathrm{per}},\quad \iota_{\mathrm{per}}\in\mathbb Z,
\end{cases}
\qquad \vartheta\in\mathbb R,
\]
and its mass and normalized fourth-power kernel
\[
\mathcal Z_{m_{\mathrm J}}
=\int_{-\pi}^{\pi}\mathcal S_{m_{\mathrm J}}(\vartheta)^4\,d\vartheta,
\qquad
\mathcal J_{m_{\mathrm J}}(\vartheta)
=\frac{\mathcal S_{m_{\mathrm J}}(\vartheta)^4}{\mathcal Z_{m_{\mathrm J}}}.
\]
The removable values in the definition are the corresponding limits.

The finite cosine identity is
\[
\mathcal S_{m_{\mathrm J}}(\vartheta)^2
=m_{\mathrm J}
 +2\sum_{j=1}^{m_{\mathrm J}-1}(m_{\mathrm J}-j)\cos(j\vartheta).
\]
For completeness, the Chebyshev polynomials of the second kind, defined by
\[
\mathsf U_0(z)=1,\qquad
\mathsf U_1(z)=2z,\qquad
\mathsf U_{j+1}(z)=2z\mathsf U_j(z)-\mathsf U_{j-1}(z)
\quad(j\ge1),
\]
satisfy
\[
\mathcal S_{m_{\mathrm J}}(\vartheta)
=\mathsf U_{m_{\mathrm J}-1}(\cos(\vartheta/2)).
\]
We first derive the polynomial identity
\[
\mathsf U_j(z)^2-\mathsf U_{j-1}(z)^2=\mathsf U_{2j}(z)
\qquad(j\ge1).
\]
Regard these real-coefficient polynomials also as polynomials on \(\mathbb C\). Surjectivity of the complex cosine allows us, for each complex argument, to choose a complex angle satisfying
\[
z_{\mathrm C}=\cos\vartheta_{\mathrm C},
\qquad z_{\mathrm C},\vartheta_{\mathrm C}\in\mathbb C.
\]
If \(\sin\vartheta_{\mathrm C}=0\), then \(\cos\vartheta_{\mathrm C}=1\) or \(-1\). Induction in the defining recurrence gives
\[
\mathsf U_j(1)=j+1,
\qquad
\mathsf U_j(-1)=(-1)^j(j+1)
\qquad(j\ge0).
\]
Thus, at either argument,
\[
\mathsf U_j(\pm1)^2-\mathsf U_{j-1}(\pm1)^2
=(j+1)^2-j^2=2j+1=\mathsf U_{2j}(\pm1)
\qquad(j\ge1).
\]
If \(\sin\vartheta_{\mathrm C}\ne0\), induction in the same recurrence, using the sine addition formula, gives
\[
\mathsf U_j(\cos\vartheta_{\mathrm C})
=\frac{\sin((j+1)\vartheta_{\mathrm C})}
       {\sin\vartheta_{\mathrm C}}
\qquad(j\ge0).
\]
Expanding the sine and cosine addition formulas and using
\(\cos^2\vartheta_{\mathrm C}+\sin^2\vartheta_{\mathrm C}=1\) yields
\[
\sin((j+1)\vartheta_{\mathrm C})^2
-\sin(j\vartheta_{\mathrm C})^2
=\sin((2j+1)\vartheta_{\mathrm C})
 \sin\vartheta_{\mathrm C}.
\]
Dividing by \(\sin^2\vartheta_{\mathrm C}\) proves
\[
\begin{aligned}
\mathsf U_j(\cos\vartheta_{\mathrm C})^2
-\mathsf U_{j-1}(\cos\vartheta_{\mathrm C})^2
&=\frac{\sin((2j+1)\vartheta_{\mathrm C})}
        {\sin\vartheta_{\mathrm C}}\\
&=\mathsf U_{2j}(\cos\vartheta_{\mathrm C}).
\end{aligned}
\]
The two cases establish the polynomial identity at every complex argument, hence also over the reals.

For the even-index cosine expansion, introduce the first-kind Chebyshev polynomials by
\[
\mathsf T^{\mathrm{Ch}}_0(z_{\mathrm{Ch}})=1,
\qquad
\mathsf T^{\mathrm{Ch}}_1(z_{\mathrm{Ch}})=z_{\mathrm{Ch}},
\qquad z_{\mathrm{Ch}}\in\mathbb R,
\]
\[
\mathsf T^{\mathrm{Ch}}_{j+1}(z_{\mathrm{Ch}})
=2z_{\mathrm{Ch}}\mathsf T^{\mathrm{Ch}}_j(z_{\mathrm{Ch}})
 -\mathsf T^{\mathrm{Ch}}_{j-1}(z_{\mathrm{Ch}})
\qquad(j\ge1).
\]
The cosine addition formula and induction give
\[
\mathsf T^{\mathrm{Ch}}_j(\cos(\vartheta/2))
=\cos(j\vartheta/2)
\qquad(j\ge0,\ \vartheta\in\mathbb R).
\]
The defining recurrences also imply
\[
\mathsf U_{j+2}(z_{\mathrm{Ch}})-\mathsf U_j(z_{\mathrm{Ch}})
=2\mathsf T^{\mathrm{Ch}}_{j+2}(z_{\mathrm{Ch}})
\qquad(j\ge0).
\]
Indeed, both sides satisfy the same second-order recurrence as sequences in \(j\), and their first two values agree:
\[
\begin{aligned}
\mathsf U_2(z_{\mathrm{Ch}})-\mathsf U_0(z_{\mathrm{Ch}})
&=4z_{\mathrm{Ch}}^2-2
 =2\mathsf T^{\mathrm{Ch}}_2(z_{\mathrm{Ch}}),\\
\mathsf U_3(z_{\mathrm{Ch}})-\mathsf U_1(z_{\mathrm{Ch}})
&=8z_{\mathrm{Ch}}^3-6z_{\mathrm{Ch}}
 =2\mathsf T^{\mathrm{Ch}}_3(z_{\mathrm{Ch}}).
\end{aligned}
\]
Consequently,
\[
\mathsf U_{2j+2}(\cos(\vartheta/2))
=\mathsf U_{2j}(\cos(\vartheta/2))
 +2\cos((j+1)\vartheta)
\qquad(j\ge0).
\]
The base case is \(\mathsf U_0=1\). If the expansion holds at \(j\), the last identity gives
\[
\begin{aligned}
\mathsf U_{2j+2}(\cos(\vartheta/2))
&=1+2\sum_{\iota_{\mathrm{cos}}=1}^{j}
 \cos(\iota_{\mathrm{cos}}\vartheta)
 +2\cos((j+1)\vartheta)\\
&=1+2\sum_{\iota_{\mathrm{cos}}=1}^{j+1}
 \cos(\iota_{\mathrm{cos}}\vartheta).
\end{aligned}
\]
Thus induction proves
\[
\mathsf U_{2j}(\cos(\vartheta/2))
=1+2\sum_{\iota_{\mathrm{cos}}=1}^{j}
 \cos(\iota_{\mathrm{cos}}\vartheta)
\qquad(j\ge0),
\]
with an empty sum at \(j=0\).
Starting with \(\mathsf U_0^2=1\) and summing these differences yields the displayed cosine identity.

Integrating that identity, and then its square, gives
\[
\int_{-\pi}^{\pi}\mathcal S_{m_{\mathrm J}}(\vartheta)^2\,d\vartheta
=2\pi m_{\mathrm J},
\]
and
\[
\begin{aligned}
\mathcal Z_{m_{\mathrm J}}
&=2\pi m_{\mathrm J}^2
  +4\pi\sum_{j=1}^{m_{\mathrm J}-1}(m_{\mathrm J}-j)^2\\
&=2\pi m_{\mathrm J}^2
  +\frac{4\pi}{6}(m_{\mathrm J}-1)m_{\mathrm J}(2m_{\mathrm J}-1)\\
&=\frac{2\pi}{3}m_{\mathrm J}(2m_{\mathrm J}^2+1)
 \ge\frac{4\pi}{3}m_{\mathrm J}^3>0.
\end{aligned}
\]
Here nonconstant cosines integrate to zero, distinct cosine modes are orthogonal, and each positive-frequency squared cosine has integral \(\pi\). In particular,
\[
\mathcal J_{m_{\mathrm J}}(\vartheta)\ge0,
\qquad
\int_{-\pi}^{\pi}\mathcal J_{m_{\mathrm J}}(\vartheta)\,d\vartheta=1.
\]
These formulas include \(m_{\mathrm J}=1\), when the sums are empty.
% lean: Causalean.Mathlib.Analysis.JacksonApproximation.jrawMass_eq

\item For \(|\vartheta|\le\pi\), the sine chord bound gives
\[
|\vartheta|\le\pi|\sin(\vartheta/2)|.
\]
Together with
\[
\mathcal S_{m_{\mathrm J}}(\vartheta)\sin(\vartheta/2)
=\sin(m_{\mathrm J}\vartheta/2),
\]
this implies
\[
\vartheta^2\mathcal S_{m_{\mathrm J}}(\vartheta)^2
\le\pi^2\sin(m_{\mathrm J}\vartheta/2)^2
\le\pi^2.
\]
The identity also holds at the removable points. Multiplying by
\(\mathcal S_{m_{\mathrm J}}(\vartheta)^2\) and integrating yields
\[
\int_{-\pi}^{\pi}
 \vartheta^2\mathcal S_{m_{\mathrm J}}(\vartheta)^4\,d\vartheta
\le
\pi^2\int_{-\pi}^{\pi}
 \mathcal S_{m_{\mathrm J}}(\vartheta)^2\,d\vartheta
=2\pi^3m_{\mathrm J}.
\]
Using the mass lower bound from the preceding step, we obtain
\[
\begin{aligned}
\int_{-\pi}^{\pi}\vartheta^2\mathcal J_{m_{\mathrm J}}(\vartheta)\,d\vartheta
&\le \frac{2\pi^3m_{\mathrm J}}{\mathcal Z_{m_{\mathrm J}}}\\
&\le \frac{2\pi^3m_{\mathrm J}}
              {(4\pi/3)m_{\mathrm J}^3}
\le \left(\frac{2\pi}{m_{\mathrm J}}\right)^2.
\end{aligned}
\]
The kernel therefore defines a probability measure
\[
d\nu_{\mathrm J}(\vartheta)
=\mathbf 1_{[-\pi,\pi]}(\vartheta)
 \mathcal J_{m_{\mathrm J}}(\vartheta)\,d\vartheta.
\]
Jensen's inequality for the concave function \(z\mapsto z^{\gamma/2}\) on \([0,\infty)\) gives
\[
\begin{aligned}
\int_{-\pi}^{\pi}
 \left(\frac{|\vartheta|}{\pi}\right)^\gamma
 \mathcal J_{m_{\mathrm J}}(\vartheta)\,d\vartheta
&=
\int_{\mathbb R}
 \left(\frac{\vartheta^2}{\pi^2}\right)^{\gamma/2}
 \,d\nu_{\mathrm J}(\vartheta)\\
&\le
\left(\int_{\mathbb R}\frac{\vartheta^2}{\pi^2}
 \,d\nu_{\mathrm J}(\vartheta)\right)^{\gamma/2}\\
&\le\left(\frac{2}{m_{\mathrm J}}\right)^\gamma.
\end{aligned}
\]
This argument applies throughout \(0<\gamma\le1\).
% lean: Causalean.Mathlib.Analysis.Approximation.Chebyshev.CosineProjection.jackson_holder_moment

\item Define the Jackson approximant on the unit interval by
\[
p_{\mathrm J}(x)
=\int_{-\pi}^{\pi}
 \bar f(\pi x-\vartheta)\mathcal J_{m_{\mathrm J}}(\vartheta)\,d\vartheta,
\qquad x\in[0,1].
\]
Normalization of the kernel and the increment bound for \(\bar f\) imply
\[
\begin{aligned}
|f(x)-p_{\mathrm J}(x)|
&=\left|\int_{-\pi}^{\pi}
  \{f(x)-\bar f(\pi x-\vartheta)\}
  \mathcal J_{m_{\mathrm J}}(\vartheta)\,d\vartheta\right|\\
&\le L_f\int_{-\pi}^{\pi}
 \left(\frac{|\vartheta|}{\pi}\right)^\gamma
 \mathcal J_{m_{\mathrm J}}(\vartheta)\,d\vartheta\\
&\le L_f\left(\frac{2}{m_{\mathrm J}}\right)^\gamma.
\end{aligned}
\]
The approximant is continuous.

It also lies in the cosine span through frequency \(2m_{\mathrm J}-2\). Indeed,
\(\mathsf U_{m_{\mathrm J}-1}^4\) is an even polynomial of degree \(4m_{\mathrm J}-4\); hence there is a polynomial \(\mathcal Q_{m_{\mathrm J}}\) such that
\[
\mathsf U_{m_{\mathrm J}-1}(z)^4
=\mathcal Q_{m_{\mathrm J}}(z^2),
\qquad
\deg\mathcal Q_{m_{\mathrm J}}\le2m_{\mathrm J}-2.
\]
Since \(\cos(\vartheta/2)^2=(1+\cos\vartheta)/2\),
\[
\mathcal J_{m_{\mathrm J}}(\vartheta)
=
\frac{\mathcal Q_{m_{\mathrm J}}((1+\cos\vartheta)/2)}
     {\mathcal Z_{m_{\mathrm J}}}.
\]
Thus the kernel is an even trigonometric polynomial. Averaging its trigonometric expansion at \(\vartheta\) and \(-\vartheta\) removes the sine terms, giving real coefficients
\(\mathfrak a_{m_{\mathrm J},j}\), \(0\le j\le2m_{\mathrm J}-2\), with
\[
\mathcal J_{m_{\mathrm J}}(\vartheta)
=\sum_{j=0}^{2m_{\mathrm J}-2}
 \mathfrak a_{m_{\mathrm J},j}\cos(j\vartheta).
\]
Reversing the convolution over one period and using the cosine subtraction formula gives
\[
\begin{aligned}
p_{\mathrm J}(x)
&=\int_{-\pi}^{\pi}
 \bar f(\vartheta)\mathcal J_{m_{\mathrm J}}(\pi x-\vartheta)\,d\vartheta\\
&=\sum_{j=0}^{2m_{\mathrm J}-2}
 \left\{\mathfrak a_{m_{\mathrm J},j}
 \int_{-\pi}^{\pi}\bar f(\vartheta)\cos(j\vartheta)\,d\vartheta\right\}
 \cos(\pi jx).
\end{aligned}
\]
The sine integrals vanish because \(\bar f\) is even. The reversal follows by a change of variable and periodicity. Since \(\varphi_0=1\) and
\(\varphi_j(x)=\sqrt2\cos(\pi jx)\) for \(j\ge1\), this is the required cosine-span representation.
% lean: Causalean.Mathlib.Analysis.Approximation.Chebyshev.CosineProjection.jacksonApproximant_mem_span

\item Now take
\[
m_{\mathrm J}=\left\lfloor\frac{k+1}{2}\right\rfloor.
\]
For every \(k\ge1\),
\[
m_{\mathrm J}\ge1,\qquad
2m_{\mathrm J}-2<k,\qquad
k\le2m_{\mathrm J}.
\]
Consequently \(p_{\mathrm J}\) belongs to the range of \(\Pi_k\), and
\[
\left(\frac{2}{m_{\mathrm J}}\right)^\gamma
\le\left(\frac4k\right)^\gamma
=4^\gamma k^{-\gamma}
\le5k^{-\gamma},
\]
where \(4^\gamma\le4\le5\). Hence
\[
|f(x)-p_{\mathrm J}(x)|\le5L_fk^{-\gamma}
\qquad(x\in[0,1]).
\]
Orthogonality of the projection gives
\[
\|f-p_{\mathrm J}\|_2^2
=\|f-\Pi_kf\|_2^2+\|\Pi_kf-p_{\mathrm J}\|_2^2,
\]
so, because \(\lambda\) has total mass one,
\[
\|f-\Pi_kf\|_2
\le\|f-p_{\mathrm J}\|_2
\le5L_fk^{-\gamma}.
\]
For \(k=1\), the chosen order is \(m_{\mathrm J}=1\) and the approximant is constant, so the same argument applies. In the finite-seminorm case these ordinary norms equal the extended norms in the statement; together with the infinite-seminorm case, this proves the approximation assertion.
% lean: CausalSmith.Stat.LogoddsLowsmoothFrontier.cosineProjection_error_le_holderSeminorm

\item Fix measurable marks \(W,V\) valued in \([0,1]\) and \(n\ge2\). Let \(O,O'\) be independent records and let \(O_1,\ldots,O_n\) be the original independent sample:
\[
(O,O')\sim P\otimes P,
\qquad
(O_1,\ldots,O_n)\sim P^{\otimes n}.
\]
Writing \(p_{ay}(x)=P(A=a,Y=y\mid X=x)\) for the specified four-cell probabilities, the conditional means are
\[
w(x)=\sum_{a,y\in\{0,1\}}p_{ay}(x)W(x,a,y),
\qquad
v(x)=\sum_{a,y\in\{0,1\}}p_{ay}(x)V(x,a,y).
\]
Their measurability and the normalization of the cells give
\[
0\le w(x),v(x)\le1,\qquad x\in[0,1],
\]
so both belong to \(L^2(\lambda)\). Disintegration and
\cref{obj:ass:uniform-design} give, for every integrable measurable \(h:[0,1]\to\mathbb R\),
\[
E_P[W(O)h(X)]=\int_0^1w(x)h(x)\,d\lambda(x),
\qquad
E_P[V(O)h(X)]=\int_0^1v(x)h(x)\,d\lambda(x).
\]

Define the symmetric record kernel
\[
q_k(o,o')
=\frac{H_k(x,x')}{2}
 \{W(o)V(o')+V(o)W(o')\},
\qquad x=X(o),\quad x'=X(o').
\]
Swapping the two indices in the complete ordered-pair sum shows that
\[
\mathbb U_{n,k}(W,V)
=\frac1{n(n-1)}\sum_{i\ne j}q_k(O_i,O_j).
\]
The finite kernel expansion and the preceding conditional-mean identities yield
\[
\int q_k(o,o')\,dP(o')
=\frac{W(o)(\Pi_kv)(X(o))+V(o)(\Pi_kw)(X(o))}{2}.
\]
Indeed,
\[
\begin{aligned}
\int H_k(x,X(o'))V(o')\,dP(o')
&=\sum_{j=0}^{k-1}\varphi_j(x)
  \int_0^1v(z)\varphi_j(z)\,d\lambda(z)\\
&=(\Pi_kv)(x),
\end{aligned}
\]
and the corresponding identity holds with \(W,w\).
% lean: CausalSmith.Stat.LogoddsLowsmoothFrontier.symmetricProjectionKernel_integral_right

\item Each distinct pair of sample coordinates has law \(P\otimes P\), and there are \(n(n-1)\) such ordered pairs. Therefore
\[
E_P\mathbb U_{n,k}(W,V)=E_Pq_k(O,O').
\]
Integrating the kernel section from the preceding step gives
\[
\begin{aligned}
E_Pq_k(O,O')
&=\frac12\{\langle w,\Pi_kv\rangle+\langle v,\Pi_kw\rangle\}\\
&=\sum_{j=0}^{k-1}
 \langle w,\varphi_j\rangle\langle v,\varphi_j\rangle\\
&=\langle\Pi_kw,\Pi_kv\rangle.
\end{aligned}
\]
Here the last equality uses the orthonormality supplied by
\cref{obj:synth_3}. All sums are finite; bounded marks and the continuous finite kernel ensure integrability.
% lean: CausalSmith.Stat.LogoddsLowsmoothFrontier.projectionStatistic_integral

\item The same finite expansion gives
\[
\langle w,\Pi_kv\rangle
=\langle v,\Pi_kw\rangle
=\langle\Pi_kw,\Pi_kv\rangle.
\]
Expanding the product of the residuals thus yields
\[
\begin{aligned}
\langle w-\Pi_kw,v-\Pi_kv\rangle
&=\langle w,v\rangle-\langle w,\Pi_kv\rangle
 -\langle v,\Pi_kw\rangle+\langle\Pi_kw,\Pi_kv\rangle\\
&=\langle w,v\rangle-\langle\Pi_kw,\Pi_kv\rangle\\
&=\langle w,v\rangle-E_P\mathbb U_{n,k}(W,V).
\end{aligned}
\]
This proves the bias identity.
% lean: CausalSmith.Stat.LogoddsLowsmoothFrontier.projectionStatistic_bias

\item For the variance argument define, for all records \(o,o'\),
\[
\begin{aligned}
\tau_k&=\langle\Pi_kw,\Pi_kv\rangle,\\
\psi_k(o)&=\int q_k(o,o')\,dP(o'),\\
\mathcal R_k(o,o')
&=q_k(o,o')-\psi_k(o)-\psi_k(o')+\tau_k.
\end{aligned}
\]
Also define the two finite-sample terms
\[
\mathbb A_{n,k}
=\frac2n\sum_{i=1}^n\{\psi_k(O_i)-\tau_k\},
\qquad
\mathbb B_{n,k}
=\frac1{n(n-1)}\sum_{i\ne j}\mathcal R_k(O_i,O_j).
\]
These quantities are square integrable: the finite cosine kernel is bounded on the compact covariate square, the marks are bounded, and all sums are finite.

The kernel mean identity gives \(E_P\psi_k(O)=\tau_k\). Hence
\[
\int\mathcal R_k(o,o')\,dP(o')
=\psi_k(o)-\psi_k(o)-\tau_k+\tau_k=0.
\]
Symmetry of \(q_k\) gives symmetry of \(\mathcal R_k\), so
\[
\int\mathcal R_k(o,o')\,dP(o)=0,
\qquad
E_P\mathcal R_k(O,O')=0.
\]
For \(i\ne j\) and \(1\le\iota_{\mathrm{sample}}\le n\),
\[
E_P\!\left[
 \{\psi_k(O_{\iota_{\mathrm{sample}}})-\tau_k\}\mathcal R_k(O_i,O_j)
\right]=0.
\]
If \(\iota_{\mathrm{sample}}=i\), integrate first over \(O_j\); if \(\iota_{\mathrm{sample}}=j\), integrate first over \(O_i\). In each case the corresponding residual section integrates to zero. If \(\iota_{\mathrm{sample}}\) differs from both indices, independence factors the expectation, and the residual mean is zero. Expanding the finite product of sums therefore gives
\[
E_P[\mathbb A_{n,k}\mathbb B_{n,k}]=0.
\]
% lean: CausalSmith.Stat.LogoddsLowsmoothFrontier.projection_hoeffding_cross

\item The explicit formula for the first projection implies
\[
\psi_k(o)^2
\le\frac{(\Pi_kv)(X(o))^2+(\Pi_kw)(X(o))^2}{2}.
\]
Indeed,
\[
\begin{aligned}
&\{W(o)(\Pi_kv)(X(o))+V(o)(\Pi_kw)(X(o))\}^2\\
&\qquad\le2\{W(o)^2(\Pi_kv)(X(o))^2
             +V(o)^2(\Pi_kw)(X(o))^2\},
\end{aligned}
\]
by the elementary square inequality, and the squared marks are at most one. Applying the residual identity with the same conditional mean in both positions gives
\[
\|\Pi_kw\|_2^2
=\|w\|_2^2-\|w-\Pi_kw\|_2^2
\le\|w\|_2^2\le1,
\]
and likewise \(\|\Pi_kv\|_2^2\le1\). Uniform design consequently yields
\[
E_P\psi_k(O)^2
\le\frac{\|\Pi_kv\|_2^2+\|\Pi_kw\|_2^2}{2}\le1.
\]
Expanding the centered square and using \(E_P\psi_k(O)=\tau_k\), we obtain
\[
E_P\{\psi_k(O)-\tau_k\}^2
=E_P\psi_k(O)^2-\tau_k^2\le1.
\]
Independence makes the cross products between distinct centered sample terms vanish. Thus
\[
\begin{aligned}
E_P\mathbb A_{n,k}^2
&=\frac4{n^2}\,
 nE_P\{\psi_k(O)-\tau_k\}^2\\
&=\frac4nE_P\{\psi_k(O)-\tau_k\}^2
\le\frac4n.
\end{aligned}
\]
% lean: CausalSmith.Stat.LogoddsLowsmoothFrontier.projectionFirst_finite_secondMoment_le

\item The mark restrictions give the pointwise inequality
\[
q_k(o,o')^2\le H_k(X(o),X(o'))^2.
\]
By the finite cosine expansion and orthonormality,
\[
\int_0^1 H_k(x,z)^2\,d\lambda(z)
=\sum_{j=0}^{k-1}\varphi_j(x)^2,
\]
and hence
\[
E_Pq_k(O,O')^2
\le\int_0^1\!\int_0^1H_k(x,z)^2\,d\lambda(z)\,d\lambda(x)
=\sum_{j=0}^{k-1}\|\varphi_j\|_2^2=k.
\]
To obtain the residual energy, note that integration of a kernel section gives
\[
E_P[q_k(O,O')\psi_k(O)]
=E_P[q_k(O,O')\psi_k(O')]
=E_P\psi_k(O)^2.
\]
Independence also gives
\[
E_P[\psi_k(O)\psi_k(O')]=\tau_k^2,
\qquad
E_Pq_k(O,O')=\tau_k.
\]
Expanding the defining square of \(\mathcal R_k\) and substituting these identities yields
\[
E_P\mathcal R_k(O,O')^2
=E_Pq_k(O,O')^2-2E_P\psi_k(O)^2+\tau_k^2.
\]
Using the centered-energy identity from the preceding step, this becomes
\[
\begin{aligned}
E_P\mathcal R_k(O,O')^2
&=E_Pq_k(O,O')^2-E_P\psi_k(O)^2
  -E_P\{\psi_k(O)-\tau_k\}^2\\
&\le E_Pq_k(O,O')^2\le k.
\end{aligned}
\]
% lean: CausalSmith.Stat.LogoddsLowsmoothFrontier.projectionResidual_energy_le_rank

\item For two ordered distinct pairs \((i,j)\) and \((i',j')\), residual degeneracy gives
\[
E_P[
 \mathcal R_k(O_i,O_j)\mathcal R_k(O_{i'},O_{j'})
]
=
\begin{cases}
E_P\mathcal R_k(O,O')^2,
 &\{i,j\}=\{i',j'\},\\
0,&\{i,j\}\ne\{i',j'\}.
\end{cases}
\]
For equal index sets, the pairs are identical or reversed, and symmetry gives the first value. For unequal index sets, one record appears in just one factor; integrating over that record first gives zero. This covers both disjoint pairs and pairs sharing one record.

Each of the \(n(n-1)\) ordered pairs has exactly two ordered pairs with the same index set. Expanding the square of the sum therefore gives
\[
E_P\left[
 \left\{\sum_{i\ne j}\mathcal R_k(O_i,O_j)\right\}^2
\right]
=2n(n-1)E_P\mathcal R_k(O,O')^2.
\]
Equivalently,
\[
nE_P\mathbb B_{n,k}^2
=\frac2{n-1}E_P\mathcal R_k(O,O')^2.
\]
Since \(n\ge2\), division by \(n\) and the residual energy bound yield
\[
E_P\mathbb B_{n,k}^2
=\frac{2E_P\mathcal R_k(O,O')^2}{n(n-1)}
\le\frac{2k}{n(n-1)}.
\]
The counting includes \(n=2\), when the two ordered pairs are reversals of one another.
% lean: CausalSmith.Stat.LogoddsLowsmoothFrontier.projectionResidual_finite_secondMoment_le

\item Finally, summing the defining residual identity over ordered distinct pairs gives
\[
\sum_{i\ne j}\mathcal R_k(O_i,O_j)
=\sum_{i\ne j}q_k(O_i,O_j)
 -2(n-1)\sum_{i=1}^n\psi_k(O_i)+n(n-1)\tau_k.
\]
Each first-coordinate term and each second-coordinate term occurs \(n-1\) times. Dividing by \(n(n-1)\) proves the pointwise decomposition
\[
\mathbb U_{n,k}(W,V)-\tau_k
=\mathbb A_{n,k}+\mathbb B_{n,k}.
\]
Since \(E_P\mathbb U_{n,k}(W,V)=\tau_k\), expansion of the square, the zero cross moment, and the two established bounds give
\[
\begin{aligned}
\operatorname{Var}_P\{\mathbb U_{n,k}(W,V)\}
&=E_P(\mathbb A_{n,k}+\mathbb B_{n,k})^2\\
&=E_P\mathbb A_{n,k}^2+E_P\mathbb B_{n,k}^2
  +2E_P[\mathbb A_{n,k}\mathbb B_{n,k}]\\
&\le\frac4n+\frac{2k}{n(n-1)}.
\end{aligned}
\]
Together with the approximation, mean, and bias identities, this proves the result.
% lean: CausalSmith.Stat.LogoddsLowsmoothFrontier.projectionStatistic_variance_le
\end{enumerate}
\end{proof}

The product-bias identity in \cref{obj:lem:projection-control} expresses the approximation error as the inner product of two projection residuals. It therefore connects the smoothness of both conditional means to the bias of their estimated product. The approximation inequality covers the Lipschitz endpoint with the same constant, and the variance bound separates a term proportional to \(n^{-1}\) from a term proportional to the projection rank divided by the number of ordered pairs. Under the known uniform design, these bounds apply directly in \(L^2(\lambda)\). Together with the denominator floor and inversion in \cref{obj:lem:observable-odds}, they supply the statistical bounds used to calibrate the coordinate error radii in \cref{obj:def:upper-handle} and establish the interval guarantee in \cref{obj:thm:finite-honest-upper}.


\section{Rational realization and the upper bound}
\label{sec:rational-realization}

We begin with the representation and primitive arithmetic contracts used by the finite rational construction.

The projection bounds in \cref{obj:lem:projection-control} supply the statistical ingredients for the coordinate error radii. This appendix establishes their finite rational realization: an explicit admissible arithmetic engine, uniform guarantees at the two prescribed precisions, and a measurable interval output. These conclusions connect the exact coordinate expressions in \cref{obj:def:coordinate-estimates} to the literal interval procedure in \cref{obj:def:upper-handle}.

The two precision levels control rank selection and final endpoint enclosure. Each primitive is evaluated by a finite rational recurrence with a public iteration bound on valid represented inputs. The guarantees establish termination and endpoint enclosure under the stated rational computation model, with public bounds on primitive iteration counts.

\subsection{Naming conventions and primitive admissibility}

Under \cref{obj:def:dyadic-name-convention}, a real input is represented by nested closed dyadic intervals containing it, with width at most \(2^{-q}\) at integer precision \(q\ge0\). An admissible Borel naming policy \(\mathsf N\) selects public exponent names before sampling and covariate names as Borel functions of the public exponents and original records. The dyadic-floor policy \(\mathsf N_0\) provides the concrete selection specified there.

The independently fixed engine \(\mathsf E\) satisfies the primitive contracts of \cref{obj:def:arithmetic-engine}. Its transcendental routines enclose the images of rational input intervals, with outward endpoint excess controlled by the dyadic tolerance \(\epsilon_q=2^{-q}\). The cosine routine uses a midpoint enclosure enlarged by half the input width and clipped to \([-1,1]\). Exact rational range operations handle arithmetic, extrema, intersections, and clipping. Domain and sign conditions ensure positive logarithm arguments, nonzero division denominators, positive exponential and power lower endpoints, and nonnegative square-root lower endpoints. Each primitive uses a finite rational recurrence with a prescribed public bound on its iteration count. In the power routine, the bound argument \(b\) denotes a positive integer base.

These requirements specify local arithmetic guarantees. The following statement supplies an explicit engine satisfying them.

\begin{lemmav}{lem:concrete-arithmetic-engine}[Admissibility of the concrete engine]
The explicit rational arithmetic engine \(\mathsf E_0\), consisting of the Machin enclosure for \(\pi\), the cosine interval routine, and the endpoint extensions of the scalar exponential, logarithm, square-root, and integer-base power routines, satisfies all primitive admissibility contracts of \cref{obj:def:arithmetic-engine}.
\end{lemmav}

\begin{proof}[Proof of \cref{obj:lem:concrete-arithmetic-engine}]
We verify the finite-computation certificate first, followed by the constant, exponential, logarithm, square-root, power, and cosine contracts. Rational quantities in real inequalities are identified with their real values.

\begin{enumerate}
\item \emph{Finite rational computation.}
For \(q\in\mathbb N_0\), define
\[
\epsilon_q=2^{-q},
\qquad
p_{\mathrm{sc}}(q)=q+3,
\qquad
e_{\mathrm{sc}}(q)=2^{-p_{\mathrm{sc}}(q)}.
\]
For rational \(\xi_{\mathrm{rat}}\) and nonnegative integer \(n_{\mathrm{int}}\), put
\[
\lfloor\xi_{\mathrm{rat}}\rfloor_+
 =\max\{0,\lfloor\xi_{\mathrm{rat}}\rfloor\},
\qquad
\lceil\xi_{\mathrm{rat}}\rceil_+
 =\max\{0,\lceil\xi_{\mathrm{rat}}\rceil\},
\qquad
\kappa(n_{\mathrm{int}})
 =\min\{\jmath\in\mathbb N_0:n_{\mathrm{int}}\le 2^\jmath\}.
\]
Thus \(\kappa(0)=0\). In the rational computation model, division by zero is assigned the value zero. This convention specifies the routines on their full input domains; the analytic contracts use the domains in \cref{obj:def:arithmetic-engine}.

For rational endpoints define
\[
\operatorname{box}(\xi_{\mathrm{lo}},\xi_{\mathrm{hi}})
 =
 [\min\{\xi_{\mathrm{lo}},\xi_{\mathrm{hi}}\},
  \max\{\xi_{\mathrm{lo}},\xi_{\mathrm{hi}}\}].
\]
For any rational interval \(\mathcal J\), write
\[
\mathcal J=[\underline{\mathcal J},\overline{\mathcal J}],
\qquad
\operatorname{wid}(\mathcal J)
 =\overline{\mathcal J}-\underline{\mathcal J}.
\]

Here are the scalar routines implemented by the programs. For \(z\in\mathbb Q\), define
\[
B_{\mathrm{ser}}(z)=\max\{1,\lceil|z|\rceil_+\},
\qquad
m_{\mathrm{ser}}(q,z)
 =\max\{2B_{\mathrm{ser}}(z)^2,q+6\},
\]
\[
\Sigma_{\exp}(q,z)
 =\sum_{\jmath=0}^{2m_{\mathrm{ser}}(q,z)-1}\frac{z^\jmath}{\jmath!},
\qquad
\Sigma_{\cos}(q,z)
 =\sum_{\jmath=0}^{m_{\mathrm{ser}}(q,z)-1}
   \frac{(-1)^\jmath z^{2\jmath}}{(2\jmath)!},
\]
and
\[
\mathcal B_{\exp}(q,z)
 =\operatorname{box}\!\left(
   \max\{3^{-B_{\mathrm{ser}}(z)},
          \Sigma_{\exp}(q,z)-e_{\mathrm{sc}}(q)\},
   \Sigma_{\exp}(q,z)+e_{\mathrm{sc}}(q)\right),
\]
\[
\mathcal B_{\cos}(q,z)
 =\operatorname{box}\!\left(
   \max\{-1,\Sigma_{\cos}(q,z)-e_{\mathrm{sc}}(q)\},
   \min\{1,\Sigma_{\cos}(q,z)+e_{\mathrm{sc}}(q)\}\right).
\]
For the logarithm, define on all rational inputs
\[
B_{\log}(z)=\max\{2,\lceil z\rceil_+,\lceil1/z\rceil_+\},
\qquad
\vartheta_{\log}(z)=\frac{z-1}{z+1},
\]
\[
c_{\log}(q,z)=q+3+\kappa(8B_{\log}(z))+2,
\qquad
n_{\log}(q,z)=B_{\log}(z)c_{\log}(q,z),
\]
\[
\Sigma_{\log}(q,z)
 =2\sum_{\jmath=0}^{n_{\log}(q,z)-1}
       \frac{\vartheta_{\log}(z)^{2\jmath+1}}{2\jmath+1},
\qquad
\mathcal B_{\log}(q,z)=
\begin{cases}
\operatorname{box}(\Sigma_{\log}(q,z)-e_{\mathrm{sc}}(q),
                   \Sigma_{\log}(q,z)+e_{\mathrm{sc}}(q)),&z>0,\\
[0,0],&z\le0.
\end{cases}
\]

For the square-root routine, set
\[
n_{\sqrt{\ }}(q,z)=\lfloor 2^{2p_{\mathrm{sc}}(q)}z\rfloor_+.
\]
For any \(n_{\mathrm{int}}\in\mathbb N_0\), start a binary recurrence with
\[
l_0=0,\qquad u_0=n_{\mathrm{int}}+1.
\]
At each iteration retain the endpoints when \(u_\jmath-l_\jmath=1\).
Otherwise, with
\[
h_\jmath=\left\lfloor\frac{l_\jmath+u_\jmath}{2}\right\rfloor,
\]
set
\[
(l_{\jmath+1},u_{\jmath+1})=
\begin{cases}
(h_\jmath,u_\jmath),&h_\jmath^2\le n_{\mathrm{int}},\\
(l_\jmath,h_\jmath),&h_\jmath^2>n_{\mathrm{int}}.
\end{cases}
\]
The invariant and width implication are
\[
l_\jmath^2\le n_{\mathrm{int}}<u_\jmath^2,\qquad l_\jmath<u_\jmath,
\]
\[
u_\jmath-l_\jmath\le2^{s_{\mathrm{rem}}+1}
\quad\Longrightarrow\quad
u_{\jmath+1}-l_{\jmath+1}\le2^{s_{\mathrm{rem}}},
\qquad s_{\mathrm{rem}}\in\mathbb N_0.
\]
Indeed, in a nonterminal iteration the midpoint lies strictly between the endpoints, and the comparison selects the endpoint preserving the squared bracket. Each possible new width is at most the old width divided by two and rounded upward. A terminal width of one also satisfies the displayed implication. Since
\[
n_{\mathrm{int}}+1\le2^{\kappa(n_{\mathrm{int}}+1)+1},
\]
induction on the remaining iterations shows that after
\(\kappa(n_{\mathrm{int}}+1)+1\) iterations the endpoints are consecutive.
Define \(a_{\mathrm{int}}(n_{\mathrm{int}})\) to be the resulting lower endpoint. Then
\[
a_{\mathrm{int}}(n_{\mathrm{int}})^2
 \le n_{\mathrm{int}}
 <(a_{\mathrm{int}}(n_{\mathrm{int}})+1)^2.
\]
These inequalities uniquely characterize the integer square root. When
\(n_{\mathrm{int}}=0\), the initial bracket is already \([0,1]\) and remains fixed. Define
\[
\mathcal B_{\sqrt{\ }}(q,z)
 =\operatorname{box}\!\left(
 \frac{a_{\mathrm{int}}(n_{\sqrt{\ }}(q,z))}{2^{p_{\mathrm{sc}}(q)}},
 \frac{a_{\mathrm{int}}(n_{\sqrt{\ }}(q,z))+1}{2^{p_{\mathrm{sc}}(q)}}\right).
\]
In particular, this formula returns
\([0,2^{-p_{\mathrm{sc}}(q)}]\) when \(z<0\).
% lean: CausalSmith.Stat.LogoddsLowsmoothFrontier.paperIntegerSqrt_spec

For \(n_{\mathrm{atan}}\in\mathbb N_0\) and \(z\in\mathbb Q\), put
\[
A_{\mathrm{atan}}(n_{\mathrm{atan}},z)
 =\sum_{\jmath=0}^{n_{\mathrm{atan}}-1}
       \frac{(-1)^\jmath z^{2\jmath+1}}{2\jmath+1},
\]
\[
\Sigma_\pi(q)
 =16A_{\mathrm{atan}}(q+7,1/5)
  -4A_{\mathrm{atan}}(q+7,1/239),
\qquad
\mathcal B_\pi(q)
 =\operatorname{box}\!\left(
   \max\{3,\Sigma_\pi(q)-e_{\mathrm{sc}}(q)\},
   \min\{4,\Sigma_\pi(q)+e_{\mathrm{sc}}(q)\}\right).
\]

For the power routine, \(b\in\mathbb N_0\) denotes the primitive integer base and
\(v_{\mathrm{pow}}\in\mathbb Q\) the scalar exponent. Define
\[
B_{\mathrm{pow}}(b)=\max\{2,b\},
\qquad
p_{\mathrm{pow}}(q,b)=q+12+3\kappa(B_{\mathrm{pow}}(b)),
\qquad
\tau_{\mathrm{pow}}(q,b)=2^{-p_{\mathrm{pow}}(q,b)},
\]
\[
\mathcal L_{\mathrm{pow}}(q,b)
 =\mathcal B_{\log}(p_{\mathrm{pow}}(q,b),b),
\qquad
\lambda_{\mathrm{pow}}(q,b)
 =\frac{\underline{\mathcal L_{\mathrm{pow}}(q,b)}
          +\overline{\mathcal L_{\mathrm{pow}}(q,b)}}2,
\]
\[
\zeta_{\mathrm{pow}}(q,b,v_{\mathrm{pow}})
 =v_{\mathrm{pow}}\lambda_{\mathrm{pow}}(q,b),
\qquad
\mathcal J_{\mathrm{pow}}(q,b,v_{\mathrm{pow}})
 =\mathcal B_{\exp}(p_{\mathrm{pow}}(q,b),
                   \zeta_{\mathrm{pow}}(q,b,v_{\mathrm{pow}})),
\]
\[
\delta_{\mathrm{pow}}(q,b)
 =96B_{\mathrm{pow}}(b)^3\tau_{\mathrm{pow}}(q,b).
\]
The scalar output is
\[
\mathcal B_{\mathrm{pow}}(q,b,v_{\mathrm{pow}})=
\begin{cases}
[1,1],&b=1,\\
\operatorname{box}\!\left(
 \max\{1/b^3,
       \underline{\mathcal J_{\mathrm{pow}}(q,b,v_{\mathrm{pow}})}
          -\delta_{\mathrm{pow}}(q,b)\},
 \min\{b^3,
       \overline{\mathcal J_{\mathrm{pow}}(q,b,v_{\mathrm{pow}})}
          +\delta_{\mathrm{pow}}(q,b)\}\right),&b\ne1.
\end{cases}
\]
The rational-division convention makes this formula total at \(b=0\).

For a rational input interval
\[
I_{\mathrm{arg}}=[a_{\mathrm{in}},b_{\mathrm{in}}],
\qquad a_{\mathrm{in}}\le b_{\mathrm{in}},
\]
the engine uses
\[
\mathsf E_0^\pi(q)=\mathcal B_\pi(q),
\qquad
\mathsf E_0^F(q,I_{\mathrm{arg}})
 =\operatorname{box}\!\left(
 \underline{\mathcal B_F(q,a_{\mathrm{in}})},
 \overline{\mathcal B_F(q,b_{\mathrm{in}})}\right),
\quad F\in\{\exp,\log,\sqrt{\ }\},
\]
\[
\mathsf E_0^{\mathrm{pow}}(q,b,I_{\mathrm{arg}})
 =\operatorname{box}\!\left(
 \underline{\mathcal B_{\mathrm{pow}}(q,b,a_{\mathrm{in}})},
 \overline{\mathcal B_{\mathrm{pow}}(q,b,b_{\mathrm{in}})}\right).
\]
For cosine, define the rational midpoint and half-width
\[
z_{\mathrm{mid}}=\frac{a_{\mathrm{in}}+b_{\mathrm{in}}}{2},
\qquad
w_{\mathrm{arg}}=\frac{b_{\mathrm{in}}-a_{\mathrm{in}}}{2},
\]
and use
\[
\mathsf E_0^{\cos}(q,I_{\mathrm{arg}})
 =\operatorname{box}\!\left(
 \max\{-1,\underline{\mathcal B_{\cos}(q,z_{\mathrm{mid}})}
                    -w_{\mathrm{arg}}\},
 \min\{1,\overline{\mathcal B_{\cos}(q,z_{\mathrm{mid}})}
                   +w_{\mathrm{arg}}\}\right).
\]

To certify finite computation, use a finite instruction list acting on rational registers, together with a program counter and a halt flag. Instructions load or copy a rational value, perform rational arithmetic, minimum or maximum, perform integer floor, ceiling, factorial, integer square root, \(\kappa\), or nonnegative integer power, compare registers and branch, jump, or halt. Each instruction has unit cost in this model. Integer arguments are obtained by the nonnegative floor convention above. Input coordinates initialize the registers, missing coordinates are zero, and a halted state remains fixed.

The term recurrences explain both the returned sums and the loop counts. Their current terms are
\[
t_{\exp,\jmath}(z)=\frac{z^\jmath}{\jmath!},
\qquad
t_{\cos,\jmath}(z)=\frac{(-1)^\jmath z^{2\jmath}}{(2\jmath)!},
\]
\[
t_{\mathrm{atan},\jmath}(z)=(-1)^\jmath z^{2\jmath+1},
\qquad
t_{\log,\jmath}(z)=\vartheta_{\log}(z)^{2\jmath+1},
\]
with updates
\[
t_{\exp,\jmath+1}
 =t_{\exp,\jmath}\frac{z}{\jmath+1},
\qquad
t_{\cos,\jmath+1}
 =t_{\cos,\jmath}\frac{-z^2}{(2\jmath+1)(2\jmath+2)},
\]
\[
t_{\mathrm{atan},\jmath+1}
 =-z^2t_{\mathrm{atan},\jmath},
\qquad
t_{\log,\jmath+1}
 =\vartheta_{\log}(z)^2t_{\log,\jmath}.
\]
Starting the exponential and cosine terms at one, and the arctangent and logarithm terms at \(z\) and \(\vartheta_{\log}(z)\), respectively, induction shows that the accumulated registers contain precisely the displayed partial sums.

For example, one joint Machin iteration uses one branch, two divisions by the current odd integer, two additions to the partial sums, two term multiplications, one odd-index increment, one counter decrement, and one jump: ten instructions. Its initialization uses fourteen instructions, and its exit uses seventeen. Thus its public bound is
\[
14+10(q+7)+17.
\]
The corresponding accounting for the other endpoint programs is
\[
\begin{array}{c|l}
\text{Primitive}&\text{Public instruction bound}\\ \hline
\pi&14+10(q+7)+17\\
\exp&
23+(14m_{\mathrm{ser}}(q,a_{\mathrm{in}})+15)
 +(3+20)+(14m_{\mathrm{ser}}(q,b_{\mathrm{in}})+15)+2\\
\cos&26+9m_{\mathrm{ser}}(q,z_{\mathrm{mid}})+20\\
\log&
28+(7n_{\log}(q,a_{\mathrm{in}})+14)
 +(3+25)+(7n_{\log}(q,b_{\mathrm{in}})+14)+2\\
\sqrt{\ }&14 .
\end{array}
\]
The exponential and logarithm iterations each use seven instructions; the cosine iteration uses nine. Their fixed blocks initialize the terms, preserve the left endpoint while evaluating the right endpoint, form and sort the scalar endpoints, and halt. Induction on the remaining loop count, using the term identities, proves the endpoint-output identities at these bounds. The logarithm exit returns zero on nonpositive scalar inputs.

The square-root program increments the precision and forms its dyadic scale in four instructions, forms the scaled square and lower endpoint in four, forms the upper endpoint in five, and halts in one. Its integer-square-root instruction agrees with the binary recurrence by the uniqueness established above. Consequently its output registers contain
\[
\frac{a_{\mathrm{int}}(n_{\sqrt{\ }}(q,a_{\mathrm{in}}))}
     {2^{p_{\mathrm{sc}}(q)}},
\qquad
\frac{a_{\mathrm{int}}(n_{\sqrt{\ }}(q,b_{\mathrm{in}}))+1}
     {2^{p_{\mathrm{sc}}(q)}}.
\]
% lean: CausalSmith.Stat.LogoddsLowsmoothFrontier.concreteEngine_sqrt_register_certificate

For power, the public bound must be computable from the original inputs. Define, for all \(q,b\in\mathbb N_0\),
\[
n_{\mathrm{pow},\log}(q,b)
 =n_{\log}(p_{\mathrm{pow}}(q,b),b).
\]
The symmetric logarithm bracket gives the exact midpoint identity
\[
\lambda_{\mathrm{pow}}(q,b)=
\begin{cases}
\Sigma_{\log}(p_{\mathrm{pow}}(q,b),b),&b>0,\\
0,&b=0.
\end{cases}
\]
Since \(b\ge0\),
\[
\left|\frac{b-1}{b+1}\right|\le1,
\qquad
\left|
 \frac{((b-1)/(b+1))^{2\jmath+1}}{2\jmath+1}
\right|\le1.
\]
The triangle inequality for the finite sum therefore gives
\[
|\Sigma_{\log}(p_{\mathrm{pow}}(q,b),b)|
 \le2n_{\mathrm{pow},\log}(q,b),
\qquad
|\zeta_{\mathrm{pow}}(q,b,v_{\mathrm{pow}})|
 \le2n_{\mathrm{pow},\log}(q,b)|v_{\mathrm{pow}}|.
\]
Define
\[
B_{\mathrm{pow},\mathrm{budget}}(q,b,v_{\mathrm{pow}})
 =\max\{1,\lceil
       2n_{\mathrm{pow},\log}(q,b)|v_{\mathrm{pow}}|
       \rceil_+\}.
\]
Monotonicity of ceiling yields
\[
B_{\mathrm{ser}}(\zeta_{\mathrm{pow}}(q,b,v_{\mathrm{pow}}))
 \le B_{\mathrm{pow},\mathrm{budget}}(q,b,v_{\mathrm{pow}}).
\]
The fused scalar program computes the logarithm midpoint, evaluates the exponential at the perturbed argument, and performs the displayed clipping. The logarithm stage is bounded by
\(45+7n_{\mathrm{pow},\log}(q,b)\) instructions. Connecting to the exponential stage uses one instruction; preparing and evaluating the exponential at a repeated scalar argument uses
\(6+78+28m_{\mathrm{ser}}(p_{\mathrm{pow}},\zeta_{\mathrm{pow}})\);
the clipping continuation is bounded by thirty-one further instructions. Hence a public scalar bound is
\[
\begin{split}
f_{\mathrm{pow},\mathrm{sc}}(q,b,v_{\mathrm{pow}})
={}&161+7n_{\mathrm{pow},\log}(q,b)\\
&+28\max\{
  2B_{\mathrm{pow},\mathrm{budget}}(q,b,v_{\mathrm{pow}})^2,
  p_{\mathrm{pow}}(q,b)+6\}.
\end{split}
\]
This bounds the actual fused run because
\[
45+1+6+78+31=161
\]
and the displayed bound on \(B_{\mathrm{ser}}(\zeta_{\mathrm{pow}})\)
bounds the exponential loop count.
% lean: CausalSmith.Stat.LogoddsLowsmoothFrontier.powerScalarFuel_le_public

For an exponent interval, put
\[
M_{\mathrm{pow},\mathrm{arg}}
 =\max\{|a_{\mathrm{in}}|,|b_{\mathrm{in}}|\},
\qquad
f_{\mathrm{pow}}(q,b,I_{\mathrm{arg}})
 =90+2f_{\mathrm{pow},\mathrm{sc}}
       (q,b,M_{\mathrm{pow},\mathrm{arg}}).
\]
The scalar bound is monotone in the absolute exponent, so it bounds both endpoint calls. When \(b=1\), the endpoint program returns \([1,1]\) and halts after nine instructions. Otherwise its prelude uses at most forty-six instructions; it runs the two scalar calls, uses forty-one connecting instructions, restores the selected endpoints in two, and halts in one. Thus its count is at most
\[
\begin{split}
&46+f_{\mathrm{pow},\mathrm{sc}}(q,b,a_{\mathrm{in}})
 +41+f_{\mathrm{pow},\mathrm{sc}}(q,b,b_{\mathrm{in}})+2+1\\
&\hspace{2cm}\le
90+2f_{\mathrm{pow},\mathrm{sc}}
       (q,b,M_{\mathrm{pow},\mathrm{arg}}).
\end{split}
\]
The case \(b=0\) uses the forty-six-instruction prelude and the totalized scalar formulas. The returned endpoints are the lower endpoint of the left scalar call and the upper endpoint of the right scalar call; normalization therefore gives exactly
\(\mathsf E_0^{\mathrm{pow}}(q,b,I_{\mathrm{arg}})\).
Running to the public bound preserves this output because the halted state is fixed.
% lean: CausalSmith.Stat.LogoddsLowsmoothFrontier.concreteEngine_power_register_certificate

Every displayed bound is itself computable by a fixed finite register program from the precision, rational input endpoints, and base when present. On arbitrary rational input lists, precision and base are interpreted by nonnegative floor and missing coordinates by zero. Thus the bounds and output identities hold on the full register-input domain. The six programs together supply finite rational recurrences for all six primitive maps. Their definitions use precisely these primitive arguments, so the choice of \(\mathsf E_0\) is fixed independently of the statistical and naming inputs in \cref{obj:def:arithmetic-engine}.
% lean: CausalSmith.Stat.LogoddsLowsmoothFrontier.concreteEngine_registers_of_other_primitives

\item \emph{The constant enclosure.}
For \(z\ge0\), the finite geometric identity
\[
(1+x_{\mathrm{arg}}^2)
 \sum_{\jmath=0}^{n_{\mathrm{atan}}-1}
       (-1)^\jmath x_{\mathrm{arg}}^{2\jmath}
 =
1-(-1)^{n_{\mathrm{atan}}}
       x_{\mathrm{arg}}^{2n_{\mathrm{atan}}}
\]
follows by induction on \(n_{\mathrm{atan}}\), starting with the empty sum. Dividing by \(1+x_{\mathrm{arg}}^2\) and integrating gives
\[
\arctan z-A_{\mathrm{atan}}(n_{\mathrm{atan}},z)
 =
\int_0^z
 \frac{(-1)^{n_{\mathrm{atan}}}
        x_{\mathrm{arg}}^{2n_{\mathrm{atan}}}}
      {1+x_{\mathrm{arg}}^2}\,dx_{\mathrm{arg}}.
\]
On \(0\le x_{\mathrm{arg}}\le z\), the integrand has absolute value at most
\(z^{2n_{\mathrm{atan}}}\). Consequently
\[
|\arctan z-A_{\mathrm{atan}}(n_{\mathrm{atan}},z)|
 \le z^{2n_{\mathrm{atan}}+1}.
\]
Machin's identity,
\[
\pi=16\arctan(1/5)-4\arctan(1/239),
\]
and the triangle inequality now imply
\[
|\pi-\Sigma_\pi(q)|
 \le16(1/5)^{2(q+7)+1}
    +4(1/239)^{2(q+7)+1}
 \le20(1/5)^{2(q+7)+1}.
\]
The remaining calibration is
\[
20(1/5)^{2(q+7)+1}\le2^{-(q+3)}=e_{\mathrm{sc}}(q).
\]
For \(q=0\), this is \(20\,5^{-15}\le2^{-3}\). Increasing \(q\) by one multiplies its left side by \(1/25\) and its right side by \(1/2\), proving it for every \(q\in\mathbb N_0\).

Together with \(3<\pi<4\), the approximation bound gives
\[
\max\{3,\Sigma_\pi(q)-e_{\mathrm{sc}}(q)\}
 \le\pi\le
\min\{4,\Sigma_\pi(q)+e_{\mathrm{sc}}(q)\}.
\]
Thus these endpoints are already ordered, their interval lies in \([3,4]\), and
\[
0\le\pi-\underline{\mathcal B_\pi(q)}
 \le2e_{\mathrm{sc}}(q)\le\epsilon_q,
\qquad
0\le\overline{\mathcal B_\pi(q)}-\pi
 \le2e_{\mathrm{sc}}(q)\le\epsilon_q.
\]
This proves the constant contract.
% lean: CausalSmith.Stat.LogoddsLowsmoothFrontier.paperPi_contract

\item \emph{The exponential contract.}
For \(z\in\mathbb Q\), the definitions give
\[
|z|\le B_{\mathrm{ser}}(z),\qquad
B_{\mathrm{ser}}(z)\ge1,\qquad
m_{\mathrm{ser}}(q,z)\ge2B_{\mathrm{ser}}(z)^2.
\]
The factorial inequality
\[
m_{\mathrm{ser}}(q,z)!\,
m_{\mathrm{ser}}(q,z)^{m_{\mathrm{ser}}(q,z)}
 \le(2m_{\mathrm{ser}}(q,z))!
\]
follows by bounding each of the last
\(m_{\mathrm{ser}}(q,z)\) factorial factors from below by
\(m_{\mathrm{ser}}(q,z)\). Since the first factorial factor on the left is at least one,
\[
\frac{|z|^{2m_{\mathrm{ser}}(q,z)}}
     {(2m_{\mathrm{ser}}(q,z))!}
 \le
 \left(\frac{B_{\mathrm{ser}}(z)^2}{m_{\mathrm{ser}}(q,z)}\right)^{
       m_{\mathrm{ser}}(q,z)}
 \le2^{-m_{\mathrm{ser}}(q,z)}.
\]
Also,
\[
\frac{|z|}{2m_{\mathrm{ser}}(q,z)+1}\le\frac12.
\]
The exponential-series remainder estimate states that, for
\(\zeta_{\mathrm{cx}}\in\mathbb C\) and \(n_{\mathrm{cut}}\in\mathbb N_0\),
\[
\frac{|\zeta_{\mathrm{cx}}|}{n_{\mathrm{cut}}+1}\le\frac12
\quad\Longrightarrow\quad
\left|
 \exp(\zeta_{\mathrm{cx}})
 -\sum_{\jmath=0}^{n_{\mathrm{cut}}-1}
       \frac{\zeta_{\mathrm{cx}}^\jmath}{\jmath!}
\right|
 \le2\frac{|\zeta_{\mathrm{cx}}|^{n_{\mathrm{cut}}}}
               {n_{\mathrm{cut}}!}.
\]
Applying it to the real argument \(z\) and using
\(m_{\mathrm{ser}}(q,z)\ge q+6\ge q+4\) gives
\[
|\exp z-\Sigma_{\exp}(q,z)|
 \le2\,2^{-m_{\mathrm{ser}}(q,z)}
 \le2\,2^{-(q+4)}
 =e_{\mathrm{sc}}(q).
\]

The positive floor lies below the exact value: since
\(-B_{\mathrm{ser}}(z)\le z\) and \(\exp(1)<3\),
\[
3^{-B_{\mathrm{ser}}(z)}
 \le(\exp1)^{-B_{\mathrm{ser}}(z)}
 =\exp(-B_{\mathrm{ser}}(z))
 \le\exp z.
\]
It follows that the two raw scalar endpoints are ordered and that
\[
0<\underline{\mathcal B_{\exp}(q,z)}
 \le\exp z\le\overline{\mathcal B_{\exp}(q,z)},
\]
\[
0\le\exp z-\underline{\mathcal B_{\exp}(q,z)}
 \le2e_{\mathrm{sc}}(q)\le\epsilon_q,
\qquad
0\le\overline{\mathcal B_{\exp}(q,z)}-\exp z
 \le2e_{\mathrm{sc}}(q)\le\epsilon_q.
\]

For the endpoint extension, exponential monotonicity gives the ordering
\[
\underline{\mathcal B_{\exp}(q,a_{\mathrm{in}})}
 \le\exp a_{\mathrm{in}}
 \le\exp b_{\mathrm{in}}
 \le\overline{\mathcal B_{\exp}(q,b_{\mathrm{in}})}.
\]
For every \(x_{\mathrm{arg}}\in I_{\mathrm{arg}}\),
\[
\underline{\mathcal B_{\exp}(q,a_{\mathrm{in}})}
 \le\exp a_{\mathrm{in}}
 \le\exp x_{\mathrm{arg}}
 \le\exp b_{\mathrm{in}}
 \le\overline{\mathcal B_{\exp}(q,b_{\mathrm{in}})}.
\]
The lower and upper endpoint excesses are therefore the scalar excesses at
\(a_{\mathrm{in}}\) and \(b_{\mathrm{in}}\), respectively, and each is at most
\(\epsilon_q\). The returned lower endpoint is positive by the scalar sign bound. This proves the exponential contract.
% lean: CausalSmith.Stat.LogoddsLowsmoothFrontier.concreteEngine_exp_contract

\item \emph{The logarithm contract.}
Let \(z>0\) be rational. The ceiling bounds imply
\[
B_{\log}(z)\ge2,\qquad
z\le B_{\log}(z),\qquad
1/z\le B_{\log}(z).
\]
Using \(1\le B_{\log}(z)z\) and clearing the positive denominators in
\(\vartheta_{\log}(z)\) yields
\[
|\vartheta_{\log}(z)|
 \le\frac{B_{\log}(z)-1}{B_{\log}(z)+1}<1.
\]
Squaring this inequality gives
\[
1-\vartheta_{\log}(z)^2
 \ge1-\left(\frac{B_{\log}(z)-1}{B_{\log}(z)+1}\right)^2
 \ge\frac1{B_{\log}(z)},
\]
where the last inequality follows by clearing denominators and using
\(B_{\log}(z)\ge2\).

To bound the numerator of the remainder, observe that
\[
|\vartheta_{\log}(z)|
 \le1-\frac1{B_{\log}(z)}
 \le\exp(-1/B_{\log}(z)).
\]
Since
\[
-\frac{2B_{\log}(z)c_{\log}(q,z)+1}{B_{\log}(z)}
 \le-c_{\log}(q,z),
\]
and \(\exp(1)\ge2\), we obtain
\[
\begin{split}
|\vartheta_{\log}(z)|^{2n_{\log}(q,z)+1}
&\le
 \exp\!\left(-\frac{2B_{\log}(z)c_{\log}(q,z)+1}
                  {B_{\log}(z)}\right)\\
&\le\exp(-c_{\log}(q,z))
 =\exp(-1)^{c_{\log}(q,z)}
 \le2^{-c_{\log}(q,z)}.
\end{split}
\]
The logarithm-series remainder estimate, applied using
\[
\frac{1+\vartheta_{\log}(z)}{1-\vartheta_{\log}(z)}=z,
\]
now gives
\[
\begin{split}
|\log z-\Sigma_{\log}(q,z)|
&\le
 \frac{2|\vartheta_{\log}(z)|^{2n_{\log}(q,z)+1}}
      {1-\vartheta_{\log}(z)^2}\\
&\le2B_{\log}(z)\,2^{-c_{\log}(q,z)}
 \le e_{\mathrm{sc}}(q).
\end{split}
\]
For the last inequality,
\(8B_{\log}(z)\le2^{\kappa(8B_{\log}(z))}\), so
\[
2B_{\log}(z)\,2^{-c_{\log}(q,z)}
 =
2^{-(q+3)}
 \frac{2B_{\log}(z)}
      {4\,2^{\kappa(8B_{\log}(z))}}
 \le2^{-(q+3)}.
\]
Thus the symmetric scalar bracket satisfies
\[
\underline{\mathcal B_{\log}(q,z)}
 \le\log z\le\overline{\mathcal B_{\log}(q,z)},
\]
\[
0\le\log z-\underline{\mathcal B_{\log}(q,z)}
 \le2e_{\mathrm{sc}}(q)\le\epsilon_q,
\qquad
0\le\overline{\mathcal B_{\log}(q,z)}-\log z
 \le2e_{\mathrm{sc}}(q)\le\epsilon_q.
\]

Now suppose \(0<a_{\mathrm{in}}\le b_{\mathrm{in}}\). Both endpoint evaluations lie in the positive logarithm domain, and monotonicity gives
\[
\underline{\mathcal B_{\log}(q,a_{\mathrm{in}})}
 \le\log a_{\mathrm{in}}
 \le\log b_{\mathrm{in}}
 \le\overline{\mathcal B_{\log}(q,b_{\mathrm{in}})}.
\]
Every \(x_{\mathrm{arg}}\in I_{\mathrm{arg}}\) is positive and satisfies
\[
\underline{\mathcal B_{\log}(q,a_{\mathrm{in}})}
 \le\log a_{\mathrm{in}}
 \le\log x_{\mathrm{arg}}
 \le\log b_{\mathrm{in}}
 \le\overline{\mathcal B_{\log}(q,b_{\mathrm{in}})}.
\]
The endpoint extension therefore contains the whole image and has outward endpoint excess at most \(\epsilon_q\).
% lean: CausalSmith.Stat.LogoddsLowsmoothFrontier.concreteEngine_log_contract

\item \emph{The square-root contract.}
For rational \(z\ge0\), the floor and integer-square-root brackets give
\[
a_{\mathrm{int}}(n_{\sqrt{\ }}(q,z))^2
 \le n_{\sqrt{\ }}(q,z)
 \le2^{2p_{\mathrm{sc}}(q)}z
 <n_{\sqrt{\ }}(q,z)+1
 \le(a_{\mathrm{int}}(n_{\sqrt{\ }}(q,z))+1)^2.
\]
Division by \(2^{2p_{\mathrm{sc}}(q)}\), followed by taking nonnegative square roots, yields
\[
0\le
\frac{a_{\mathrm{int}}(n_{\sqrt{\ }}(q,z))}
     {2^{p_{\mathrm{sc}}(q)}}
 \le\sqrt z\le
\frac{a_{\mathrm{int}}(n_{\sqrt{\ }}(q,z))+1}
     {2^{p_{\mathrm{sc}}(q)}}.
\]
The endpoints are ordered and their difference is exactly
\[
\operatorname{wid}(\mathcal B_{\sqrt{\ }}(q,z))
 =2^{-p_{\mathrm{sc}}(q)}
 =e_{\mathrm{sc}}(q)\le\epsilon_q.
\]
Containment consequently gives both scalar excess bounds:
\[
0\le\sqrt z-\underline{\mathcal B_{\sqrt{\ }}(q,z)}
 \le e_{\mathrm{sc}}(q),
\qquad
0\le\overline{\mathcal B_{\sqrt{\ }}(q,z)}-\sqrt z
 \le e_{\mathrm{sc}}(q).
\]

Suppose \(0\le a_{\mathrm{in}}\le b_{\mathrm{in}}\). The endpoint evaluations satisfy
\[
\underline{\mathcal B_{\sqrt{\ }}(q,a_{\mathrm{in}})}
 \le\sqrt{a_{\mathrm{in}}}
 \le\sqrt{b_{\mathrm{in}}}
 \le\overline{\mathcal B_{\sqrt{\ }}(q,b_{\mathrm{in}})}.
\]
For \(x_{\mathrm{arg}}\in I_{\mathrm{arg}}\), square-root monotonicity gives
\[
\underline{\mathcal B_{\sqrt{\ }}(q,a_{\mathrm{in}})}
 \le\sqrt{x_{\mathrm{arg}}}
 \le\overline{\mathcal B_{\sqrt{\ }}(q,b_{\mathrm{in}})}.
\]
The two endpoint excesses are bounded by the corresponding scalar widths, hence by \(\epsilon_q\), and the returned lower endpoint is nonnegative. The argument includes \(a_{\mathrm{in}}=0\) and degenerate intervals.
% lean: CausalSmith.Stat.LogoddsLowsmoothFrontier.concreteEngine_sqrt_contract

\item \emph{The power contract.}
Let \(b\ge1\) and \(-3\le v_{\mathrm{pow}}\le3\). If \(b=1\), the scalar interval
\([1,1]\) has exact containment, zero excess, and positive lower endpoint.

For \(b\ne1\), use the quantities defining the scalar routine above. The precision calibration is
\[
193B_{\mathrm{pow}}(b)^3\tau_{\mathrm{pow}}(q,b)
 \le\epsilon_q,
\qquad
\tau_{\mathrm{pow}}(q,b)\le\frac1{4096}.
\]
Indeed,
\[
B_{\mathrm{pow}}(b)\le2^{\kappa(B_{\mathrm{pow}}(b))}
\quad\Longrightarrow\quad
B_{\mathrm{pow}}(b)^3
 \le2^{3\kappa(B_{\mathrm{pow}}(b))},
\]
and therefore
\[
193B_{\mathrm{pow}}(b)^3\tau_{\mathrm{pow}}(q,b)
 \le\frac{193}{4096}\,2^{-q}\le\epsilon_q.
\]
The second bound follows from \(q,\kappa(B_{\mathrm{pow}}(b))\ge0\).

Exponent monotonicity for \(b\ge1\) gives the exact clip range
\[
\frac1{b^3}=b^{-3}\le b^{v_{\mathrm{pow}}}\le b^3.
\]
The scalar logarithm contract, applied at precision \(p_{\mathrm{pow}}(q,b)\), implies
\[
|\lambda_{\mathrm{pow}}(q,b)-\log b|
 \le\tau_{\mathrm{pow}}(q,b),
\]
because each endpoint of \(\mathcal L_{\mathrm{pow}}(q,b)\) is within that tolerance of \(\log b\).

Define the exact argument and the argument perturbation by
\[
\eta_{\mathrm{pow}}(b,v_{\mathrm{pow}})
 =v_{\mathrm{pow}}\log b,
\qquad
t_{\mathrm{pow}}(q,b,v_{\mathrm{pow}})
 =\zeta_{\mathrm{pow}}(q,b,v_{\mathrm{pow}})
  -\eta_{\mathrm{pow}}(b,v_{\mathrm{pow}}).
\]
Then
\[
|t_{\mathrm{pow}}(q,b,v_{\mathrm{pow}})|
 \le3\tau_{\mathrm{pow}}(q,b)\le1.
\]
The local exponential inequality
\[
|\exp t-1|\le2|t|\qquad (t\in[-1,1])
\]
and the identities
\[
\exp(\eta_{\mathrm{pow}}(b,v_{\mathrm{pow}}))
 =b^{v_{\mathrm{pow}}},
\]
\[
\exp(\zeta_{\mathrm{pow}})
 -\exp(\eta_{\mathrm{pow}})
 =
\exp(\eta_{\mathrm{pow}})\bigl(\exp(t_{\mathrm{pow}})-1\bigr)
\]
yield
\[
\begin{split}
|\exp(\zeta_{\mathrm{pow}})-b^{v_{\mathrm{pow}}}|
&\le B_{\mathrm{pow}}(b)^3\,2|t_{\mathrm{pow}}|\\
&\le6B_{\mathrm{pow}}(b)^3\tau_{\mathrm{pow}}(q,b)\\
&\le96B_{\mathrm{pow}}(b)^3\tau_{\mathrm{pow}}(q,b)
 =\delta_{\mathrm{pow}}(q,b).
\end{split}
\]
Here the bound on \(\exp(\eta_{\mathrm{pow}})\) follows from
\(b^{v_{\mathrm{pow}}}\le b^3\le B_{\mathrm{pow}}(b)^3\).
% lean: CausalSmith.Stat.LogoddsLowsmoothFrontier.power_midpoint_value_error

The exponential bracket \(\mathcal J_{\mathrm{pow}}\) contains
\(\exp(\zeta_{\mathrm{pow}})\), with each endpoint excess at most
\(\tau_{\mathrm{pow}}\). Hence
\[
\underline{\mathcal J_{\mathrm{pow}}}-\delta_{\mathrm{pow}}
 \le b^{v_{\mathrm{pow}}}
 \le\overline{\mathcal J_{\mathrm{pow}}}+\delta_{\mathrm{pow}}.
\]
Combining this with the exact clip range proves that the raw endpoints of
\(\mathcal B_{\mathrm{pow}}\) are ordered and contain \(b^{v_{\mathrm{pow}}}\).
Its lower endpoint is at least \(1/b^3>0\). Since clipping increases the lower endpoint and decreases the upper endpoint, both outward excesses are at most
\[
\tau_{\mathrm{pow}}+2\delta_{\mathrm{pow}}
 =(1+192B_{\mathrm{pow}}(b)^3)\tau_{\mathrm{pow}}
 \le193B_{\mathrm{pow}}(b)^3\tau_{\mathrm{pow}}
 \le\epsilon_q,
\]
where \(B_{\mathrm{pow}}(b)^3\ge1\). Thus
\[
0<\underline{\mathcal B_{\mathrm{pow}}(q,b,v_{\mathrm{pow}})}
 \le b^{v_{\mathrm{pow}}}
 \le\overline{\mathcal B_{\mathrm{pow}}(q,b,v_{\mathrm{pow}})},
\]
with the required scalar endpoint excesses.

Now let
\[
-3\le a_{\mathrm{in}}\le b_{\mathrm{in}}\le3.
\]
Both endpoint exponents satisfy the scalar hypotheses. Because
\(v_{\mathrm{pow}}\mapsto b^{v_{\mathrm{pow}}}\) is monotone,
\[
\underline{\mathcal B_{\mathrm{pow}}(q,b,a_{\mathrm{in}})}
 \le b^{a_{\mathrm{in}}}
 \le b^{b_{\mathrm{in}}}
 \le\overline{\mathcal B_{\mathrm{pow}}(q,b,b_{\mathrm{in}})}.
\]
For every \(x_{\mathrm{arg}}\in I_{\mathrm{arg}}\),
\[
\underline{\mathcal B_{\mathrm{pow}}(q,b,a_{\mathrm{in}})}
 \le b^{x_{\mathrm{arg}}}
 \le\overline{\mathcal B_{\mathrm{pow}}(q,b,b_{\mathrm{in}})}.
\]
The endpoint extension inherits the two scalar excess bounds and the positive lower endpoint. This includes \(b=1\), the exponent endpoints \(-3,3\), and degenerate intervals.
% lean: CausalSmith.Stat.LogoddsLowsmoothFrontier.concreteEngine_power_contract

\item \emph{The cosine contract.}
Let \(\mathrm i\in\mathbb C\) be defined by
\[
\operatorname{Re}\mathrm i=0,\qquad
\operatorname{Im}\mathrm i=1.
\]
Pairing the even and odd terms gives
\[
\operatorname{Re}\left(
 \sum_{\jmath=0}^{2m_{\mathrm{ser}}(q,z)-1}
       \frac{(z\mathrm i)^\jmath}{\jmath!}\right)
 =\Sigma_{\cos}(q,z).
\]
To verify this identity, start with the empty sum and append successive pairs. In each pair,
\[
(z\mathrm i)^{2\jmath}=(-1)^\jmath z^{2\jmath},
\qquad
\operatorname{Re}(z\mathrm i)^{2\jmath+1}=0,
\]
so the real part increases by exactly the next cosine term.

Since \(|z\mathrm i|=|z|\) and
\(\operatorname{Re}\exp(z\mathrm i)=\cos z\), the complex exponential remainder estimate and the bounds on its first omitted term give
\[
\begin{split}
|\cos z-\Sigma_{\cos}(q,z)|
&\le
 \left|\exp(z\mathrm i)
  -\sum_{\jmath=0}^{2m_{\mathrm{ser}}(q,z)-1}
       \frac{(z\mathrm i)^\jmath}{\jmath!}\right|\\
&\le2\frac{|z|^{2m_{\mathrm{ser}}(q,z)}}
               {(2m_{\mathrm{ser}}(q,z))!}
 \le e_{\mathrm{sc}}(q).
\end{split}
\]
Together with \(-1\le\cos z\le1\), this implies that the clipped scalar endpoints are ordered and satisfy
\[
-1\le\underline{\mathcal B_{\cos}(q,z)}
 \le\cos z\le\overline{\mathcal B_{\cos}(q,z)}\le1,
\]
\[
0\le\cos z-\underline{\mathcal B_{\cos}(q,z)}
 \le2e_{\mathrm{sc}}(q)\le\epsilon_q,
\qquad
0\le\overline{\mathcal B_{\cos}(q,z)}-\cos z
 \le2e_{\mathrm{sc}}(q)\le\epsilon_q.
\]

For the interval routine, define its midpoint enclosure endpoints
\[
a_{\cos,\mathrm{mid}}
 =\underline{\mathcal B_{\cos}(q,z_{\mathrm{mid}})},
\qquad
b_{\cos,\mathrm{mid}}
 =\overline{\mathcal B_{\cos}(q,z_{\mathrm{mid}})}.
\]
They obey
\[
a_{\cos,\mathrm{mid}}\le\cos z_{\mathrm{mid}}
 \le b_{\cos,\mathrm{mid}},
\qquad
\cos z_{\mathrm{mid}}-a_{\cos,\mathrm{mid}}\le\epsilon_q,
\qquad
b_{\cos,\mathrm{mid}}-\cos z_{\mathrm{mid}}\le\epsilon_q,
\]
\[
-1\le a_{\cos,\mathrm{mid}}\le b_{\cos,\mathrm{mid}}\le1.
\]
Because \(w_{\mathrm{arg}}\ge0\), these inequalities also give
\[
\max\{-1,a_{\cos,\mathrm{mid}}-w_{\mathrm{arg}}\}
 \le
\min\{1,b_{\cos,\mathrm{mid}}+w_{\mathrm{arg}}\}.
\]
Thus the interval routine has exactly the prescribed endpoints
\[
\underline{\mathsf E_0^{\cos}(q,I_{\mathrm{arg}})}
 =\max\{-1,a_{\cos,\mathrm{mid}}-w_{\mathrm{arg}}\},
\qquad
\overline{\mathsf E_0^{\cos}(q,I_{\mathrm{arg}})}
 =\min\{1,b_{\cos,\mathrm{mid}}+w_{\mathrm{arg}}\}.
\]

For every \(x_{\mathrm{arg}}\in I_{\mathrm{arg}}\),
\[
|x_{\mathrm{arg}}-z_{\mathrm{mid}}|\le w_{\mathrm{arg}}.
\]
The cosine Lipschitz inequality therefore yields
\[
|\cos x_{\mathrm{arg}}-\cos z_{\mathrm{mid}}|
 \le|x_{\mathrm{arg}}-z_{\mathrm{mid}}|
 \le w_{\mathrm{arg}}.
\]
Combining this with the midpoint enclosure and the global cosine range proves
\[
\max\{-1,a_{\cos,\mathrm{mid}}-w_{\mathrm{arg}}\}
 \le\cos x_{\mathrm{arg}}
 \le\min\{1,b_{\cos,\mathrm{mid}}+w_{\mathrm{arg}}\}.
\]
Finally, the endpoint formulas imply
\[
\begin{split}
\operatorname{wid}(\mathsf E_0^{\cos}(q,I_{\mathrm{arg}}))
&\le b_{\cos,\mathrm{mid}}-a_{\cos,\mathrm{mid}}
       +2w_{\mathrm{arg}}\\
&\le2\epsilon_q+2w_{\mathrm{arg}}
 =(b_{\mathrm{in}}-a_{\mathrm{in}})+2\epsilon_q.
\end{split}
\]
This establishes the midpoint excess, enlargement, clipping, image enclosure, and width requirements, including zero-width input intervals.
% lean: CausalSmith.Stat.LogoddsLowsmoothFrontier.concreteEngine_cos_contract
\end{enumerate}

The finite rational recurrences and the six primitive contracts establish admissibility under \cref{obj:def:arithmetic-engine}.
% lean: CausalSmith.Stat.LogoddsLowsmoothFrontier.concreteEngine_admissible
\end{proof}

The witness in \cref{obj:lem:concrete-arithmetic-engine} establishes existence at the primitive level. Its finite rational routines provide one concrete choice for the interval construction. The precision guarantees below apply uniformly to every admissible engine, including this witness.

\subsection{Uniform realization at fixed precision}

Rank selection follows \cref{obj:def:resolutions}: the public exponent intervals are intersected with the prescribed exponent boxes, and the ceilings of the upper target-enclosure endpoints determine the two integer ranks. These ranks are fixed before sampling. Endpoint evaluation then uses precisely those retained integers in the coordinate estimates, variance envelope, error radii, and precision budget of \cref{obj:def:upper-handle}. At that step, the queried exponent and covariate intervals enter the endpoint expression tree in their supplied form.

The next statement establishes admissibility of the concrete naming policy and derives the rank, endpoint, and measurability guarantees from the primitive contracts.

\begin{lemmav}{lem:fixed-budget-realization}[Fixed-budget realization and measurability]
The dyadic-floor naming policy \(\mathsf N_0\) is admissible in the sense of \cref{obj:def:dyadic-name-convention}. Moreover, suppose that:
\begin{itemize}
\item \textbf{(Engine and names.)} \(\mathsf E\) is an admissible rational interval-arithmetic engine in the sense of \cref{obj:def:arithmetic-engine}, and \(\mathsf N\) is an admissible Borel naming policy in the sense of \cref{obj:def:dyadic-name-convention}.
\item \textbf{(Sample size.)} \(n\) is an integer with \(n\ge2\).
\item \textbf{(Public exponents.)} The real smoothness exponents satisfy
\[
0<\beta<\frac14,\qquad \beta<\alpha\le1.
\]
\end{itemize}
For every record tuple \(\mathcal O\in([0,1]\times\{0,1\}\times\{0,1\})^n\), let \(k_C,k_S\) be the ranks selected according to \cref{obj:def:resolutions}, with target resolutions
\[
R_C=n^{2/(2\alpha+2\beta+1)},\qquad
R_S=n^{2/(4\beta+1)}.
\]
The prescribed public budgets are
\[
q_{\mathrm{rank}}(n)=7+3\lceil\log_2 n\rceil,\qquad
q_{\mathrm{end}}(n,k_C,k_S)
=44+\lceil\log_2(n+1)\rceil
 +4\lceil\log_2(K+1)\rceil,\qquad
K=\max\{k_C,k_S\}.
\]
At the rank budget, both rational target enclosures \([a_C,b_C]\) and \([a_S,b_S]\) are successfully returned and satisfy
\[
b_C-a_C\le1,\qquad b_S-a_S\le1,
\]
\[
R_C\le k_C<R_C+2\le3R_C,\qquad
R_S\le k_S<R_S+2\le3R_S.
\]
Here \(b_C,b_S\) denote the upper endpoints of the resolution enclosures.

For the endpoint construction, use \(b_C,b_S\) instead to denote the coordinate error radii
\[
b_C=8k_C^{-(\alpha+\beta)}+\sqrt{20V_n(k_C)},\qquad
b_S=25k_S^{-2\beta}+\sqrt{20V_n(k_S)},
\]
where \(V_n(k)\) is the public variance envelope at rank \(k\). Let \(\widehat C_n,\widehat S_n\) be the numerator and denominator coordinate estimates, and set
\[
c_\pm=\widehat C_n\pm b_C,\qquad
s_\pm=\operatorname{clip}_{[s_0,1]}(\widehat S_n\pm b_S),
\]
where \(s_0\) is the denominator floor. With \(F(c,d)\) denoting the clipped logarithmic inversion of the coordinate ratio, the ideal lower and upper endpoints are
\[
L=\min\{F(c_-,s_-),F(c_-,s_+)\},\qquad
R=\max\{F(c_+,s_-),F(c_+,s_+)\}.
\]
At the endpoint budget, rational enclosures \([L_-,L_+]\) and \([R_-,R_+]\) are successfully returned with
\[
L\in[L_-,L_+],\qquad R\in[R_-,R_+],\qquad
L_+-L_-\le\frac1n,\qquad R_+-R_-\le\frac1n.
\]
The literal output satisfies
\[
[L,R]\subseteq
I^{\mathrm{up}}_{n,\mathsf E,\mathsf N}(\alpha,\beta)(\mathcal O),
\qquad
\operatorname{length}\!\left(
I^{\mathrm{up}}_{n,\mathsf E,\mathsf N}(\alpha,\beta)(\mathcal O)
\right)
\le R-L+\frac2n.
\]
It is a nonempty closed connected interval with rational endpoints in \([-1/2,1/2]\), and its output map, as a function of the record tuple, is Borel measurable.
\end{lemmav}

\begin{proof}[Proof of \cref{obj:lem:fixed-budget-realization}]
\begin{enumerate}
\item \textbf{Admissibility of the dyadic-floor policy.}
For a real input \(v_{\mathrm{name}}\in\mathbb R\) and an integer precision \(q\ge0\), define
\[
m_q(v_{\mathrm{name}})
=\lfloor 2^qv_{\mathrm{name}}\rfloor,
\qquad
\mathcal J_q(v_{\mathrm{name}})
=
\left[
\frac{m_q(v_{\mathrm{name}})}{2^q},
\frac{m_q(v_{\mathrm{name}})+1}{2^q}
\right].
\]
The floor inequalities give
\[
m_q(v_{\mathrm{name}})
\le 2^qv_{\mathrm{name}}
<m_q(v_{\mathrm{name}})+1,
\]
and therefore
\[
v_{\mathrm{name}}\in\mathcal J_q(v_{\mathrm{name}}),
\qquad
\operatorname{wid}\bigl(\mathcal J_q(v_{\mathrm{name}})\bigr)=2^{-q},
\]
where throughout the proof
\[
\operatorname{wid}([u_{\mathrm{lo}},u_{\mathrm{hi}}])
=u_{\mathrm{hi}}-u_{\mathrm{lo}}.
\]
Moreover, multiplying the floor inequalities by two yields
\[
2m_q(v_{\mathrm{name}})
\le m_{q+1}(v_{\mathrm{name}})
<2\bigl(m_q(v_{\mathrm{name}})+1\bigr).
\]
Since the middle term is an integer, these inequalities imply
\[
\mathcal J_{q+1}(v_{\mathrm{name}})
\subseteq\mathcal J_q(v_{\mathrm{name}}).
\]
The conclusions include dyadic inputs, which may equal the lower endpoint of a queried interval. At each fixed precision, both rational endpoints are Borel functions of the real input because the floor function is Borel. Composing these functions with the public exponent coordinates and the observed covariate coordinates gives precisely the validity and Borel requirements in \cref{obj:def:dyadic-name-convention}. Thus \(\mathsf N_0\) is admissible.
% lean: CausalSmith.Stat.LogoddsLowsmoothFrontier.dyadicFloorPolicy_admissible

\item \textbf{Resolution enclosures at the rank budget.}
Fix an admissible engine \(\mathsf E\), an admissible policy \(\mathsf N\), \(n\ge2\), exponents in the stated domain, and a record tuple \(\mathcal O\). Set
\[
q_{\mathrm r}=q_{\mathrm{rank}}(n),
\qquad
\epsilon_{\mathrm r}=2^{-q_{\mathrm r}}.
\]
Let the exponent-name intervals after the intersections prescribed in \cref{obj:def:resolutions} be
\[
\mathcal J_{\alpha,\mathrm r}
=[\alpha_{\mathrm r,-},\alpha_{\mathrm r,+}],
\qquad
\mathcal J_{\beta,\mathrm r}
=[\beta_{\mathrm r,-},\beta_{\mathrm r,+}].
\]
Validity of the names and the exponent domain ensure that both intersections succeed and that
\[
\begin{gathered}
\alpha\in\mathcal J_{\alpha,\mathrm r}\subseteq[0,1],
\qquad
\beta\in\mathcal J_{\beta,\mathrm r}\subseteq[0,1/4],\\
\operatorname{wid}(\mathcal J_{\alpha,\mathrm r})\le\epsilon_{\mathrm r},
\qquad
\operatorname{wid}(\mathcal J_{\beta,\mathrm r})\le\epsilon_{\mathrm r}.
\end{gathered}
\]
Define the rational denominator and power-exponent boxes by exact interval operations:
\[
\begin{aligned}
\mathcal D_{C,\mathrm r}
&=2\mathcal J_{\alpha,\mathrm r}
  +2\mathcal J_{\beta,\mathrm r}+[1,1],
&
\mathcal D_{S,\mathrm r}
&=4\mathcal J_{\beta,\mathrm r}+[1,1],\\
\mathcal P_{C,\mathrm r}
&=[2,2]/\mathcal D_{C,\mathrm r},
&
\mathcal P_{S,\mathrm r}
&=[2,2]/\mathcal D_{S,\mathrm r}.
\end{aligned}
\]
Both denominator boxes have lower endpoints at least one and widths at most \(4\epsilon_{\mathrm r}\). For any denominator box
\(\mathcal D_{\mathrm{aux}}=[d_{\mathrm{lo}},d_{\mathrm{hi}}]\) with
\(1\le d_{\mathrm{lo}}\le d_{\mathrm{hi}}\), exact division gives
\[
[2,2]/\mathcal D_{\mathrm{aux}}
=
[2/d_{\mathrm{hi}},2/d_{\mathrm{lo}}],
\]
and
\[
\operatorname{wid}([2,2]/\mathcal D_{\mathrm{aux}})
=
\frac{2(d_{\mathrm{hi}}-d_{\mathrm{lo}})}
     {d_{\mathrm{lo}}d_{\mathrm{hi}}}
\le2\operatorname{wid}(\mathcal D_{\mathrm{aux}}).
\]
Consequently, both power-exponent boxes lie in \([0,2]\), contain their respective exact exponents, and have width at most \(8\epsilon_{\mathrm r}\).

For a power exponent \(\tau\in[0,2]\),
\[
\left|\frac{d}{d\tau}n^\tau\right|
=n^\tau\log n
\le n^2n=n^3,
\]
using \(n\ge1\) and \(\log n\le n\). The mean-value bound and the outward endpoint errors required by \cref{obj:def:arithmetic-engine} therefore show that the returned target boxes satisfy
\[
R_C\in[a_C,b_C],\qquad R_S\in[a_S,b_S],
\]
\[
b_C-a_C,\ b_S-a_S
\le8n^3\epsilon_{\mathrm r}+2\epsilon_{\mathrm r}.
\]
To pay for this width, the integer logarithmic ceiling in \cref{obj:def:resolutions} gives
\[
n\le2^{\lceil\log_2n\rceil},
\qquad
128n^3\epsilon_{\mathrm r}
=
\frac{n^3}{2^{3\lceil\log_2n\rceil}}
\le1.
\]
Since \(n^3\ge1\), it follows that
\[
8n^3\epsilon_{\mathrm r}+2\epsilon_{\mathrm r}
\le10n^3\epsilon_{\mathrm r}
\le\frac{10}{128}\le1.
\]
All rank-step intersections and divisions have thus succeeded, the power inputs satisfy the primitive domain requirements, and both returned boxes have the required unit-width bound.
% lean: CausalSmith.Stat.LogoddsLowsmoothFrontier.rankBoxes_rank_budget_contract

\item \textbf{Bounds for the selected ranks.}
The exponent domain makes both target exponents positive. Hence
\[
R_C\ge1,\qquad R_S\ge1.
\]
In particular \(b_C,b_S\ge1\), so the nonnegative ceilings in \cref{obj:def:resolutions} are their ordinary integer ceilings. Containment and the unit-width bounds give
\[
R_C\le b_C\le a_C+1\le R_C+1,
\qquad
R_S\le b_S\le a_S+1\le R_S+1.
\]
Using the strict upper bound for an integer ceiling,
\[
\begin{aligned}
R_C&\le k_C=\lceil b_C\rceil<b_C+1\le R_C+2\le3R_C,\\
R_S&\le k_S=\lceil b_S\rceil<b_S+1\le R_S+2\le3R_S.
\end{aligned}
\]
The last inequalities follow from \(R_C,R_S\ge1\). These ranks depend on the public exponent names and are fixed before sampling. In particular,
\[
k_C\ge1,\qquad k_S\ge1.
\]
From this point onward, \(b_C,b_S\) denote the coordinate error radii in \cref{obj:def:upper-handle}.
% lean: CausalSmith.Stat.LogoddsLowsmoothFrontier.enclosure_ceiling_rank_bounds

\item \textbf{Coordinate enclosures at the endpoint budget.}
Retain the selected ranks and set
\[
K=\max\{k_C,k_S\},
\qquad
q_{\mathrm e}=q_{\mathrm{end}}(n,k_C,k_S),
\qquad
\epsilon_{\mathrm e}=2^{-q_{\mathrm e}}.
\]
At this precision, write the raw name responses as
\[
\begin{gathered}
\mathcal B_\alpha
=[\alpha_{\mathrm{name},-},\alpha_{\mathrm{name},+}],
\qquad
\mathcal B_\beta
=[\beta_{\mathrm{name},-},\beta_{\mathrm{name},+}],\\
\mathcal B_{X_i}
=[X_{i,\mathrm{name},-},X_{i,\mathrm{name},+}],
\qquad 1\le i\le n.
\end{gathered}
\]
These boxes contain their named inputs and have widths at most \(\epsilon_{\mathrm e}\). Since \(q_{\mathrm e}\ge2\), one has
\[
\epsilon_{\mathrm e}\le\frac14.
\]
Containment and the width bound place each raw covariate box in
\[
\mathcal B_{X_i}\subseteq[-\epsilon_{\mathrm e},1+\epsilon_{\mathrm e}].
\]
Indeed, each endpoint is within one box width of the contained covariate.

For any displayed scalar expression \(G\), denote its computed enclosure in the prescribed tree at precision \(q_{\mathrm e}\) by
\[
\mathcal B[G]=[G_{\mathrm{box},-},G_{\mathrm{box},+}].
\]
This notation refers to the supplied raw responses, the fixed engine, and the exact rational range operations.

We use two elementary properties of these operations. For closed rational boxes \(\mathcal J_1,\mathcal J_2\), define
\[
M_{\mathcal J_\iota}
=\max\{|u|:u\in\mathcal J_\iota\},
\qquad \iota\in\{1,2\}.
\]
Inserting a cross product when subtracting two products, and then taking the extrema among endpoint products, gives
\[
\operatorname{wid}(\mathcal J_1\mathcal J_2)
\le
M_{\mathcal J_2}\operatorname{wid}(\mathcal J_1)
+
M_{\mathcal J_1}\operatorname{wid}(\mathcal J_2).
\]
Widths add under addition and subtraction. The endpoint formulas for minimum and maximum also give
\[
\begin{aligned}
\operatorname{wid}(\min(\mathcal J_1,\mathcal J_2))
&\le\max\{\operatorname{wid}(\mathcal J_1),
          \operatorname{wid}(\mathcal J_2)\},\\
\operatorname{wid}(\max(\mathcal J_1,\mathcal J_2))
&\le\max\{\operatorname{wid}(\mathcal J_1),
          \operatorname{wid}(\mathcal J_2)\}.
\end{aligned}
\]
For example, each upper endpoint is at most its corresponding lower endpoint plus the larger width; taking minima or maxima preserves that inequality.

We now evaluate the basis and kernel from \cref{obj:synth_3}. Define the primitive constant boxes
\[
\mathcal B_\pi=\mathcal B[\pi],
\qquad
\mathcal B_{\sqrt2}=\mathcal B[\sqrt2].
\]
Their contracts give
\[
\mathcal B_\pi\subseteq[3,4],
\qquad
\operatorname{wid}(\mathcal B_\pi)\le2\epsilon_{\mathrm e},
\]
\[
\mathcal B_{\sqrt2}\subseteq[0,2],
\qquad
\operatorname{wid}(\mathcal B_{\sqrt2})\le2\epsilon_{\mathrm e}.
\]
For the second inclusion, \(\sqrt2\le3/2\) follows by squaring, and the upper endpoint error is at most \(\epsilon_{\mathrm e}\le1/4\).

For a basis index \(j\ge1\), the intermediate argument boxes are
\[
\mathcal B_{\pi j}=\mathcal B_\pi[j,j],
\qquad
\mathcal B_{\pi jX_i}=\mathcal B_{\pi j}\mathcal B_{X_i}.
\]
They satisfy
\[
\operatorname{wid}(\mathcal B_{\pi j})\le2j\epsilon_{\mathrm e},
\qquad
M_{\mathcal B_{\pi j}}\le4j,
\qquad
M_{\mathcal B_{X_i}}\le\frac54.
\]
Thus the product-width inequality yields
\[
\operatorname{wid}(\mathcal B_{\pi jX_i})
\le
\frac54(2j\epsilon_{\mathrm e})+4j\epsilon_{\mathrm e}
=\frac{13}{2}j\epsilon_{\mathrm e}.
\]
The cosine contract returns a box in \([-1,1]\) whose width is at most
\[
\operatorname{wid}\bigl(\mathcal B[\cos(\pi jX_i)]\bigr)
\le\left(\frac{13}{2}j+2\right)\epsilon_{\mathrm e}.
\]
Multiplication by the square-root box therefore gives
\[
M_{\mathcal B[\varphi_j(X_i)]}\le2,
\]
\[
\operatorname{wid}\bigl(\mathcal B[\varphi_j(X_i)]\bigr)
\le
2\epsilon_{\mathrm e}
+2\left(\frac{13}{2}j+2\right)\epsilon_{\mathrm e}
=(13j+6)\epsilon_{\mathrm e}
\le20(j+1)\epsilon_{\mathrm e}.
\]
For \(j=0\), the basis expression is the exact constant one, so these magnitude and width bounds also hold.

Each computed basis-product term consequently has width at most
\(80(j+1)\epsilon_{\mathrm e}\). Summing the terms in the kernel, for either retained rank \(k\in\{k_C,k_S\}\), gives
\[
\begin{aligned}
\operatorname{wid}\bigl(\mathcal B[H_k(X_i,X_{i'})]\bigr)
&\le\sum_{j=0}^{k-1}80(j+1)\epsilon_{\mathrm e}\\
&\le\sum_{j=0}^{k-1}80k\epsilon_{\mathrm e}
=80k^2\epsilon_{\mathrm e}
\le160k^2\epsilon_{\mathrm e}.
\end{aligned}
\]
At \(k=1\), the sole basis-product term is exact.

The binary marks used in \cref{obj:def:coordinate-estimates} are exact rational numbers in \([0,1]\). Multiplication by their products preserves the displayed kernel-width bound. There are exactly
\[
\#\{(i,i'):1\le i,i'\le n,\ i\ne i'\}=n(n-1)
\]
ordered pairs. Hence the positive rational normalization in \cref{obj:def:projection-statistic} cancels this factor in the sum of widths:
\[
\operatorname{wid}\bigl(\mathcal B[\mathbb U_{n,k}(W,V)]\bigr)
\le
\frac{n(n-1)\,160k^2\epsilon_{\mathrm e}}{n(n-1)}
=160k^2\epsilon_{\mathrm e}
\]
for the marks used here. The empirical mean in \(\widehat C_n\) is exact. Soundness of the operations therefore supplies successful coordinate boxes with
\[
\begin{gathered}
\widehat C_n\in\mathcal B[\widehat C_n],
\qquad
\widehat S_n\in\mathcal B[\widehat S_n],\\
\operatorname{wid}\bigl(\mathcal B[\widehat C_n]\bigr)
\le160k_C^2\epsilon_{\mathrm e},
\qquad
\operatorname{wid}\bigl(\mathcal B[\widehat S_n]\bigr)
\le160k_S^2\epsilon_{\mathrm e}.
\end{gathered}
\]
% lean: CausalSmith.Stat.LogoddsLowsmoothFrontier.endpoint_coordinates_budget_contract

\item \textbf{Error-radius enclosures.}
The endpoint budget gives
\[
K\le2^{\lceil\log_2(K+1)\rceil},
\qquad
q_{\mathrm e}\ge4+\lceil\log_2(K+1)\rceil,
\]
and therefore
\[
K\epsilon_{\mathrm e}\le\frac1{16}.
\]
Validity of the raw exponent responses gives
\[
\mathcal B_\alpha\subseteq[-\epsilon_{\mathrm e},1+\epsilon_{\mathrm e}],
\qquad
\mathcal B_\beta\subseteq[-\epsilon_{\mathrm e},1/4+\epsilon_{\mathrm e}].
\]
Define their bias-exponent boxes by
\[
\mathcal B_{C,\mathrm{pow}}=[0,0]-(\mathcal B_\alpha+\mathcal B_\beta),
\qquad
\mathcal B_{S,\mathrm{pow}}=[-2,-2]\mathcal B_\beta.
\]
They contain \(-(\alpha+\beta)\) and \(-2\beta\), respectively, and satisfy
\[
\begin{gathered}
\mathcal B_{C,\mathrm{pow}}
\subseteq[-5/4-2\epsilon_{\mathrm e},2\epsilon_{\mathrm e}]
\subseteq[-3,3],\\
\mathcal B_{S,\mathrm{pow}}
\subseteq[-1/2-2\epsilon_{\mathrm e},2\epsilon_{\mathrm e}]
\subseteq[-3,3],\\
\operatorname{wid}(\mathcal B_{C,\mathrm{pow}}),
\ \operatorname{wid}(\mathcal B_{S,\mathrm{pow}})
\le2\epsilon_{\mathrm e}.
\end{gathered}
\]

For either retained rank \(k\), the elementary inequality
\(\log u\le u-1\) for \(u>0\) gives
\[
\log k=\log(k/2)+\log2
\le(k/2-1)+1=k/2.
\]
Applying the same inequality at \(u=1/2\) gives \(\log2\ge1/2\).
Thus, for every exponent \(\tau\) in the corresponding raw power-exponent box,
\[
\tau\log k
\le2\epsilon_{\mathrm e}\log k
\le k\epsilon_{\mathrm e}
\le\frac1{16}\le\log2.
\]
It follows that
\[
k^\tau\le2,
\qquad
\left|\frac{d}{d\tau}k^\tau\right|
=k^\tau\log k\le k.
\]
These inequalities include \(k=1\), for which the derivative is zero.
The mean-value bound and the primitive endpoint errors now imply
\[
\operatorname{wid}\bigl(\mathcal B[k^\tau]\bigr)
\le k\operatorname{wid}(\mathcal B_{\mathrm{pow}})
    +2\epsilon_{\mathrm e}
\le(2k+2)\epsilon_{\mathrm e},
\]
with \(\mathcal B_{\mathrm{pow}}\) the corresponding bias-exponent box.

The variance quantity
\[
V_n(k)=\frac{10}{n}+\frac{4k}{n(n-1)}
\]
is an exactly evaluated positive rational number. Its square-root enclosure therefore satisfies
\[
\operatorname{wid}\bigl(\mathcal B[\sqrt{20V_n(k)}]\bigr)
\le2\epsilon_{\mathrm e}.
\]
Using the constants in the radius expressions, we obtain successful enclosures containing \(b_C,b_S\), with
\[
\begin{aligned}
\operatorname{wid}(\mathcal B[b_C])
&\le\bigl(8(2k_C+2)+2\bigr)\epsilon_{\mathrm e}
=(16k_C+18)\epsilon_{\mathrm e}
\le34K\epsilon_{\mathrm e},\\
\operatorname{wid}(\mathcal B[b_S])
&\le\bigl(25(2k_S+2)+2\bigr)\epsilon_{\mathrm e}
=(50k_S+52)\epsilon_{\mathrm e}
\le102K\epsilon_{\mathrm e}.
\end{aligned}
\]
The last inequalities use \(k_C,k_S\le K\) and \(K\ge1\).
% lean: CausalSmith.Stat.LogoddsLowsmoothFrontier.endpoint_radii_budget_contract

\item \textbf{Denominator-corner enclosures.}
Define the computed denominator boxes by the prescribed exact clipping tree:
\[
\mathcal B[s_\pm]
=
\max\left(
[s_0,s_0],
\min\left([1,1],
\mathcal B[\widehat S_n]\pm\mathcal B[b_S]\right)
\right),
\qquad s_0=\frac1{512}.
\]
Here the minus sign denotes interval subtraction. Soundness supplies containment of \(s_\pm\), and the endpoint formulas give
\[
\mathcal B[s_\pm]\subseteq[1/512,1].
\]
Before clipping, the widths are bounded by
\[
160k_S^2\epsilon_{\mathrm e}+102K\epsilon_{\mathrm e}
\le262K^2\epsilon_{\mathrm e}
\le300K^2\epsilon_{\mathrm e},
\]
because \(k_S\le K\) and \(K\le K^2\). Clipping by point intervals preserves this bound by the minimum and maximum width inequalities. Consequently,
\[
s_\pm\in\mathcal B[s_\pm],
\qquad
\operatorname{wid}(\mathcal B[s_\pm])
\le300K^2\epsilon_{\mathrm e}.
\]
In particular, every subsequent denominator interval has a positive rational lower endpoint.
% lean: CausalSmith.Stat.LogoddsLowsmoothFrontier.endpoint_denominator_corners

\item \textbf{Numerator-corner widths and magnitudes.}
Define
\[
\mathcal B[c_\pm]
=\mathcal B[\widehat C_n]\pm\mathcal B[b_C].
\]
These boxes contain \(c_\pm\), and
\[
\operatorname{wid}(\mathcal B[c_\pm])
\le160k_C^2\epsilon_{\mathrm e}+34K\epsilon_{\mathrm e}
\le194K^2\epsilon_{\mathrm e}
\le300K^2\epsilon_{\mathrm e}.
\]

We also require bounds on their endpoints. The exact basis satisfies
\[
|\varphi_j(x)|\le\sqrt2
\qquad(j\ge0,\ x\in[0,1]),
\]
including \(j=0\). Therefore
\[
|H_k(x,z)|
\le\sum_{j=0}^{k-1}|\varphi_j(x)\varphi_j(z)|
\le2k.
\]
The ordered-pair normalization and marks of absolute value at most one give
\[
|\mathbb U_{n,k}(W,V)|\le2k.
\]
The empirical treated-outcome mean has absolute value at most one, so
\[
|\widehat C_n|\le1+2k_C.
\]
Since \(n\ge2\),
\[
V_n(k)\le5+2k,
\qquad
20V_n(k)\le100+40k\le(10+7k)^2.
\]
Taking square roots yields
\[
\sqrt{20V_n(k)}\le10+7k.
\]
Moreover, \(k_C\ge1\) and \(\alpha+\beta>0\) give
\[
k_C^{-(\alpha+\beta)}\le1,
\qquad
0\le b_C\le8+10+7k_C\le25k_C.
\]
Thus
\[
|c_\pm|
\le|\widehat C_n|+b_C
\le1+27k_C\le1+27K.
\]

To control the computed endpoint overhang, the two logarithmic ceilings in the endpoint budget give
\[
n+1\le2^{\lceil\log_2(n+1)\rceil},
\qquad
K+1\le2^{\lceil\log_2(K+1)\rceil},
\]
hence
\[
2^{44}(n+1)(K+1)^4\epsilon_{\mathrm e}\le1.
\]
Comparing coefficients with this product shows
\[
\epsilon_{\mathrm e}\le\frac1{16},
\qquad
300K^2\epsilon_{\mathrm e}\le1,
\qquad
2^{37}(K+1)^3\epsilon_{\mathrm e}\le\frac1n.
\]
For the last two comparisons, use
\[
K^2\le(K+1)^4,
\qquad
300K^2\le2^{44}(n+1)(K+1)^4,
\]
\[
2^{37}n(K+1)^3
\le2^{44}(n+1)(K+1)^4.
\]

A box containing \(c_\pm\) and having width at most one has every point within distance one of \(c_\pm\). It follows that
\[
M_{\mathcal B[c_\pm]}
\le1+27K+1\le100(K+1).
\]
This supplies both the width and magnitude bounds at the actual numerator inputs to inversion.
% lean: CausalSmith.Stat.LogoddsLowsmoothFrontier.endpoint_numerator_corners

\item \textbf{Enclosures of the four inversion corners.}
For signs \(\sigma_c,\sigma_s\in\{-,+\}\), define
\[
\mathcal B_{\mathrm{quot},\sigma_c\sigma_s}
=\mathcal B[c_{\sigma_c}]/\mathcal B[s_{\sigma_s}],
\qquad
\mathcal B_{\mathrm{odds},\sigma_c\sigma_s}
=[1,1]+\mathcal B_{\mathrm{quot},\sigma_c\sigma_s}.
\]
For a denominator box
\(\mathcal J_{\mathrm{den}}=[d_{\mathrm{lo}},d_{\mathrm{hi}}]\)
with \(d_{\mathrm{lo}}\ge1/512\), its reciprocal range satisfies
\[
\mathcal J_{\mathrm{den}}^{-1}
=[1/d_{\mathrm{hi}},1/d_{\mathrm{lo}}],
\qquad
M_{\mathcal J_{\mathrm{den}}^{-1}}\le512,
\]
\[
\operatorname{wid}(\mathcal J_{\mathrm{den}}^{-1})
=
\frac{d_{\mathrm{hi}}-d_{\mathrm{lo}}}
     {d_{\mathrm{lo}}d_{\mathrm{hi}}}
\le512^2\operatorname{wid}(\mathcal J_{\mathrm{den}}).
\]
Exact division is multiplication by this reciprocal range. The product-width bound therefore gives containment of the true ratio and
\[
\begin{aligned}
\operatorname{wid}(\mathcal B_{\mathrm{quot},\sigma_c\sigma_s})
\le{}&
512\operatorname{wid}(\mathcal B[c_{\sigma_c}])\\
&+100(K+1)512^2
 \operatorname{wid}(\mathcal B[s_{\sigma_s}]).
\end{aligned}
\]

Define the exponential boundary boxes and the clipped odds boxes by
\[
\mathcal B_{\exp,-}=\mathcal B[\exp(-1/2)],
\qquad
\mathcal B_{\exp,+}=\mathcal B[\exp(1/2)],
\]
\[
\mathcal B_{\mathrm{clip},\sigma_c\sigma_s}
=
\max\left(
\mathcal B_{\exp,-},
\min\left(\mathcal B_{\exp,+},
          \mathcal B_{\mathrm{odds},\sigma_c\sigma_s}\right)
\right).
\]
The exponential series, bounded termwise by the geometric series, gives
\[
\exp(1/4)\le\sum_{j=0}^{\infty}(1/4)^j=\frac43,
\]
and consequently
\[
\exp(1/2)=\exp(1/4)^2\le\frac{16}{9},
\qquad
\exp(-1/2)\ge\frac9{16}.
\]
The primitive endpoint errors and \(\epsilon_{\mathrm e}\le1/16\) imply
\[
\inf\mathcal B_{\exp,-}\ge\frac9{16}-\epsilon_{\mathrm e}\ge\frac12,
\qquad
\sup\mathcal B_{\exp,+}\le\frac{16}{9}+\epsilon_{\mathrm e}\le2,
\]
\[
\operatorname{wid}(\mathcal B_{\exp,-}),
\ \operatorname{wid}(\mathcal B_{\exp,+})
\le2\epsilon_{\mathrm e}.
\]
Thus every clipped odds box has lower endpoint at least \(1/2\), contains the exact clipped odds, and satisfies
\[
\operatorname{wid}(\mathcal B_{\mathrm{clip},\sigma_c\sigma_s})
\le
\operatorname{wid}(\mathcal B_{\mathrm{quot},\sigma_c\sigma_s})
+2\epsilon_{\mathrm e}.
\]
Here the added exact constant one preserves width, and the minimum and maximum width inequalities bound both clipping operations.

Define the inversion-corner boxes by
\[
\mathcal B_{F,\sigma_c\sigma_s}
=\mathcal B[F(c_{\sigma_c},s_{\sigma_s})].
\]
The logarithm derivative is at most two on every argument interval with lower endpoint at least \(1/2\). Its primitive contract therefore supplies successful boxes containing the four exact corner values, with
\[
\begin{aligned}
\operatorname{wid}(\mathcal B_{F,\sigma_c\sigma_s})
\le{}&
2\left(
512\operatorname{wid}(\mathcal B[c_{\sigma_c}])
+100(K+1)512^2
 \operatorname{wid}(\mathcal B[s_{\sigma_s}])
\right)
+6\epsilon_{\mathrm e}\\
\le{}&
\bigl(307200K^2
+15728640000(K+1)K^2+6\bigr)\epsilon_{\mathrm e}\\
\le{}&2^{37}(K+1)^3\epsilon_{\mathrm e}
\le\frac1n.
\end{aligned}
\]
For the coefficient comparison, \(K\ge0\) implies
\[
K^2\le(K+1)^3,\qquad
(K+1)K^2\le(K+1)^3,\qquad
1\le(K+1)^3,
\]
and
\[
307200+15728640000+6\le2^{37}.
\]
This establishes the required precision for each of the four corner evaluations.
% lean: CausalSmith.Stat.LogoddsLowsmoothFrontier.endpoint_inversion_budget_contract

\item \textbf{The two endpoint boxes.}
The literal minimum and maximum operations return
\[
\mathcal B_L
=\min(\mathcal B_{F,--},\mathcal B_{F,-+})
=[L_-,L_+],
\]
\[
\mathcal B_R
=\max(\mathcal B_{F,+-},\mathcal B_{F,++})
=[R_-,R_+].
\]
Monotonicity of minimum and maximum in their arguments preserves containment, giving
\[
L\in[L_-,L_+],
\qquad
R\in[R_-,R_+].
\]
Their widths satisfy
\[
L_+-L_-
\le\max\{
\operatorname{wid}(\mathcal B_{F,--}),
\operatorname{wid}(\mathcal B_{F,-+})\}
\le\frac1n,
\]
\[
R_+-R_-
\le\max\{
\operatorname{wid}(\mathcal B_{F,+-}),
\operatorname{wid}(\mathcal B_{F,++})\}
\le\frac1n.
\]
All raw input queries and all intervening evaluations have succeeded. Every power input lies in its required domain, every square-root input is nonnegative, every division has a positive denominator interval, and every logarithm input has a positive lower endpoint. The finite sums at the retained finite ranks and the primitive termination contracts in \cref{obj:def:arithmetic-engine} therefore give successful return of both boxes at the single endpoint precision.
% lean: CausalSmith.Stat.LogoddsLowsmoothFrontier.endpoint_min_max_width

\item \textbf{Outward enclosure and rational interval shape.}
The exact clipped logarithmic map in \cref{obj:def:upper-handle} satisfies
\[
-\frac12\le F(c,d)\le\frac12
\]
at the endpoint inputs: its logarithm argument lies between
\(\exp(-1/2)\) and \(\exp(1/2)\). For each positive denominator, division, clipping, and logarithm are nondecreasing in the numerator. Since \(b_C\ge0\) and \(s_-\ge s_0>0\),
\[
L
\le F(c_-,s_-)
\le F(c_+,s_-)
\le R.
\]
Consequently,
\[
-\frac12\le L\le R\le\frac12.
\]
Define the literal rational output endpoints by
\[
l_{\mathrm{out}}=\max\{-1/2,L_-\},
\qquad
u_{\mathrm{out}}=\min\{1/2,R_+\}.
\]
Containment of the exact endpoints and their effect bounds give
\[
l_{\mathrm{out}}\le L\le R\le u_{\mathrm{out}}.
\]
Thus the intersection in the output construction succeeds and the literal output is
\[
I^{\mathrm{up}}_{n,\mathsf E,\mathsf N}(\alpha,\beta)(\mathcal O)
=[l_{\mathrm{out}},u_{\mathrm{out}}],
\qquad
[L,R]\subseteq[l_{\mathrm{out}},u_{\mathrm{out}}].
\]
The endpoint-width bounds give
\[
0\le L-l_{\mathrm{out}}
\le L-L_-\le L_+-L_-\le\frac1n,
\]
\[
0\le u_{\mathrm{out}}-R
\le R_+-R\le R_+-R_-\le\frac1n.
\]
Adding these inequalities yields
\[
u_{\mathrm{out}}-l_{\mathrm{out}}
\le R-L+\frac2n.
\]
Finally,
\[
l_{\mathrm{out}},u_{\mathrm{out}}\in\mathbb Q,
\qquad
-\frac12\le l_{\mathrm{out}}\le u_{\mathrm{out}}\le\frac12.
\]
The output is therefore a nonempty closed connected interval with the asserted rational endpoint shape.
% lean: CausalSmith.Stat.LogoddsLowsmoothFrontier.upperOutput_enclosure_guarantees

\item \textbf{Borel measurability.}
Keep the engine, policy, sample size, and public exponents fixed. The ranks and endpoint precision are then fixed integers. Define the finite response code and its countable range by
\[
\mathfrak z(\mathcal O)
=
\left(
(A_i,Y_i),
(X_{i,\mathrm{name},-},X_{i,\mathrm{name},+})
\right)_{i=1}^n,
\qquad
\mathcal Z
=
\bigl(\{0,1\}^2\times\mathbb Q^2\bigr)^n.
\]
The binary coordinates and the fixed-precision name responses are Borel, so \(\mathfrak z\) is Borel into the discrete countable space \(\mathcal Z\). The public exponent responses are fixed, and the remaining computation uses this code through finite rational operations. Hence
\[
\mathfrak z(\mathcal O)=\mathfrak z(\mathcal O')
\quad\Longrightarrow\quad
I^{\mathrm{up}}_{n,\mathsf E,\mathsf N}(\alpha,\beta)(\mathcal O)
=
I^{\mathrm{up}}_{n,\mathsf E,\mathsf N}(\alpha,\beta)(\mathcal O').
\]
For any set \(\mathcal T\) of output intervals, define
\[
\mathcal Z_{\mathcal T}
=
\left\{
z\in\mathcal Z:
\exists\mathcal O,\ 
\mathfrak z(\mathcal O)=z,\
I^{\mathrm{up}}_{n,\mathsf E,\mathsf N}(\alpha,\beta)(\mathcal O)
\in\mathcal T
\right\}.
\]
The preceding implication gives the fiber decomposition
\[
\left\{
\mathcal O:
I^{\mathrm{up}}_{n,\mathsf E,\mathsf N}(\alpha,\beta)(\mathcal O)
\in\mathcal T
\right\}
=
\bigcup_{z\in\mathcal Z_{\mathcal T}}
\mathfrak z^{-1}(\{z\}).
\]
This is a countable union of Borel sets. In particular, inverse images of Borel sets of intervals are Borel, establishing the claimed measurability of the output map.
% lean: CausalSmith.Stat.LogoddsLowsmoothFrontier.upperOutput_regular
\end{enumerate}
\end{proof}

The rank precision \leanref{sym:q_{\mathrm{rank}}}{\ensuremath{q_{\mathrm{rank}}(n)}} gives an additive rank excess below two for each coordinate. The endpoint precision \leanref{sym:q_{\mathrm{end}}}{\ensuremath{q_{\mathrm{end}}(n,k_C,k_S)}} depends on the sample size and the larger retained rank. The same selected integers therefore govern both the exact statistics and their rational evaluation. Together with the finite recurrence contracts, successful evaluation at these prescribed precisions establishes termination of the construction in \cref{obj:def:upper-handle}.

In the inversion, \(R\) denotes the ideal upper effect endpoint. The rational output uses the outward endpoints \(L_-,R_+\), followed by clipping to the effect envelope. Its samplewise containment of the ideal interval and additional length bound of \(2/n\) hold for every record tuple. The Borel conclusion concerns the resulting map from the original observed records, under each fixed admissible naming policy.

\subsection{Connection to honesty and expected length}

The rational interval \leanref{sym:I^{\mathrm{up}}}{\ensuremath{I^{\mathrm{up}}_{n,\mathsf E,\mathsf N}(\alpha,\beta)}} in \cref{obj:def:upper-handle} combines these realization guarantees with the statistical bounds. The bias and variance bounds in \cref{obj:lem:projection-control} calibrate its coordinate error radii; the denominator floor and odds identity in \cref{obj:lem:observable-odds} supply the joint inversion. \Cref{obj:lem:fixed-budget-realization} establishes measurable outward enclosure of the resulting ideal interval with a controlled length increment.

Under the engine, naming-policy, and exponent conditions of \cref{obj:thm:finite-honest-upper}, the literal procedure has global \(90\%\) coverage over the ambient model for every integer sample size \(n\ge1\). On every effect-radius slice, that result bounds its unconditional expected length by an absolute constant times the diagnostic envelope in \cref{obj:def:diagnostic-envelope}. Numerator uncertainty contributes directly to this envelope, while denominator uncertainty contributes at the scale of the effect radius. At \(n=1\), the specified full effect interval completes the procedure.

Different admissible engines and valid naming policies may select different ranks and return different rational endpoints. For each choice, the realization bounds apply to its own retained ranks and ideal endpoints. The coverage and expected-length guarantees in \cref{obj:thm:finite-honest-upper} share one absolute constant across all such choices. In particular, \cref{obj:lem:concrete-arithmetic-engine,obj:lem:fixed-budget-realization} supplies the explicit engine and dyadic-floor policy for the concrete attaining sequence defined in \cref{obj:def:frontier-handle}.


\section{Exact calibration and same-model alternatives}

The calibrated alternatives are specified by the following maps and continuous sign families.

The lower-bound constructions use exact calibration to combine a homogeneous conditional log-odds coefficient with continuous nuisance functions. We first describe the local equations and their smooth neighborhoods, then state the support guarantees for the finite sign families in \cref{obj:def:continuous-calibrated-priors}. The calibration maps are defined in \cref{obj:def:calibration-maps}.

The covariance adjustment \leanref{sym:\mathfrak c}{\ensuremath{\mathfrak c(t,\xi,\upsilon)}} prescribes the odds ratio of a binary table with treatment margin \leanref{sym:\xi}{\ensuremath{\xi}} and outcome margin \leanref{sym:\upsilon}{\ensuremath{\upsilon}}. Its normalized extension \leanref{sym:\mathfrak B}{\ensuremath{\mathfrak B(t,\xi,\upsilon)}} extends the effect-normalized adjustment smoothly through zero effect. The mixed equation uses the signed propensity amplitude \leanref{sym:\eta}{\ensuremath{\eta}} and signed outcome amplitude \leanref{sym:\zeta}{\ensuremath{\zeta}}; the fair equation uses the signed control-risk amplitude \leanref{sym:\delta}{\ensuremath{\delta}}. In the latter equation, the residual \leanref{sym:\mathfrak H}{\ensuremath{\mathfrak H(t,\delta,\xi,u)}} smoothly extends the matching discrepancy normalized by the effect and squared amplitude. The within-cell coordinate \leanref{sym:u}{\ensuremath{u}} ranges over \([0,1]\).

The scalar implicit-function theorem supplies the local graph representation needed for these calibration equations. Its conditions concern an open parameter domain, differentiability, and a nonzero derivative with respect to the scalar being calibrated.

\begin{lemmav}{lem:scalar-implicit-function}[Scalar implicit function theorem]
Let \(d\) be a nonnegative integer, let \(F:\mathbb R^d\times\mathbb R\to\mathbb R\), and let \(D\subseteq\mathbb R^d\times\mathbb R\). Fix \(\mathbf b_0\in\mathbb R^d\) and \(\chi_0\in\mathbb R\). Suppose that:
\begin{itemize}
\item \textbf{(Open domain.)} \(D\) is open and \((\mathbf b_0,\chi_0)\in D\).
\item \textbf{(Smoothness.)} \(F\) is continuously differentiable on \(D\).
\item \textbf{(Zero.)} \(F(\mathbf b_0,\chi_0)=0\).
\item \textbf{(Nonzero partial derivative.)} \(\partial_\chi F(\mathbf b_0,\chi_0)\ne0\).
\end{itemize}
Then there exist open sets \(U\subseteq\mathbb R^d\) and \(V\subseteq\mathbb R\), and a function \(q:\mathbb R^d\to\mathbb R\), such that
\[
\mathbf b_0\in U,\qquad \chi_0\in V,\qquad U\times V\subseteq D.
\]
The function \(q\) is continuously differentiable on \(U\), satisfies \(q(\mathbf b_0)=\chi_0\), and has \(q(\mathbf b)\in V\) for every \(\mathbf b\in U\). Moreover, for every \(\mathbf b\in U\) and every \(\chi\in V\),
\[
F(\mathbf b,\chi)=0
\quad\Longleftrightarrow\quad
\chi=q(\mathbf b).
\]
\end{lemmav}

\begin{proof}[Proof of \cref{obj:lem:scalar-implicit-function}]
\begin{enumerate}
\item Since \(D\) is open and contains \((\mathbf b_0,\chi_0)\), the smoothness hypothesis gives continuous differentiability at this point. Define the continuous linear map
\[
L_\chi:\mathbb R\to\mathbb R,
\qquad
L_\chi(t)=DF(\mathbf b_0,\chi_0)[(\mathbf 0,t)].
\]
The chain rule for the map \(\chi\mapsto(\mathbf b_0,\chi)\) gives
\[
\left.\frac{d}{d\chi}F(\mathbf b_0,\chi)\right|_{\chi=\chi_0}
=L_\chi(1)\ne0.
\]
By linearity,
\[
L_\chi(t)=tL_\chi(1)
\qquad(t\in\mathbb R).
\]
Define
\[
J_\chi:\mathbb R\to\mathbb R,
\qquad
J_\chi(t)=\frac{t}{L_\chi(1)}.
\]
This map is continuous and linear, and the preceding identity gives
\[
L_\chi(J_\chi(t))=t,
\qquad
J_\chi(L_\chi(t))=t
\qquad(t\in\mathbb R).
\]
Thus \(L_\chi\) is invertible. % lean: CausalSmith.Stat.LogoddsLowsmoothFrontier.scalar_implicit_function

\item The general \(C^1\) implicit-function theorem for Banach-space products, applied with scalar-coordinate derivative \(L_\chi\), supplies a function
\[
q:\mathbb R^d\to\mathbb R
\]
such that
\[
q(\mathbf b_0)=\chi_0,
\qquad
q\text{ is continuously differentiable at }\mathbf b_0.
\]
Its local graph identity, together with \(F(\mathbf b_0,\chi_0)=0\), gives
\[
F(\mathbf b,\chi)=0
\quad\Longleftrightarrow\quad
q(\mathbf b)=\chi
\]
throughout a neighborhood of \((\mathbf b_0,\chi_0)\). % lean: CausalSmith.Stat.LogoddsLowsmoothFrontier.scalar_implicit_function

\item Intersect this neighborhood with \(D\). The product topology supplies neighborhoods \(U_1\subseteq\mathbb R^d\) of \(\mathbf b_0\) and \(V_1\subseteq\mathbb R\) of \(\chi_0\) satisfying
\[
U_1\times V_1
\subseteq
D\cap
\bigl\{(\mathbf b,\chi):
F(\mathbf b,\chi)=0
\Longleftrightarrow q(\mathbf b)=\chi\bigr\}.
\]
Choose open neighborhoods \(U_2\subseteq\mathbb R^d\) and \(V\subseteq\mathbb R\) with
\[
\mathbf b_0\in U_2\subseteq U_1,
\qquad
\chi_0\in V\subseteq V_1.
\]
Continuous differentiability of \(q\) at \(\mathbf b_0\) also supplies a neighborhood \(S\subseteq\mathbb R^d\) such that
\[
\mathbf b_0\in\operatorname{int}(S),
\qquad
q\text{ is continuously differentiable on }S.
\]
% lean: CausalSmith.Stat.LogoddsLowsmoothFrontier.scalar_implicit_function

\item Continuity of \(q\) at \(\mathbf b_0\), the identity \(q(\mathbf b_0)=\chi_0\), and openness of \(V\) imply
\[
\mathbf b_0\in\operatorname{int}\bigl(q^{-1}(V)\bigr),
\qquad
q^{-1}(V)=\{\mathbf b\in\mathbb R^d:q(\mathbf b)\in V\}.
\]
Hence we may choose an open set \(U\subseteq\mathbb R^d\) satisfying
\[
\mathbf b_0\in U
\subseteq U_2\cap S\cap q^{-1}(V).
\]
The inclusions established above yield
\[
U\times V\subseteq U_1\times V_1\subseteq D,
\]
while restriction to \(U\subseteq S\) gives continuous differentiability of \(q\) on \(U\), and
\[
q(\mathbf b)\in V
\qquad(\mathbf b\in U).
\]
Finally, for every \(\mathbf b\in U\) and \(\chi\in V\), the graph identity on \(U_1\times V_1\) gives
\[
F(\mathbf b,\chi)=0
\quad\Longleftrightarrow\quad
q(\mathbf b)=\chi
\quad\Longleftrightarrow\quad
\chi=q(\mathbf b).
\]
These are the required properties. % lean: CausalSmith.Stat.LogoddsLowsmoothFrontier.scalar_implicit_function
\end{enumerate}
\end{proof}

The open-domain condition makes the calibrated scalar a differentiable function of the remaining parameters, including at parameter values that will lie on the boundary of the eventual support region. The local uniqueness conclusion identifies the zero set with that function throughout the product neighborhood. For the two calibration equations, \cref{obj:lem:exact-calibrations} supplies fixed brackets and common open neighborhoods for the resulting branches.

Specifically, the mixed neighborhood \(U\) contains a closed region with signed amplitude radius \(\varepsilon_{\mathrm{mix}}\), and the fair neighborhood \(V\) contains a closed region with signed amplitude radius \(\varepsilon_{\mathrm{fair}}\) and effect coordinate in \([0,1/4]\). Both amplitude radii exceed the calibration radius \(\varepsilon\) used for support. A further neighborhood \(W\) allows the fair margin to vary freely within a centered bracket of half-width \(b\); here \(b\) denotes a root-bracket width. The next statement establishes these neighborhoods together with the exact matching identities and derivative bounds.

\begin{lemmav}{lem:exact-calibrations}[Exact smooth calibration branches]
There exist constants \(\varepsilon>0\) and \(M>0\) with the following properties. Expectations below are over two independent fair signs, and
\[
Z_u=s_1\cos(\pi u/2)+s_2\sin(\pi u/2),
\qquad
E F(s_1,s_2)=\frac14\sum_{s_1,s_2\in\{-1,1\}}F(s_1,s_2).
\]
The covariance adjustment \(\mathfrak c\) and its normalized extension \(\mathfrak B\) are given by
\[
\begin{gathered}
d=e^t-1,\qquad
v_*=\xi(1-\xi)\upsilon(1-\upsilon),\qquad
L_*=1+d\{\xi(1-\upsilon)+(1-\xi)\upsilon\},\\
D(t)=\int_0^1 e^{rt}\,dr,\\
\mathfrak c(t,\xi,\upsilon)
=\frac{2d v_*}{L_*+\sqrt{L_*^2-4d^2v_*}},
\qquad
\mathfrak B(t,\xi,\upsilon)
=\frac{2D(t)v_*}{L_*+\sqrt{L_*^2-4d^2v_*}}.
\end{gathered}
\]
The mixed root \(\mathfrak q(\eta,\zeta,u)\) is selected from the roots in \((3/4,5/4)\) of
\[
16E\mathfrak B
 \bigl(32\eta\zeta,\tfrac12-\eta qZ_u,\tfrac12+\zeta Z_u\bigr)-q=0,
\]
with value \(1\) when that bracket contains no root. The fair root
\(\mathfrak p(t,\delta,u)\) is selected from the roots in \((3/10,1/2)\) of
\(\mathfrak H(t,\delta,p,u)=0\), with value \(2/5\) when that bracket contains no root. Here \(\mathfrak H\) is the smoothly extended normalized fair-matching residual, whose defining identity away from the axes is displayed below.

\begin{itemize}
\item \textbf{(Positive odds tables.)}
For every \(t,\xi,\upsilon\in\mathbb R\) satisfying
\[
|t|\le\varepsilon,\qquad
|\xi-\tfrac12|\le\varepsilon,\qquad
|\upsilon-\tfrac12|\le\varepsilon,
\]
put \(c=\mathfrak c(t,\xi,\upsilon)\). The four cells
\[
\begin{aligned}
p_{11}&=\xi\upsilon+c,&
p_{10}&=\xi(1-\upsilon)-c,\\
p_{01}&=(1-\xi)\upsilon-c,&
p_{00}&=(1-\xi)(1-\upsilon)+c
\end{aligned}
\]
are strictly positive and satisfy
\[
p_{11}+p_{10}=\xi,\qquad
p_{11}+p_{01}=\upsilon,\qquad
p_{11}+p_{10}+p_{01}+p_{00}=1,\qquad
\frac{p_{11}p_{00}}{p_{10}p_{01}}=e^t.
\]
Moreover,
\[
\mathfrak B(t,\xi,\upsilon)=\frac{\mathfrak c(t,\xi,\upsilon)}{t}
\quad(t\ne0),\qquad
\mathfrak B(0,\xi,\upsilon)=\xi(1-\xi)\upsilon(1-\upsilon).
\]

\item \textbf{(Signed mixed calibration.)}
For every \(|\eta|\le\varepsilon\), \(|\zeta|\le\varepsilon\), and \(u\in[0,1]\), put
\(q=\mathfrak q(\eta,\zeta,u)\). Then
\[
\frac34<q<\frac54,\qquad
16E\mathfrak B
 \bigl(32\eta\zeta,\tfrac12-\eta qZ_u,\tfrac12+\zeta Z_u\bigr)-q=0,
\]
and
\[
\begin{gathered}
\mathfrak q(\eta,\zeta,0)=\mathfrak q(\eta,\zeta,1),\\
E\mathfrak c
 \bigl(32\eta\zeta,\tfrac12-\eta qZ_u,\tfrac12+\zeta Z_u\bigr)
 =2\eta\zeta q,\\
\mathfrak q(-\eta,\zeta,u)=q,\qquad
\mathfrak q(\eta,-\zeta,u)=q,\\
\mathfrak q(0,\zeta,u)=1-4\zeta^2,\qquad
\mathfrak q(\eta,0,u)+4\eta^2\mathfrak q(\eta,0,u)^2=1.
\end{gathered}
\]

\item \textbf{(Fair calibration and separation.)}
For every \(t\in[0,1/4]\), \(|\delta|\le\varepsilon\), and \(u\in[0,1]\), put
\(p=\mathfrak p(t,\delta,u)\). Then
\[
\begin{gathered}
\mathfrak p(t,\delta,0)=\mathfrak p(t,\delta,1)=\frac25,\qquad
\mathfrak p(t,0,u)=\mathfrak p(0,\delta,u)=\frac25,\\
\mathfrak p(t,-\delta,u)=p,\qquad
\mathfrak H(t,\delta,p,u)=0,\\
E\mathsf G(t,p+\delta Z_u)=\mathsf G(\mathsf T(t,\delta),p).
\end{gathered}
\]
For \(t\delta\ne0\), the residual at the selected root agrees with the unnormalized fair-matching numerator divided by \(t\delta^2\). Furthermore,
\[
|p-\tfrac25|+|\partial_u\mathfrak p(t,\delta,u)|\le M\delta^2,
\]
and
\[
3t\delta^2\le t-\mathsf T(t,\delta)\le6t\delta^2,\qquad
\frac t2\le\mathsf T(t,\delta)\le t.
\]

\item \textbf{(Smooth neighborhoods and derivative bounds.)}
Both roots and every fully substituted local four-cell formula of the mixed and fair families are \(C^\infty\) on their respective closed parameter regions
\[
\begin{gathered}
|\eta|\le\varepsilon,\quad |\zeta|\le\varepsilon,\quad u\in[0,1],\\
t\in[0,1/4],\quad |\delta|\le\varepsilon,\quad u\in[0,1].
\end{gathered}
\]
For every integer \(m\ge0\), there is a constant \(B>0\) such that the norm of the \(m\)th total derivative of each root and each fully substituted local four-cell formula is at most \(B\) throughout its respective closed parameter region, uniformly over all family indices, auxiliary signs, and binary cell indices.

There also exist signed amplitude radii
\(\varepsilon_{\mathrm{mix}}>\varepsilon\) and
\(\varepsilon_{\mathrm{fair}}>\varepsilon\), a fixed fair-root bracket half-width
\(0<b<1/10\), and open sets \(U,V\subset\mathbb R^3\) containing, respectively,
\[
\begin{gathered}
\{(\eta,\zeta,u):|\eta|\le\varepsilon_{\mathrm{mix}},
                       \ |\zeta|\le\varepsilon_{\mathrm{mix}},\ u\in[0,1]\},\\
\{(t,\delta,u):t\in[0,1/4],\
                  |\delta|\le\varepsilon_{\mathrm{fair}},\ u\in[0,1]\}.
\end{gathered}
\]
Throughout \(U\), the selected mixed branch is the unique zero in its fixed bracket \((3/4,5/4)\); throughout \(V\), the selected fair branch is the unique zero in the fixed centered bracket \((2/5-b,2/5+b)\). Both roots and every fully substituted local four-cell formula are \(C^\infty\) on the corresponding open set.

Finally, there is an open set \(W\subset\mathbb R^4\) containing
\[
\{(t,\delta,\xi,u):t\in[0,1/4],\
 |\delta|\le\varepsilon_{\mathrm{fair}},\
 \xi\in(2/5-b,2/5+b),\ u\in[0,1]\}
\]
on which \(\mathfrak H\) is \(C^\infty\). Throughout this displayed core, for every margin \(\xi\) in the centered bracket and whenever \(t\delta\ne0\),
\[
\mathfrak H(t,\delta,\xi,u)
=\frac{E\mathsf G(t,\xi+\delta Z_u)-\mathsf G(\mathsf T(t,\delta),\xi)}
       {t\delta^2}.
\]
\end{itemize}
\end{lemmav}

\begin{proof}[Proof of \cref{obj:lem:exact-calibrations}]
We use the maps in \cref{obj:def:calibration-maps}, with the root-selection and fallback conventions in the statement. For \(r_{\mathrm{amp}}\ge0\), write
\[
\begin{aligned}
\mathcal P_{\mathrm{mix}}(r_{\mathrm{amp}})
&=[-r_{\mathrm{amp}},r_{\mathrm{amp}}]^2\times[0,1],\\
\mathcal P_{\mathrm{fair}}(r_{\mathrm{amp}})
&=[0,1/4]\times[-r_{\mathrm{amp}},r_{\mathrm{amp}}]\times[0,1].
\end{aligned}
\]
These parameter regions are compact. We write \(\mathrm D^m f\) for the \(m\)th total derivative of a parameter function \(f\), with \(\mathrm D^0f=f\).

\begin{enumerate}
\item \emph{Uniform existence and uniqueness of the mixed root.}
Define the mixed residual, wherever its branch denominators are positive, by
\[
\mathcal F_{\mathrm{mix}}(\eta,\zeta,u,\chi)
=
16E\mathfrak B
\bigl(32\eta\zeta,\tfrac12-\eta\chi Z_u,
                    \tfrac12+\zeta Z_u\bigr)-\chi .
\]
The integral defining \(D\) is smooth, and the fundamental theorem of calculus gives
\[
tD(t)=e^t-1,\qquad D(0)=1.
\]
At zero amplitudes the branch radicand is \(1\), its denominator is \(2\), and
\[
\mathfrak B(0,\tfrac12,\tfrac12)=\frac1{16}.
\]
Consequently, for every \(u,\chi\in\mathbb R\),
\[
\mathcal F_{\mathrm{mix}}(0,0,u,\chi)=1-\chi,
\qquad
\partial_\chi\mathcal F_{\mathrm{mix}}(0,0,u,\chi)=-1.
\]
The radicands and denominators remain positive near these points, so the residual is smooth there.

Compactness of \([0,1]\times[7/8,9/8]\), continuity, and the finite number of sign pairs give a positive signed amplitude radius on which all branch denominators are positive and
\[
\left|
\mathcal F_{\mathrm{mix}}(\eta,\zeta,u,\chi)-(1-\chi)
\right|<\frac1{16}
\quad
\left(u\in[0,1],\ \chi\in[7/8,9/8]\right).
\]
In particular,
\[
\mathcal F_{\mathrm{mix}}(\eta,\zeta,u,7/8)>\frac1{16},
\qquad
\mathcal F_{\mathrm{mix}}(\eta,\zeta,u,9/8)<-\frac1{16}.
\]
The intermediate value theorem supplies a zero in \([7/8,9/8]\), hence in \((3/4,5/4)\).

Separately, continuity of the root partial derivative and compactness of
\([0,1]\times[3/4,5/4]\) give another positive signed amplitude radius on which
\[
\partial_\chi\mathcal F_{\mathrm{mix}}(\eta,\zeta,u,\chi)<-\frac12
\quad
\left(u\in[0,1],\ \chi\in[3/4,5/4]\right).
\]
Taking the minimum of these two radii gives
\(\varepsilon_{\mathrm{root}}>0\) such that
\[
\forall(\eta,\zeta,u)\in
\mathcal P_{\mathrm{mix}}(\varepsilon_{\mathrm{root}}),
\qquad
\exists!\chi\in(3/4,5/4):
\mathcal F_{\mathrm{mix}}(\eta,\zeta,u,\chi)=0.
\]
Indeed, the derivative inequality makes the residual strictly decreasing throughout the selection bracket.
% lean: CausalSmith.Stat.LogoddsLowsmoothFrontier.mixedEquation_uniform_existsUnique

\item \emph{Smoothness and derivative bounds for the mixed root and cells.}
At each \((0,0,u)\), the root is \(1\) and its root partial derivative is \(-1\). Apply \cref{obj:lem:scalar-implicit-function} on an open neighborhood where the residual is smooth. The resulting local graph stays in \((3/4,5/4)\) after shrinking its parameter neighborhood. The negative root derivative persists throughout that bracket nearby, so the selected root agrees with this graph.

The same graph is smooth to every finite order. To see the regularity upgrade explicitly, implicit differentiation gives
\[
\mathrm D\mathfrak q(\eta,\zeta,u)
=
-\left.
\frac{
\bigl(\partial_\eta\mathcal F_{\mathrm{mix}},
      \partial_\zeta\mathcal F_{\mathrm{mix}},
      \partial_u\mathcal F_{\mathrm{mix}}\bigr)}
{\partial_\chi\mathcal F_{\mathrm{mix}}}
\right|_{\chi=\mathfrak q(\eta,\zeta,u)} .
\]
On a smaller neighborhood the denominator remains nonzero. If the graph is \(C^k\), this formula makes its first derivative \(C^k\), and therefore makes the graph \(C^{k+1}\). Starting from the \(C^1\) graph supplied by the cited result proves smoothness to every finite order. Compactness of the spatial interval then gives
\(\varepsilon_{\mathrm q}>0\) such that \(\mathfrak q\) is smooth on a neighborhood of every point of
\(\mathcal P_{\mathrm{mix}}(\varepsilon_{\mathrm q})\).

For clarity, define the fully substituted mixed cells. For a family index
\(\iota\in\{0,1\}\), a sign pair
\(\boldsymbol{s}=(s_1,s_2)\in\{-1,1\}^2\), and
\(a,y\in\{0,1\}\), put
\[
\begin{aligned}
e_{\iota}^{\mathrm{loc}}(\eta,\zeta,u,\boldsymbol{s})
&=\frac12+(1-2\iota)\eta\mathfrak q(\eta,\zeta,u)Z_u,\\
m^{\mathrm{loc}}(\zeta,u,\boldsymbol{s})
&=\frac12+\zeta Z_u,\\
c_{\iota}^{\mathrm{loc}}(\eta,\zeta,u,\boldsymbol{s})
&=
\begin{cases}
0,&\iota=0,\\
\mathfrak c(32\eta\zeta,e_1^{\mathrm{loc}},m^{\mathrm{loc}}),
&\iota=1,
\end{cases}\\
\bigl(
\mathcal K^{\mathrm{mix}}_{\iota,\boldsymbol{s},1,1},
\mathcal K^{\mathrm{mix}}_{\iota,\boldsymbol{s},1,0},
\mathcal K^{\mathrm{mix}}_{\iota,\boldsymbol{s},0,1},
\mathcal K^{\mathrm{mix}}_{\iota,\boldsymbol{s},0,0}
\bigr)
&=
\bigl(
e_{\iota}^{\mathrm{loc}}m^{\mathrm{loc}}+c_{\iota}^{\mathrm{loc}},
e_{\iota}^{\mathrm{loc}}(1-m^{\mathrm{loc}})-c_{\iota}^{\mathrm{loc}},
(1-e_{\iota}^{\mathrm{loc}})m^{\mathrm{loc}}-c_{\iota}^{\mathrm{loc}},
(1-e_{\iota}^{\mathrm{loc}})(1-m^{\mathrm{loc}})
+c_{\iota}^{\mathrm{loc}}
\bigr).
\end{aligned}
\]
All arguments on the last line are the displayed parameters and sign pair. At zero amplitudes, the covariance formula is evaluated at
\((0,1/2,1/2)\), where its radicand and denominator are positive. The chain rule therefore makes every displayed cell smooth near each \((0,0,u)\). Compactness in \(u\) and finiteness of the indices give
\(\varepsilon_{\mathrm{cell}}>0\) on whose closed mixed region all these cells are smooth on pointwise neighborhoods.

For each fixed \(m\ge0\), continuity of the total derivatives on the corresponding compact regions gives positive constants satisfying
\[
\|\mathrm D^m\mathfrak q\|\le B_{\mathrm q,m}
\quad\hbox{on }\mathcal P_{\mathrm{mix}}(\varepsilon_{\mathrm q}),
\qquad
\|\mathrm D^m\mathcal K^{\mathrm{mix}}_j\|
\le C^{\mathrm{mix}}_{j,m}
\quad\hbox{on }\mathcal P_{\mathrm{mix}}(\varepsilon_{\mathrm{cell}}).
\]
Here the finite index set is
\[
\mathcal I_{\mathrm{aux}}
=\{0,1\}\times\{-1,1\}^2\times\{0,1\}^2,
\qquad
j=(\iota,\boldsymbol{s},a,y).
\]
Thus the positive constant
\[
B_{\mathrm{cell},m}
=\sum_{j\in\mathcal I_{\mathrm{aux}}}C^{\mathrm{mix}}_{j,m}
\]
bounds the \(m\)th derivative of every mixed cell uniformly over its indices.
% lean: CausalSmith.Stat.LogoddsLowsmoothFrontier.mixedRoot_uniform_smooth_derivative_bounds
% lean: CausalSmith.Stat.LogoddsLowsmoothFrontier.localMixedCell_uniform_smooth_derivative_bounds

\item \emph{The smooth fair residual.}
In the fair calculations use the baseline risk and its normalization
\[
p_*=\frac25,\qquad B_*=1+dp_*,
\qquad d=e^t-1,
\]
and define the unnormalized residual by
\[
\mathcal R_{\mathrm{fair}}(t,\delta,\xi,u)
=
E\mathsf G(t,\xi+\delta Z_u)
-\mathsf G(\mathsf T(t,\delta),\xi).
\]
We first check the domains needed for its derivatives.

For \(0\le t\le1/4\), the elementary exponential bounds
\[
1+t\le e^t\le\frac1{1-t}
\]
give
\[
t\le d\le\frac43t,\qquad
0\le d\le\frac13,\qquad
1\le B_*\le\frac43.
\]
If \(|\delta|\le1/100\), then
\[
\delta^2\le\frac1{10000},
\qquad
d\delta^2\le\frac1{30000},
\]
so both logarithm arguments in \(\mathsf T\) are strictly positive.

The following calculation also supplies the comparator bounds used to justify smooth composition. Define
\[
\mathcal A_{\mathrm{gap}}=\frac{d\delta^2}{B_*}.
\]
The preceding inequalities imply
\[
0\le\mathcal A_{\mathrm{gap}}\le\frac1{1000},
\qquad
\frac34t\delta^2\le\mathcal A_{\mathrm{gap}}
\le\frac43t\delta^2.
\]
The comparator formula becomes
\[
t-\mathsf T(t,\delta)
=
-\log\!\left(1-\frac52\mathcal A_{\mathrm{gap}}\right)
+\log\!\left(1+\frac53\mathcal A_{\mathrm{gap}}\right).
\]
For \(0\le x_{\log},y_{\log}\le1/100\), the inequalities
\(\log z\le z-1\) and \(1-z^{-1}\le\log z\), valid for \(z>0\), yield
\[
x_{\log}+\frac{100}{101}y_{\log}
\le-\log(1-x_{\log})+\log(1+y_{\log})
\le\frac{100}{99}x_{\log}+y_{\log}.
\]
Indeed,
\[
x_{\log}\le-\log(1-x_{\log})
\le\frac{x_{\log}}{1-x_{\log}}
\le\frac{100}{99}x_{\log},
\]
and
\[
\frac{100}{101}y_{\log}
\le\frac{y_{\log}}{1+y_{\log}}
\le\log(1+y_{\log})\le y_{\log}.
\]
Apply these inequalities with
\(x_{\log}=\frac52\mathcal A_{\mathrm{gap}}\) and
\(y_{\log}=\frac53\mathcal A_{\mathrm{gap}}\). Since
\[
\frac34\left(\frac52+\frac{100}{101}\frac53\right)>3,
\qquad
\frac43\left(\frac{100}{99}\frac52+\frac53\right)<6,
\]
we obtain, including \(t=0\) and \(\delta=0\),
\[
3t\delta^2\le t-\mathsf T(t,\delta)\le6t\delta^2,
\qquad
\frac t2\le\mathsf T(t,\delta)\le t.
\]
The last pair follows from \(6\delta^2\le1/2\).

For every real \(u\), the two sign terms give \(|Z_u|\le2\). Thus, for
\(\xi\in[3/10,1/2]\) and \(|\delta|\le1/100\),
\[
|\delta Z_u|\le\frac1{50},
\qquad
\xi+\delta Z_u\ge\frac3{10}-\frac1{50}>0.
\]
Together with \(t\ge0\) and \(\mathsf T(t,\delta)\ge0\), this makes the risk-shift denominators positive. All these strict positivity conditions persist on ambient neighborhoods. Hence \(\mathcal R_{\mathrm{fair}}\) is jointly smooth near every point of the indicated parameter range.

Sign symmetry and the dependence of \(\mathsf T\) on \(\delta^2\) give
\[
\begin{gathered}
\mathcal R_{\mathrm{fair}}(t,-\delta,\xi,u)
=\mathcal R_{\mathrm{fair}}(t,\delta,\xi,u),\\
\mathcal R_{\mathrm{fair}}(t,0,\xi,u)=0,
\qquad
\partial_\delta\mathcal R_{\mathrm{fair}}(t,0,\xi,u)=0,
\qquad
\mathcal R_{\mathrm{fair}}(0,\delta,\xi,u)=0.
\end{gathered}
\]
The last equality uses \(\mathsf T(0,\delta)=0\),
\(\mathsf G(0,\xi)=\xi\), and \(EZ_u=0\).

The normalized extension is the double Taylor integral
\[
\mathfrak H(t,\delta,\xi,u)
=
\int_0^1\int_0^1
(1-\omega)\,
\partial_t\partial_\delta^2
\mathcal R_{\mathrm{fair}}(\tau t,\omega\delta,\xi,u)
\,d\omega\,d\tau.
\]
Here \(\tau,\omega\in[0,1]\) are integration variables. Their compact range permits one ambient neighborhood on which the integrand is smooth for all integration points. Differentiation under the two integrals then proves joint smoothness of \(\mathfrak H\) on that neighborhood, at every finite order.

The fundamental theorem of calculus, using
\(\partial_\delta^2\mathcal R_{\mathrm{fair}}(0,\delta,\xi,u)=0\), gives
\[
t\int_0^1
\partial_t\partial_\delta^2
\mathcal R_{\mathrm{fair}}(\tau t,\omega\delta,\xi,u)\,d\tau
=
\partial_\delta^2
\mathcal R_{\mathrm{fair}}(t,\omega\delta,\xi,u).
\]
The second-order integral Taylor formula gives
\[
\delta^2\int_0^1
(1-\omega)\,
\partial_\delta^2
\mathcal R_{\mathrm{fair}}(t,\omega\delta,\xi,u)\,d\omega
=
\mathcal R_{\mathrm{fair}}(t,\delta,\xi,u).
\]
Interchanging the two integrals therefore proves the undivided identity
\[
t\delta^2\mathfrak H(t,\delta,\xi,u)
=\mathcal R_{\mathrm{fair}}(t,\delta,\xi,u)
\]
throughout
\(t\in[0,1/4]\), \(|\delta|\le1/100\),
\(\xi\in[3/10,1/2]\), and \(u\in[0,1]\).
For \(t\delta\ne0\), division gives the quotient in the statement.
% lean: CausalSmith.Stat.LogoddsLowsmoothFrontier.comparatorEffect_gap_bounds
% lean: CausalSmith.Stat.LogoddsLowsmoothFrontier.comparatorEffect_range_bounds
% lean: CausalSmith.Stat.LogoddsLowsmoothFrontier.fairEquation_mul_parameters
% lean: CausalSmith.Stat.LogoddsLowsmoothFrontier.fairEquation_eq_div

\item \emph{Fair endpoint factorization and the removable axes.}
At either spatial endpoint the law of \(Z_u\) is one fair sign. Exponentiating the comparator formula gives
\[
e^{\mathsf T(t,\delta)}
=e^t
\frac{1-d\delta^2/(p_*B_*)}
     {1+d\delta^2/((1-p_*)B_*)},
\]
or, equivalently,
\[
e^{\mathsf T(t,\delta)}
\left(\frac35B_*+d\delta^2\right)
=
e^t\left(\frac35B_*-\frac32d\delta^2\right).
\]
Clearing the positive risk denominators using this identity proves
\[
\frac{\mathsf G(t,p_*+\delta)+\mathsf G(t,p_*-\delta)}2
=\mathsf G(\mathsf T(t,\delta),p_*).
\]
It follows first for \(t\delta\ne0\) that
\(\mathfrak H(t,\delta,p_*,u)=0\) at \(u=0,1\).
For each nonzero \(\delta\), continuity in \(t\) extends this equality to \(t=0\); continuity in \(\delta\) then extends it to \(\delta=0\). Thus
\[
\mathfrak H(t,\delta,p_*,0)
=\mathfrak H(t,\delta,p_*,1)=0
\quad
\left(t\in[0,1/4],\ |\delta|\le1/100\right).
\]

To identify every endpoint zero, define
\[
\begin{aligned}
\mathcal S_{\mathrm{end}}(d,\delta,\xi)
&=-25d\delta^2+25d\xi^2-15d\xi-6d+25\xi-15,\\
\mathcal J_{\mathrm{end}}(t,\delta,\xi)
&=\bigl(1+d(\xi+\delta)\bigr)
  \bigl(1+d(\xi-\delta)\bigr)
  \bigl(1+(e^{\mathsf T(t,\delta)}-1)\xi\bigr)
  \left(\frac35B_*+d\delta^2\right).
\end{aligned}
\]
The second expression is positive on the fair range
\(t\in[0,1/4]\), \(|\delta|\le1/100\),
\(\xi\in[3/10,1/2]\). Directly clearing the same denominators and substituting the corrected exponential identity above gives
\[
\mathcal R_{\mathrm{fair}}(t,\delta,\xi,0)
\mathcal J_{\mathrm{end}}(t,\delta,\xi)
=
-\frac{d\delta^2e^t}{50}
(5\xi-2)\mathcal S_{\mathrm{end}}(d,\delta,\xi).
\]
For \(\delta\ne0\), the undivided Taylor identity cancels the amplitude square:
\[
t\mathfrak H(t,\delta,\xi,0)
\mathcal J_{\mathrm{end}}(t,\delta,\xi)
=
-\frac{de^t}{50}
(5\xi-2)\mathcal S_{\mathrm{end}}(d,\delta,\xi).
\]
Continuity in \(\delta\) extends this identity to zero amplitude, where it reads
\[
t\mathfrak H(t,0,\xi,0)(1+d\xi)^3\frac35B_*
=
-\frac{de^t}{50}
(5\xi-2)\mathcal S_{\mathrm{end}}(d,0,\xi).
\]
For \(d\ge0\) and \(\xi\in[3/10,1/2]\), the inequality
\(\xi^2\le\xi/2\) gives
\[
\mathcal S_{\mathrm{end}}(d,\delta,\xi)
\le-\frac52d\xi-6d+25\xi-15
\le-\frac52<0.
\]
Thus, when \(t>0\), the zero-amplitude endpoint residual has the sign of \(5\xi-2\). At either endpoint, when also \(\delta\ne0\), the factored unnormalized identity shows that every matching root in the selection bracket equals \(p_*\).

We also need position independence at zero amplitude. The sign moments are
\[
EZ_u=0,\qquad EZ_u^2
=\cos^2(\pi u/2)+\sin^2(\pi u/2)=1.
\]
Differentiating the averaged risk twice in amplitude gives
\[
\left.
\partial_\delta^2 E\mathsf G(t,\xi+\delta Z_u)
\right|_{\delta=0}
=
-\frac{2e^td}{(1+d\xi)^3}.
\]
The comparator term has no spatial argument, so
\[
\partial_\delta^2
\mathcal R_{\mathrm{fair}}(t,0,\xi,u)
=
\partial_\delta^2
\mathcal R_{\mathrm{fair}}(t,0,\xi,0).
\]
Their effect derivatives agree as well, including at the endpoints of the effect interval by smoothness and the one-sided derivative there. Substitution in the double integral proves
\[
\mathfrak H(t,0,\xi,u)=\mathfrak H(t,0,\xi,0).
\]

For zero effect, differentiation of the explicit comparator gives
\[
\partial_t\mathsf T(0,\delta)=1-\frac{\delta^2}{6/25},
\qquad
\partial_t\mathsf G(0,\xi)=\xi(1-\xi).
\]
Using the two sign moments,
\[
\partial_t\mathcal R_{\mathrm{fair}}(0,\delta,\xi,u)
=
\delta^2\left(-1+\frac{\xi(1-\xi)}{6/25}\right).
\]
Differentiate the undivided Taylor identity at \(t=0\). For \(\delta\ne0\), cancellation of \(\delta^2\) gives the first identity below; continuity in \(\delta\) supplies the same identity at \(\delta=0\):
\[
\mathfrak H(0,\delta,\xi,u)
=-1+\frac{\xi(1-\xi)}{6/25},
\qquad
\mathfrak H(0,\delta,\xi,u)=0
\ \Longleftrightarrow\
(\xi-\tfrac25)(\xi-\tfrac35)=0.
\]
Within \((3/10,1/2)\), the zero is therefore \(p_*\).
Combining these facts yields
\[
\mathfrak H(t,0,p_*,u)=0,
\qquad
\mathfrak H(t,0,\xi,u)=0,\quad \xi\in(3/10,1/2)
\ \Longrightarrow\ \xi=p_*.
\]
Here \(t=0\) uses the explicit quadratic formula, and \(t>0\) uses the negative residual factor.
% lean: CausalSmith.Stat.LogoddsLowsmoothFrontier.fairEquation_endpoint_factor_zero_amplitude
% lean: CausalSmith.Stat.LogoddsLowsmoothFrontier.fairEquation_zero_amplitude_unique
% lean: CausalSmith.Stat.LogoddsLowsmoothFrontier.fairEquation_zero_effect

\item \emph{Uniform fair-root existence.}
The preceding factorization shows opposite bracket signs at zero amplitude when \(t>0\). At \(t=0\), the explicit extension gives
\[
\mathfrak H(0,0,3/10,u)=-\frac18,
\qquad
\mathfrak H(0,0,1/2,u)=\frac1{24}.
\]
Hence throughout the compact effect and spatial ranges,
\[
\mathfrak H(t,0,3/10,u)<0,
\qquad
\mathfrak H(t,0,1/2,u)>0.
\]
Joint continuity and compactness preserve these strict inequalities for
\(|\delta|\le\varepsilon_{\mathrm{exist}}\), for some
\[
0<\varepsilon_{\mathrm{exist}}\le\frac1{100}.
\]
Continuity in the margin and the intermediate value theorem then give
\[
\forall(t,\delta,u)\in
\mathcal P_{\mathrm{fair}}(\varepsilon_{\mathrm{exist}}),
\qquad
\exists\xi\in(3/10,1/2):
\mathfrak H(t,\delta,\xi,u)=0.
\]
The strict endpoint signs put the zero in the open bracket.
% lean: CausalSmith.Stat.LogoddsLowsmoothFrontier.fairEquation_uniform_exists

\item \emph{Smoothness of the selected fair root.}
Differentiate the zero-amplitude factorization at \(\xi=p_*\), where the residual itself vanishes. This gives
\[
t\,\partial_\xi\mathfrak H(t,0,p_*,u)
B_*^3\frac35B_*
=
\frac{de^t(8d+5)}{10}.
\]
For \(t>0\), all factors except the derivative on the left are positive, and the right side is positive. At zero effect, the explicit extension gives
\[
\partial_\xi\mathfrak H(0,0,p_*,u)=\frac56.
\]
Therefore the center slope is positive throughout the full effect and spatial ranges.

Continuity of this partial derivative and compactness give a radius
\(r_{\mathrm{center}}\) satisfying
\[
0<r_{\mathrm{center}}\le\frac1{100},
\qquad
\partial_\xi\mathfrak H(t,\delta,\xi,u)>0
\]
whenever
\[
t\in[0,1/4],\quad
|\delta|\le r_{\mathrm{center}},\quad
|\xi-p_*|\le r_{\mathrm{center}},\quad
u\in[0,1].
\]
The zero-amplitude uniqueness from the preceding steps also confines all nearby roots. Indeed, the set
\[
\left\{
(t,u,\xi):
t\in[0,1/4],\ u\in[0,1],\
\xi\in[3/10,1/2],\
|\xi-p_*|\ge r_{\mathrm{center}}
\right\}
\]
is compact, and \(\mathfrak H(t,0,\xi,u)\) is nonzero on it. Continuity therefore gives
\(0<\varepsilon_{\mathrm{conf}}\le1/100\) such that
\[
\mathfrak H(t,\delta,\xi,u)=0,\quad
\xi\in[3/10,1/2],\quad
|\delta|\le\varepsilon_{\mathrm{conf}}
\quad\Longrightarrow\quad
|\xi-p_*|<r_{\mathrm{center}}.
\]
If the displayed compact set is empty, the implication follows directly from its defining condition.

Take the minimum of
\(r_{\mathrm{center}},\varepsilon_{\mathrm{conf}}\), and
\(\varepsilon_{\mathrm{exist}}\). Every selected root on the resulting closed fair region is a genuine bracket zero, lies in the center neighborhood, and has nonzero margin derivative. Moreover, these controls give uniqueness on ambient parameter neighborhoods: positivity persists on the compact center bracket, and nonvanishing persists on the compact portion of the larger bracket outside it. Any two nearby bracket zeros must therefore lie in the center bracket, where the residual is strictly increasing.

Apply \cref{obj:lem:scalar-implicit-function} at each selected zero. Nearby bracket uniqueness identifies the selected root with the local graph. The first-derivative identity is
\[
\mathrm D\mathfrak p(t,\delta,u)
=
-\left.
\frac{
\bigl(\partial_t\mathfrak H,
      \partial_\delta\mathfrak H,
      \partial_u\mathfrak H\bigr)}
{\partial_\xi\mathfrak H}
\right|_{\xi=\mathfrak p(t,\delta,u)} .
\]
The regularity upgrade used for the mixed graph applies again. Thus there is
\[
0<\varepsilon_{\mathrm{smooth}}\le\frac1{100}
\]
such that \(\mathfrak p\) is smooth on a neighborhood of every point of
\(\mathcal P_{\mathrm{fair}}(\varepsilon_{\mathrm{smooth}})\).
% lean: CausalSmith.Stat.LogoddsLowsmoothFrontier.fairRoot_uniform_contDiffAt

\item \emph{A preliminary common radius and the quadratic fair bound.}
Define
\[
\varepsilon_{\mathrm{pre}}
=
\min\{\varepsilon_{\mathrm{smooth}},
      \varepsilon_{\mathrm q},
      \varepsilon_{\mathrm{cell}}\},
\qquad
\varepsilon_0
=
\min\{\varepsilon_{\mathrm{pre}},
      \varepsilon_{\mathrm{exist}}\}.
\]
Then
\[
0<\varepsilon_0\le\frac1{100},
\qquad
\varepsilon_0\le\varepsilon_{\mathrm q},
\qquad
\varepsilon_0\le\varepsilon_{\mathrm{cell}},
\]
and \(\mathfrak p\) is smooth on pointwise neighborhoods of its closed region.

The endpoint and axis computations identify the selected roots there:
\[
\mathfrak p(t,\delta,0)=p_*,
\qquad
\mathfrak p(t,0,u)=p_*,
\qquad
\mathfrak p(0,\delta,u)=p_*.
\]
For the first equality, split first on \(t=0\), then on \(\delta=0\); the zero-effect and zero-amplitude uniqueness computations settle these cases. When both are nonzero, endpoint matching and the negative residual factor force the selected bracket root to be \(p_*\).

The numerator is even in \(\delta\). Its second amplitude derivative is consequently even, so its effect derivative and the double integral satisfy
\[
\mathfrak H(t,-\delta,\xi,u)=\mathfrak H(t,\delta,\xi,u).
\]
The selectors for these identical equations have the same value, including their common fallback convention. Hence
\[
\mathfrak p(t,-\delta,u)=\mathfrak p(t,\delta,u).
\]

Smoothness and compactness give positive bounds for the two partial derivatives needed below. More explicitly, define
\[
\begin{aligned}
C_{\delta\delta}
&=
\max_{\mathcal P_{\mathrm{fair}}(\varepsilon_0)}
|\partial_\delta^2\mathfrak p|,\\
C_{\delta\delta u}
&=
\max_{\mathcal P_{\mathrm{fair}}(\varepsilon_0)}
|\partial_\delta^2\partial_u\mathfrak p|,\\
B_{\mathrm{quad}}
&=\max\{1,C_{\delta\delta},C_{\delta\delta u}\},
\qquad M=2B_{\mathrm{quad}}>0.
\end{aligned}
\]
Both maxima exist because the indicated partial derivatives are continuous on the compact region.

For a smooth even function \(\varphi\), its derivative at zero vanishes, and the integral Taylor formula gives
\[
\varphi(\delta)-\varphi(0)
=
\delta^2\int_0^1(1-\omega)\varphi''(\omega\delta)\,d\omega.
\]
If \(|\varphi''|\le B_{\mathrm{quad}}\) on the segment joining \(0\) to \(\delta\), then
\[
|\varphi(\delta)-\varphi(0)|
\le B_{\mathrm{quad}}\delta^2.
\]
Every amplitude on that segment has absolute value at most
\(|\delta|\le\varepsilon_0\).

Apply this first to
\(\varphi(\delta')=\mathfrak p(t,\delta',u)\), and then to
\(\varphi(\delta')=\partial_u\mathfrak p(t,\delta',u)\), with \(t,u\) fixed in their prescribed ranges. Both functions are smooth and even. Their zero-amplitude values are respectively \(p_*\) and \(0\). For the latter value, constancy of
\(\mathfrak p(t,0,u)=p_*\) on \([0,1]\) forces the spatial derivative to vanish also at \(u=0,1\): the available two-sided derivative equals the zero one-sided derivative from the interval. Therefore
\[
|\mathfrak p(t,\delta,u)-p_*|
\le B_{\mathrm{quad}}\delta^2,
\qquad
|\partial_u\mathfrak p(t,\delta,u)|
\le B_{\mathrm{quad}}\delta^2,
\]
and hence
\[
|\mathfrak p(t,\delta,u)-p_*|
+|\partial_u\mathfrak p(t,\delta,u)|
\le M\delta^2
\quad\hbox{on }\mathcal P_{\mathrm{fair}}(\varepsilon_0).
\]
% lean: CausalSmith.Stat.LogoddsLowsmoothFrontier.fairRoot_endpoint_center
% lean: CausalSmith.Stat.LogoddsLowsmoothFrontier.fairRoot_quadratic_bound_of_smooth

\item \emph{Fair-cell regularity and common derivative bounds.}
For the same auxiliary indices as before, define
\[
\begin{aligned}
\mu^{\mathrm{loc}}_{0,\iota}
&=
\begin{cases}
\mathfrak p(t,\delta,u),&\iota=0,\\
\mathfrak p(t,\delta,u)+\delta Z_u,&\iota=1,
\end{cases}\\
t^{\mathrm{loc}}_\iota
&=
\begin{cases}
\mathsf T(t,\delta),&\iota=0,\\
t,&\iota=1,
\end{cases}\\
\mu^{\mathrm{loc}}_{1,\iota}
&=\mathsf G(t^{\mathrm{loc}}_\iota,\mu^{\mathrm{loc}}_{0,\iota}),\\
\bigl(
\mathcal K^{\mathrm{fair}}_{\iota,\boldsymbol{s},1,1},
\mathcal K^{\mathrm{fair}}_{\iota,\boldsymbol{s},1,0},
\mathcal K^{\mathrm{fair}}_{\iota,\boldsymbol{s},0,1},
\mathcal K^{\mathrm{fair}}_{\iota,\boldsymbol{s},0,0}
\bigr)
&=
\frac12\bigl(
\mu^{\mathrm{loc}}_{1,\iota},
1-\mu^{\mathrm{loc}}_{1,\iota},
\mu^{\mathrm{loc}}_{0,\iota},
1-\mu^{\mathrm{loc}}_{0,\iota}
\bigr).
\end{aligned}
\]
These are functions of \((t,\delta,u)\), with the signs evaluated in \(Z_u\).

On every input the fair selector belongs to \((3/10,1/2)\): it either selects a root there or takes the fallback value \(p_*\), which is also in that bracket. Consequently, when \(|\delta|\le1/100\),
\[
\frac14\le\mu^{\mathrm{loc}}_{0,\iota}\le\frac35.
\]
Indeed, \(|\delta Z_u|\le1/50\), and adding this bound to either side of the selection bracket proves the displayed bounds. The local effect is nonnegative by the comparator range estimate. Hence
\[
1+\bigl(e^{t^{\mathrm{loc}}_\iota}-1\bigr)
\mu^{\mathrm{loc}}_{0,\iota}>0.
\]
The chain rule, root smoothness, and positivity of the logarithm arguments now prove smoothness of every fully substituted fair cell on neighborhoods of the closed fair region.

For each \(m\ge0\), compactness gives positive bounds
\[
\|\mathrm D^m\mathfrak p\|\le B_{\mathrm p,m},
\qquad
\|\mathrm D^m\mathcal K^{\mathrm{fair}}_j\|
\le C^{\mathrm{fair}}_{j,m}
\quad\hbox{on }\mathcal P_{\mathrm{fair}}(\varepsilon_0).
\]
Set
\[
\begin{aligned}
B_{\mathrm{fair},m}
&=B_{\mathrm p,m}
  +\sum_{j\in\mathcal I_{\mathrm{aux}}}C^{\mathrm{fair}}_{j,m},\\
B_{\mathrm{rem},m}
&=\max\{B_{\mathrm{cell},m},B_{\mathrm{fair},m}\},\\
B_m&=\max\{B_{\mathrm q,m},B_{\mathrm{rem},m}\}>0.
\end{aligned}
\]
Restriction to the smaller mixed region
\(\mathcal P_{\mathrm{mix}}(\varepsilon_0)\) preserves its earlier bounds. Thus
\[
\begin{array}{ll}
\|\mathrm D^m\mathfrak q\|\le B_m,\quad
\|\mathrm D^m\mathcal K^{\mathrm{mix}}_j\|\le B_m
&\text{on }\mathcal P_{\mathrm{mix}}(\varepsilon_0),\\[2mm]
\|\mathrm D^m\mathfrak p\|\le B_m,\quad
\|\mathrm D^m\mathcal K^{\mathrm{fair}}_j\|\le B_m
&\text{on }\mathcal P_{\mathrm{fair}}(\varepsilon_0).
\end{array}
\]
All these functions are \(C^\infty\) on their respective closed regions.
% lean: CausalSmith.Stat.LogoddsLowsmoothFrontier.localFairCell_smooth_derivative_bounds_of_root
% lean: CausalSmith.Stat.LogoddsLowsmoothFrontier.exact_calibrations

\item \emph{Fixed open branch neighborhoods.}
We construct open sets for the actual selectors and cells. Define
\[
\begin{aligned}
\mathcal A_{\mathrm{mix}}
=\bigl\{(\eta,\zeta,u)\in\mathbb R^3:\;&
\mathfrak q(\eta,\zeta,u)
\text{ is the unique zero of }
\mathcal F_{\mathrm{mix}}(\eta,\zeta,u,\cdot)
\text{ in }(3/4,5/4),\\[-1mm]
&\mathfrak q\text{ and every }\mathcal K^{\mathrm{mix}}_j
\text{ are smooth at }(\eta,\zeta,u)\bigr\},\\
U&=\operatorname{int}\mathcal A_{\mathrm{mix}}.
\end{aligned}
\]
At every \((0,0,u)\), the local implicit graph through \(1\), its bracket range, and the persistent negative root derivative give the unique selected zero nearby. The root and all mixed cells are smooth on those same neighborhoods. Therefore
\[
(0,0,u)\in U\qquad(u\in[0,1]).
\]
Openness and compactness of this spatial segment give
\(\varepsilon_{\mathrm{mix}}>0\) such that
\[
\mathcal P_{\mathrm{mix}}(\varepsilon_{\mathrm{mix}})\subset U.
\]
By the defining conditions, the mixed selector is the unique bracket zero throughout \(U\), and it and all mixed cells are smooth there.

For the fair branch fix
\[
b=\frac1{20},\qquad 0<b<\frac1{10},
\]
and define
\[
\begin{aligned}
\mathcal A_{\mathrm{fair}}
=\bigl\{(t,\delta,u)\in\mathbb R^3:\;&
\mathfrak p(t,\delta,u)\in(p_*-b,p_*+b)\cap(3/10,1/2),\\
&\mathfrak p(t,\delta,u)
\text{ is the unique zero of }
\mathfrak H(t,\delta,\cdot,u)
\text{ in }(p_*-b,p_*+b),\\
&\mathfrak p\text{ and every }\mathcal K^{\mathrm{fair}}_j
\text{ are smooth at }(t,\delta,u)\bigr\},\\
V&=\operatorname{int}\mathcal A_{\mathrm{fair}}.
\end{aligned}
\]
At each \((t,0,u)\) in the full effect and spatial ranges, the local implicit graph starts at \(p_*\), and therefore stays in
\((p_*-b,p_*+b)\) nearby. The ambient uniqueness argument from the center derivative and compact root exclusion identifies it with the selector from the larger bracket. Smoothness of the residual holds on an ambient open set; the derivative upgrade applies to this same graph at every nearby point. The outer cell formulas have positive denominators and logarithm arguments at zero amplitude, so all fair cells are smooth nearby as well. Consequently,
\[
(t,0,u)\in V
\qquad(t\in[0,1/4],\ u\in[0,1]).
\]
Compactness gives a positive amplitude radius for which the corresponding closed fair region lies in \(V\). Reduce that radius to obtain
\[
0<\varepsilon_{\mathrm{fair}}\le\frac1{100},
\qquad
\mathcal P_{\mathrm{fair}}(\varepsilon_{\mathrm{fair}})\subset V.
\]

For the residual itself define
\[
W=
\operatorname{int}
\left\{
(t,\delta,\xi,u)\in\mathbb R^4:
\mathfrak H\text{ is smooth at }(t,\delta,\xi,u)
\right\}.
\]
The compact integration ranges in its Taylor representation give a common neighborhood of smoothness around every point of
\[
t\in[0,1/4],\quad |\delta|\le1/100,\quad
\xi\in[3/10,1/2],\quad u\in[0,1].
\]
Since
\((p_*-b,p_*+b)\subset[3/10,1/2]\), the required centered fair core with amplitude radius \(\varepsilon_{\mathrm{fair}}\) is contained in \(W\). The residual is smooth on \(W\), and its quotient identity on this core follows from the undivided Taylor identity.

Finally define
\[
\varepsilon_{\mathrm{branch}}
=\frac12\min\{\varepsilon_{\mathrm{mix}},
             \varepsilon_{\mathrm{fair}}\}>0.
\]
Then
\[
\varepsilon_{\mathrm{branch}}<\varepsilon_{\mathrm{mix}},
\qquad
\varepsilon_{\mathrm{branch}}<\varepsilon_{\mathrm{fair}},
\]
which supplies the strict separation between the eventual closed-region radius and both open-neighborhood radii.
% lean: CausalSmith.Stat.LogoddsLowsmoothFrontier.calibration_branch_neighbourhood_exists

\item \emph{The final common radius.}
Set
\[
\varepsilon
=
\min\left\{
\varepsilon_0,\varepsilon_{\mathrm{root}},
\frac1{100},\varepsilon_{\mathrm{branch}}
\right\}>0.
\]
In particular,
\[
\varepsilon\le\varepsilon_0,\qquad
\varepsilon\le\varepsilon_{\mathrm{root}},\qquad
\varepsilon\le\frac1{100}\le\frac18,
\]
and
\[
\varepsilon<\varepsilon_{\mathrm{mix}},
\qquad
\varepsilon<\varepsilon_{\mathrm{fair}}.
\]
The inclusions
\[
\mathcal P_{\mathrm{mix}}(\varepsilon)
\subset\mathcal P_{\mathrm{mix}}(\varepsilon_0),
\qquad
\mathcal P_{\mathrm{fair}}(\varepsilon)
\subset\mathcal P_{\mathrm{fair}}(\varepsilon_0)
\]
preserve smoothness, every derivative bound \(B_m\), and the quadratic estimate with the same positive \(M\). The open sets and the larger signed radii remain as constructed.
% lean: CausalSmith.Stat.LogoddsLowsmoothFrontier.exact_calibrations

\item \emph{Positive odds tables.}
Suppose
\[
|t|\le\varepsilon,\qquad
|\xi-\tfrac12|\le\varepsilon,\qquad
|\upsilon-\tfrac12|\le\varepsilon.
\]
Since \(\varepsilon\le1/8\), the standard estimate
\(|e^t-1|\le2|t|\) for \(|t|\le1\) gives
\[
|d|\le\frac14,\qquad
\frac38\le\xi,\upsilon,1-\xi,1-\upsilon\le\frac58.
\]
Write the branch radicand as
\[
\Delta_{\mathrm{cov}}=L_*^2-4d^2v_*.
\]
The elementary variance and margin-product bounds give
\[
0\le v_*\le\frac1{16},
\qquad
0\le\xi(1-\upsilon)+(1-\xi)\upsilon\le1,
\]
and therefore
\[
L_*\ge\frac34,\qquad
d^2\le\frac1{16},\qquad
\Delta_{\mathrm{cov}}\ge0,\qquad
L_*+\sqrt{\Delta_{\mathrm{cov}}}\ge\frac34.
\]
The covariance formula consequently satisfies
\[
|\mathfrak c(t,\xi,\upsilon)|
=
\frac{2|d|v_*}{L_*+\sqrt{\Delta_{\mathrm{cov}}}}
\le
\frac{2(1/4)(1/16)}{3/4}
=\frac1{24}.
\]
Every unadjusted cell product is at least \(9/64\). Thus each adjusted cell is at least
\[
\frac9{64}-\frac1{24}>0.
\]
Direct addition gives both prescribed margins and total mass one.

Put \(c=\mathfrak c(t,\xi,\upsilon)\). Squaring the square root and clearing its positive branch denominator gives
\[
dc^2-L_*c+dv_*=0.
\]
Expanding the four cells now yields
\[
p_{11}p_{00}-e^t p_{10}p_{01}
=-(dc^2-L_*c+dv_*)=0.
\]
Since \(p_{10}p_{01}>0\), division proves the stated odds ratio.

Finally, \(d=tD(t)\) in the branch formulas gives
\[
\mathfrak c(t,\xi,\upsilon)=t\mathfrak B(t,\xi,\upsilon).
\]
Thus \(\mathfrak B=\mathfrak c/t\) when \(t\ne0\). At \(t=0\), the branch denominator is \(2\) and \(D(0)=1\), so
\[
\mathfrak B(0,\xi,\upsilon)
=\xi(1-\xi)\upsilon(1-\upsilon).
\]
% lean: CausalSmith.Stat.LogoddsLowsmoothFrontier.covarianceBranch_table_positive
% lean: CausalSmith.Stat.LogoddsLowsmoothFrontier.covarianceBranch_odds_ratio
% lean: CausalSmith.Stat.LogoddsLowsmoothFrontier.normalizedBranch_eq_div
% lean: CausalSmith.Stat.LogoddsLowsmoothFrontier.normalizedBranch_zero

\item \emph{The signed mixed identities.}
For \((\eta,\zeta,u)\in\mathcal P_{\mathrm{mix}}(\varepsilon)\), the established bracket-root existence makes the selector a genuine root:
\[
\frac34<\mathfrak q(\eta,\zeta,u)<\frac54,
\qquad
\mathcal F_{\mathrm{mix}}
(\eta,\zeta,u,\mathfrak q(\eta,\zeta,u))=0.
\]
The two endpoint sign laws give identical residual functions and therefore identical selected values:
\[
\mathfrak q(\eta,\zeta,0)=\mathfrak q(\eta,\zeta,1).
\]
Using \(\mathfrak c=t\mathfrak B\), multiplication of the normalized root equation gives
\[
\begin{aligned}
&E\mathfrak c
\bigl(32\eta\zeta,
\tfrac12-\eta\mathfrak q(\eta,\zeta,u)Z_u,
\tfrac12+\zeta Z_u\bigr)\\
&\hspace{15mm}
=32\eta\zeta\,E\mathfrak B
\bigl(32\eta\zeta,
\tfrac12-\eta\mathfrak q(\eta,\zeta,u)Z_u,
\tfrac12+\zeta Z_u\bigr)
=2\eta\zeta\mathfrak q(\eta,\zeta,u).
\end{aligned}
\]
This multiplication also covers either zero-amplitude axis.

The separate sign symmetries follow algebraically from the normalized branch. First,
\[
D(-t)=\frac{D(t)}{e^t};
\]
for \(t\ne0\) this follows from \(tD(t)=e^t-1\), and at \(t=0\) it follows from \(D(0)=1\). Under
\((t,\xi,\upsilon)\mapsto(-t,1-\xi,\upsilon)\), the branch coefficients transform as
\[
d\mapsto-\frac d{e^t},
\qquad
v_*\mapsto v_*,
\qquad
L_*\mapsto\frac{L_*}{e^t}.
\]
The square root and the whole denominator are consequently divided by the positive number \(e^t\), just as the normalized numerator is. This proves
\[
\mathfrak B(-t,1-\xi,\upsilon)
=\mathfrak B(t,\xi,\upsilon).
\]
Exchanging the two margins gives
\[
\mathfrak B(-t,\xi,1-\upsilon)
=\mathfrak B(t,\xi,\upsilon).
\]
Substitution in the mixed residual gives, with the root argument fixed,
\[
\mathcal F_{\mathrm{mix}}(-\eta,\zeta,u,\chi)
=\mathcal F_{\mathrm{mix}}(\eta,\zeta,u,\chi)
=\mathcal F_{\mathrm{mix}}(\eta,-\zeta,u,\chi).
\]
Identical residuals give identical selectors, so the two asserted root symmetries follow.

On the axes, the zero-effect normalized formula and \(EZ_u^2=1\) reduce the residuals to
\[
\mathcal F_{\mathrm{mix}}(0,\zeta,u,\chi)
=1-4\zeta^2-\chi,
\qquad
\mathcal F_{\mathrm{mix}}(\eta,0,u,\chi)
=1-4\eta^2\chi^2-\chi.
\]
When \(|\zeta|\le1/8\), the candidate \(1-4\zeta^2\) lies strictly in \((3/4,5/4)\); the affine equation forces every selected zero to equal it. Hence
\[
\mathfrak q(0,\zeta,u)=1-4\zeta^2.
\]
For the other axis, \(\eta^2\le1/64\) gives
\[
1-4\eta^2(7/8)^2-\frac78\ge0,
\qquad
1-4\eta^2-1\le0.
\]
The intermediate value theorem supplies a zero in \([7/8,1]\), hence in the selection bracket. Substitution of the selected zero into the reduced equation proves
\[
\mathfrak q(\eta,0,u)
+4\eta^2\mathfrak q(\eta,0,u)^2=1.
\]
These arguments include \(\eta=0\) and \(\zeta=0\).
% lean: CausalSmith.Stat.LogoddsLowsmoothFrontier.mixedRoot_spec
% lean: CausalSmith.Stat.LogoddsLowsmoothFrontier.mixedEquation_covariance_matching
% lean: CausalSmith.Stat.LogoddsLowsmoothFrontier.mixedRoot_even_left
% lean: CausalSmith.Stat.LogoddsLowsmoothFrontier.mixedRoot_even_right
% lean: CausalSmith.Stat.LogoddsLowsmoothFrontier.mixedRoot_zero_left
% lean: CausalSmith.Stat.LogoddsLowsmoothFrontier.mixedRoot_zero_right

\item \emph{The fair identities and separation.}
Fix
\[
t\in[0,1/4],\qquad |\delta|\le\varepsilon,\qquad u\in[0,1],
\]
and put \(p=\mathfrak p(t,\delta,u)\). The comparator calculation already gives
\[
3t\delta^2\le t-\mathsf T(t,\delta)\le6t\delta^2,
\qquad
\frac t2\le\mathsf T(t,\delta)\le t.
\]
The preliminary radius construction supplies a bracket root and the identities
\[
\mathfrak p(t,\delta,0)=p_*,
\qquad
\mathfrak p(t,0,u)=p_*,
\qquad
\mathfrak p(0,\delta,u)=p_*.
\]
The endpoint sign laws give identical fair residuals at \(u=0,1\), hence
\[
\mathfrak p(t,\delta,0)=\mathfrak p(t,\delta,1)=p_*.
\]
Evenness supplies \(\mathfrak p(t,-\delta,u)=p\), and the previously obtained quadratic bound gives
\[
|p-p_*|+|\partial_u\mathfrak p(t,\delta,u)|
\le M\delta^2.
\]

To verify the normalized equation on the entire range, first consider \(\delta=0\). The zero-amplitude center is a genuine bracket zero. If \(\delta\ne0\) and \(u\in\{0,1\}\), the endpoint center is a genuine bracket zero by the continuity argument through both axes. In the remaining case, uniform bracket-root existence and the selector definition give the root equation directly. Thus in all cases
\[
\mathfrak H(t,\delta,p,u)=0.
\]
The selector lies in \((3/10,1/2)\), so the fair Taylor identity applies at this margin. For \(t\delta\ne0\), it gives
\[
\mathfrak H(t,\delta,p,u)
=
\frac{\mathcal R_{\mathrm{fair}}(t,\delta,p,u)}{t\delta^2}.
\]
For all \(t,\delta\), including the axes, its undivided form gives
\[
\mathcal R_{\mathrm{fair}}(t,\delta,p,u)
=t\delta^2\mathfrak H(t,\delta,p,u)=0.
\]
By the definition of the numerator, this is exactly
\[
E\mathsf G(t,p+\delta Z_u)
=\mathsf G(\mathsf T(t,\delta),p).
\]
Together with the closed-region regularity, derivative bounds, and open neighborhoods already constructed, these are all the asserted properties of the positive constants \(\varepsilon\) and \(M\).
% lean: CausalSmith.Stat.LogoddsLowsmoothFrontier.fair_matching_of_equation
% lean: CausalSmith.Stat.LogoddsLowsmoothFrontier.exact_calibrations
\end{enumerate}
\end{proof}

The mixed calibration fixes the mean covariance adjustment exactly while allowing both margins to vary with the signs. The fair calibration instead matches the mean shifted risk, expressed through the logistic risk shift \leanref{sym:\mathsf G}{\ensuremath{\mathsf G(t,\xi)}} from \cref{obj:def:calibration-maps}, to a deterministic comparator. Its effect separation is between \(3t\delta^2\) and \(6t\delta^2\), while the comparator effect remains between \(t/2\) and \(t\). These identities specify the separation and matching simultaneously on the stated parameter regions.

The smoothness conclusion concerns the fully substituted cell probabilities as well as the calibration roots. Thus the derivative bounds incorporate the dependence of each cell on its calibrated margin. The open neighborhoods also include zero amplitudes and zero effect, where the normalized equations have their smooth extensions. In the fair family, both the displacement of the calibrated margin from \(2/5\) and its within-cell derivative are bounded by a constant times the squared amplitude.

We now apply these local maps to the shared-endpoint construction in \cref{obj:def:continuous-calibrated-priors}. A finite endpoint sign vector \leanref{sym:\sigma}{\ensuremath{\sigma}} determines the trigonometric field on the equal partition, with adjacent cells sharing the sign at their common endpoint. The endpoint identities in \cref{obj:lem:exact-calibrations} make the calibrated margins compatible with this continuous field. The following statement certifies model membership for every sign realization, together with the exact singleton matching and uniform amplitude derivatives needed for the sample mixtures.

\begin{lemmav}{lem:calibrated-support}[Uniform calibrated support and smoothness]
There exist jointly chosen absolute constants \(\varepsilon\), the calibration radius, and \(M\), the derivative bound, satisfying
\[
0<\varepsilon\le \frac{1}{100},
\qquad M\ge 1,
\]
with the following properties for the calibrated families of \cref{obj:def:continuous-calibrated-priors}. Throughout, \(k\ge1\) is an integer and \(\sigma\in\{-1,1\}^{k+1}\) is a sign vector.

\begin{itemize}
\item \textbf{(Mixed membership and effects.)}
For every \((\alpha,\beta)\in\mathcal D\), the prescribed domain of public smoothness exponents, set
\[
\eta=\varepsilon k^{-\alpha},
\qquad
\zeta=\varepsilon k^{-\beta}.
\]
The calibrated mixed laws \(P_{0,\sigma}\) and \(P_{1,\sigma}\) at these amplitudes belong to \(\mathcal M(\alpha,\beta)\) as defined in \cref{obj:def:model}, and their homogeneous conditional log-odds coefficients satisfy
\[
\theta(P_{0,\sigma})=0,
\qquad
\theta(P_{1,\sigma})=32\eta\zeta.
\]
For each treatment label, outcome label, and covariate \(x\), the sums over all sign vectors of the corresponding conditional cell probabilities under the two mixed laws are equal.

\item \textbf{(Fair membership and effects.)}
For every \((\alpha,\beta)\in\mathcal D\), set
\[
\delta=\varepsilon k^{-\beta}.
\]
For every \(t\in[0,1/4]\), the fair random law \(P_{t,\sigma}\) and the deterministic fair comparator \(P_{*,t,\delta}\) belong to \(\mathcal M(\alpha,\beta)\), and
\[
\theta(P_{t,\sigma})=t,
\qquad
\theta(P_{*,t,\delta})=\mathsf T(t,\delta),
\]
where \(\mathsf T(t,\delta)\) is the calibrated comparator effect defined in \cref{obj:def:calibration-maps}. For each treatment label, outcome label, and covariate \(x\), the uniform average over the \(2^{k+1}\) sign vectors of the conditional cell probability under \(P_{t,\sigma}\) equals the corresponding comparator cell probability.

\item \textbf{(Signed validity and derivative bounds.)}
For all signed amplitudes satisfying
\[
|\eta|\le\varepsilon,
\qquad
|\zeta|\le\varepsilon,
\qquad
|\delta|\le\varepsilon,
\]
the conditional four-cell tables of both mixed branches are valid probability tables, as are those of both fair branches for every \(t\in[0,1/4]\). For every sign vector, branch, treatment label, outcome label, and covariate \(x\), their one-record densities relative to uniform four-label measure lie in \([1/2,2]\). The mixed densities, viewed as functions of \((\eta,\zeta)\), and the fair densities, viewed as functions of \(\delta\) with \(t\) fixed, are twice continuously differentiable at every point of the corresponding closed amplitude box. The norms of their first and second total amplitude derivatives are at most \(M\), uniformly over all these indices and amplitudes. These bounds include the mixed second derivative in \(\eta\) and \(\zeta\) and the deterministic comparator's derivatives in \(\delta\).

\item \textbf{(Common smooth neighborhoods.)}
There exists an open subset of \(\mathbb R^2\) containing
\[
\{(\eta,\zeta):|\eta|\le\varepsilon,\ |\zeta|\le\varepsilon\}
\]
on which every mixed density is infinitely continuously differentiable in its amplitude arguments. There also exists an open subset of \(\mathbb R\) containing
\[
[-\varepsilon,\varepsilon]
\]
on which every fair density is infinitely continuously differentiable in \(\delta\). Each neighborhood is common to every integer \(k\ge1\), sign vector, branch, treatment label, outcome label, and covariate \(x\); the fair neighborhood is also common to every \(t\in[0,1/4]\).

\item \textbf{(Signed singleton matching.)}
For every \(|\eta|\le\varepsilon\) and \(|\zeta|\le\varepsilon\), and each treatment label, outcome label, and covariate \(x\), the sums over all sign vectors of the corresponding conditional cell probabilities under the two mixed laws are equal. For every \(|\delta|\le\varepsilon\) and \(t\in[0,1/4]\), the uniform average over all \(2^{k+1}\) sign vectors of each conditional cell probability under the fair random law equals the corresponding cell probability under the deterministic fair comparator.
\end{itemize}
\end{lemmav}

\begin{proof}[Proof of \cref{obj:lem:calibrated-support}]
We use the calibration maps and families of
\cref{obj:def:calibration-maps,obj:def:continuous-calibrated-priors}.
All finite coordinate spaces below carry the maximum norm, and total derivatives carry the induced multilinear operator norm.

\begin{enumerate}
\item \textit{Local formulas and smoothness radii.}
By \cref{obj:lem:exact-calibrations}, choose a calibration radius
\(\varepsilon_{\mathrm{cal}}>0\) on which the stated smoothness, endpoint, and matching identities hold.

For a local sign pair \(\mathbf s=(s_1,s_2)\in\{-1,1\}^2\), write
\[
Z_u(\mathbf s)
=s_1\cos(\pi u/2)+s_2\sin(\pi u/2).
\]
For \(a,y\in\{0,1\}\), define the local mixed cells by
\[
\begin{gathered}
e^{\mathrm{loc}}_{b,\mathbf s}(\eta,\zeta,u)
=\frac12+(1-2b)\eta\mathfrak q(\eta,\zeta,u)Z_u(\mathbf s),
\qquad b\in\{0,1\},\\
m^{\mathrm{loc}}_{\mathbf s}(\zeta,u)
=\frac12+\zeta Z_u(\mathbf s),\\
c^{\mathrm{loc}}_{0,\mathbf s}(\eta,\zeta,u)=0,\qquad
c^{\mathrm{loc}}_{1,\mathbf s}(\eta,\zeta,u)
=\mathfrak c\!\left(
32\eta\zeta,e^{\mathrm{loc}}_{1,\mathbf s}(\eta,\zeta,u),
m^{\mathrm{loc}}_{\mathbf s}(\zeta,u)\right),
\end{gathered}
\]
and
\[
\begin{aligned}
\pi^{\mathrm{mix}}_{b,\mathbf s,11}
&=e^{\mathrm{loc}}_{b,\mathbf s}m^{\mathrm{loc}}_{\mathbf s}
  +c^{\mathrm{loc}}_{b,\mathbf s},&
\pi^{\mathrm{mix}}_{b,\mathbf s,10}
&=e^{\mathrm{loc}}_{b,\mathbf s}(1-m^{\mathrm{loc}}_{\mathbf s})
  -c^{\mathrm{loc}}_{b,\mathbf s},\\
\pi^{\mathrm{mix}}_{b,\mathbf s,01}
&=(1-e^{\mathrm{loc}}_{b,\mathbf s})m^{\mathrm{loc}}_{\mathbf s}
  -c^{\mathrm{loc}}_{b,\mathbf s},&
\pi^{\mathrm{mix}}_{b,\mathbf s,00}
&=(1-e^{\mathrm{loc}}_{b,\mathbf s})(1-m^{\mathrm{loc}}_{\mathbf s})
  +c^{\mathrm{loc}}_{b,\mathbf s}.
\end{aligned}
\]
The arguments of these cells are \((\eta,\zeta,u)\).

Use \(\mathrm r\) and \(\mathrm c\) for the random and comparator fair branches, respectively. Their local risks and cells are
\[
\begin{aligned}
\mu^{\mathrm{loc}}_{0,\mathrm r}(t,\delta,u,\mathbf s)
&=\mathfrak p(t,\delta,u)+\delta Z_u(\mathbf s),&
\mu^{\mathrm{loc}}_{1,\mathrm r}
&=\mathsf G(t,\mu^{\mathrm{loc}}_{0,\mathrm r}),\\
\mu^{\mathrm{loc}}_{0,\mathrm c}(t,\delta,u,\mathbf s)
&=\mathfrak p(t,\delta,u),&
\mu^{\mathrm{loc}}_{1,\mathrm c}
&=\mathsf G(\mathsf T(t,\delta),\mu^{\mathrm{loc}}_{0,\mathrm c}),\\
\pi^{\mathrm{fair}}_{\iota,\mathbf s,ay}(t,\delta,u)
&=\frac12
\begin{cases}
\mu^{\mathrm{loc}}_{a,\iota}(t,\delta,u,\mathbf s),&y=1,\\
1-\mu^{\mathrm{loc}}_{a,\iota}(t,\delta,u,\mathbf s),&y=0,
\end{cases}
&&\iota\in\{\mathrm r,\mathrm c\}.
\end{aligned}
\]
The comparator expressions are independent of \(\mathbf s\).

The smooth-neighborhood conclusion of
\cref{obj:lem:exact-calibrations} supplies positive radii
\(r_{\mathrm{mix},\infty}\) and \(r_{\mathrm{fair},\infty}\) on whose closed parameter regions all these local formulas are smooth at every point. Set
\[
\rho=\frac12\min\{r_{\mathrm{mix},\infty},
                         r_{\mathrm{fair},\infty}\}>0.
\]
Thus, whenever \(0<\varepsilon\le\rho\), the local mixed formulas are smooth at every
\[
|\eta|<2\varepsilon,\qquad |\zeta|<2\varepsilon,\qquad u\in[0,1],
\]
and the local fair formulas are smooth at every
\[
t\in[0,1/4],\qquad |\delta|<2\varepsilon,\qquad u\in[0,1].
\]

For preliminary estimates, interpret a supplied spatial cell profile
\(\varpi=(\varpi_{ay})_{a,y\in\{0,1\}}\) through the convention
\[
\varpi^\dagger_{ay}(x)=
\begin{cases}
\varpi_{ay}(x),
&\text{if all cells are continuous, }
  0<\varpi_{ay}(x)<1,\ \sum_{a,y}\varpi_{ay}(x)=1
  \text{ for every }x,\\
1/4,&\text{otherwise}.
\end{cases}
\]
In either case the covariate law is \(\lambda\).
We will establish the first case throughout the required neighborhoods.
% lean: CausalSmith.Stat.LogoddsLowsmoothFrontier.calibrated_density_smoothness_of_valid

\item \textit{A uniform bound for amplitude derivatives.}
The local derivative bounds supplied by
\cref{obj:lem:exact-calibrations} give positive constants
\(B_{\mathrm{mix},1},B_{\mathrm{mix},2},
 B_{\mathrm{fair},1},B_{\mathrm{fair},2}\)
and positive radii \(r_{\mathrm{mix},D},r_{\mathrm{fair},D}\) such that, for \(m=1,2\),
\[
\begin{aligned}
\left\|D^m\pi^{\mathrm{mix}}_{b,\mathbf s,ay}
                   (\eta,\zeta,u)\right\|
&\le B_{\mathrm{mix},m},
&&|\eta|,|\zeta|\le r_{\mathrm{mix},D},\quad u\in[0,1],\\
\left\|D^m\pi^{\mathrm{fair}}_{\iota,\mathbf s,ay}
                   (t,\delta,u)\right\|
&\le B_{\mathrm{fair},m},
&&t\in[0,1/4],\quad |\delta|\le r_{\mathrm{fair},D},
  \quad u\in[0,1].
\end{aligned}
\]
These are total derivatives in the three displayed arguments, uniformly over the finite indices. Define
\[
r_D=\min\{r_{\mathrm{mix},D},r_{\mathrm{fair},D}\},
\qquad
M=4\bigl(B_{\mathrm{mix},1}+B_{\mathrm{mix},2}
        +B_{\mathrm{fair},1}+B_{\mathrm{fair},2}\bigr)+1.
\]
In particular, \(r_D>0\) and \(M\ge1\).

The amplitude slices have linear parts
\[
L_{\mathrm{mix}}(\eta,\zeta)=(\eta,\zeta,0),
\qquad
L_{\mathrm{fair}}(\delta)=(0,\delta,0),
\qquad
\|L_{\mathrm{mix}}\|=\|L_{\mathrm{fair}}\|=1.
\]
For a smooth function, the \(m\)th derivative of its composition with an affine map is its \(m\)th derivative evaluated on the linear part in each argument. Consequently, fixing \(u\), or fixing \(t,u\), gives
\[
\begin{aligned}
\left\|D^m_{(\eta,\zeta)}
      \bigl[4\pi^{\mathrm{mix}}_{b,\mathbf s,ay}(\eta,\zeta,u)\bigr]\right\|
&\le4B_{\mathrm{mix},m}\le M,\\
\left\|D^m_\delta
      \bigl[4\pi^{\mathrm{fair}}_{\iota,\mathbf s,ay}(t,\delta,u)\bigr]\right\|
&\le4B_{\mathrm{fair},m}\le M,
\qquad m=1,2.
\end{aligned}
\]
The factor four comes from density relative to the uniform measure on the four labels. These bounds involve the fully substituted formulas, including the comparator effect.
% lean: CausalSmith.Stat.LogoddsLowsmoothFrontier.calibrated_local_density_derivative_envelopes

\item \textit{Fair prognosis estimates.}
By \cref{obj:lem:exact-calibrations}, choose
\(\varepsilon_{\mathrm f}>0\) and \(B_{\mathrm p}>0\) such that
\[
\left|\mathfrak p(t,\delta,u)-\frac25\right|
+\left|\partial_u\mathfrak p(t,\delta,u)\right|
\le B_{\mathrm p}\delta^2
\]
for \(t\in[0,1/4]\), \(|\delta|\le\varepsilon_{\mathrm f}\), and \(u\in[0,1]\).
Put
\[
r_E=r_O
=\min\left\{\varepsilon_{\mathrm f},\frac1{100},
                    \frac1{100(B_{\mathrm p}+1)}\right\},
\qquad
r_L=\min\{r_E,r_{\mathrm{fair},\infty}\}.
\]
All three radii are positive, and
\[
B_{\mathrm p}r_E\le\frac1{100}.
\]
Hence, for \(|\delta|\le r_E\),
\[
\left|\mathfrak p(t,\delta,u)-\frac25\right|
\le B_{\mathrm p}\delta^2
\le B_{\mathrm p}r_E^2
\le\frac1{100}.
\]
Since \(|Z_u(\mathbf s)|\le2\), both fair control risks lie in
\([1/3,2/3]\): the random risk differs from \(2/5\) by at most
\(1/100+2/100\), and the comparator risk by at most \(1/100\).

We record the logit estimates used here and below. For
\(\xi_1,\xi_2\in[1/3,2/3]\), the mean value theorem applied separately to the logarithms of the risks and their complements gives
\[
|\log \xi_1-\log \xi_2|\le3|\xi_1-\xi_2|,
\qquad
|\log(1-\xi_1)-\log(1-\xi_2)|\le3|\xi_1-\xi_2|.
\]
Thus
\[
|\operatorname{logit}(\xi_1)-\operatorname{logit}(\xi_2)|
\le6|\xi_1-\xi_2|.
\]
Also, with
\[
\omega_\xi=\frac{\xi}{1-\xi},
\qquad \xi\in[1/3,2/3],
\]
we have \(\omega_\xi\in[1/2,2]\). The elementary inequalities
\[
1-\omega_\xi^{-1}\le\log\omega_\xi\le\omega_\xi-1
\]
therefore imply
\[
|\operatorname{logit}(\xi)|\le1.
\]
This proves the fair prognosis envelope on \(|\delta|\le r_E\).

For a valid fair profile, comparison with the fixed risk \(2/5\) yields
\[
\left|\nu(P,x)-\operatorname{logit}\!\left(\frac25\right)\right|
\le6(B_{\mathrm p}\delta^2+2|\delta|).
\]
The triangle inequality consequently gives
\[
|\nu(P,x)-\nu(P,z)|
\le12B_{\mathrm p}\delta^2+24|\delta|.
\]
For an invalid preliminary profile, its replacement has prognosis logit zero, so the same oscillation inequality holds.

For any auxiliary radius \(0<\varepsilon_{\mathrm{amp}}\le r_O\), any
\(\beta\ge0\), and any \(k\ge1\), set
\[
\delta_{\mathrm{amp}}
=\varepsilon_{\mathrm{amp}}k^{-\beta}.
\]
Because \(0<k^{-\beta}\le1\),
\[
0\le\delta_{\mathrm{amp}}\le\varepsilon_{\mathrm{amp}}\le r_O,
\qquad
\delta_{\mathrm{amp}}^2\le r_O|\delta_{\mathrm{amp}}|.
\]
It follows that
\[
\begin{aligned}
12B_{\mathrm p}\delta_{\mathrm{amp}}^2
   +24|\delta_{\mathrm{amp}}|
&\le(12B_{\mathrm p}r_O+24)|\delta_{\mathrm{amp}}|\\
&\le26|\delta_{\mathrm{amp}}|
\le2k^{-\beta},
\end{aligned}
\]
where the last inequality uses \(26\varepsilon_{\mathrm{amp}}\le2\).

For the linear estimate, define the within-cell coordinate by
\[
j_k(x)=\min\{\lfloor kx\rfloor,k-1\},
\qquad
u_k(x)=kx-j_k(x),
\qquad x\in[0,1].
\]
Here \(0\le j_k(x)<k\) and \(u_k(x)\in[0,1]\): below the last cell this follows from the floor inequalities, and in the last cell it follows from
\(k-1\le kx\le k\). In particular this convention includes \(x=1\) and \(k=1\).

Within a cell, the elementary sine and cosine difference bounds give
\[
|Z_u(\mathbf s)-Z_v(\mathbf s)|
\le\pi|u-v|\le4|u-v|,
\qquad u,v\in[0,1].
\]
Adjacent profiles agree at the shared sign. Subdividing a segment at its cell boundaries therefore gives
\[
|\mathsf Z_{k,\sigma}(x)-\mathsf Z_{k,\sigma}(z)|
\le4k|x-z|.
\]
Similarly, smoothness of the fair root, its derivative bound, and
\(\mathfrak p(t,\delta,0)=\mathfrak p(t,\delta,1)=2/5\) give
\[
|\mathfrak p(t,\delta,u_k(x))
 -\mathfrak p(t,\delta,u_k(z))|
\le B_{\mathrm p}\delta^2k|x-z|.
\]
Thus each fair control risk satisfies
\[
|\mu_0(x)-\mu_0(z)|
\le(B_{\mathrm p}\delta^2+4|\delta|)k|x-z|.
\]
Applying the logit difference bound, or using the zero logit in the replacement case, proves
\[
|\nu(P,x)-\nu(P,z)|
\le6(B_{\mathrm p}\delta^2+4|\delta|)k|x-z|.
\]
For \(0<\varepsilon_{\mathrm{amp}}\le r_L\),
\[
6(B_{\mathrm p}\delta_{\mathrm{amp}}^2
   +4|\delta_{\mathrm{amp}}|)
\le25|\delta_{\mathrm{amp}}|,
\qquad
25\varepsilon_{\mathrm{amp}}\le1.
\]
Since \(k^{1-\beta}=k^{-\beta}k\), we obtain
\[
|\nu(P,x)-\nu(P,z)|
\le k^{1-\beta}|x-z|.
\]
These estimates apply to both fair branches uniformly over
\(t\in[0,1/4]\).
% lean: CausalSmith.Stat.LogoddsLowsmoothFrontier.fairCells_uniform_prognosis_envelope
% lean: CausalSmith.Stat.LogoddsLowsmoothFrontier.fairCells_uniform_prognosis_oscillation
% lean: CausalSmith.Stat.LogoddsLowsmoothFrontier.fairCells_uniform_prognosis_linear

\item \textit{Mixed propensity and alternative-prognosis estimates.}
The mixed-root smoothness and derivative bounds in
\cref{obj:lem:exact-calibrations} supply a positive radius \(r_{\mathrm q}\) and a constant \(B_{\mathrm q}>0\) such that
\[
|\partial_u\mathfrak q(\eta,\zeta,u)|\le B_{\mathrm q}
\]
for \(|\eta|,|\zeta|\le r_{\mathrm q}\), \(u\in[0,1]\).
Indeed, the spatial derivative is the total first derivative evaluated on the unit vector in the spatial coordinate.

The endpoint identity for \(\mathfrak q\), followed by the mean value theorem within each cell and subdivision across cells, gives
\[
|\mathfrak q(\eta,\zeta,u_k(x))
 -\mathfrak q(\eta,\zeta,u_k(z))|
\le B_{\mathrm q}k|x-z|.
\]
Together with \(|\mathfrak q|\le5/4\), \(|\mathsf Z_{k,\sigma}|\le2\), and the field estimate above, the product decomposition gives
\[
\begin{aligned}
&|\mathfrak q(\eta,\zeta,u_k(x))\mathsf Z_{k,\sigma}(x)
 -\mathfrak q(\eta,\zeta,u_k(z))\mathsf Z_{k,\sigma}(z)|\\
&\quad\le
\frac54|\mathsf Z_{k,\sigma}(x)-\mathsf Z_{k,\sigma}(z)|
+2|\mathfrak q(\eta,\zeta,u_k(x))
      -\mathfrak q(\eta,\zeta,u_k(z))|\\
&\quad\le(5+2B_{\mathrm q})k|x-z|.
\end{aligned}
\]
For \(|\eta|,|\zeta|\le1/100\), the mixed margins satisfy
\[
|e_b(x)-1/2|\le\frac52|\eta|\le\frac1{40},
\qquad
|m_\sigma(x)-1/2|\le2|\zeta|\le\frac1{40}.
\]
Their propensity margins consequently belong to \([1/3,2/3]\). For a valid profile, summing the treated cells identifies its propensity with \(e_b\), and hence
\[
|g(P_{b,\sigma},x)-g(P_{b,\sigma},z)|
\le6|\eta|(5+2B_{\mathrm q})k|x-z|.
\]
The replacement profile has propensity logit zero, giving the same bound. Define
\[
r_P=\min\left\{r_{\mathrm q},\frac1{100},
                  \frac1{6(5+2B_{\mathrm q})}\right\}>0.
\]
For
\[
\eta_{\mathrm{amp}}
=\varepsilon_{\mathrm{amp}}k^{-\alpha},
\qquad
\zeta_{\mathrm{amp}}
=\varepsilon_{\mathrm{amp}}k^{-\beta},
\]
with \(0<\varepsilon_{\mathrm{amp}}\le r_P\) and
\(\alpha,\beta\ge0\), we therefore have
\[
|g(P_{b,\sigma},x)-g(P_{b,\sigma},z)|
\le k^{1-\alpha}|x-z|.
\]

For the alternative prognosis, consider the control-logit map
\[
\mathcal F_{\mathrm{ctrl}}(\vartheta,\xi,\upsilon)
=\operatorname{logit}\!\left(
\upsilon-\frac{\mathfrak c(\vartheta,\xi,\upsilon)}{1-\xi}
\right).
\]
It is smooth near \((0,1/2,1/2)\): the covariance branch is smooth there, the denominator equals \(1/2\) at the center, and the argument of the logit equals \(1/2\). A \(C^1\) function is locally Lipschitz. Choose a neighborhood
\(\mathcal U_{\mathrm{ctrl}}\), a Lipschitz constant
\(K_{\mathrm{ctrl}}\ge0\), and \(h_{\mathrm{ctrl}}>0\) such that
\[
\left\{(\vartheta,\xi,\upsilon):
\max\{|\vartheta|,|\xi-1/2|,|\upsilon-1/2|\}
<h_{\mathrm{ctrl}}\right\}
\subseteq\mathcal U_{\mathrm{ctrl}}.
\]
Set
\[
r_{\mathrm{ctrl}}=h_{\mathrm{ctrl}}/2,
\qquad C_{\mathrm{ctrl}}=K_{\mathrm{ctrl}}+1.
\]
Comparing points with the same first coordinate gives
\[
\begin{aligned}
&|\mathcal F_{\mathrm{ctrl}}(\vartheta,\xi_1,\upsilon_1)
 -\mathcal F_{\mathrm{ctrl}}(\vartheta,\xi_2,\upsilon_2)|\\
&\qquad\le C_{\mathrm{ctrl}}
       (|\xi_1-\xi_2|+|\upsilon_1-\upsilon_2|)
\end{aligned}
\]
whenever
\[
|\vartheta|,\ |\xi_1-1/2|,\ |\xi_2-1/2|,
\ |\upsilon_1-1/2|,\ |\upsilon_2-1/2|
\le r_{\mathrm{ctrl}}.
\]

Define
\[
r_N=\min\left\{r_{\mathrm q},\frac1{100},
 \frac{r_{\mathrm{ctrl}}}{3},
 \frac1{C_{\mathrm{ctrl}}(9+2B_{\mathrm q})}\right\}>0.
\]
For \((\alpha,\beta)\in\mathcal D\), we have
\(0<\beta<\alpha\le1\). At the rate-scaled amplitudes with
\(0<\varepsilon_{\mathrm{amp}}\le r_N\),
\[
0\le\eta_{\mathrm{amp}}\le\zeta_{\mathrm{amp}}
\le r_N.
\]
Moreover,
\[
32|\eta_{\mathrm{amp}}\zeta_{\mathrm{amp}}|
\le\frac{32}{100}r_N\le r_{\mathrm{ctrl}},
\qquad
\frac52|\eta_{\mathrm{amp}}|\le r_{\mathrm{ctrl}},
\qquad
2|\zeta_{\mathrm{amp}}|\le r_{\mathrm{ctrl}}.
\]
The control-logit comparison therefore applies to the alternative margins. The two margin differences obey
\[
\begin{aligned}
|e_1(x)-e_1(z)|
&\le|\eta_{\mathrm{amp}}|(5+2B_{\mathrm q})k|x-z|,\\
|m_\sigma(x)-m_\sigma(z)|
&\le4|\zeta_{\mathrm{amp}}|k|x-z|.
\end{aligned}
\]
For a valid table, its control risk is
\[
\mu_0(x)=m_\sigma(x)-\frac{c_1(x)}{1-e_1(x)}.
\]
Using \(\eta_{\mathrm{amp}}\le\zeta_{\mathrm{amp}}\) gives
\[
\begin{aligned}
|\nu(P_{1,\sigma},x)-\nu(P_{1,\sigma},z)|
&\le C_{\mathrm{ctrl}}(9+2B_{\mathrm q})
       \zeta_{\mathrm{amp}}k|x-z|\\
&\le k^{1-\beta}|x-z|.
\end{aligned}
\]
The replacement case again has zero prognosis logit.
% lean: CausalSmith.Stat.LogoddsLowsmoothFrontier.mixedCells_uniform_propensity_linear
% lean: CausalSmith.Stat.LogoddsLowsmoothFrontier.mixedCells_uniform_alternative_prognosis_linear

\item \textit{One radius and the remaining mixed spatial bounds.}
Choose
\[
R_{\mathrm{support}}
=\min\left\{\varepsilon_{\mathrm{cal}},\rho,r_D,r_E,r_O,r_L,
                         r_P,r_N,\frac1{100}\right\},
\qquad
\varepsilon=\frac12R_{\mathrm{support}}.
\]
Every entry in the minimum is positive. Therefore
\[
0<\varepsilon\le\frac1{100},
\qquad
2\varepsilon\le\varepsilon_{\mathrm{cal}},
\qquad
2\varepsilon\le\frac1{100},
\qquad
\varepsilon\le\rho,r_D,r_E,r_O,r_L,r_P,r_N.
\]

We supply the remaining mixed estimates at
\[
\eta=\varepsilon k^{-\alpha},
\qquad
\zeta=\varepsilon k^{-\beta}.
\]
For every nonnegative exponent \(\gamma\),
\[
0<k^{-\gamma}\le1,
\qquad
|\varepsilon k^{-\gamma}|=\varepsilon k^{-\gamma}\le\varepsilon.
\]
Thus \(|\eta|,|\zeta|\le1/100\).

For the null branch the control risk equals \(m_\sigma\) whenever the profile is valid. The logit and field bounds give
\[
|\nu(P_{0,\sigma},x)-\nu(P_{0,\sigma},z)|
\le24|\zeta|k|x-z|
\le k^{1-\beta}|x-z|,
\]
because \(24\varepsilon\le1\). In the replacement case the left side is zero.

For oscillation, comparison of the propensity margin with \(1/2\) gives
\[
|g(P_{b,\sigma},x)|
\le6|e_b(x)-1/2|\le15|\eta|,
\]
and hence
\[
|g(P_{b,\sigma},x)-g(P_{b,\sigma},z)|
\le30|\eta|
\le2k^{-\alpha}.
\]

We also need a covariance estimate for the mixed prognosis. For
\[
|\vartheta|,\ |\xi-1/2|,\ |\upsilon-1/2|\le\frac18,
\]
define
\[
\begin{gathered}
d_\vartheta=e^\vartheta-1,\qquad
v_{\xi,\upsilon}=\xi(1-\xi)\upsilon(1-\upsilon),\\
L_{\vartheta,\xi,\upsilon}
=1+d_\vartheta\{\xi(1-\upsilon)+(1-\xi)\upsilon\},\\
S_{\vartheta,\xi,\upsilon}
=\sqrt{L_{\vartheta,\xi,\upsilon}^2
              -4d_\vartheta^2v_{\xi,\upsilon}},\qquad
D_{\vartheta,\xi,\upsilon}
=L_{\vartheta,\xi,\upsilon}+S_{\vartheta,\xi,\upsilon}.
\end{gathered}
\]
The elementary exponential bound and the margin bounds give
\[
|d_\vartheta|\le2|\vartheta|\le\frac14,\qquad
0\le v_{\xi,\upsilon}\le\frac1{16},\qquad
0\le\xi(1-\upsilon)+(1-\xi)\upsilon\le1.
\]
Consequently,
\[
L_{\vartheta,\xi,\upsilon}\ge\frac34,\qquad
L_{\vartheta,\xi,\upsilon}^2
       -4d_\vartheta^2v_{\xi,\upsilon}\ge0,
\qquad
D_{\vartheta,\xi,\upsilon}\ge\frac34.
\]
The rationalized covariance formula now yields
\[
|\mathfrak c(\vartheta,\xi,\upsilon)|
D_{\vartheta,\xi,\upsilon}
=2|d_\vartheta|v_{\xi,\upsilon},
\]
so
\[
|\mathfrak c(\vartheta,\xi,\upsilon)|
\le\frac1{24},
\qquad
|\mathfrak c(\vartheta,\xi,\upsilon)|
\le|\vartheta|.
\]
The first inequality uses \(|d_\vartheta|\le1/4\); the second uses
\(|d_\vartheta|\le2|\vartheta|\).

For the mixed alternative,
\[
|\vartheta|=|32\eta\zeta|\le\frac18,
\]
and its margins lie within \(1/40\) of \(1/2\). Thus, for either branch,
\[
|c_b(x)|\le32|\eta||\zeta|,
\qquad |c_b(x)|\le\frac1{24}.
\]
Summing the control cells gives \(1-e_b\), so their risk satisfies
\[
\mu_0(x)=m_\sigma(x)-\frac{c_b(x)}{1-e_b(x)}.
\]
The bounds
\[
\frac{19}{40}\le e_b,m_\sigma\le\frac{21}{40},
\qquad |c_b|\le\frac1{24}
\]
imply
\[
\frac13(1-e_b)\le(1-e_b)m_\sigma-c_b
                 \le\frac23(1-e_b),
\]
by multiplication and the displayed rational bounds. Hence
\(\mu_0(x)\in[1/3,2/3]\). Since \(1-e_b\ge19/40>1/3\),
\[
|\mu_0(x)-1/2|
\le|m_\sigma(x)-1/2|+3|c_b(x)|
\le2|\zeta|+96|\eta||\zeta|.
\]
Applying the logit comparison with \(1/2\), and using \(|\eta|\le1/100\), gives
\[
|\nu(P_{b,\sigma},x)|
\le6(2|\zeta|+96|\eta||\zeta|)
\le18|\zeta|.
\]
Therefore
\[
|\nu(P_{b,\sigma},x)-\nu(P_{b,\sigma},z)|
\le36|\zeta|\le2k^{-\beta}.
\]
All these bounds also hold for the zero logits of a replacement profile.

Combining these estimates with the preceding linear estimates proves, for both mixed branches,
\[
\begin{aligned}
|g(P_{b,\sigma},x)-g(P_{b,\sigma},z)|
&\le k^{1-\alpha}|x-z|,&
|g(P_{b,\sigma},x)-g(P_{b,\sigma},z)|
&\le2k^{-\alpha},\\
|\nu(P_{b,\sigma},x)-\nu(P_{b,\sigma},z)|
&\le k^{1-\beta}|x-z|,&
|\nu(P_{b,\sigma},x)-\nu(P_{b,\sigma},z)|
&\le2k^{-\beta}.
\end{aligned}
\]
% lean: CausalSmith.Stat.LogoddsLowsmoothFrontier.mixedCells_uniform_null_prognosis_linear
% lean: CausalSmith.Stat.LogoddsLowsmoothFrontier.mixedCells_uniform_propensity_oscillation
% lean: CausalSmith.Stat.LogoddsLowsmoothFrontier.mixedCells_uniform_prognosis_oscillation

\item \textit{Continuity and validity on the doubled boxes.}
Fix \(k\ge1\) and \(\sigma\). On cell \(j\), use the local pair
\((\sigma_j,\sigma_{j+1})\) and coordinate \(u=kx-j\).
For
\[
|\eta|,|\zeta|\le2\varepsilon,
\qquad
t\in[0,1/4],\quad |\delta|\le2\varepsilon,
\]
the local cell formulas are continuous in \(u\), by
\cref{obj:lem:exact-calibrations} and
\(2\varepsilon\le\varepsilon_{\mathrm{cal}}\).

At a shared endpoint, the two sign fields take the same value:
\[
Z_1(\sigma_j,\sigma_{j+1})
=\sigma_{j+1}
=Z_0(\sigma_{j+1},\sigma_{j+2}).
\]
The endpoint identities
\[
\mathfrak q(\eta,\zeta,1)=\mathfrak q(\eta,\zeta,0),
\qquad
\mathfrak p(t,\delta,1)=\mathfrak p(t,\delta,0)=\frac25
\]
therefore make every mixed and fair cell agree across adjoining cells. Gluing the finitely many continuous cell profiles proves continuity on \([0,1]\), including its endpoints. When \(k=1\), there is one continuous profile and no interior boundary to glue.

Because \(2\varepsilon\le1/100\), the raw mixed cells satisfy
\[
\frac{19}{40}\le e_b,m_\sigma,1-e_b,1-m_\sigma\le\frac{21}{40},
\qquad |c_b|\le\frac1{24}.
\]
Each product entering a mixed cell lies in
\([361/1600,441/1600]\). Accordingly,
\[
4\left(\frac{361}{1600}-\frac1{24}\right)
\le4\pi^{\mathrm{mix}}_{b,\mathbf s,ay}
\le4\left(\frac{441}{1600}+\frac1{24}\right),
\]
which implies
\[
\frac12\le4\pi^{\mathrm{mix}}_{b,\mathbf s,ay}\le2.
\]

For the fair cells, the selected root belongs to \((3/10,1/2)\); this includes its default value \(2/5\). Since
\[
|\delta Z_u(\mathbf s)|\le\frac1{50},
\]
both fair control risks belong to \([1/4,3/5]\).
For \(t\in[0,1/4]\), the elementary exponential estimate gives
\[
0\le e^t-1\le\frac43t\le\frac13.
\]
The comparator effect lies between \(t/2\) and \(t\), by
\cref{obj:lem:exact-calibrations}. Thus the effect used in either fair branch lies in \([0,1/4]\). For
\(\vartheta\in[0,1/4]\) and \(\xi\in[1/4,3/5]\),
\[
\mathsf G(\vartheta,\xi)
=\frac{e^\vartheta\xi}{1+(e^\vartheta-1)\xi},
\qquad
\frac14\le\mathsf G(\vartheta,\xi)\le\frac34.
\]
To check the last bounds, the denominator is positive; multiplication by it reduces the lower bound to
\[
(1+e^\vartheta-1)\xi
\ge\frac14\{1+(e^\vartheta-1)\xi\},
\]
using \(\xi\ge1/4\), and the upper bound follows from
\(\xi\le3/5\) and \(e^\vartheta-1\le1/3\).
Consequently both fair arm risks and their complements are at least \(1/4\), and
\[
\frac12\le4\pi^{\mathrm{fair}}_{\iota,\mathbf s,ay}\le2.
\]

Finally, direct addition of either four-cell formula gives
\[
\sum_{a,y}\pi^{\mathrm{mix}}_{b,\mathbf s,ay}=1,
\qquad
\sum_{a,y}\pi^{\mathrm{fair}}_{\iota,\mathbf s,ay}=1.
\]
Every cell lies in \([1/8,1/2]\), so the cells are strictly interior. Together with continuity, this proves validity throughout both doubled closed boxes. The constructed laws therefore retain their supplied cells there.
% lean: CausalSmith.Stat.LogoddsLowsmoothFrontier.mixedCells_valid_of_continuous
% lean: CausalSmith.Stat.LogoddsLowsmoothFrontier.fairCells_valid_of_continuous

\item \textit{Transfer to actual densities and common open neighborhoods.}
Let \(\tau\) be the uniform probability measure on the four treatment--outcome labels:
\[
\tau(\{(a,y)\})=\frac14,\qquad a,y\in\{0,1\}.
\]
Define the actual one-record densities, relative to \(\lambda\otimes\tau\), by
\[
\begin{aligned}
\mathsf d^{\mathrm{mix}}_{b,k,\sigma,a,y,x}(\eta,\zeta)
&=4P_{b,\sigma}(A=a,Y=y\mid X=x),\\
\mathsf d^{\mathrm{fair}}_{\mathrm r,k,\sigma,t,a,y,x}(\delta)
&=4P_{t,\sigma}(A=a,Y=y\mid X=x),\\
\mathsf d^{\mathrm{fair}}_{\mathrm c,k,\sigma,t,a,y,x}(\delta)
&=4P_{*,t,\delta}(A=a,Y=y\mid X=x).
\end{aligned}
\]
The amplitudes of the laws on the right are those on the left.

Fix \(x\), and define its local sign pair by
\[
\mathbf s_{k,\sigma}(x)
=(\sigma_{j_k(x)},\sigma_{j_k(x)+1}).
\]
Validity on the doubled boxes identifies the actual densities there with the local slices:
\[
\begin{aligned}
\mathsf d^{\mathrm{mix}}_{b,k,\sigma,a,y,x}(\eta,\zeta)
&=4\pi^{\mathrm{mix}}_{b,\mathbf s_{k,\sigma}(x),ay}
                    (\eta,\zeta,u_k(x)),\\
\mathsf d^{\mathrm{fair}}_{\iota,k,\sigma,t,a,y,x}(\delta)
&=4\pi^{\mathrm{fair}}_{\iota,\mathbf s_{k,\sigma}(x),ay}
                    (t,\delta,u_k(x)).
\end{aligned}
\]
For the comparator the chosen pair is immaterial.

At a point of the smaller closed box, continuity of absolute value supplies a neighborhood lying inside the doubled box. Thus these identities hold on a neighborhood of every such point, and differentiation transfers the local bounds to the actual densities. Since \(\varepsilon\le r_D\),
\[
\begin{aligned}
\left\|D^m\mathsf d^{\mathrm{mix}}_{b,k,\sigma,a,y,x}
                            (\eta,\zeta)\right\|&\le M,
&&|\eta|,|\zeta|\le\varepsilon,\\
\left\|D^m\mathsf d^{\mathrm{fair}}_{\iota,k,\sigma,t,a,y,x}
                            (\delta)\right\|&\le M,
&&|\delta|\le\varepsilon,\quad t\in[0,1/4],
\end{aligned}
\qquad m=1,2.
\]

Moreover, set
\[
U_{\mathrm{mix}}
=\{(\eta,\zeta)\in\mathbb R^2:
                   |\eta|<2\varepsilon,\ |\zeta|<2\varepsilon\},
\qquad
U_{\mathrm{fair}}
=(-2\varepsilon,2\varepsilon).
\]
These sets are open and contain the required closed boxes. At every point of them, nearby validity identifies the actual density with a smooth local slice. Because \(\varepsilon\le\rho\), every mixed density is \(C^\infty\) on \(U_{\mathrm{mix}}\), and every fair density is \(C^\infty\) on \(U_{\mathrm{fair}}\), uniformly over the indicated indices and over \(t\in[0,1/4]\). In particular, they are \(C^2\) at every point of the smaller closed boxes.

For clarity, with
\[
\mathbf e_\eta=(1,0),\qquad \mathbf e_\zeta=(0,1),
\]
the total second-derivative bound includes
\[
\left|\partial_\eta\partial_\zeta
       \mathsf d^{\mathrm{mix}}(\eta,\zeta)\right|
=\left|D^2\mathsf d^{\mathrm{mix}}(\eta,\zeta)
                  [\mathbf e_\eta,\mathbf e_\zeta]\right|
\le M.
\]
It also applies to both first and second \(\delta\)-derivatives of the deterministic comparator.
% lean: CausalSmith.Stat.LogoddsLowsmoothFrontier.calibrated_density_derivative_envelopes_of_valid
% lean: CausalSmith.Stat.LogoddsLowsmoothFrontier.calibrated_density_smoothness_of_valid

\item \textit{Mixed model membership and effects.}
Fix \((\alpha,\beta)\in\mathcal D\), \(k\ge1\), and the rate-scaled amplitudes
\[
\eta=\varepsilon k^{-\alpha},\qquad
\zeta=\varepsilon k^{-\beta}.
\]
They lie in the valid smaller box. The constructed covariate marginal is uniform: for every Borel set \(\mathcal B_X\subseteq[0,1]\),
\[
P(X\in\mathcal B_X)
=\int_{\mathcal B_X}\sum_{a,y}p_{ay}(x)\,d\lambda(x)
=\lambda(\mathcal B_X).
\]

We identify the homogeneous effects from the positive cell odds. For the covariance formula introduced above, the identity
\[
S_{\vartheta,\xi,\upsilon}^2
=L_{\vartheta,\xi,\upsilon}^2
 -4d_\vartheta^2v_{\xi,\upsilon}
\]
and
\[
\mathfrak c(\vartheta,\xi,\upsilon)
=\frac{2d_\vartheta v_{\xi,\upsilon}}
       {D_{\vartheta,\xi,\upsilon}}
\]
give, by multiplication by the positive denominator,
\[
d_\vartheta\mathfrak c(\vartheta,\xi,\upsilon)^2
-L_{\vartheta,\xi,\upsilon}\mathfrak c(\vartheta,\xi,\upsilon)
+d_\vartheta v_{\xi,\upsilon}=0.
\]
For a table with margins \(\xi,\upsilon\) and covariance
\(\mathfrak c(\vartheta,\xi,\upsilon)\), expansion of its four cells then gives
\[
p_{11}p_{00}-e^\vartheta p_{10}p_{01}
=-d_\vartheta\mathfrak c(\vartheta,\xi,\upsilon)^2
 +L_{\vartheta,\xi,\upsilon}\mathfrak c(\vartheta,\xi,\upsilon)
 -d_\vartheta v_{\xi,\upsilon}=0.
\]
This computation includes \(\vartheta=0\). The null table, with covariance zero, directly satisfies
\(p_{11}p_{00}=p_{10}p_{01}\). Therefore the mixed branches satisfy
\[
\frac{p_{11}p_{00}}{p_{10}p_{01}}
=
\begin{cases}
1,&b=0,\\
e^{32\eta\zeta},&b=1.
\end{cases}
\]
All four cells are positive. Dividing by the positive treatment margins gives
\[
\mu_0=\frac{p_{01}}{p_{00}+p_{01}},
\qquad
\mu_1=\frac{p_{11}}{p_{10}+p_{11}},
\qquad
\frac{\mu_1/(1-\mu_1)}{\mu_0/(1-\mu_0)}
=\frac{p_{11}p_{00}}{p_{10}p_{01}}.
\]
Taking logarithms, and evaluating at covariate zero to identify the scalar effect, proves
\[
\theta(P_{0,\sigma})=0,\qquad
\theta(P_{1,\sigma})=32\eta\zeta,
\qquad
\mu_1(P_{b,\sigma},x)
=\ell\{\nu(P_{b,\sigma},x)+\theta(P_{b,\sigma})\}.
\]
The effect envelope follows from
\[
|32\eta\zeta|\le32\varepsilon^2\le\frac12.
\]
The propensity margin and the control risk lie in \([1/3,2/3]\), as established above, so
\[
\|g(P_{b,\sigma},\cdot)\|_\infty\le1,
\qquad
\|\nu(P_{b,\sigma},\cdot)\|_\infty\le1.
\]

To convert the spatial bounds into the required Hölder bounds, let
\(f:[0,1]\to\mathbb R\), \(0<\gamma\le1\), satisfy
\[
|f(x)-f(z)|\le k^{1-\gamma}|x-z|,
\qquad
|f(x)-f(z)|\le2k^{-\gamma}.
\]
Set
\[
d_{xz}=|x-z|\ge0.
\]
If \(kd_{xz}\le1\), then
\[
|f(x)-f(z)|
\le k^{-\gamma}(kd_{xz})
\le k^{-\gamma}(kd_{xz})^\gamma
=d_{xz}^\gamma.
\]
If \(kd_{xz}\ge1\), then
\[
|f(x)-f(z)|
\le2k^{-\gamma}
\le2k^{-\gamma}(kd_{xz})^\gamma
=2d_{xz}^\gamma.
\]
The first case includes \(x=z\), and both cases include \(\gamma=1\). Thus
\[
[f]_\gamma\le2.
\]
Apply this argument to \(g\) with \(\gamma=\alpha\) and to \(\nu\) with
\(\gamma=\beta\). Together with uniform design, homogeneity, and the three envelopes, these are precisely the conditions of
\cref{obj:def:model}. Hence both mixed laws belong to
\(\mathcal M(\alpha,\beta)\).
% lean: CausalSmith.Stat.LogoddsLowsmoothFrontier.mixedCells_model_bounds_of_spatial
% lean: CausalSmith.Stat.LogoddsLowsmoothFrontier.mixedCells_homogeneous_effect
% lean: CausalSmith.Stat.LogoddsLowsmoothFrontier.calibrated_model_of_bounds

\item \textit{Fair model membership and effects.}
Set
\[
\delta=\varepsilon k^{-\beta},
\qquad t\in[0,1/4].
\]
Both fair branches are valid. Their treatment probability is exactly \(1/2\), because the two cells in either treatment arm sum to \(1/2\). Therefore
\[
g(P_{t,\sigma},x)=g(P_{*,t,\delta},x)=0.
\]
This supplies the propensity envelope and both propensity spatial bounds.

The choices \(\varepsilon\le r_E,r_O,r_L\) and the fair estimates above give, for both branches,
\[
\|\nu(P,\cdot)\|_\infty\le1,\qquad
|\nu(P,x)-\nu(P,z)|\le k^{1-\beta}|x-z|,
\qquad
|\nu(P,x)-\nu(P,z)|\le2k^{-\beta}.
\]
The two-scale calculation in the preceding step gives
\[
[\nu(P,\cdot)]_\beta\le2.
\]

For any interior risk \(\xi\in(0,1)\), the positive denominator in the shift formula yields
\[
\mathsf G(\vartheta,\xi)(1-\xi)
=e^\vartheta\{1-\mathsf G(\vartheta,\xi)\}\xi.
\]
Applying this identity to the control risks of the two fair branches gives their positive cell cross-odds relations, and hence
\[
\theta(P_{t,\sigma})=t,\qquad
\theta(P_{*,t,\delta})=\mathsf T(t,\delta),
\qquad
\mu_1(P,x)=\ell\{\nu(P,x)+\theta(P)\}.
\]
By \cref{obj:lem:exact-calibrations},
\[
0\le\frac t2\le\mathsf T(t,\delta)\le t\le\frac14.
\]
Thus both effects satisfy the envelope \( |\theta(P)|\le1/2 \).
Uniform design follows by the same cell-normalization integral as above.
All conditions of \cref{obj:def:model} now hold for every random fair law and for its deterministic comparator.
% lean: CausalSmith.Stat.LogoddsLowsmoothFrontier.fairCells_model_bounds_of_prognosis
% lean: CausalSmith.Stat.LogoddsLowsmoothFrontier.fairCells_homogeneous_effect
% lean: CausalSmith.Stat.LogoddsLowsmoothFrontier.calibrated_model_of_bounds

\item \textit{Singleton identities at the rate-scaled amplitudes.}
Fix \(x\in[0,1]\), and put
\[
u=u_k(x),\qquad j=j_k(x).
\]
The indices \(j\) and \(j+1\) are distinct members of \(\{0,\ldots,k\}\), including when \(k=1\). For any function \(\Phi\) of their two signs, counting the \(2^{k-1}\) extensions of each pair gives
\[
2^{-(k+1)}\sum_{\sigma\in\{-1,1\}^{k+1}}
                 \Phi(\sigma_j,\sigma_{j+1})
=\frac14\sum_{s_1,s_2\in\{-1,1\}}\Phi(s_1,s_2).
\]
Consequently the full sign average of each cell is its local two-sign average.

Write \(E_{\mathbf s}\) for this uniform two-sign average:
\[
E_{\mathbf s}\Phi(\mathbf s)
=\frac14\sum_{\mathbf s\in\{-1,1\}^2}\Phi(\mathbf s).
\]
Symmetry and the trigonometric identity give
\[
E_{\mathbf s}Z_u=0,\qquad
E_{\mathbf s}Z_u^2
=\cos^2(\pi u/2)+\sin^2(\pi u/2)=1.
\]
The mixed calibration identity in
\cref{obj:lem:exact-calibrations} supplies
\[
E_{\mathbf s}\mathfrak c\!\left(
32\eta\zeta,\frac12-\eta\mathfrak q(\eta,\zeta,u)Z_u,
                 \frac12+\zeta Z_u\right)
=2\eta\zeta\mathfrak q(\eta,\zeta,u).
\]
Expanding the four local cells, using these moments and this identity, yields
\[
\begin{aligned}
E_{\mathbf s}\pi^{\mathrm{mix}}_{0,\mathbf s,11}
=E_{\mathbf s}\pi^{\mathrm{mix}}_{1,\mathbf s,11}
&=\frac14+\eta\zeta\mathfrak q(\eta,\zeta,u),\\
E_{\mathbf s}\pi^{\mathrm{mix}}_{0,\mathbf s,00}
=E_{\mathbf s}\pi^{\mathrm{mix}}_{1,\mathbf s,00}
&=\frac14+\eta\zeta\mathfrak q(\eta,\zeta,u),\\
E_{\mathbf s}\pi^{\mathrm{mix}}_{0,\mathbf s,10}
=E_{\mathbf s}\pi^{\mathrm{mix}}_{1,\mathbf s,10}
&=\frac14-\eta\zeta\mathfrak q(\eta,\zeta,u),\\
E_{\mathbf s}\pi^{\mathrm{mix}}_{0,\mathbf s,01}
=E_{\mathbf s}\pi^{\mathrm{mix}}_{1,\mathbf s,01}
&=\frac14-\eta\zeta\mathfrak q(\eta,\zeta,u).
\end{aligned}
\]
The cells of the actual laws equal these supplied cells by validity. Multiplying the equal full sign averages by \(2^{k+1}\) proves the asserted mixed sum identities.

For the fair family,
\cref{obj:lem:exact-calibrations} gives
\[
E_{\mathbf s}\mathsf G\!\left(
 t,\mathfrak p(t,\delta,u)+\delta Z_u\right)
=\mathsf G\!\left(\mathsf T(t,\delta),
                       \mathfrak p(t,\delta,u)\right).
\]
The control-risk average is \(\mathfrak p(t,\delta,u)\), by
\(E_{\mathbf s}Z_u=0\). Since treatment is fair, the four averaged cells are therefore
\[
\begin{aligned}
E_{\mathbf s}\pi^{\mathrm{fair}}_{\mathrm r,\mathbf s,01}
&=\frac12\mathfrak p(t,\delta,u),&
E_{\mathbf s}\pi^{\mathrm{fair}}_{\mathrm r,\mathbf s,00}
&=\frac12\{1-\mathfrak p(t,\delta,u)\},\\
E_{\mathbf s}\pi^{\mathrm{fair}}_{\mathrm r,\mathbf s,11}
&=\frac12\mathsf G(\mathsf T(t,\delta),\mathfrak p(t,\delta,u)),&
E_{\mathbf s}\pi^{\mathrm{fair}}_{\mathrm r,\mathbf s,10}
&=\frac12\{1-\mathsf G(\mathsf T(t,\delta),\mathfrak p(t,\delta,u))\}.
\end{aligned}
\]
These are exactly the comparator cells. Counting sign extensions and using validity proves the required full uniform-average identity. The calibration identities invoked here apply also at \(t=0\), at zero amplitudes, and at \(u=0,1\).
% lean: CausalSmith.Stat.LogoddsLowsmoothFrontier.mixedLaw_singleton_matching
% lean: CausalSmith.Stat.LogoddsLowsmoothFrontier.fairLaw_singleton_matching

\item \textit{Signed matching and completion.}
Now allow arbitrary signed amplitudes satisfying
\[
|\eta|\le\varepsilon,\qquad
|\zeta|\le\varepsilon,\qquad
|\delta|\le\varepsilon.
\]
They lie in the calibration region because
\(\varepsilon\le\varepsilon_{\mathrm{cal}}\), and their tables are valid by the doubled-box argument. The preceding local moment calculation and the matching identities of
\cref{obj:lem:exact-calibrations} apply on this whole signed region. Thus, for every \(a,y,x\),
\[
\sum_\sigma P_{0,\sigma}(A=a,Y=y\mid X=x)
=\sum_\sigma P_{1,\sigma}(A=a,Y=y\mid X=x),
\]
and, for every \(t\in[0,1/4]\),
\[
2^{-(k+1)}
\sum_\sigma P_{t,\sigma}(A=a,Y=y\mid X=x)
=P_{*,t,\delta}(A=a,Y=y\mid X=x).
\]

Finally, the raw cell bounds established on the doubled boxes transfer directly to the actual densities on the smaller boxes:
\[
\frac12\le
\mathsf d^{\mathrm{mix}}_{b,k,\sigma,a,y,x}(\eta,\zeta)
\le2,
\qquad
\frac12\le
\mathsf d^{\mathrm{fair}}_{\iota,k,\sigma,t,a,y,x}(\delta)
\le2.
\]
Together with the first- and second-derivative bounds, the common smooth open neighborhoods, and the rate-scaled membership and effect calculations, these prove every assertion. The radius \(\varepsilon\) and bound \(M\) were selected from absolute local calibration constants; their choices are independent of \(k\), the public exponents, the sign vector, the labels, the covariate, and the fair effect \(t\).
% lean: CausalSmith.Stat.LogoddsLowsmoothFrontier.SingletonNeighbourhood
% lean: CausalSmith.Stat.LogoddsLowsmoothFrontier.calibrated_support
\end{enumerate}
\end{proof}

Membership in \cref{obj:def:model} certifies the prescribed envelopes and Hölder restrictions for the native propensity logit \leanref{sym:g}{\ensuremath{g(P,\cdot)}} and prognosis logit \(\nu(P,\cdot)\), together with the homogeneous logistic relation. This certification applies to every sign realization at the stated amplitude scales and throughout the prescribed domain of public exponents. Continuity across partition endpoints and the homogeneous scalar coefficient therefore hold for each law entering the finite priors in \cref{obj:def:continuous-calibrated-priors}.

Exact singleton matching concerns every conditional cell probability at every covariate value. For the mixed families, the two uniform sign averages agree despite their distinct homogeneous coefficients. For the fair family, the uniform sign average agrees with the deterministic comparator at its calibrated coefficient. The matching identities continue to hold throughout the signed amplitude boxes, providing exact cancellations alongside the first and second amplitude derivative bounds.

Finally, the density bounds and common smooth neighborhoods in \cref{obj:lem:calibrated-support} are uniform in the partition size, sign vector, labels, and covariate. They furnish a joint regularity certificate for the actual one-record densities, including the mixed cross derivative and the comparator's amplitude derivatives. Together with the effect identities, these properties provide same-model alternatives with controlled separation for the original-record lower-bound argument.


\section{Original-record lower bounds and sharp comparison}

The calibrated families in \cref{obj:def:continuous-calibrated-priors} and their joint support certificate in \cref{obj:lem:calibrated-support} provide the alternatives for the expected-length lower bounds. The comparisons concern sample laws on the original record space \leanref{sym:\Omega}{\ensuremath{\Omega}}, including the observed covariates. We first fix the normalization of the distance used to compare these laws.

\begin{definitionv}{synth_6}[Squared Hellinger distance for original-record sample laws]
For probability laws \(Q,Q'\) on the original-record sample space \(\Omega^n\), define their unhalved squared Hellinger distance by
\[
\mathcal H^2(Q,Q')
=
\int_{\Omega^n}
\left(
\sqrt{\frac{dQ}{dR}}
-
\sqrt{\frac{dQ'}{dR}}
\right)^2\,dR,
\]
where \(R\) is any common dominating measure. The value is independent of the choice of \(R\) and lies in \([0,2]\). For the continuous calibrated mixtures, take \(R\) to be the original-record sample reference measure. This normalization assigns coefficient one to the displayed integral.
\end{definitionv}

The normalization in \cref{obj:synth_6} fixes the numerical constants in the mixture comparisons. Both comparisons use the radius and envelope constant from the same joint certificate in \cref{obj:lem:calibrated-support}. The following bounds apply when the average number of observations per partition cell satisfies the stipulated occupancy threshold.

\begin{lemmav}{lem:original-record-mixtures}[Original-record mixture bounds]
Let \(\varepsilon\) and \(M\) be the radius and envelope constant selected from the fixed joint calibrated-support certificate, with \(0<\varepsilon\le 1/100\) and \(M\ge 1\). Define the positive constants
\[
\kappa=\frac{1}{96\exp(1)},
\qquad
C_{\mathrm{mix}}=96\exp(2)\,8^2M^4.
\]
For sample laws \(Q,Q'\), use the unhalved squared Hellinger distance
\[
\mathcal H^2(Q,Q')
=
\int
\left(
\sqrt{\frac{dQ}{dR}}-
\sqrt{\frac{dQ'}{dR}}
\right)^2\,dR,
\]
where \(R\) is the original-record sample reference measure.

For every choice satisfying
\begin{itemize}
\item \textbf{(Sample and partition sizes.)} \(n,k\) are integers with \(n\ge 1\) and \(k\ge 1\).
\item \textbf{(Occupancy.)} \(n/k\le\kappa\).
\item \textbf{(Signed amplitudes.)} \(\eta,\zeta,\delta\in\mathbb R\) satisfy
\[
|\eta|\le\varepsilon,\qquad
|\zeta|\le\varepsilon,\qquad
|\delta|\le\varepsilon,
\]
\end{itemize}
let \(Q_b\), for \(b\in\{0,1\}\), and \(Q_t\) be the uniform finite mixtures of the calibrated mixed and fair original-record sample laws, respectively:
\[
Q_b=\mathbb E_\sigma\!\left[P_{b,\sigma}^{\otimes n}\right],
\qquad
Q_t=\mathbb E_\sigma\!\left[P_{t,\sigma}^{\otimes n}\right],
\]
where the expectation averages uniformly over the finite sign vectors. Let \(P_{*,t,\delta}\) be the deterministic fair comparator obtained from the fair construction at partition size \(k\), center \(t\), and amplitude \(\delta\), using the constant numerical sign vector with every entry equal to \( -1 \). Then
\[
\mathcal H^2(Q_0,Q_1)
\le
C_{\mathrm{mix}}\frac{n^2}{k}\eta^2\zeta^2,
\]
and, for every \(t\in[0,1/4]\),
\[
\mathcal H^2\!\left(Q_t,P_{*,t,\delta}^{\otimes n}\right)
\le
C_{\mathrm{mix}}\frac{n^2}{k}\delta^4.
\]
\end{lemmav}

\begin{proof}[Proof of \cref{obj:lem:original-record-mixtures}]
\begin{enumerate}
\item Fix the joint certificate \((\varepsilon,M)\) from
\cref{obj:lem:calibrated-support}, and fix admissible sample sizes and
amplitudes. Since \(\exp(1),\exp(2)>0\) and \(M\ge1\),
\[
\kappa>0,\qquad C_{\mathrm{mix}}>0.
\]
Write
\[
\mathcal I_n=\{1,\ldots,n\},\qquad
\mathscr L=\{0,1\}^2,\qquad
\tau=\frac14\operatorname{count}_{\mathscr L}.
\]
Under the identification of an ordered record sample with its covariate
and label vectors, the reference measure is
\[
R=\lambda^{\otimes n}\otimes\tau^{\otimes n}.
\]
For a law \(P\) in the calibrated families, define its conditional label
density by
\[
\mathsf f_P(x,l)
=4P\{(A,Y)=l\mid X=x\},
\qquad x\in[0,1],\quad l\in\mathscr L.
\]
The support certificate supplies
\[
\frac12\le \mathsf f_P(x,l)\le2,
\qquad
\frac14\sum_{l\in\mathscr L}\mathsf f_P(x,l)=1
\]
throughout the prescribed signed amplitude domains.
% lean: CausalSmith.Stat.LogoddsLowsmoothFrontier.calib_constants_spec

\item We first bound the covariate-dependent component weight. For
\(\mathbf x=(x_i)_{i\in\mathcal I_n}\in[0,1]^n\), define
\[
j_k(x)=\min\{k-1,\lfloor kx\rfloor\},
\qquad
\mathscr G_{k,\mathbf x}:
\quad
i\sim i'
\ \Longleftrightarrow\
i\ne i'\ \text{ and }\ |j_k(x_i)-j_k(x_{i'})|\le1.
\]
Thus the cells are half-open, with \(x=1\) assigned to the last cell.
Let
\[
\mathscr V_{\mathbf x}
=\{\text{vertex sets of connected components of }\mathscr G_{k,\mathbf x}\},
\qquad
\mathscr B_{\mathbf x}=\mathscr V_{\mathbf x}\cup\{\varnothing\}.
\]
For every \(\mathcal V\in\mathscr B_{\mathbf x}\), put
\[
m_{\mathcal V}=|\mathcal V|,
\qquad
\mathsf w_{\mathcal V}
=
\begin{cases}
m_{\mathcal V}^4\,8^{m_{\mathcal V}},&m_{\mathcal V}\ge2,\\
0,&m_{\mathcal V}\le1,
\end{cases}
\]
and define
\[
\mathsf W_{\mathrm{occ}}(\mathbf x)
=\sum_{\mathcal V\in\mathscr B_{\mathbf x}}\mathsf w_{\mathcal V},
\qquad
\overline{\mathsf W}_{\mathrm{occ}}
=\int\mathsf W_{\mathrm{occ}}(\mathbf x)\,
d\lambda^{\otimes n}(\mathbf x).
\]
Counting components by their vertex sets gives the identity
\[
\mathsf W_{\mathrm{occ}}(\mathbf x)
=\sum_{m_{\mathrm{occ}}=2}^{n+1}
m_{\mathrm{occ}}^4\,8^{m_{\mathrm{occ}}}
\sum_{\substack{\mathcal I_{\mathrm{rec}}\subseteq\mathcal I_n\\
|\mathcal I_{\mathrm{rec}}|=m_{\mathrm{occ}}}}
\mathbf1_{\{\mathcal I_{\mathrm{rec}}\in\mathscr V_{\mathbf x}\}}.
\]
The term of size \(n+1\) is zero.

Here is the component probability estimate used to integrate this
identity. Fix a nonempty set \(\mathcal I_{\mathrm{rec}}\subseteq\mathcal I_n\)
and a root \(i_{\mathrm{root}}\in\mathcal I_{\mathrm{rec}}\). Define the set
of parent arrays by
\[
\mathscr P_{\mathrm{rec}}
=\left\{
\pi_{\mathrm{par}}:
\mathcal I_{\mathrm{rec}}\setminus\{i_{\mathrm{root}}\}
\longrightarrow\mathcal I_{\mathrm{rec}}
\right\}.
\]
An array is valid when there exists a height function satisfying
\[
\mathsf h_{\mathrm{par}}:\mathcal I_{\mathrm{rec}}\to\mathbb N_0,
\qquad
\mathsf h_{\mathrm{par}}(i_{\mathrm{root}})=0,
\qquad
\mathsf h_{\mathrm{par}}(\pi_{\mathrm{par}}(i))
<\mathsf h_{\mathrm{par}}(i)
\quad(i\ne i_{\mathrm{root}}).
\]
For a valid array, define its edge event by
\[
\mathscr A_{\pi_{\mathrm{par}}}
=
\left\{\mathbf x:
|j_k(x_i)-j_k(x_{\pi_{\mathrm{par}}(i)})|\le1
\text{ for every }i\ne i_{\mathrm{root}}\right\}.
\]
If \(\mathcal I_{\mathrm{rec}}\) is a component, choose each nonroot
vertex's parent one step closer to the root in the induced graph.
Graph distance then supplies the required height function. Consequently,
\[
\mathbf1_{\{\mathcal I_{\mathrm{rec}}\in\mathscr V_{\mathbf x}\}}
\le
\sum_{\substack{\pi_{\mathrm{par}}\in\mathscr P_{\mathrm{rec}}\\
\pi_{\mathrm{par}}\ \mathrm{valid}}}
\mathbf1_{\mathscr A_{\pi_{\mathrm{par}}}}(\mathbf x),
\qquad
|\mathscr P_{\mathrm{rec}}|
=|\mathcal I_{\mathrm{rec}}|^{|\mathcal I_{\mathrm{rec}}|-1}.
\]

For any fixed parent covariate \(x_{\mathrm{par}}\in[0,1]\),
\[
\left\{x_{\mathrm{child}}:
|j_k(x_{\mathrm{child}})-j_k(x_{\mathrm{par}})|\le1\right\}
\subseteq
\left[
\frac{j_k(x_{\mathrm{par}})-1}{k},
\frac{j_k(x_{\mathrm{par}})+2}{k}
\right]\cap[0,1].
\]
Indeed, a point in cell \(j_k(x_{\mathrm{child}})\) lies between that
cell's two endpoints, and the two cell numbers differ by at most one.
The displayed interval has length at most \(3/k\); hence
\[
\lambda\left\{x_{\mathrm{child}}:
|j_k(x_{\mathrm{child}})-j_k(x_{\mathrm{par}})|\le1\right\}
\le\frac3k.
\]
In a valid array, a nonroot vertex of maximal height has no children.
Integrating its covariate first removes one edge and costs at most
\(3/k\). Repeating this operation leaves the unrestricted root, whose
mass is one. Thus
\[
\lambda^{\otimes n}(\mathscr A_{\pi_{\mathrm{par}}})
\le
\left(\frac3k\right)^{|\mathcal I_{\mathrm{rec}}|-1},
\]
including the one-vertex case, where the edge product is empty. Summing
over parent arrays yields
\[
\int
\mathbf1_{\{\mathcal I_{\mathrm{rec}}\in\mathscr V_{\mathbf x}\}}\,
d\lambda^{\otimes n}(\mathbf x)
\le
|\mathcal I_{\mathrm{rec}}|^{|\mathcal I_{\mathrm{rec}}|-1}
\left(\frac3k\right)^{|\mathcal I_{\mathrm{rec}}|-1}.
\]
Therefore
\[
\overline{\mathsf W}_{\mathrm{occ}}
\le
\sum_{m_{\mathrm{occ}}=2}^{n+1}
m_{\mathrm{occ}}^4\,8^{m_{\mathrm{occ}}}
\binom n{m_{\mathrm{occ}}}
m_{\mathrm{occ}}^{m_{\mathrm{occ}}-1}
\left(\frac3k\right)^{m_{\mathrm{occ}}-1}.
\]

To evaluate this finite envelope, use the nonnegative exponential series:
for every integer \(m_{\mathrm{occ}}\ge1\),
\[
\frac{m_{\mathrm{occ}}^{m_{\mathrm{occ}}}}{m_{\mathrm{occ}}!}
\le \exp(m_{\mathrm{occ}})
=\exp(1)^{m_{\mathrm{occ}}},
\qquad
\frac{m_{\mathrm{occ}}^{m_{\mathrm{occ}}-1}}
{m_{\mathrm{occ}}!}
\le\frac{\exp(1)^{m_{\mathrm{occ}}}}{m_{\mathrm{occ}}}.
\]
Also,
\[
\binom n{m_{\mathrm{occ}}}
\le\frac{n^{m_{\mathrm{occ}}}}{m_{\mathrm{occ}}!},
\]
because each factor in the falling factorial is at most \(n\), with
the binomial coefficient equal to zero when \(m_{\mathrm{occ}}>n\).
Set
\[
q_{\mathrm{occ}}=24\exp(1)\frac nk.
\]
The occupancy hypothesis gives \(0\le q_{\mathrm{occ}}\le1/4\), and the
preceding two estimates imply, term by term,
\[
\begin{aligned}
&m_{\mathrm{occ}}^4\,8^{m_{\mathrm{occ}}}
\binom n{m_{\mathrm{occ}}}
m_{\mathrm{occ}}^{m_{\mathrm{occ}}-1}
\left(\frac3k\right)^{m_{\mathrm{occ}}-1}\\
&\hspace{2em}\le
8\exp(1)n\,m_{\mathrm{occ}}^3
q_{\mathrm{occ}}^{m_{\mathrm{occ}}-1}.
\end{aligned}
\]

For completeness, the finite cubic sum has the required numerical bound
by an elementary induction. Define
\[
\mathsf A_{\mathrm{tail}}(h)
=2h^3+12h^2+30h+32,
\qquad h\in\mathbb N_0.
\]
For \(h,N_{\mathrm{tail}}\in\mathbb N_0\),
\[
\sum_{h_{\mathrm{sum}}=0}^{N_{\mathrm{tail}}-1}
(h+h_{\mathrm{sum}}+2)^3q_{\mathrm{occ}}^{h_{\mathrm{sum}}}
\le \mathsf A_{\mathrm{tail}}(h).
\]
The empty sum satisfies this inequality. For the induction step, split
off the first term and apply the induction hypothesis at \(h+1\);
the resulting bound follows from
\[
\begin{aligned}
(h+2)^3+q_{\mathrm{occ}}\mathsf A_{\mathrm{tail}}(h+1)
&\le (h+2)^3+\frac14\mathsf A_{\mathrm{tail}}(h+1),\\
\mathsf A_{\mathrm{tail}}(h)
-\left((h+2)^3+\frac14\mathsf A_{\mathrm{tail}}(h+1)\right)
&=\frac12h^3+\frac32h^2+3h+5\ge0.
\end{aligned}
\]
At \(h=0\), multiplication by \(q_{\mathrm{occ}}\) gives
\[
\sum_{h_{\mathrm{sum}}=0}^{n-1}
(h_{\mathrm{sum}}+2)^3q_{\mathrm{occ}}^{h_{\mathrm{sum}}+1}
\le32q_{\mathrm{occ}}.
\]
Combining the estimates,
\[
\overline{\mathsf W}_{\mathrm{occ}}
\le8\exp(1)n\,32q_{\mathrm{occ}}
=96\exp(2)\,8^2\frac{n^2}{k}.
\]
In particular, for every \(\omega_{\mathrm{amp}}\in\mathbb R\),
\[
M^4\omega_{\mathrm{amp}}^2\overline{\mathsf W}_{\mathrm{occ}}
\le
C_{\mathrm{mix}}\frac{n^2}{k}\omega_{\mathrm{amp}}^2.
\]
The weights are integrable finite sums of measurable component
indicators.
% lean: CausalSmith.Stat.LogoddsLowsmoothFrontier.mixture_weighted_component_integral_bound

\item We now prove the mixed inequality, first treating its degenerate
cases. If \(n=1\), signed singleton matching in
\cref{obj:lem:calibrated-support} gives
\[
\mathbb E_\sigma\mathsf f_{P_{0,\sigma}}(x,l)
=\mathbb E_\sigma\mathsf f_{P_{1,\sigma}}(x,l),
\qquad
\mathcal H^2(Q_0,Q_1)=0.
\]
If \(\eta=0\), the two four-cell tables in
\cref{obj:def:continuous-calibrated-priors} agree for each sign vector:
the propensity is \(1/2\), the effect is zero, and the covariance
adjustment at zero effect is zero. Thus \(Q_0=Q_1\).
If \(\zeta=0\), reversing all signs gives instead the draw-wise identity
\[
P_{0,\sigma}(\eta,0)=P_{1,-\sigma}(\eta,0).
\]
Here the outcome margin is \(1/2\), the covariance adjustment is zero,
and sign reversal exchanges the two propensity margins. Since
\(\sigma\mapsto-\sigma\) permutes the uniform finite prior, again
\(Q_0=Q_1\). In both amplitude cases the Hellinger distance and its
claimed upper bound are zero. Henceforth, for the mixed comparison,
take \(n\ge2\), \(\eta\ne0\), and \(\zeta\ne0\).
% lean: CausalSmith.Stat.LogoddsLowsmoothFrontier.mixed_originalRecordHellinger_one
% lean: CausalSmith.Stat.LogoddsLowsmoothFrontier.mixedMixture_zero_left
% lean: CausalSmith.Stat.LogoddsLowsmoothFrontier.mixedMixture_zero_right

\item Fix the full covariate vector \(\mathbf x\). We describe the
conditional component decomposition explicitly. For each nonempty
\(\mathcal V\in\mathscr B_{\mathbf x}\), and for the empty block, define
\[
\mathscr S_{\mathcal V}
=\bigcup_{i\in\mathcal V}\{j_k(x_i),j_k(x_i)+1\},
\qquad
\mathscr S_{\varnothing}
=\{0,\ldots,k\}\setminus
\bigcup_{\mathcal V\in\mathscr V_{\mathbf x}}\mathscr S_{\mathcal V}.
\]
Distinct nonempty components have disjoint endpoint-sign sets: sharing
an endpoint would make the corresponding cell numbers equal or
adjacent. These sets therefore partition the \(k+1\) endpoint signs.
For an assignment
\(\sigma_{\mathcal V}\in\{-1,1\}^{\mathscr S_{\mathcal V}}\), use the
extension
\[
\widehat\sigma_{\mathcal V}(s)
=
\begin{cases}
\sigma_{\mathcal V}(s),&s\in\mathscr S_{\mathcal V},\\
-1,&s\notin\mathscr S_{\mathcal V}.
\end{cases}
\]
Every record in \(\mathcal V\) depends only on the assigned signs.
Each assignment has exactly \(2^{k+1-|\mathscr S_{\mathcal V}|}\)
completions to a full sign vector, so restricting the uniform prior
preserves each block average.

For \(b\in\{0,1\}\), define the full and component conditional densities
by
\[
\begin{aligned}
\mathsf L_b^{\mathrm{mix}}(\mathbf x,\boldsymbol l;\eta,\zeta)
&=\mathbb E_\sigma
\prod_{i\in\mathcal I_n}
\mathsf f_{P_{b,\sigma}}(x_i,l_i),\\
\mathsf L_{b,\mathcal V}^{\mathrm{mix}}
(\boldsymbol l_{\mathcal V};\eta,\zeta)
&=\mathbb E_{\sigma_{\mathcal V}}
\prod_{i\in\mathcal V}
\mathsf f_{P_{b,\widehat\sigma_{\mathcal V}}}(x_i,l_i),
\qquad
\boldsymbol l_{\mathcal V}\in\mathscr L^{\mathcal V}.
\end{aligned}
\]
The covariates remain fixed in the component notation. The empty
component density is one. Independence of the disjoint sign groups
and finite distributivity give
\[
\mathsf L_b^{\mathrm{mix}}(\mathbf x,\boldsymbol l;\eta,\zeta)
=
\prod_{\mathcal V\in\mathscr B_{\mathbf x}}
\mathsf L_{b,\mathcal V}^{\mathrm{mix}}
(\boldsymbol l|_{\mathcal V};\eta,\zeta).
\]

To make the Hellinger tensor bound transparent, define component label
masses and affinities by
\[
\mathsf p_{b,\mathcal V}^{\mathrm{mix}}(\boldsymbol l_{\mathcal V})
=4^{-m_{\mathcal V}}
\mathsf L_{b,\mathcal V}^{\mathrm{mix}}
(\boldsymbol l_{\mathcal V};\eta,\zeta),
\qquad
\mathfrak a_{\mathcal V}^{\mathrm{mix}}
=
\sum_{\boldsymbol l_{\mathcal V}\in\mathscr L^{\mathcal V}}
\sqrt{\mathsf p_{0,\mathcal V}^{\mathrm{mix}}
(\boldsymbol l_{\mathcal V})
\mathsf p_{1,\mathcal V}^{\mathrm{mix}}
(\boldsymbol l_{\mathcal V})}.
\]
Normalization of each conditional four-cell table yields
\[
\begin{aligned}
\sum_{\boldsymbol l_{\mathcal V}}
\mathsf p_{b,\mathcal V}^{\mathrm{mix}}(\boldsymbol l_{\mathcal V})
&=
\mathbb E_{\sigma_{\mathcal V}}
\prod_{i\in\mathcal V}
\left(\frac14\sum_{l_i\in\mathscr L}
\mathsf f_{P_{b,\widehat\sigma_{\mathcal V}}}(x_i,l_i)\right)
=1,\\
\sum_{\boldsymbol l_{\mathcal V}}
\left(\sqrt{\mathsf p_{0,\mathcal V}^{\mathrm{mix}}}
-\sqrt{\mathsf p_{1,\mathcal V}^{\mathrm{mix}}}\right)^2
&=2(1-\mathfrak a_{\mathcal V}^{\mathrm{mix}}).
\end{aligned}
\]
All masses are nonnegative. The last equality, together with
nonnegativity of the squared sum, shows
\(0\le\mathfrak a_{\mathcal V}^{\mathrm{mix}}\le1\).
Regrouping labels by their component restrictions factors the full
affinity into the product of these affinities. The elementary
inequality
\[
1-\prod_{\mathcal V}\mathfrak a_{\mathcal V}^{\mathrm{mix}}
\le
\sum_{\mathcal V}(1-\mathfrak a_{\mathcal V}^{\mathrm{mix}})
\]
follows by induction from
\(1-aa'=(1-a)+a(1-a')\le(1-a)+(1-a')\) for
\(a,a'\in[0,1]\).

Define the conditional discrepancy by
\[
\mathsf h_{\mathrm{mix}}(\mathbf x)
=
\int_{\mathscr L^{\mathcal I_n}}
\left(
\sqrt{\mathsf L_0^{\mathrm{mix}}(\mathbf x,\boldsymbol l;\eta,\zeta)}
-\sqrt{\mathsf L_1^{\mathrm{mix}}(\mathbf x,\boldsymbol l;\eta,\zeta)}
\right)^2\,d\tau^{\otimes n}(\boldsymbol l).
\]
Since each label vector has reference mass \(4^{-n}\), scaling the
densities by \(4^{-n}\) turns this integral exactly into the squared
root discrepancy of the corresponding label masses. The affinity
calculation therefore proves
\[
\mathsf h_{\mathrm{mix}}(\mathbf x)
\le
\sum_{\mathcal V\in\mathscr B_{\mathbf x}}
4^{-m_{\mathcal V}}
\sum_{\boldsymbol l_{\mathcal V}\in\mathscr L^{\mathcal V}}
\left(
\sqrt{\mathsf L_{0,\mathcal V}^{\mathrm{mix}}}
-\sqrt{\mathsf L_{1,\mathcal V}^{\mathrm{mix}}}
\right)^2.
\]
Finally, the common uniform covariate marginal and the displayed
product densities give, by integration over covariates and labels,
\[
\mathcal H^2(Q_0,Q_1)
=\int\mathsf h_{\mathrm{mix}}(\mathbf x)\,
d\lambda^{\otimes n}(\mathbf x).
\]
% lean: CausalSmith.Stat.LogoddsLowsmoothFrontier.mixed_conditional_label_hellinger_le_components
% lean: CausalSmith.Stat.LogoddsLowsmoothFrontier.mixed_originalRecordHellinger_conditional

\item We bound each mixed component's pointwise square-root discrepancy.
If \(m_{\mathcal V}=0\), both densities are one. If
\(m_{\mathcal V}=1\), restriction of the uniform prior and signed
singleton matching in \cref{obj:lem:calibrated-support} give equality
of the two densities. Thus, for \(m_{\mathcal V}\le1\),
\[
\left(
\sqrt{\mathsf L_{0,\mathcal V}^{\mathrm{mix}}}
-\sqrt{\mathsf L_{1,\mathcal V}^{\mathrm{mix}}}
\right)^2=0=M^4(\eta\zeta)^2\mathsf w_{\mathcal V}.
\]

Suppose \(m_{\mathcal V}\ge2\), and fix its labels. Define, throughout
the signed amplitude square,
\[
\mathsf D_{\mathcal V}^{\mathrm{mix}}(\eta_0,\zeta_0)
=
\mathsf L_{0,\mathcal V}^{\mathrm{mix}}
(\boldsymbol l_{\mathcal V};\eta_0,\zeta_0)
-
\mathsf L_{1,\mathcal V}^{\mathrm{mix}}
(\boldsymbol l_{\mathcal V};\eta_0,\zeta_0),
\qquad
|\eta_0|,|\zeta_0|\le\varepsilon.
\]
For a fixed sign assignment, put
\[
\mathsf f_{b,i}^{\mathrm{mix}}(\eta_0,\zeta_0)
=\mathsf f_{P_{b,\widehat\sigma_{\mathcal V}}}(x_i,l_i),
\qquad
\mathsf F_{b,\mathcal J}^{\mathrm{mix}}(\eta_0,\zeta_0)
=\prod_{i\in\mathcal J}\mathsf f_{b,i}^{\mathrm{mix}}(\eta_0,\zeta_0),
\quad \mathcal J\subseteq\mathcal V.
\]
Here each law on the right is evaluated at \((\eta_0,\zeta_0)\).
The derivative and common-neighborhood assertions of
\cref{obj:lem:calibrated-support} imply
\[
|\mathsf f_{b,i}^{\mathrm{mix}}|\le2,\qquad
|\partial_{\eta_0}\mathsf f_{b,i}^{\mathrm{mix}}|,
|\partial_{\zeta_0}\mathsf f_{b,i}^{\mathrm{mix}}|,
|\partial_{\zeta_0}\partial_{\eta_0}\mathsf f_{b,i}^{\mathrm{mix}}|
\le M.
\]
The partial derivative bounds follow by evaluating the total
derivatives on unit coordinate vectors.

Product differentiation gives, with \(m_{\mathcal J}=|\mathcal J|\),
\[
\begin{aligned}
|\mathsf F_{b,\mathcal J}^{\mathrm{mix}}|&\le2^{m_{\mathcal J}},\\
|\partial_{\eta_0}\mathsf F_{b,\mathcal J}^{\mathrm{mix}}|,
|\partial_{\zeta_0}\mathsf F_{b,\mathcal J}^{\mathrm{mix}}|
&\le M m_{\mathcal J}2^{m_{\mathcal J}},\\
|\partial_{\zeta_0}\partial_{\eta_0}
\mathsf F_{b,\mathcal J}^{\mathrm{mix}}|
&\le M^2m_{\mathcal J}^2 2^{m_{\mathcal J}}.
\end{aligned}
\]
These estimates can be proved together by induction on
\(\mathcal J\). For the empty product its value is one and all its
derivatives vanish. On adjoining \(i_{\mathrm{new}}\), the first
derivative bound is controlled by
\[
M2^{m_{\mathcal J}}+
2M m_{\mathcal J}2^{m_{\mathcal J}}
\le
M(m_{\mathcal J}+1)2^{m_{\mathcal J}+1}.
\]
The mixed product rule has four terms:
\[
\begin{aligned}
\partial_{\zeta_0}\partial_{\eta_0}
\left(\mathsf f_{b,i_{\mathrm{new}}}^{\mathrm{mix}}
\mathsf F_{b,\mathcal J}^{\mathrm{mix}}\right)
={}&
(\partial_{\zeta_0}\partial_{\eta_0}
\mathsf f_{b,i_{\mathrm{new}}}^{\mathrm{mix}})
\mathsf F_{b,\mathcal J}^{\mathrm{mix}}
+
(\partial_{\eta_0}\mathsf f_{b,i_{\mathrm{new}}}^{\mathrm{mix}})
(\partial_{\zeta_0}\mathsf F_{b,\mathcal J}^{\mathrm{mix}})\\
&+
(\partial_{\zeta_0}\mathsf f_{b,i_{\mathrm{new}}}^{\mathrm{mix}})
(\partial_{\eta_0}\mathsf F_{b,\mathcal J}^{\mathrm{mix}})
+
\mathsf f_{b,i_{\mathrm{new}}}^{\mathrm{mix}}
(\partial_{\zeta_0}\partial_{\eta_0}
\mathsf F_{b,\mathcal J}^{\mathrm{mix}}).
\end{aligned}
\]
Their absolute sum is at most
\[
\left(M+2M^2m_{\mathcal J}
+2M^2m_{\mathcal J}^2\right)2^{m_{\mathcal J}}
\le
M^2(m_{\mathcal J}+1)^2 2^{m_{\mathcal J}+1},
\]
using \(M\ge1\). Differentiation commutes with the finite sign
averages, so subtraction of the two averages yields
\[
|\partial_{\zeta_0}\partial_{\eta_0}
\mathsf D_{\mathcal V}^{\mathrm{mix}}(\eta_0,\zeta_0)|
\le
2M^2m_{\mathcal V}^2 2^{m_{\mathcal V}}.
\]
The one-record density floor also gives
\[
\mathsf L_{0,\mathcal V}^{\mathrm{mix}},
\mathsf L_{1,\mathcal V}^{\mathrm{mix}}
\ge2^{-m_{\mathcal V}}.
\]

The same table identities used in the amplitude cases, now applied
to a block average, imply
\[
\mathsf D_{\mathcal V}^{\mathrm{mix}}(0,\zeta_0)=0,\qquad
\mathsf D_{\mathcal V}^{\mathrm{mix}}(\eta_0,0)=0,\qquad
\partial_{\eta_0}\mathsf D_{\mathcal V}^{\mathrm{mix}}(\eta_0,0)=0.
\]
For the second identity, reverse the component signs, a permutation
of their uniform prior. The two closed segments joining zero to
\(\eta\) and \(\zeta\) lie inside the certificate's signed square.
The mean value bound along the \(\zeta_0\)-segment therefore gives
\[
|\partial_{\eta_0}\mathsf D_{\mathcal V}^{\mathrm{mix}}(\eta_0,\zeta)|
\le
2M^2m_{\mathcal V}^2 2^{m_{\mathcal V}}|\zeta|
\]
for every \(\eta_0\) between zero and \(\eta\). Applying the mean value
bound along that \(\eta_0\)-segment next gives
\[
|\mathsf D_{\mathcal V}^{\mathrm{mix}}(\eta,\zeta)|
\le
2M^2m_{\mathcal V}^2 2^{m_{\mathcal V}}|\eta\zeta|.
\]
Both estimates apply to either orientation of the segments.

The positive density floors justify the square-root identity and
bound
\[
\begin{aligned}
\left(
\sqrt{\mathsf L_{0,\mathcal V}^{\mathrm{mix}}}
-\sqrt{\mathsf L_{1,\mathcal V}^{\mathrm{mix}}}
\right)^2
&=
\frac{(\mathsf D_{\mathcal V}^{\mathrm{mix}}(\eta,\zeta))^2}
{\left(
\sqrt{\mathsf L_{0,\mathcal V}^{\mathrm{mix}}}
+\sqrt{\mathsf L_{1,\mathcal V}^{\mathrm{mix}}}
\right)^2}\\
&\le
\frac{4M^4m_{\mathcal V}^4 4^{m_{\mathcal V}}(\eta\zeta)^2}
{4\,2^{-m_{\mathcal V}}}\\
&=M^4(\eta\zeta)^2\mathsf w_{\mathcal V}.
\end{aligned}
\]
This proves the pointwise component bound for every block, including
the unused-sign block.
% lean: CausalSmith.Stat.LogoddsLowsmoothFrontier.mixedDensity_small_block_matching
% lean: CausalSmith.Stat.LogoddsLowsmoothFrontier.mixed_component_hellinger_from_support

\item Averaging the pointwise bound over a component's labels preserves
it, since there are \(4^{m_{\mathcal V}}\) label vectors:
\[
4^{-m_{\mathcal V}}
\sum_{\boldsymbol l_{\mathcal V}}
\left(
\sqrt{\mathsf L_{0,\mathcal V}^{\mathrm{mix}}}
-\sqrt{\mathsf L_{1,\mathcal V}^{\mathrm{mix}}}
\right)^2
\le M^4(\eta\zeta)^2\mathsf w_{\mathcal V}.
\]
Summing over components and integrating over the original covariates
now yields
\[
\begin{aligned}
\mathcal H^2(Q_0,Q_1)
&\le
\int M^4(\eta\zeta)^2\mathsf W_{\mathrm{occ}}(\mathbf x)\,
d\lambda^{\otimes n}(\mathbf x)\\
&=M^4(\eta\zeta)^2\overline{\mathsf W}_{\mathrm{occ}}\\
&\le C_{\mathrm{mix}}\frac{n^2}{k}\eta^2\zeta^2.
\end{aligned}
\]
The conditional discrepancy is nonnegative, and the integrable
component weight supplies the integrable upper bound.
% lean: CausalSmith.Stat.LogoddsLowsmoothFrontier.conditionalComponent_mass_hellinger_le_density_bound
% lean: CausalSmith.Stat.LogoddsLowsmoothFrontier.original_record_mixtures

\item Fix \(t\in[0,1/4]\) for the fair comparison. If \(n=1\), signed
singleton matching in \cref{obj:lem:calibrated-support} gives
\[
\mathbb E_\sigma\mathsf f_{P_{t,\sigma}}(x,l)
=\mathsf f_{P_{*,t,\delta}}(x,l),
\qquad
\mathcal H^2(Q_t,P_{*,t,\delta})=0.
\]
If \(\delta=0\), the formulas in
\cref{obj:def:calibration-maps,obj:def:continuous-calibrated-priors}
give \(\mathsf T(t,0)=t\); the random and comparator tables then agree
for every sign draw. Thus
\[
Q_t=P_{*,t,0}^{\otimes n},
\qquad
\mathcal H^2(Q_t,P_{*,t,0}^{\otimes n})=0.
\]
These cases satisfy the desired inequality. Take henceforth \(n\ge2\)
and \(\delta\ne0\).
% lean: CausalSmith.Stat.LogoddsLowsmoothFrontier.fair_originalRecordHellinger_one
% lean: CausalSmith.Stat.LogoddsLowsmoothFrontier.fairMixture_zero_amplitude

\item First establish a reduction for any auxiliary amplitude
\(\delta_{\mathrm{aux}}\in[0,\varepsilon]\). With the covariates fixed,
define
\[
\begin{aligned}
\mathsf L_t^{\mathrm{fair}}(\mathbf x,\boldsymbol l;\delta_{\mathrm{aux}})
&=\mathbb E_\sigma\prod_{i\in\mathcal I_n}
\mathsf f_{P_{t,\sigma}}(x_i,l_i),\\
\mathsf L_*^{\mathrm{fair}}(\mathbf x,\boldsymbol l;\delta_{\mathrm{aux}})
&=\prod_{i\in\mathcal I_n}
\mathsf f_{P_{*,t,\delta_{\mathrm{aux}}}}(x_i,l_i),\\
\mathsf L_{t,\mathcal V}^{\mathrm{fair}}
(\boldsymbol l_{\mathcal V};\delta_{\mathrm{aux}})
&=\mathbb E_{\sigma_{\mathcal V}}\prod_{i\in\mathcal V}
\mathsf f_{P_{t,\widehat\sigma_{\mathcal V}}}(x_i,l_i),\\
\mathsf L_{*,\mathcal V}^{\mathrm{fair}}
(\boldsymbol l_{\mathcal V};\delta_{\mathrm{aux}})
&=\prod_{i\in\mathcal V}
\mathsf f_{P_{*,t,\delta_{\mathrm{aux}}}}(x_i,l_i).
\end{aligned}
\]
The random laws in these displays are evaluated at
\(\delta_{\mathrm{aux}}\). Averaging a sign-independent product is
exactly that product:
\[
\mathbb E_\sigma
\prod_{i\in\mathcal I_n}
\mathsf f_{P_{*,t,\delta_{\mathrm{aux}}}}(x_i,l_i)
=
\prod_{i\in\mathcal I_n}
\mathsf f_{P_{*,t,\delta_{\mathrm{aux}}}}(x_i,l_i).
\]
Restriction of the random prior to each component and finite
factorization consequently give
\[
\mathsf L_t^{\mathrm{fair}}
=\prod_{\mathcal V\in\mathscr B_{\mathbf x}}
\mathsf L_{t,\mathcal V}^{\mathrm{fair}},
\qquad
\mathsf L_*^{\mathrm{fair}}
=\prod_{\mathcal V\in\mathscr B_{\mathbf x}}
\mathsf L_{*,\mathcal V}^{\mathrm{fair}}.
\]
For each of the two component densities, summing its scaled label
mass gives
\[
4^{-m_{\mathcal V}}
\sum_{\boldsymbol l_{\mathcal V}}
\mathsf L_{t,\mathcal V}^{\mathrm{fair}}
=
4^{-m_{\mathcal V}}
\sum_{\boldsymbol l_{\mathcal V}}
\mathsf L_{*,\mathcal V}^{\mathrm{fair}}
=1.
\]
For the random density this follows by commuting the finite sign
average with the label sum; in both cases the label sum factors into
normalized one-record sums.

Define the component affinities and conditional discrepancy by
\[
\begin{aligned}
\mathfrak a_{\mathcal V}^{\mathrm{fair}}
&=
4^{-m_{\mathcal V}}\sum_{\boldsymbol l_{\mathcal V}}
\sqrt{\mathsf L_{t,\mathcal V}^{\mathrm{fair}}
\mathsf L_{*,\mathcal V}^{\mathrm{fair}}},\\
\mathsf h_{\mathrm{fair}}(\mathbf x;\delta_{\mathrm{aux}})
&=
\int
\left(
\sqrt{\mathsf L_t^{\mathrm{fair}}}
-\sqrt{\mathsf L_*^{\mathrm{fair}}}
\right)^2\,d\tau^{\otimes n}.
\end{aligned}
\]
Expanding the squared difference and using normalization gives
\(0\le\mathfrak a_{\mathcal V}^{\mathrm{fair}}\le1\) and
\[
\begin{aligned}
\mathsf h_{\mathrm{fair}}(\mathbf x;\delta_{\mathrm{aux}})
&=2\left(1-\prod_{\mathcal V}
\mathfrak a_{\mathcal V}^{\mathrm{fair}}\right)\\
&\le
2\sum_{\mathcal V}(1-\mathfrak a_{\mathcal V}^{\mathrm{fair}})\\
&=
\sum_{\mathcal V\in\mathscr B_{\mathbf x}}
4^{-m_{\mathcal V}}\sum_{\boldsymbol l_{\mathcal V}}
\left(
\sqrt{\mathsf L_{t,\mathcal V}^{\mathrm{fair}}}
-\sqrt{\mathsf L_{*,\mathcal V}^{\mathrm{fair}}}
\right)^2.
\end{aligned}
\]
Here the affinity factors by finite distributivity, and the inequality
is the finite-product inequality for numbers in \([0,1]\).
The common covariate marginal also gives
\[
\mathcal H^2(Q_t,P_{*,t,\delta_{\mathrm{aux}}}^{\otimes n})
=
\int\mathsf h_{\mathrm{fair}}(\mathbf x;\delta_{\mathrm{aux}})
\,d\lambda^{\otimes n}(\mathbf x),
\]
where \(Q_t\) in this display is evaluated at the auxiliary amplitude.
% lean: CausalSmith.Stat.LogoddsLowsmoothFrontier.fair_conditional_label_hellinger_le_components
% lean: CausalSmith.Stat.LogoddsLowsmoothFrontier.fair_originalRecordHellinger_conditional

\item For \(m_{\mathcal V}\le1\), the two fair component densities agree:
the empty products are one, and the singleton identity follows from
signed singleton matching in \cref{obj:lem:calibrated-support} and
restriction of the uniform sign prior. Hence their square-root
discrepancy is zero.

For \(m_{\mathcal V}\ge2\), fix the component labels and define, on the
whole signed interval,
\[
\mathsf D_{\mathcal V}^{\mathrm{fair}}(\delta_0)
=
\mathsf L_{t,\mathcal V}^{\mathrm{fair}}
(\boldsymbol l_{\mathcal V};\delta_0)
-
\mathsf L_{*,\mathcal V}^{\mathrm{fair}}
(\boldsymbol l_{\mathcal V};\delta_0),
\qquad |\delta_0|\le\varepsilon.
\]
For the derivative calculation, define the one-record factors and
their partial products by
\[
\begin{aligned}
\mathsf f_{t,i}^{\mathrm{fair}}(\delta_0)
&=\mathsf f_{P_{t,\widehat\sigma_{\mathcal V}}}(x_i,l_i),&
\mathsf f_{*,i}^{\mathrm{fair}}(\delta_0)
&=\mathsf f_{P_{*,t,\delta_0}}(x_i,l_i),\\
\mathsf F_{t,\mathcal J}^{\mathrm{fair}}(\delta_0)
&=\prod_{i\in\mathcal J}\mathsf f_{t,i}^{\mathrm{fair}}(\delta_0),&
\mathsf F_{*,\mathcal J}^{\mathrm{fair}}(\delta_0)
&=\prod_{i\in\mathcal J}\mathsf f_{*,i}^{\mathrm{fair}}(\delta_0),
\qquad \mathcal J\subseteq\mathcal V.
\end{aligned}
\]
The random law on the right is evaluated at \(\delta_0\). By
\cref{obj:lem:calibrated-support}, both kinds of factors are twice
continuously differentiable and satisfy
\[
|\mathsf f_{t,i}^{\mathrm{fair}}|,
|\mathsf f_{*,i}^{\mathrm{fair}}|\le2,\qquad
|(\mathsf f_{t,i}^{\mathrm{fair}})'|,
|(\mathsf f_{*,i}^{\mathrm{fair}})'|,
|(\mathsf f_{t,i}^{\mathrm{fair}})''|,
|(\mathsf f_{*,i}^{\mathrm{fair}})''|\le M.
\]
For either kind of partial product, the product rule and induction
give
\[
|\mathsf F|\le2^{m_{\mathcal J}},\qquad
|\mathsf F'|\le M m_{\mathcal J}2^{m_{\mathcal J}},\qquad
|\mathsf F''|\le M^2m_{\mathcal J}^2 2^{m_{\mathcal J}}.
\]
To see the second-derivative induction explicitly, adjoining a
factor gives
\[
(\mathsf f\,\mathsf F)''
=\mathsf f\,\mathsf F''+2\mathsf f'\mathsf F'
+\mathsf f''\mathsf F,
\]
whose absolute value is at most
\[
\left(2M^2m_{\mathcal J}^2+2M^2m_{\mathcal J}+M\right)
2^{m_{\mathcal J}}
\le
M^2(m_{\mathcal J}+1)^2 2^{m_{\mathcal J}+1}.
\]
Here \(\mathsf f\) denotes the adjoined one-record factor and
\(\mathsf F\) the product over \(\mathcal J\), of either kind. The
first-derivative induction uses
\[
M2^{m_{\mathcal J}}+
2M m_{\mathcal J}2^{m_{\mathcal J}}
\le M(m_{\mathcal J}+1)2^{m_{\mathcal J}+1}.
\]
The empty product again has value one and zero derivatives.
Finite averaging and subtraction therefore imply, for
\(0\le\delta_0\le\delta_{\mathrm{aux}}\),
\[
|(\mathsf D_{\mathcal V}^{\mathrm{fair}})''(\delta_0)|
\le2M^2m_{\mathcal V}^2 2^{m_{\mathcal V}}.
\]
The density floors give
\[
\mathsf L_{t,\mathcal V}^{\mathrm{fair}},
\mathsf L_{*,\mathcal V}^{\mathrm{fair}}
\ge2^{-m_{\mathcal V}}.
\]

At zero amplitude the random and comparator tables coincide, so
\[
\mathsf D_{\mathcal V}^{\mathrm{fair}}(0)=0.
\]
The fair root is even in its amplitude by
\cref{obj:lem:exact-calibrations}; the comparator effect in
\cref{obj:def:calibration-maps} depends on the amplitude through its
square. Reversing the component signs consequently proves
\[
\mathsf D_{\mathcal V}^{\mathrm{fair}}(-\delta_0)
=\mathsf D_{\mathcal V}^{\mathrm{fair}}(\delta_0).
\]
Differentiating this equality at zero, using the supplied smoothness,
gives
\[
(\mathsf D_{\mathcal V}^{\mathrm{fair}})'(0)
=-(\mathsf D_{\mathcal V}^{\mathrm{fair}})'(0)=0.
\]
If \(\delta_{\mathrm{aux}}=0\), the difference is already zero. If
\(\delta_{\mathrm{aux}}>0\), Taylor's theorem with Lagrange remainder
provides \(\delta_{\mathrm{mid}}\in(0,\delta_{\mathrm{aux}})\) such that
\[
\mathsf D_{\mathcal V}^{\mathrm{fair}}(\delta_{\mathrm{aux}})
=
\frac{\delta_{\mathrm{aux}}^2}{2}
(\mathsf D_{\mathcal V}^{\mathrm{fair}})''(\delta_{\mathrm{mid}}).
\]
Thus in either case,
\[
|\mathsf D_{\mathcal V}^{\mathrm{fair}}(\delta_{\mathrm{aux}})|
\le
M^2m_{\mathcal V}^2 2^{m_{\mathcal V}}\delta_{\mathrm{aux}}^2
\le
2M^2m_{\mathcal V}^2 2^{m_{\mathcal V}}\delta_{\mathrm{aux}}^2.
\]
Using this last bound and the two positive floors,
\[
\begin{aligned}
\left(
\sqrt{\mathsf L_{t,\mathcal V}^{\mathrm{fair}}}
-\sqrt{\mathsf L_{*,\mathcal V}^{\mathrm{fair}}}
\right)^2
&=
\frac{(\mathsf D_{\mathcal V}^{\mathrm{fair}}
(\delta_{\mathrm{aux}}))^2}
{\left(
\sqrt{\mathsf L_{t,\mathcal V}^{\mathrm{fair}}}
+\sqrt{\mathsf L_{*,\mathcal V}^{\mathrm{fair}}}
\right)^2}\\
&\le
\frac{4M^4m_{\mathcal V}^4 4^{m_{\mathcal V}}
\delta_{\mathrm{aux}}^4}{4\,2^{-m_{\mathcal V}}}\\
&=M^4\delta_{\mathrm{aux}}^4\mathsf w_{\mathcal V}.
\end{aligned}
\]
Together with the small-component cases, this estimate holds for every
\(\mathcal V\in\mathscr B_{\mathbf x}\).
% lean: CausalSmith.Stat.LogoddsLowsmoothFrontier.fairDensity_small_block_matching
% lean: CausalSmith.Stat.LogoddsLowsmoothFrontier.fair_component_hellinger_from_support

\item Summing the fair pointwise bound over component labels with their
reference masses, then summing over components, gives
\[
\mathsf h_{\mathrm{fair}}(\mathbf x;\delta_{\mathrm{aux}})
\le
M^4(\delta_{\mathrm{aux}}^2)^2
\mathsf W_{\mathrm{occ}}(\mathbf x).
\]
Integration over the original covariates therefore establishes the
nonnegative-amplitude reduction
\[
\mathcal H^2(Q_t,P_{*,t,\delta_{\mathrm{aux}}}^{\otimes n})
\le
M^4(\delta_{\mathrm{aux}}^2)^2
\overline{\mathsf W}_{\mathrm{occ}},
\qquad 0\le\delta_{\mathrm{aux}}\le\varepsilon.
\]
As before, the conditional discrepancy is nonnegative and the
component weight is integrable.
% lean: CausalSmith.Stat.LogoddsLowsmoothFrontier.original_record_mixtures

\item Finally, restore the signed amplitude. The fair-root evenness and
the square dependence of the comparator effect give, throughout the
signed certificate interval,
\[
P_{t,\sigma}(-\delta)=P_{t,-\sigma}(\delta),
\qquad
P_{*,t,-\delta}=P_{*,t,\delta}.
\]
The notation in the first identity displays the amplitude argument of
the random fair law. Reindexing the finite sign average by
\(\sigma\mapsto-\sigma\) shows that the fair sample mixture is even
as well. Thus, if \(\delta\ge0\) use \(|\delta|=\delta\), and if
\(\delta<0\) use \(|\delta|=-\delta\), to obtain
\[
\mathcal H^2(Q_t,P_{*,t,\delta}^{\otimes n})
=
\left.
\mathcal H^2(Q_t,P_{*,t,\delta_{\mathrm{aux}}}^{\otimes n})
\right|_{\delta_{\mathrm{aux}}=|\delta|}.
\]
The mixture on the right is evaluated at \(|\delta|\).
Since \(0\le|\delta|\le\varepsilon\), the nonnegative-amplitude
reduction applies, and \(|\delta|^2=\delta^2\) gives
\[
\begin{aligned}
\mathcal H^2(Q_t,P_{*,t,\delta}^{\otimes n})
&\le M^4(\delta^2)^2\overline{\mathsf W}_{\mathrm{occ}}\\
&\le C_{\mathrm{mix}}\frac{n^2}{k}(\delta^2)^2\\
&=C_{\mathrm{mix}}\frac{n^2}{k}\delta^4.
\end{aligned}
\]
Every estimate used the same fixed certificate and was uniform over
\(t\in[0,1/4]\), proving both asserted bounds.
% lean: CausalSmith.Stat.LogoddsLowsmoothFrontier.fairMixture_even
% lean: CausalSmith.Stat.LogoddsLowsmoothFrontier.fairComparator_even
% lean: CausalSmith.Stat.LogoddsLowsmoothFrontier.original_record_mixtures
\end{enumerate}
\end{proof}

The two bounds in \cref{obj:lem:original-record-mixtures} quantify the statistical proximity of the same-model alternatives certified by \cref{obj:lem:calibrated-support}. The mixed comparison depends on the squared product of the propensity and outcome amplitudes; the fair comparison depends on the fourth power of its outcome amplitude and is uniform over the stated centers. Both retain the factor \(n^2/k\) and use a common constant determined by the joint density envelope. Exact singleton matching and the signed derivative bounds in \cref{obj:lem:calibrated-support} supply the calibration properties underlying these comparisons.

Expected length is evaluated on the radius slice in \cref{obj:def:radius-model}, while coverage is required throughout the ambient model by \cref{obj:def:honest-procedures}. Consequently, alternatives in the ambient model constrain interval length at laws in an evaluation slice. The parametric comparison records this implication uniformly over all radii, including zero.

\begin{lemmav}{lem:parametric-null-floor}[Parametric lower bound on length]
Let \(\alpha,\beta,r\in\mathbb R\) and \(n\in\mathbb N\) satisfy:
\begin{itemize}
\item \textbf{(Exponent domain.)} \((\alpha,\beta)\in\mathcal D\), the public exponent domain given by
\[
\mathcal D=\{(\alpha,\beta)\in\mathbb R^2:0<\beta<1/4,\ \beta<\alpha\le1\}.
\]
\item \textbf{(Sample size.)} \(n\ge1\).
\item \textbf{(Radius.)} \(0\le r\le1/2\).
\end{itemize}
Then the minimax expected interval length \(J_n(\alpha,\beta;r)\) defined in \cref{obj:def:length-objective} satisfies
\[
J_n(\alpha,\beta;r)\ge \frac{3}{16}n^{-1/2}.
\]
\end{lemmav}

\begin{proof}[Proof of \cref{obj:lem:parametric-null-floor}]
\begin{enumerate}
\item Fix the parameters in the statement. For every \(\tau\in\mathbb R\), define the observed law \(P_{\mathrm{par},\tau}\) by
\[
P_{\mathrm{par},\tau}(dx,\{(a,y)\})
=
\lambda(dx)
\begin{cases}
1/4,&a=0,\ y\in\{0,1\},\\
\frac12\{1-\ell(\tau)\},&a=1,\ y=0,\\
\frac12\ell(\tau),&a=1,\ y=1,
\end{cases}
\qquad
x\in[0,1],\quad a,y\in\{0,1\},
\]
where \(\ell\) is the logistic map,
\[
\ell(\tau)=\frac{1}{1+\exp(-\tau)}.
\]
The four conditional cell probabilities are constant in \(x\), strictly positive, and sum to one. Thus this construction has continuous interior cells and uniform covariate marginal \(\lambda\), which has full support.

Summing the treated cells gives treatment probability \(1/2\); dividing each success cell by its arm probability gives control risk \(1/2\) and treated risk \(\ell(\tau)\). Moreover,
\[
\frac{\ell(\tau)}{1-\ell(\tau)}=\exp(\tau),
\qquad
\operatorname{logit}\{\ell(\tau)\}=\tau,
\qquad
\operatorname{logit}(1/2)=0.
\]
Consequently, for every \(x\in[0,1]\),
\[
\begin{aligned}
e(P_{\mathrm{par},\tau},x)&=\frac12,\\
g(P_{\mathrm{par},\tau},x)&=0,
&
\nu(P_{\mathrm{par},\tau},x)&=0,\\
\theta(P_{\mathrm{par},\tau})&=\tau,
&
\mu_1(P_{\mathrm{par},\tau},x)
&=\ell\{\nu(P_{\mathrm{par},\tau},x)+\theta(P_{\mathrm{par},\tau})\}.
\end{aligned}
\]
The nuisance envelopes and seminorms therefore satisfy
\[
\|g(P_{\mathrm{par},\tau},\cdot)\|_\infty
=
\|\nu(P_{\mathrm{par},\tau},\cdot)\|_\infty
=
[g(P_{\mathrm{par},\tau},\cdot)]_\alpha
=
[\nu(P_{\mathrm{par},\tau},\cdot)]_\beta
=0.
\]
Together with the uniform design and homogeneous relation, these identities verify every condition of \cref{obj:def:model} whenever \(|\tau|\le1/2\).

Define the comparison amplitude and laws by
\[
\tau_n:=\frac14 n^{-1/2},
\qquad
P_{\mathrm{null}}:=P_{\mathrm{par},0},
\qquad
P_{\mathrm{shift}}:=P_{\mathrm{par},\tau_n}.
\]
Since \(n\ge1\),
\[
0\le n^{-1/2}\le1,
\qquad
0\le\tau_n\le\frac14\le\frac12.
\]
Hence, using \cref{obj:def:radius-model},
\[
P_{\mathrm{null}}\in\mathcal M_r(\alpha,\beta),
\qquad
P_{\mathrm{shift}}\in\mathcal M(\alpha,\beta),
\qquad
\theta(P_{\mathrm{null}})=0,
\qquad
\theta(P_{\mathrm{shift}})=\tau_n.
\]
Indeed, the radius condition at the null is \(0\le r\), including when \(r=0\).
% lean: CausalSmith.Stat.LogoddsLowsmoothFrontier.parametric_comparison

\item All four null cells have probability \(1/4\). The likelihood ratio of \(P_{\mathrm{par},\tau}\) relative to \(P_{\mathrm{null}}\) is therefore
\[
\Lambda_\tau(x,a,y)
:=
\frac{dP_{\mathrm{par},\tau}}{dP_{\mathrm{null}}}(x,a,y)
=
\begin{cases}
1,&a=0,\\
2\{1-\ell(\tau)\},&a=1,\ y=0,\\
2\ell(\tau),&a=1,\ y=1,
\end{cases}
\qquad \tau\in\mathbb R.
\]
This bounded density establishes absolute continuity and integrability of its squared deviation. Define its one-record chi-square divergence by
\[
\chi^2_\tau
:=
\int_\Omega(\Lambda_\tau-1)^2\,dP_{\mathrm{null}}.
\]
Summing over the four null cells gives
\[
\begin{aligned}
\chi^2_\tau
&=
\frac14\left[
\{2\ell(\tau)-1\}^2+
\{2(1-\ell(\tau))-1\}^2
\right]\\
&=2\{\ell(\tau)-1/2\}^2.
\end{aligned}
\]
For every real \(\tau\), differentiation yields
\[
0\le\ell'(\tau)
=\ell(\tau)\{1-\ell(\tau)\}
=\frac14-\{\ell(\tau)-1/2\}^2
\le\frac14.
\]
The mean-value bound, together with \(\ell(0)=1/2\), consequently gives
\[
|\ell(\tau)-1/2|\le\frac{|\tau|}{4},
\qquad
\chi^2_\tau
\le2\left(\frac{|\tau|}{4}\right)^2
=\frac{\tau^2}{8}.
\]
% lean: CausalSmith.Stat.LogoddsLowsmoothFrontier.parametricLaw_chiSq_le

\item The amplitude has the exact information budget
\[
\frac{n\tau_n^2}{8}=\frac1{128},
\qquad
n\chi^2_{\tau_n}\le\frac1{128}.
\]
Since the exponential lies above its tangent at zero,
\[
(1+\chi^2_{\tau_n})^n
\le
\{\exp(\chi^2_{\tau_n})\}^n
=
\exp(n\chi^2_{\tau_n})
\le
\exp(1/128).
\]
Comparison of the exponential and geometric series gives the required numerical bound:
\[
\exp(1/128)
=
\sum_{j=0}^{\infty}\frac{(1/128)^j}{j!}
\le
\sum_{j=0}^{\infty}(1/128)^j
=
\frac{128}{127}.
\]

For the record tuple
\[
\mathcal O=(\mathcal O_1,\ldots,\mathcal O_n)\in\Omega^n,
\]
define the product likelihood ratio and its chi-square divergence by
\[
\Lambda_{\tau,n}(\mathcal O)
:=
\prod_{j=1}^n\Lambda_\tau(\mathcal O_j),
\qquad
\chi^2_{\tau,n}
:=
\int_{\Omega^n}(\Lambda_{\tau,n}-1)^2
\,dP_{\mathrm{null}}^{\otimes n},
\qquad \tau\in\mathbb R.
\]
The product likelihood ratio is bounded and is the density of
\(P_{\mathrm{par},\tau}^{\otimes n}\) relative to
\(P_{\mathrm{null}}^{\otimes n}\). Because each likelihood ratio integrates to one, independence implies
\[
\begin{aligned}
1+\chi^2_{\tau,n}
&=
\int_{\Omega^n}\Lambda_{\tau,n}^2
\,dP_{\mathrm{null}}^{\otimes n}\\
&=
\left(\int_\Omega\Lambda_\tau^2\,dP_{\mathrm{null}}\right)^n
=
(1+\chi^2_\tau)^n.
\end{aligned}
\]
Applying this identity at \(\tau=\tau_n\) to the preceding exponential bound yields
\[
\chi^2_{\tau_n,n}
\le\frac{128}{127}-1
=\frac1{127}.
\]
% lean: CausalSmith.Stat.LogoddsLowsmoothFrontier.parametricLaw_iid_tv

\item Total variation is the supremum of the absolute differences of measurable-event probabilities. For probability laws with a density, it equals half the integral of the absolute density difference. Cauchy--Schwarz therefore gives
\[
\begin{aligned}
\operatorname{TV}
\bigl(P_{\mathrm{null}}^{\otimes n},
      P_{\mathrm{shift}}^{\otimes n}\bigr)
&=
\frac12
\int_{\Omega^n}|\Lambda_{\tau_n,n}-1|
\,dP_{\mathrm{null}}^{\otimes n}\\
&\le
\frac12\sqrt{\chi^2_{\tau_n,n}}
\le
\frac12\sqrt{\frac1{127}}
\le
\frac1{20},
\end{aligned}
\]
where the last inequality follows from \(1/127\le1/100\).

Define the full sample-and-seed experiment laws by
\[
\mathbb P_{\mathrm{null}}
:=
P_{\mathrm{null}}^{\otimes n}\otimes\lambda,
\qquad
\mathbb P_{\mathrm{shift}}
:=
P_{\mathrm{shift}}^{\otimes n}\otimes\lambda.
\]
Their likelihood ratio is still \(\Lambda_{\tau_n,n}(\mathcal O)\). Integrating over the common independent seed, whose law has mass one, gives
\[
\operatorname{TV}
(\mathbb P_{\mathrm{null}},\mathbb P_{\mathrm{shift}})
=
\operatorname{TV}
\bigl(P_{\mathrm{null}}^{\otimes n},
      P_{\mathrm{shift}}^{\otimes n}\bigr)
\le\frac1{20}.
\]
The comparison effects are ordered:
\[
\theta(P_{\mathrm{null}})
=0
\le\tau_n
=\theta(P_{\mathrm{shift}}).
\]
% lean: CausalSmith.Stat.LogoddsLowsmoothFrontier.parametricLaw_experiment_tv

\item To bound the infimum in \cref{obj:def:length-objective}, first observe that the honest-procedure class is nonempty. Define
\[
I_{\mathrm{full}}(\mathcal O,U):=[-1/2,1/2],
\qquad
(\mathcal O,U)\in\Omega^n\times[0,1].
\]
By \cref{obj:ass:effect-envelope}, for every
\(P\in\mathcal M(\alpha,\beta)\),
\[
(P^{\otimes n}\otimes\lambda)
\{\theta(P)\in I_{\mathrm{full}}(\mathcal O,U)\}=1.
\]
The constant procedure is measurable and hence belongs to the class in
\cref{obj:def:honest-procedures}.

Now fix any
\[
I\in\mathcal A_n(\alpha,\beta).
\]
For
\(\omega=(\mathcal O,U)\in\Omega^n\times[0,1]\), define the measurable coverage events
\[
\begin{aligned}
\mathcal E_{\mathrm{null}}
&:=
\{(\mathcal O,U):0\in I(\mathcal O,U)\},\\
\mathcal E_{\mathrm{shift}}
&:=
\{(\mathcal O,U):\tau_n\in I(\mathcal O,U)\}.
\end{aligned}
\]
Both comparison laws belong to the ambient model, so
\cref{obj:ass:global-coverage} gives
\[
\mathbb P_{\mathrm{null}}(\mathcal E_{\mathrm{null}})
\ge\frac9{10},
\qquad
\mathbb P_{\mathrm{shift}}(\mathcal E_{\mathrm{shift}})
\ge\frac9{10}.
\]
% lean: CausalSmith.Stat.LogoddsLowsmoothFrontier.twoPoint_lengthObjective_lower

\item The total variation bound transfers the second coverage event to the null experiment:
\[
\begin{aligned}
\mathbb P_{\mathrm{null}}(\mathcal E_{\mathrm{shift}})
&\ge
\mathbb P_{\mathrm{shift}}(\mathcal E_{\mathrm{shift}})
-
\operatorname{TV}
(\mathbb P_{\mathrm{null}},\mathbb P_{\mathrm{shift}})\\
&\ge\frac9{10}-\frac1{20}.
\end{aligned}
\]
Using the union--intersection identity and the fact that the union has probability at most one,
\[
\begin{aligned}
\mathbb P_{\mathrm{null}}
(\mathcal E_{\mathrm{null}}\cap\mathcal E_{\mathrm{shift}})
&=
\mathbb P_{\mathrm{null}}(\mathcal E_{\mathrm{null}})
+
\mathbb P_{\mathrm{null}}(\mathcal E_{\mathrm{shift}})
-
\mathbb P_{\mathrm{null}}
(\mathcal E_{\mathrm{null}}\cup\mathcal E_{\mathrm{shift}})\\
&\ge
\frac9{10}+\frac9{10}-\frac1{20}-1
=\frac34.
\end{aligned}
\]

Write \(I(\omega)=I(\mathcal O,U)\). On the intersection event, membership of both effects implies
\[
\inf I(\omega)\le0\le\tau_n\le\sup I(\omega),
\qquad
|I(\omega)|
=
\sup I(\omega)-\inf I(\omega)
\ge\tau_n.
\]
These weak endpoint inequalities apply to every choice of endpoint inclusion. Off the intersection event, interval length is nonnegative. Thus
\[
|I(\omega)|\ge
\begin{cases}
\tau_n,
&\omega\in\mathcal E_{\mathrm{null}}\cap\mathcal E_{\mathrm{shift}},\\
0,
&\omega\notin\mathcal E_{\mathrm{null}}\cap\mathcal E_{\mathrm{shift}}.
\end{cases}
\]
Admissible intervals lie in the scalar effect region, so
\[
0\le |I(\omega)|\le1.
\]
Their measurable lengths are therefore integrable. Integrating the pointwise bound yields
\[
\mathbb E_{\mathbb P_{\mathrm{null}}}|I(\mathcal O,U)|
\ge
\tau_n\,
\mathbb P_{\mathrm{null}}
(\mathcal E_{\mathrm{null}}\cap\mathcal E_{\mathrm{shift}})
\ge\frac34\tau_n.
\]
% lean: CausalSmith.Stat.LogoddsLowsmoothFrontier.twoPoint_expectedLength_lower

\item Since \(P_{\mathrm{null}}\in\mathcal M_r(\alpha,\beta)\), its expected length is bounded above by the supremum over that slice:
\[
\sup_{P\in\mathcal M_r(\alpha,\beta)}
\mathbb E_{P^{\otimes n}\otimes\lambda}|I(\mathcal O,U)|
\ge
\mathbb E_{\mathbb P_{\mathrm{null}}}|I(\mathcal O,U)|
\ge\frac34\tau_n.
\]
This supremum is finite because every expected length is at most one. The bound holds for every honest procedure, and that class is nonempty. Taking its infimum as in \cref{obj:def:length-objective} gives
\[
J_n(\alpha,\beta;r)
\ge
\frac34\{\theta(P_{\mathrm{shift}})-\theta(P_{\mathrm{null}})\}
=
\frac34\tau_n
=
\frac3{16}n^{-1/2}.
\]
% lean: CausalSmith.Stat.LogoddsLowsmoothFrontier.parametric_null_floor
\end{enumerate}
\end{proof}

At zero radius, \cref{obj:lem:parametric-null-floor} gives a positive uncertainty floor for expected length under the null slice. The coverage requirement still ranges over the ambient model, so uncertainty about neighboring effects enters the criterion even when length is evaluated at zero effect. The statement also covers procedures using the independent randomization seed \leanref{sym:U}{\ensuremath{U}}, since these procedures belong to the class defining the minimax objective in \cref{obj:def:length-objective}.

Combining this parametric floor with the mixed and fair comparisons gives the lower side of the sharp expected-length profile. The upper side is attained by the literal interval output \leanref{sym:I^{\mathrm{up}}}{\ensuremath{I^{\mathrm{up}}_{n,\mathsf E,\mathsf N}(\alpha,\beta)}} defined in \cref{obj:def:upper-handle}, with its coverage and length guarantees from \cref{obj:thm:finite-honest-upper}. The resulting comparison holds throughout the public exponent domain \leanref{sym:\mathcal D}{\ensuremath{\mathcal D}} and for every admissible engine and naming policy.

\begin{theoremv}{thm:sharp-honest-frontier}[Sharp honest interval frontier]
There exist positive finite absolute constants \(c,C\) with the following guarantees. Write
\[
\mathcal D=\{(\alpha,\beta)\in\mathbb R^2:0<\beta<1/4,\ \beta<\alpha\le1\},
\]
and define the numerator and denominator exponents and the expected-length profile by
\[
\mathsf a(\alpha,\beta)
=\min\left\{\frac12,\frac{2(\alpha+\beta)}{2\alpha+2\beta+1}\right\},
\qquad
\mathsf b(\beta)=\frac{4\beta}{4\beta+1},
\qquad
\rho_n(\alpha,\beta;r)
=n^{-\mathsf a(\alpha,\beta)}+r n^{-\mathsf b(\beta)}.
\]

For every admissible rational interval-arithmetic engine \(\mathsf E\) and every admissible Borel policy \(\mathsf N\) satisfying the quantitative contracts of \cref{obj:def:arithmetic-engine,obj:def:dyadic-name-convention} and selecting valid names for the public exponents and observed covariates, consider parameters satisfying:
\begin{itemize}
\item \textbf{(Exponent domain.)} \((\alpha,\beta)\in\mathcal D\).
\item \textbf{(Sample size.)} \(n\in\mathbb N\) and \(n\ge1\).
\item \textbf{(Radius.)} \(0\le r\le1/2\).
\end{itemize}
Let \(I^\star_{n,\mathsf E,\mathsf N}(\alpha,\beta)
=I^{\mathrm{up}}_{n,\mathsf E,\mathsf N}(\alpha,\beta)\) denote the interval returned by the fixed-budget upper construction. With \(\mathcal A_n(\alpha,\beta)\) as in \cref{obj:def:honest-procedures}, define the minimax expected length by
\[
J_n(\alpha,\beta;r)
=
\inf_{I\in\mathcal A_n(\alpha,\beta)}
\sup_{P\in\mathcal M_r(\alpha,\beta)}
\mathbb E_{P^{\otimes n}\otimes\lambda}
\bigl|I(\mathcal O,U)\bigr|,
\]
where \(\mathcal M_r(\alpha,\beta)\) is the prescribed effect-radius slice, \(\mathcal O\) is the sample of original records, and \(U\) is an independent seed with uniform law \(\lambda\) on the unit interval. Then
\[
c\rho_n(\alpha,\beta;r)
\le J_n(\alpha,\beta;r)
\le
\sup_{P\in\mathcal M_r(\alpha,\beta)}
\mathbb E_{P^{\otimes n}}
\bigl|I^\star_{n,\mathsf E,\mathsf N}(\alpha,\beta)\bigr|
\le C\rho_n(\alpha,\beta;r).
\]
The same constants apply to every admissible engine and naming policy.

For every such engine and policy, every \((\alpha,\beta)\in\mathcal D\), and every \(n\ge1\), the actual name-dependent output is a deterministic, radius-independent Borel interval with global \(90\%\) coverage over \(\mathcal M(\alpha,\beta)\), the prescribed ambient model. It is a nonempty closed connected interval contained in \([-1/2,1/2]\), with rational endpoints, returned after the two prescribed finite precision budgets and finite rational arithmetic. Its inputs consist of the fixed engine, the public exponents and their valid public names, and the \(n\) original records and their valid covariate names. Fixing the explicit admissible engine \(\mathsf E_0\) and dyadic-floor naming policy \(\mathsf N_0\) gives the concrete attaining interval
\[
I_n^\star
=I^{\mathrm{up}}_{n,\mathsf E_0,\mathsf N_0},
\qquad
c_\star=c,\qquad C_\star=C,\qquad n_0=1,
\]
where \(c_\star,C_\star\) are the comparison constants and \(n_0\) is the sample-size threshold.

Moreover, for every \((\alpha,\beta)\in\mathcal D\), every \(0\le r\le1/2\), every \(n\ge1\), and every \(I\in\mathcal A_n(\alpha,\beta)\), including procedures using independent randomization,
\[
c\rho_n(\alpha,\beta;r)
\le
\sup_{P\in\mathcal M_r(\alpha,\beta)}
\mathbb E_{P^{\otimes n}\otimes\lambda}
\bigl|I(\mathcal O,U)\bigr|.
\]

For every \((\alpha,\beta)\in\mathcal D\), the exponent difference is positive and satisfies
\[
\mathsf a(\alpha,\beta)-\mathsf b(\beta)>0,
\qquad
\mathsf a(\alpha,\beta)-\mathsf b(\beta)
=
\begin{cases}
\displaystyle
\frac{2(\alpha-\beta)}
{(2\alpha+2\beta+1)(4\beta+1)},
&\alpha+\beta\le1/2,\\[6pt]
\displaystyle
\frac{1-4\beta}{2(4\beta+1)},
&\alpha+\beta\ge1/2.
\end{cases}
\]
At \(\alpha+\beta=1/2\), the two displayed expressions agree.

Finally, for every nonempty compact set \(K\subseteq\mathcal D\), the same constants and threshold satisfy the compact-uniform conditions
\[
0<\inf_{(\alpha,\beta)\in K}c,
\qquad
\sup_{(\alpha,\beta)\in K}C<\infty,
\qquad
\sup_{(\alpha,\beta)\in K}1=1.
\]
\end{theoremv}

The two terms in \cref{obj:thm:sharp-honest-frontier} distinguish uncertainty in the numerator coordinate from effect-scaled uncertainty in the denominator coordinate. The mixed alternatives supply the radius-independent lower contribution, together with the parametric floor when the numerator exponent reaches \(1/2\); the fair alternatives supply the contribution proportional to the evaluation radius. Thus the null slice retains the numerator rate, whereas every fixed positive radius eventually makes the slower denominator rate dominant. The positive exponent difference stated in the theorem establishes this ordering throughout the prescribed domain. A single globally honest interval attains the comparison simultaneously across radii.

The following synthesis collects the statistical comparison, the finite-budget realization from \cref{obj:lem:fixed-budget-realization}, and the concrete engine supplied by \cref{obj:lem:concrete-arithmetic-engine}. It identifies the attaining procedure with the actual output for each admissible engine and naming policy.

For the concrete sequence below, the inputs are the sample size, public exponent representations, and represented original records, together with the independently fixed arithmetic engine and naming policy. The same radius-independent sequence is evaluated over each radius subclass.

\begin{theoremv}{thm:full-honest-frontier-solution}[Full honest frontier solution]
There exist absolute constants \(c,C\in\mathbb R\), with \(0<c\) and \(0<C\), for which the following assertions hold. Write
\[
\mathcal D=\{(\alpha,\beta)\in\mathbb R^2:0<\beta<1/4,\ \beta<\alpha\le1\},
\qquad
\mathsf a(\alpha,\beta)=
\min\left\{\frac12,\frac{2(\alpha+\beta)}{2\alpha+2\beta+1}\right\},
\qquad
\mathsf b(\beta)=\frac{4\beta}{4\beta+1},
\]
and define the expected-length profile by
\[
\rho_n(\alpha,\beta;r)
=d_n(\alpha,\beta;r)
=n^{-\mathsf a(\alpha,\beta)}
+r n^{-\mathsf b(\beta)}.
\]

For every admissible rational interval-arithmetic engine \(\mathsf E\), every admissible Borel naming policy \(\mathsf N\), both satisfying the quantitative contracts of \cref{obj:def:arithmetic-engine,obj:def:dyadic-name-convention}, every \((\alpha,\beta)\in\mathcal D\), and every integer \(n\ge1\), take
\[
I^\star_{n,\mathsf E,\mathsf N}(\alpha,\beta)
=I^{\mathrm{up}}_{n,\mathsf E,\mathsf N}(\alpha,\beta),
\]
the literal output of the fixed-budget represented upper construction. This procedure belongs to the globally \(90\%\)-honest Borel class \(\mathcal A_n(\alpha,\beta)\) of \cref{obj:def:honest-procedures}. It receives neither the radius nor unknown nuisance functions and terminates on every valid named input within the prescribed rank and endpoint budgets of \cref{obj:def:resolutions,obj:def:upper-handle,obj:lem:fixed-budget-realization}. For every sample and randomization seed its output is a closed interval \([l,u]\) with
\[
l,u\in\mathbb Q,
\qquad -\frac12\le l\le u\le\frac12.
\]
For every \(r\in[0,1/2]\), the same constants satisfy
\[
c\rho_n(\alpha,\beta;r)
\le J_n(\alpha,\beta;r)
\le
\sup_{P\in\mathcal M_r(\alpha,\beta)}
E_P\!\left[
\left|I^\star_{n,\mathsf E,\mathsf N}(\alpha,\beta)(\mathcal O,U)\right|
\right]
\le C\rho_n(\alpha,\beta;r).
\]
Here \(J_n(\alpha,\beta;r)\) is the minimax expected length of \cref{obj:def:length-objective}, and the expectation uses \(n\) original records drawn independently from \(P\), together with the independent randomization seed \(U\).

In particular, the explicit concrete engine \(\mathsf E_0\) and the dyadic-floor naming policy \(\mathsf N_0\) give the sequence
\[
I_n^\star(\alpha,\beta)
=I^{\mathrm{up}}_{n,\mathsf E_0,\mathsf N_0}(\alpha,\beta).
\]
For every \((\alpha,\beta)\in\mathcal D\) and \(n\ge1\), this sequence is globally honest, has the rational closed ordered output described above at every input, and satisfies the same displayed comparison for every \(r\in[0,1/2]\).

Moreover, for every \((\alpha,\beta)\in\mathcal D\), every \(r\in[0,1/2]\), every integer \(n\ge1\), and every \(I\in\mathcal A_n(\alpha,\beta)\), including procedures using independent randomization,
\[
c\rho_n(\alpha,\beta;r)
\le
\sup_{P\in\mathcal M_r(\alpha,\beta)}
E_P\!\left[|I(\mathcal O,U)|\right].
\]

The exponent difference is positive throughout \(\mathcal D\) and obeys
\[
\mathsf a(\alpha,\beta)-\mathsf b(\beta)=
\begin{cases}
\displaystyle
\frac{2(\alpha-\beta)}
{(2\alpha+2\beta+1)(4\beta+1)},
&\alpha+\beta\le1/2,\\[6pt]
\displaystyle
\frac{1-4\beta}{2(4\beta+1)},
&\alpha+\beta\ge1/2.
\end{cases}
\]
At \(\alpha+\beta=1/2\), the two displayed expressions agree.

Finally, select the lower and upper comparison constants and sample-size threshold as
\[
c_\star(\alpha,\beta)=c,\qquad
C_\star(\alpha,\beta)=C,\qquad
n_0(\alpha,\beta)=1.
\]
For every nonempty compact \(\mathcal K\subset\mathcal D\), these choices satisfy
\[
0<
\inf_{(\alpha,\beta)\in\mathcal K}c_\star(\alpha,\beta),
\qquad
\sup_{(\alpha,\beta)\in\mathcal K}C_\star(\alpha,\beta)<\infty,
\qquad
\sup_{(\alpha,\beta)\in\mathcal K}n_0(\alpha,\beta)=1.
\]
\end{theoremv}

The statistical and computational guarantees in \cref{obj:thm:full-honest-frontier-solution} concern the same interval: its rational endpoint construction preserves global coverage and the sharp expected-length comparison at every sample size. The original-record lower bound covers the entire honest class, including randomized procedures, while each admissible implementation attains the comparison with the same absolute constants. Selecting these constants and the threshold \(n_0=1\) also gives uniform comparisons on every nonempty compact subset \leanref{sym:\mathcal K}{\ensuremath{\mathcal K}} of the public exponent domain.


\section{Verification note}

The supplied theorem and lemma statements and their proofs are machine-checked in Lean 4 under their stated assumptions. The frozen formal layer covers identification and projection control in \cref{obj:lem:observable-odds,obj:lem:projection-control}; rational arithmetic and finite-budget realization in \cref{obj:lem:concrete-arithmetic-engine,obj:lem:fixed-budget-realization}; local implicit-function arguments, exact calibration, and model support in \cref{obj:lem:scalar-implicit-function,obj:lem:exact-calibrations,obj:lem:calibrated-support}; and original-record mixture bounds and the parametric floor in \cref{obj:lem:original-record-mixtures,obj:lem:parametric-null-floor}. The resulting coverage, expected-length comparison, and attaining construction are consolidated in \cref{obj:thm:finite-honest-upper,obj:thm:sharp-honest-frontier,obj:thm:full-honest-frontier-solution,obj:thm:full-honest-frontier-answer}.

The model restrictions defining \cref{obj:def:model} and the primitive contracts stated in the corresponding results are hypotheses of these implications. The global coverage criterion defining \cref{obj:def:honest-procedures} specifies the admissible statistical procedures. Admissibility of the concrete arithmetic engine, finite termination and endpoint enclosure for the represented construction, global coverage, and the sharp comparison for the length objective in \cref{obj:def:length-objective} are established conclusions. The recorded external-dependency sets for these declarations are empty.

For the represented construction, valid names satisfying \cref{obj:def:dyadic-name-convention} supply enclosing dyadic intervals for the public exponents and observed covariates. Together with the rational primitive contracts in \cref{obj:def:arithmetic-engine}, they support the rank selection and endpoint realization in \cref{obj:lem:fixed-budget-realization}. The explicit arithmetic witness in \cref{obj:lem:concrete-arithmetic-engine} supplies an admissible engine, and its combination with the dyadic-floor naming policy yields the concrete attaining interval identified in \cref{obj:thm:full-honest-frontier-solution}.

The formal source tree is
\texttt{CausalSmith/CausalSmith/Stat/STAT\_LogoddsLowsmoothFrontier\_Research/},
with its aggregate module at the corresponding \texttt{.lean} path. The
declaration-to-paper mapping is
\texttt{CausalSmith/doc/presentation/stat\_logodds\_lowsmooth\_frontier\_v1/presentation\_crosswalk.json};
it records repository source revision
\texttt{1efc342d4a68f9860ec0aa7692071d90a9230018}. In a checkout of that revision,
run
\begin{verbatim}
cd CausalSmith
lake build CausalSmith.Stat.STAT_LogoddsLowsmoothFrontier_Research
\end{verbatim}
The Lean toolchain is pinned in \texttt{CausalSmith/lean-toolchain}, and package
dependencies are pinned in \texttt{CausalSmith/lake-manifest.json} under the
project configuration \texttt{CausalSmith/lakefile.toml}. The same presentation
directory contains \texttt{lean\_snippets.json}, \texttt{references.bib},
\texttt{paper\_macros.tex}, and the crosswalk needed to reproduce the manuscript
links and inspect the exact declarations rendered in each numbered environment.


\section{Proofs of the main results}\label{sec:deferred-proofs}

\begin{proof}[Proof of \cref{obj:thm:full-honest-frontier-answer}]
\begin{enumerate}
\item Choose the constants \(c,C\) supplied by
\cref{obj:thm:full-honest-frontier-solution}, with
\[
0<c,\qquad 0<C.
\]
Its universal construction guarantee applies simultaneously to every admissible engine and naming policy, every \((\alpha,\beta)\in\mathcal D\), every integer \(n\ge1\), and every \(r\in[0,1/2]\), with these same constants. Its compact-uniform guarantee, elbow identities, and converse for every globally honest procedure, including independently randomized procedures, give the corresponding assertions of the present theorem. In particular, its choices
\[
c_\star=c,\qquad C_\star=C,\qquad n_0=1
\]
give, for every nonempty compact \(K\subseteq\mathcal D\),
\[
0<\inf_{(\alpha,\beta)\in K}c_\star,\qquad
\sup_{(\alpha,\beta)\in K}C_\star<\infty,\qquad
\sup_{(\alpha,\beta)\in K}n_0=1.
\]
% lean: CausalSmith.Stat.LogoddsLowsmoothFrontier.full_honest_frontier_solution

\item Fix \((\alpha,\beta)\in\mathcal D\) and an integer \(n\ge1\). The concrete sequence in
\cref{obj:thm:full-honest-frontier-solution} is exactly
\[
I_n^\star(\alpha,\beta)
=I^{\mathrm{up}}_{n,\mathsf E_0,\mathsf N_0}(\alpha,\beta),
\]
as specified in \cref{obj:def:frontier-handle}. That guarantee yields
\[
I_n^\star(\alpha,\beta)\in\mathcal A_n(\alpha,\beta).
\]
For every sample and seed, it also gives rational endpoints \(l,u\) such that
\[
I_n^\star(\alpha,\beta)(\mathcal O,U)=[l,u],
\qquad
-\frac12\le l\le u\le\frac12.
\]
% lean: CausalSmith.Stat.LogoddsLowsmoothFrontier.full_honest_frontier_solution

\item Fix \(r\in[0,1/2]\). Expanding the diagnostic envelope of
\cref{obj:def:diagnostic-envelope} and its identification with the rate profile in
\cref{obj:def:frontier-handle} gives
\[
\rho_n(\alpha,\beta;r)=d_n(\alpha,\beta;r)
=
n^{-\min\{1/2,\,2(\alpha+\beta)/(2\alpha+2\beta+1)\}}
+r\,n^{-4\beta/(4\beta+1)}.
\]
Thus the concrete comparison supplied by
\cref{obj:thm:full-honest-frontier-solution}, expressed with this explicit profile, is
\[
c\,\rho_n(\alpha,\beta;r)
\le J_n(\alpha,\beta;r)
\le
\sup_{P\in\mathcal M_r(\alpha,\beta)}
E_{P^{\otimes n}\otimes\lambda}
\bigl|I_n^\star(\alpha,\beta)(\mathcal O,U)\bigr|
\le C\,\rho_n(\alpha,\beta;r).
\]
Together with the assertions retained in the first step, this establishes all the stated guarantees with the same constants.
% lean: CausalSmith.Stat.LogoddsLowsmoothFrontier.full_honest_frontier_answer
\end{enumerate}
\end{proof}

\begin{proof}[Proof of \cref{obj:thm:finite-honest-upper}]
We first establish the bounds for \(n\ge2\). Fix an admissible arithmetic engine \(\mathsf E\), an admissible naming policy \(\mathsf N\), and public exponents satisfying \(0<\beta<1/4\) and \(\beta<\alpha\le1\).

\begin{enumerate}
\item Define the retained ranks, smoothness sum, and rate exponents by
\[
\begin{gathered}
k_C=k_C(n,\alpha,\beta;\mathsf E,\mathsf N),
\qquad
k_S=k_S(n,\beta;\mathsf E,\mathsf N),
\qquad s=\alpha+\beta,\\
\mathsf a(\alpha,\beta)
=\min\left\{\frac12,\frac{2s}{2s+1}\right\},
\qquad
\mathsf b(\beta)=\frac{4\beta}{4\beta+1}.
\end{gathered}
\]
These are the ranks and exponents of
\cref{obj:def:resolutions,obj:def:diagnostic-envelope}.
For \(k\ge1\), put
\[
V_n(k)=\frac{10}{n}+\frac{4k}{n(n-1)},
\qquad
b_C=8k_C^{-s}+\sqrt{20V_n(k_C)},
\qquad
b_S=25k_S^{-2\beta}+\sqrt{20V_n(k_S)}.
\]
Here \(b_C,b_S\) are the coordinate error radii, as in
\cref{obj:def:upper-handle}.

By \cref{obj:lem:fixed-budget-realization}, with target resolutions
\[
R_C=n^{2/(2s+1)},
\qquad R_S=n^{2/(4\beta+1)},
\]
the ranks satisfy
\[
1\le R_C\le k_C\le3R_C,
\qquad
1\le R_S\le k_S\le3R_S.
\]
In particular, both ranks are positive and are fixed before sampling.

For either pair
\[
(\gamma_{\mathrm{bal}},k_{\mathrm{bal}})
\in\{(s,k_C),(2\beta,k_S)\},
\]
the lower rank bound and \(\gamma_{\mathrm{bal}}>0\) give
\[
k_{\mathrm{bal}}^{-\gamma_{\mathrm{bal}}}
\le
n^{-2\gamma_{\mathrm{bal}}/(2\gamma_{\mathrm{bal}}+1)}.
\]
Also \(n(n-1)\ge n^2/2>0\), so the upper rank bound yields
\[
20V_n(k_{\mathrm{bal}})
\le
\frac{200}{n}
+
480\,\frac{n^{2/(2\gamma_{\mathrm{bal}}+1)}}{n^2}
=
200(n^{-1/2})^2
+
480\left(
n^{-2\gamma_{\mathrm{bal}}/(2\gamma_{\mathrm{bal}}+1)}
\right)^2.
\]
Both powers are nonnegative. The last expression is bounded by
\[
\left[
40\left(
n^{-1/2}
+
n^{-2\gamma_{\mathrm{bal}}/(2\gamma_{\mathrm{bal}}+1)}
\right)
\right]^2,
\]
because expansion of this square gives coefficients \(1600\) on both squared terms and a nonnegative cross term. Taking square roots therefore gives
\[
\sqrt{20V_n(k_{\mathrm{bal}})}
\le
40\left(
n^{-1/2}
+
n^{-2\gamma_{\mathrm{bal}}/(2\gamma_{\mathrm{bal}}+1)}
\right).
\]
Each power on the right is at most the power with exponent
\(-\min\{1/2,2\gamma_{\mathrm{bal}}/(2\gamma_{\mathrm{bal}}+1)\}\).
Since \(0<\beta<1/4\),
\[
\frac{4\beta}{4\beta+1}<\frac12.
\]
Consequently,
\[
\begin{aligned}
b_C
&\le88\,n^{-\mathsf a(\alpha,\beta)}
\le128\,n^{-\mathsf a(\alpha,\beta)},\\
b_S
&\le105\,n^{-\mathsf b(\beta)}
\le128\,n^{-\mathsf b(\beta)}.
\end{aligned}
\]
The exponents are positive, and \(\mathsf a(\alpha,\beta)\le1/2\le1\). Thus
\[
n^{-\mathsf a(\alpha,\beta)}\le1,
\qquad
n^{-\mathsf b(\beta)}\le1,
\qquad
\frac1n\le n^{-\mathsf a(\alpha,\beta)}.
\]
% lean: CausalSmith.Stat.LogoddsLowsmoothFrontier.coordinate_radii_power_bounds
% lean: CausalSmith.Stat.LogoddsLowsmoothFrontier.diagnostic_power_bounds

\item With the denominator floor and sensitivity coefficients defined by
\[
s_0=\frac1{512},
\qquad
\mathsf L_{\mathrm{sens}}=\frac{2\exp(1/2)}{s_0},
\qquad
\mathsf H_{\mathrm{sens}}=\frac{2\exp(1/2)}{s_0^2},
\]
choose
\[
K_0
=
128\mathsf L_{\mathrm{sens}}
+32768\mathsf H_{\mathrm{sens}}
+128\mathsf H_{\mathrm{sens}}\exp(1/2)+2.
\]
All its summands are nonnegative, so \(K_0\ge1>0\).
For every \(r\in[0,1/2]\), the preceding bounds imply
\[
\begin{aligned}
&\mathsf L_{\mathrm{sens}}b_C
+\mathsf H_{\mathrm{sens}}b_S
   \{\exp(1/2)r+2b_C\}
+\frac2n\\
&\quad\le
128\mathsf L_{\mathrm{sens}}n^{-\mathsf a(\alpha,\beta)}
+128\mathsf H_{\mathrm{sens}}n^{-\mathsf b(\beta)}
 \{\exp(1/2)r+256n^{-\mathsf a(\alpha,\beta)}\}
+2n^{-\mathsf a(\alpha,\beta)}\\
&\quad\le
(128\mathsf L_{\mathrm{sens}}
 +32768\mathsf H_{\mathrm{sens}}+2)
 n^{-\mathsf a(\alpha,\beta)}
+
128\mathsf H_{\mathrm{sens}}\exp(1/2)\,
r n^{-\mathsf b(\beta)}\\
&\quad\le
K_0\left\{
n^{-\mathsf a(\alpha,\beta)}
+r n^{-\mathsf b(\beta)}
\right\}
=
K_0d_n(\alpha,\beta;r).
\end{aligned}
\]
The second inequality uses
\(n^{-\mathsf b(\beta)}n^{-\mathsf a(\alpha,\beta)}
\le n^{-\mathsf a(\alpha,\beta)}\).
The last inequality follows by adding the nonnegative terms that complete the common coefficient \(K_0\).
This choice depends solely on the fixed denominator floor.
% lean: CausalSmith.Stat.LogoddsLowsmoothFrontier.radius_sensitivity_rate_assembly
% lean: CausalSmith.Stat.LogoddsLowsmoothFrontier.upper_radii_uniform_rate

\item Fix \(P\in\mathcal M(\alpha,\beta)\). Write
\[
\mathbb P_{P,n}=P^{\otimes n}\otimes\lambda,
\qquad
\mathbb E_{P,n}[\cdot]=\int(\cdot)\,d\mathbb P_{P,n},
\]
where \(\lambda\) is the uniform law of the independent seed. Variances below are under this same probability law.

We first control the unconditional absolute mean of the numerator estimate. Define, for \(x\in[0,1]\),
\[
\begin{gathered}
e(P,x)=\ell\{g(P,x)\},\\
\mu_0(P,x)=\ell\{\nu(P,x)\},
\qquad
\mu_1(P,x)=\ell\{\nu(P,x)+\theta(P)\},\\
m(P,x)=E_P[Y\mid X=x]
=(1-e(P,x))\mu_0(P,x)+e(P,x)\mu_1(P,x).
\end{gathered}
\]
The homogeneous relation and native-logit bounds are those of
\cref{obj:def:model}. The logistic derivative satisfies
\[
0<\ell'(t_{\ell})
=\ell(t_{\ell})\{1-\ell(t_{\ell})\}
\le\frac14
\qquad(t_{\ell}\in\mathbb R).
\]
The derivative bound on the interval between any two real arguments gives
\[
|\ell(t_{\ell})-\ell(t_{\ell}')|
\le\frac14|t_{\ell}-t_{\ell}'|
\qquad(t_{\ell},t_{\ell}'\in\mathbb R).
\]
Hence, for \(x,z\in[0,1]\),
\[
\begin{aligned}
|e(P,x)-e(P,z)|
&\le\frac12|x-z|^\alpha
\le\frac12|x-z|^\beta,\\
|\mu_0(P,x)-\mu_0(P,z)|
&\le\frac12|x-z|^\beta,\\
|\mu_1(P,x)-\mu_1(P,z)|
&\le\frac12|x-z|^\beta,\\
|\mu_1(P,z)-\mu_0(P,z)|
&\le\frac14|\theta(P)|\le\frac18.
\end{aligned}
\]
The comparison of the two covariate powers uses
\(|x-z|\le1\) and \(\beta<\alpha\); it also holds at \(x=z\), since both exponents are positive.

The marginal mean has the decomposition
\[
\begin{aligned}
m(P,x)-m(P,z)
={}&(1-e(P,x))\{\mu_0(P,x)-\mu_0(P,z)\}\\
&+e(P,x)\{\mu_1(P,x)-\mu_1(P,z)\}\\
&+\{e(P,x)-e(P,z)\}
  \{\mu_1(P,z)-\mu_0(P,z)\}.
\end{aligned}
\]
Since \(0<e(P,x)<1\), the preceding moduli imply
\[
\begin{aligned}
|m(P,x)-m(P,z)|
&\le
\{1-e(P,x)\}\frac12|x-z|^\beta
+e(P,x)\frac12|x-z|^\beta
+\frac12|x-z|^\beta\frac18\\
&=\frac9{16}|x-z|^\beta.
\end{aligned}
\]
Thus
\[
[e(P,\cdot)]_\alpha\le\frac12,
\qquad
[m(P,\cdot)]_\beta\le\frac9{16}.
\]
These functions are continuous: the propensity and marginal mean are sums of continuous conditional cells.

Define the empirical treated-outcome mean by
\[
\overline{AY}_n=\frac1n\sum_{i=1}^n A_iY_i.
\]
The coordinate definition in \cref{obj:def:coordinate-estimates} gives
\[
\widehat C_n
=\overline{AY}_n-\mathbb U_{n,k_C}(A,Y),
\qquad
\mathbb E_{P,n}\overline{AY}_n=E_P[AY].
\]
Using the projection expectation and residual identities from
\cref{obj:lem:projection-control}, we obtain
\[
\mathbb E_{P,n}\widehat C_n-C(P)
=
\left\langle
e(P,\cdot)-\Pi_{k_C}e(P,\cdot),
m(P,\cdot)-\Pi_{k_C}m(P,\cdot)
\right\rangle_{L^2(\lambda)}.
\]
Cauchy--Schwarz and the approximation bound in the same cited statement yield
\[
\begin{aligned}
|\mathbb E_{P,n}\widehat C_n-C(P)|
&\le
\|e(P,\cdot)-\Pi_{k_C}e(P,\cdot)\|_2
\|m(P,\cdot)-\Pi_{k_C}m(P,\cdot)\|_2\\
&\le
25\left(\frac12\right)\left(\frac9{16}\right)
k_C^{-(\alpha+\beta)}
=\frac{225}{32}k_C^{-s}
\le8k_C^{-s}.
\end{aligned}
\]
Here \(0<\alpha\le1\), \(0<\beta<1\), and \(k_C\ge1\), so the approximation bound applies throughout the stated range.

Independence of the records and \(0\le AY\le1\) give
\[
\operatorname{Var}_{P,n}(\overline{AY}_n)
=\frac{\operatorname{Var}_P(AY)}n
\le\frac1n.
\]
Also, \cref{obj:lem:projection-control} supplies
\[
\operatorname{Var}_{P,n}\{\mathbb U_{n,k_C}(A,Y)\}
\le\frac4n+\frac{2k_C}{n(n-1)}.
\]
The variance identities for a sum and a difference give
\[
\begin{aligned}
\operatorname{Var}_{P,n}(\widehat C_n)
={}&2\operatorname{Var}_{P,n}(\overline{AY}_n)
+2\operatorname{Var}_{P,n}\{\mathbb U_{n,k_C}(A,Y)\}\\
&-\operatorname{Var}_{P,n}
 \{\overline{AY}_n+\mathbb U_{n,k_C}(A,Y)\}\\
\le{}&\frac2n+2\left\{\frac4n+\frac{2k_C}{n(n-1)}\right\}
=V_n(k_C),
\end{aligned}
\]
where the inequality uses nonnegativity of the last variance.
The finite-rank statistics are square integrable, so all these identities and subsequent expectations are well defined.

Cauchy--Schwarz on the probability space gives
\[
\mathbb E_{P,n}
\left|\widehat C_n-\mathbb E_{P,n}\widehat C_n\right|
\le
\sqrt{\operatorname{Var}_{P,n}(\widehat C_n)}.
\]
Consequently,
\[
\begin{aligned}
\mathbb E_{P,n}|\widehat C_n|
&\le
|\mathbb E_{P,n}\widehat C_n|
+\sqrt{\operatorname{Var}_{P,n}(\widehat C_n)}\\
&\le
|C(P)|+8k_C^{-s}+\sqrt{V_n(k_C)}
\le |C(P)|+b_C.
\end{aligned}
\]
For the last step, \(V_n(k_C)>0\) and
\(\sqrt{V_n(k_C)}\le\sqrt{20V_n(k_C)}\).

By \cref{obj:lem:observable-odds},
\[
C(P)=\{\exp(\theta(P))-1\}S(P),
\qquad S(P)\ge s_0.
\]
The product defining \(S(P)\) is at most one pointwise, so
\[
S(P)
=\int p_{10}(P,x)p_{01}(P,x)\,\lambda(dx)\le1,
\]
where the two conditional off-diagonal probabilities are
\[
p_{10}(P,x)=P(A=1,Y=0\mid X=x),
\qquad
p_{01}(P,x)=P(A=0,Y=1\mid X=x).
\]
Since \(|\theta(P)|\le1/2\), the exponential derivative is bounded by
\(\exp(1/2)\) on the segment between \(0\) and \(\theta(P)\). Therefore
\[
|\exp(\theta(P))-1|
\le\exp(1/2)|\theta(P)|,
\qquad
|C(P)|\le\exp(1/2)|\theta(P)|.
\]
In particular, for \(P\in\mathcal M_r(\alpha,\beta)\),
\[
\mathbb E_{P,n}|\widehat C_n|
\le\exp(1/2)r+b_C.
\]
% lean: CausalSmith.Stat.LogoddsLowsmoothFrontier.cHat_integral_abs_le
% lean: CausalSmith.Stat.LogoddsLowsmoothFrontier.covarianceCoordinate_abs_le_effect

\item We next bound the interval length on every sample. Use the confidence endpoints and clipped inversion of \cref{obj:def:upper-handle}:
\[
\begin{gathered}
c_\pm=\widehat C_n\pm b_C,
\qquad
s_\pm=\operatorname{clip}_{[s_0,1]}(\widehat S_n\pm b_S),\\
F(c_{\mathrm{rect}},s_{\mathrm{rect}})
=
\log\!\left(
\operatorname{clip}_{[\exp(-1/2),\exp(1/2)]}
\left(1+\frac{c_{\mathrm{rect}}}{s_{\mathrm{rect}}}\right)
\right)
\quad
(c_{\mathrm{rect}}\in\mathbb R,\ s_{\mathrm{rect}}>0),\\
L=\min\{F(c_-,s_-),F(c_-,s_+)\},
\qquad
R=\max\{F(c_+,s_-),F(c_+,s_+)\}.
\end{gathered}
\]
Clipping is a composition of minimum and maximum with fixed constants, each of which contracts absolute differences. Thus
\[
s_0\le s_-,s_+,
\qquad
|s_--s_+|
\le|(\widehat S_n-b_S)-(\widehat S_n+b_S)|
=2b_S.
\]

For real numbers satisfying
\(\exp(-1/2)\le v_{\log}\le w_{\log}\), the inequality
\(\log t\le t-1\) for \(t>0\) gives
\[
\log w_{\log}-\log v_{\log}
=
\log\frac{w_{\log}}{v_{\log}}
\le
\frac{w_{\log}-v_{\log}}{v_{\log}}
\le
\exp(1/2)(w_{\log}-v_{\log}).
\]
Interchanging the two arguments handles the reverse order. Applying this bound to the clipped arguments, and using contraction of the exponential-range clipping, gives
\[
\left|
F(c_{\mathrm{rect}},s_{\mathrm{rect}})
-
F(c_{\mathrm{rect}}',s_{\mathrm{rect}}')
\right|
\le
\exp(1/2)
\left|
\frac{c_{\mathrm{rect}}}{s_{\mathrm{rect}}}
-
\frac{c_{\mathrm{rect}}'}{s_{\mathrm{rect}}'}
\right|
\]
for all real numerators and positive denominators.

For the rectangle points with
\[
c_{\mathrm{rect}},c_{\mathrm{rect}}'\in[c_-,c_+],
\qquad
s_{\mathrm{rect}},s_{\mathrm{rect}}'\in\{s_-,s_+\},
\]
we have
\[
|c_{\mathrm{rect}}-c_{\mathrm{rect}}'|\le2b_C,
\qquad
|c_{\mathrm{rect}}'|\le|\widehat C_n|+b_C,
\qquad
|s_{\mathrm{rect}}-s_{\mathrm{rect}}'|\le2b_S.
\]
Inserting the cross quotient with numerator
\(c_{\mathrm{rect}}'\) and denominator \(s_{\mathrm{rect}}\) yields
\[
\begin{aligned}
\left|
\frac{c_{\mathrm{rect}}}{s_{\mathrm{rect}}}
-
\frac{c_{\mathrm{rect}}'}{s_{\mathrm{rect}}'}
\right|
&\le
\frac{|c_{\mathrm{rect}}-c_{\mathrm{rect}}'|}
     {s_{\mathrm{rect}}}
+
\frac{|c_{\mathrm{rect}}'|
      |s_{\mathrm{rect}}-s_{\mathrm{rect}}'|}
     {s_{\mathrm{rect}}s_{\mathrm{rect}}'}\\
&\le
\frac{2b_C}{s_0}
+
\frac{2b_S(|\widehat C_n|+b_C)}{s_0^2}.
\end{aligned}
\]
It follows that every pairing of an upper-numerator corner and a lower-numerator corner satisfies the same difference bound. Choosing the corners attaining the maximum \(R\) and minimum \(L\), including either choice in a tie, gives
\[
R-L
\le
\mathsf L_{\mathrm{sens}}b_C
+
\mathsf H_{\mathrm{sens}}b_S(|\widehat C_n|+b_C).
\]
By \cref{obj:lem:fixed-budget-realization}, the literal rational output has additional length at most \(2/n\). Hence, pointwise,
\[
\left|
I^{\mathrm{up}}_{n,\mathsf E,\mathsf N}(\alpha,\beta)(\mathcal O,U)
\right|
\le
\mathsf L_{\mathrm{sens}}b_C
+
\mathsf H_{\mathrm{sens}}b_S(|\widehat C_n|+b_C)
+\frac2n.
\]
% lean: CausalSmith.Stat.LogoddsLowsmoothFrontier.inversion_four_corner_width_le
% lean: CausalSmith.Stat.LogoddsLowsmoothFrontier.upperOutput_length_le

\item For \(P\in\mathcal M_r(\alpha,\beta)\), integrate the pointwise bound. The ranks and radii are fixed before sampling, and the numerator estimate is integrable. Therefore
\[
\begin{aligned}
&\mathbb E_{P,n}
\left|
I^{\mathrm{up}}_{n,\mathsf E,\mathsf N}(\alpha,\beta)(\mathcal O,U)
\right|\\
&\quad\le
\mathsf L_{\mathrm{sens}}b_C
+
\mathsf H_{\mathrm{sens}}b_S
 \{\mathbb E_{P,n}|\widehat C_n|+b_C\}
+\frac2n\\
&\quad\le
\mathsf L_{\mathrm{sens}}b_C
+
\mathsf H_{\mathrm{sens}}b_S
 \{\exp(1/2)r+2b_C\}
+\frac2n\\
&\quad\le K_0d_n(\alpha,\beta;r).
\end{aligned}
\]
The last inequality is the deterministic calibration established above.
% lean: CausalSmith.Stat.LogoddsLowsmoothFrontier.upperInterval_expectedLength_le_radii

\item To establish honesty, take any \(P\in\mathcal M(\alpha,\beta)\).
The numerator bias and variance bounds have already been established. For the denominator, the conditional probabilities defined above satisfy
\[
p_{10}(P,x)=e(P,x)\{1-\mu_1(P,x)\},
\qquad
p_{01}(P,x)=\{1-e(P,x)\}\mu_0(P,x).
\]
All their factors lie in \([0,1]\). Adding and subtracting the product with the first factor evaluated at \(x\) and the second at \(z\) gives
\[
\begin{aligned}
|p_{10}(P,x)-p_{10}(P,z)|
&\le |e(P,x)-e(P,z)|
   +|\mu_1(P,x)-\mu_1(P,z)|\\
&\le |x-z|^\beta,\\
|p_{01}(P,x)-p_{01}(P,z)|
&\le |e(P,x)-e(P,z)|
   +|\mu_0(P,x)-\mu_0(P,z)|\\
&\le |x-z|^\beta.
\end{aligned}
\]
Thus both continuous functions have Hölder seminorm at exponent \(\beta\) bounded by one.

The denominator marks \(A(1-Y)\) and \((1-A)Y\) have conditional means \(p_{10}(P,\cdot)\) and \(p_{01}(P,\cdot)\).
The expectation and residual identities in \cref{obj:lem:projection-control} therefore give
\[
\mathbb E_{P,n}\widehat S_n-S(P)
=
-\left\langle
p_{10}(P,\cdot)-\Pi_{k_S}p_{10}(P,\cdot),
p_{01}(P,\cdot)-\Pi_{k_S}p_{01}(P,\cdot)
\right\rangle_{L^2(\lambda)}.
\]
Cauchy--Schwarz and the cited approximation bound imply
\[
\begin{aligned}
|\mathbb E_{P,n}\widehat S_n-S(P)|
&\le
\|p_{10}(P,\cdot)-\Pi_{k_S}p_{10}(P,\cdot)\|_2
\|p_{01}(P,\cdot)-\Pi_{k_S}p_{01}(P,\cdot)\|_2\\
&\le25k_S^{-2\beta}.
\end{aligned}
\]
The variance bound in the same cited statement applies to these bounded marks:
\[
\operatorname{Var}_{P,n}(\widehat S_n)
\le\frac4n+\frac{2k_S}{n(n-1)}
\le V_n(k_S).
\]
Together, the four coordinate moment bounds are
\[
\begin{gathered}
|\mathbb E_{P,n}\widehat C_n-C(P)|\le8k_C^{-s},
\qquad
|\mathbb E_{P,n}\widehat S_n-S(P)|\le25k_S^{-2\beta},\\
\operatorname{Var}_{P,n}(\widehat C_n)\le V_n(k_C),
\qquad
\operatorname{Var}_{P,n}(\widehat S_n)\le V_n(k_S).
\end{gathered}
\]
% lean: CausalSmith.Stat.LogoddsLowsmoothFrontier.coordinateEstimates_bias_identities
% lean: CausalSmith.Stat.LogoddsLowsmoothFrontier.numerator_projection_bias_bound
% lean: CausalSmith.Stat.LogoddsLowsmoothFrontier.denominator_projection_bias_bound
% lean: CausalSmith.Stat.LogoddsLowsmoothFrontier.coordinateEstimates_variance_bounds

\item Define the coordinate failure events by
\[
\begin{aligned}
\mathcal B_C
&=\{(\mathcal O,U):|\widehat C_n-C(P)|>b_C\},\\
\mathcal B_S
&=\{(\mathcal O,U):|\widehat S_n-S(P)|>b_S\}.
\end{aligned}
\]
The triangle inequality and the bias bounds imply
\[
\begin{aligned}
\mathcal B_C
&\subseteq
\left\{
|\widehat C_n-\mathbb E_{P,n}\widehat C_n|
\ge\sqrt{20V_n(k_C)}
\right\},\\
\mathcal B_S
&\subseteq
\left\{
|\widehat S_n-\mathbb E_{P,n}\widehat S_n|
\ge\sqrt{20V_n(k_S)}
\right\}.
\end{aligned}
\]
Both variance radii are strictly positive. Integrating each squared centered coordinate over its displayed event gives
\[
\begin{aligned}
20V_n(k_C)\,
\mathbb P_{P,n}
\left\{
|\widehat C_n-\mathbb E_{P,n}\widehat C_n|
\ge\sqrt{20V_n(k_C)}
\right\}
&\le\operatorname{Var}_{P,n}(\widehat C_n),\\
20V_n(k_S)\,
\mathbb P_{P,n}
\left\{
|\widehat S_n-\mathbb E_{P,n}\widehat S_n|
\ge\sqrt{20V_n(k_S)}
\right\}
&\le\operatorname{Var}_{P,n}(\widehat S_n).
\end{aligned}
\]
Consequently,
\[
\mathbb P_{P,n}(\mathcal B_C)\le\frac1{20},
\qquad
\mathbb P_{P,n}(\mathcal B_S)\le\frac1{20},
\qquad
\mathbb P_{P,n}(\mathcal B_C\cup\mathcal B_S)\le\frac1{10}.
\]

Outside their union,
\[
C(P)\in[c_-,c_+],
\qquad
\widehat S_n-b_S\le S(P)\le\widehat S_n+b_S.
\]
We have already shown \(s_0\le S(P)\le1\).
Clipping is nondecreasing and fixes every point of \([s_0,1]\), so
\[
s_-\le S(P)\le s_+.
\]
For a fixed positive denominator, \(F\) is nondecreasing in its numerator.
For a fixed real numerator \(c_{\mathrm{rect}}\), its behavior as the positive denominator varies follows from
\[
\begin{cases}
\displaystyle
\frac{c_{\mathrm{rect}}}{s_+}
\le\frac{c_{\mathrm{rect}}}{s_{\mathrm{rect}}}
\le\frac{c_{\mathrm{rect}}}{s_-},
& c_{\mathrm{rect}}\ge0,\\[6pt]
\displaystyle
\frac{c_{\mathrm{rect}}}{s_-}
\le\frac{c_{\mathrm{rect}}}{s_{\mathrm{rect}}}
\le\frac{c_{\mathrm{rect}}}{s_+},
& c_{\mathrm{rect}}<0,
\end{cases}
\qquad
s_{\mathrm{rect}}\in[s_-,s_+].
\]
The clipping and logarithm preserve these inequalities. Applying them first to \(c_-\), then to \(c_+\), and using numerator monotonicity yields
\[
\begin{aligned}
L
&\le F(c_-,S(P))
\le F(C(P),S(P))\\
&\le F(c_+,S(P))
\le R.
\end{aligned}
\]
By \cref{obj:lem:observable-odds} and the effect envelope,
\[
1+\frac{C(P)}{S(P)}
=\exp(\theta(P))
\in[\exp(-1/2),\exp(1/2)],
\qquad
F(C(P),S(P))=\theta(P).
\]
Thus \(\theta(P)\in[L,R]\) outside \(\mathcal B_C\cup\mathcal B_S\).
The containment supplied by \cref{obj:lem:fixed-budget-realization} gives
\[
[L,R]\subseteq
I^{\mathrm{up}}_{n,\mathsf E,\mathsf N}(\alpha,\beta)(\mathcal O,U)
\]
on every sample. Hence
\[
\begin{aligned}
1
&\le
\mathbb P_{P,n}
\left\{
\theta(P)\in
I^{\mathrm{up}}_{n,\mathsf E,\mathsf N}(\alpha,\beta)(\mathcal O,U)
\right\}
+\mathbb P_{P,n}(\mathcal B_C\cup\mathcal B_S)\\
&\le
\mathbb P_{P,n}
\left\{
\theta(P)\in
I^{\mathrm{up}}_{n,\mathsf E,\mathsf N}(\alpha,\beta)(\mathcal O,U)
\right\}
+\frac1{10}.
\end{aligned}
\]
The coverage probability is therefore at least \(9/10\).
The Borel output property from the same cited statement gives measurable coverage events and interval length. Since \(P\) was arbitrary in the ambient model, this proves
\[
I^{\mathrm{up}}_{n,\mathsf E,\mathsf N}(\alpha,\beta)
\in\mathcal A_n(\alpha,\beta).
\]
% lean: CausalSmith.Stat.LogoddsLowsmoothFrontier.coordinate_error_probability_le
% lean: CausalSmith.Stat.LogoddsLowsmoothFrontier.upperInterval_coverage_of_coordinate_moments

\item Now let \(n\ge1\). If \(n=1\), the first branch of
\cref{obj:def:upper-handle} returns
\[
I^{\mathrm{up}}_{1,\mathsf E,\mathsf N}(\alpha,\beta)(\mathcal O,U)
=[-1/2,1/2]
\]
at every input. The effect envelope gives coverage one, and its length is one. For every \(r\in[0,1/2]\),
\[
d_1(\alpha,\beta;r)=1+r,
\qquad
\mathbb E_{P,1}
\left|
I^{\mathrm{up}}_{1,\mathsf E,\mathsf N}(\alpha,\beta)(\mathcal O,U)
\right|
=1
\le K_0(1+r).
\]
If \(n\ne1\), the integer assumption \(n\ge1\) implies \(n\ge2\), and the honesty and expected-length bounds established above apply. Thus the same positive constant \(K_0\) gives both assertions for every admissible engine, policy, exponent pair, sample size, radius, and law in the theorem.
% lean: CausalSmith.Stat.LogoddsLowsmoothFrontier.upperInterval_one_honest
% lean: CausalSmith.Stat.LogoddsLowsmoothFrontier.upperInterval_one_expectedLength_bound

\item Finally, every branch of the literal output construction either takes a nonempty intersection of a rational endpoint box with the full effect box or returns the full effect box. Each resulting box has ordered rational endpoints and is contained in \([-1/2,1/2]\); the output includes both endpoints. Consequently, for every sample and seed there exist \(l,u\in\mathbb Q\) such that
\[
I^{\mathrm{up}}_{n,\mathsf E,\mathsf N}(\alpha,\beta)(\mathcal O,U)
=[l,u],
\qquad
-\frac12\le l\le u\le\frac12.
\]
For \(n=1\), the full-box branch gives the asserted equality with
\([-1/2,1/2]\) at every sample and seed.
% lean: CausalSmith.Stat.LogoddsLowsmoothFrontier.upperOutput_rational_shape
% lean: CausalSmith.Stat.LogoddsLowsmoothFrontier.upperInterval_one_set
\end{enumerate}
\end{proof}

\begin{proof}[Proof of \cref{obj:thm:sharp-honest-frontier}]
\begin{enumerate}
\item Fix the joint calibration constants \(\varepsilon,M\) supplied by
\cref{obj:lem:calibrated-support}, and use the corresponding positive constants
\(\kappa,C_{\mathrm{mix}}\) from \cref{obj:lem:original-record-mixtures}. Define
\[
L_{\mathrm{mix}}
=
\left\lceil
\max\{1,\kappa^{-1},100C_{\mathrm{mix}}\varepsilon^4\}
\right\rceil,
\qquad
d_{\mathrm{ceil}}=(L_{\mathrm{mix}}+1)^{-1/2}.
\]
These are finite absolute constants, and the ceiling gives
\[
L_{\mathrm{mix}}\ge1,\qquad
\frac1{L_{\mathrm{mix}}}\le\kappa,\qquad
\frac{C_{\mathrm{mix}}\varepsilon^4}{L_{\mathrm{mix}}}
\le\frac1{100},\qquad
d_{\mathrm{ceil}}>0.
\]
For the argument below, write the worst expected length of an admissible procedure as
\[
\mathcal W_n(\alpha,\beta;r,I)
=
\sup_{P\in\mathcal M_r(\alpha,\beta)}
\mathbb E_{P^{\otimes n}\otimes\lambda}
|I(\mathcal O,U)|.
\]
We work throughout with \((\alpha,\beta)\in\mathcal D\), \(n\ge1\),
\(r\in[0,1/2]\), and \(I\in\mathcal A_n(\alpha,\beta)\), unless a narrower
range is specified.
% lean: CausalSmith.Stat.LogoddsLowsmoothFrontier.frontier_mixture_lower_rates

\item We first establish the rank estimates used in both mixture constructions.
For every real \(\gamma_{\mathrm{rank}}\in(0,1/2]\), define
\[
k_{\mathrm{low}}(\gamma_{\mathrm{rank}},n)
=
\left\lceil
L_{\mathrm{mix}}n^{2/(2\gamma_{\mathrm{rank}}+1)}
\right\rceil.
\]
Since \(n^{2/(2\gamma_{\mathrm{rank}}+1)}\ge1\), the ceiling inequalities yield
\[
1\le k_{\mathrm{low}}(\gamma_{\mathrm{rank}},n),\qquad
L_{\mathrm{mix}}n^{2/(2\gamma_{\mathrm{rank}}+1)}
\le k_{\mathrm{low}}(\gamma_{\mathrm{rank}},n)
\le
(L_{\mathrm{mix}}+1)n^{2/(2\gamma_{\mathrm{rank}}+1)}.
\]
The exponent \(2/(2\gamma_{\mathrm{rank}}+1)\) is at least one, so
\[
L_{\mathrm{mix}}n
\le
L_{\mathrm{mix}}n^{2/(2\gamma_{\mathrm{rank}}+1)}
\le k_{\mathrm{low}}(\gamma_{\mathrm{rank}},n),
\qquad
\frac{n}{k_{\mathrm{low}}(\gamma_{\mathrm{rank}},n)}
\le\frac1{L_{\mathrm{mix}}}.
\]
Moreover, \(L_{\mathrm{mix}}^{2\gamma_{\mathrm{rank}}+1}\ge L_{\mathrm{mix}}\).
Raising the lower rank bound to the positive power \(2\gamma_{\mathrm{rank}}+1\) therefore gives
\[
k_{\mathrm{low}}(\gamma_{\mathrm{rank}},n)^{2\gamma_{\mathrm{rank}}+1}
\ge
L_{\mathrm{mix}}^{2\gamma_{\mathrm{rank}}+1}n^2
\ge L_{\mathrm{mix}}n^2,
\]
and hence
\[
n^2k_{\mathrm{low}}(\gamma_{\mathrm{rank}},n)^{-(2\gamma_{\mathrm{rank}}+1)}
\le\frac1{L_{\mathrm{mix}}}.
\]
Finally, raising the upper rank bound to the negative power
\(-\gamma_{\mathrm{rank}}\), and using
\(\gamma_{\mathrm{rank}}\le1/2\), gives
\[
\begin{aligned}
k_{\mathrm{low}}(\gamma_{\mathrm{rank}},n)^{-\gamma_{\mathrm{rank}}}
&\ge
(L_{\mathrm{mix}}+1)^{-\gamma_{\mathrm{rank}}}
n^{-2\gamma_{\mathrm{rank}}/(2\gamma_{\mathrm{rank}}+1)}\\
&\ge
d_{\mathrm{ceil}}\,
n^{-2\gamma_{\mathrm{rank}}/(2\gamma_{\mathrm{rank}}+1)}.
\end{aligned}
\]
All these estimates include \(n=1\).
% lean: CausalSmith.Stat.LogoddsLowsmoothFrontier.lower_rank_ceiling_bounds
% lean: CausalSmith.Stat.LogoddsLowsmoothFrontier.lower_rank_budget
% lean: CausalSmith.Stat.LogoddsLowsmoothFrontier.lower_rank_separation

\item Suppose first that \(\alpha+\beta\le1/2\). Define the mixed rank and amplitudes by
\[
k_{\mathrm{mixed}}=k_{\mathrm{low}}(\alpha+\beta,n),
\qquad
\eta=\varepsilon k_{\mathrm{mixed}}^{-\alpha},
\qquad
\zeta=\varepsilon k_{\mathrm{mixed}}^{-\beta}.
\]
The domain conditions imply \(0<\alpha+\beta\le1/2\) and
\(\alpha,\beta>0\). Since \(k_{\mathrm{mixed}}\ge1\),
\[
0<\eta\le\varepsilon,\qquad 0<\zeta\le\varepsilon,
\qquad
\eta\zeta=\varepsilon^2k_{\mathrm{mixed}}^{-(\alpha+\beta)}.
\]
The rank estimates and the calibration of \(L_{\mathrm{mix}}\) give
\[
\frac{n}{k_{\mathrm{mixed}}}\le\kappa
\]
and, by the laws of real powers,
\[
\begin{aligned}
C_{\mathrm{mix}}\frac{n^2}{k_{\mathrm{mixed}}}\eta^2\zeta^2
&=
C_{\mathrm{mix}}\varepsilon^4
n^2k_{\mathrm{mixed}}^{-(2\alpha+2\beta+1)}\\
&\le
\frac{C_{\mathrm{mix}}\varepsilon^4}{L_{\mathrm{mix}}}
\le\frac1{100}.
\end{aligned}
\]
Thus \cref{obj:lem:original-record-mixtures}, applied to the mixed priors of
\cref{obj:def:continuous-calibrated-priors}, supplies
\[
\mathcal H^2(Q_0,Q_1)\le\frac1{100}.
\]

We give the testing-to-length argument needed here and for the fair priors.
Let \(Q_{\mathrm{eval}}\) and \(Q_{\mathrm{alt}}\) be uniform finite mixtures
of original-record sample laws. More precisely, for
\(\iota_{\mathrm{prior}}\in\{\mathrm{eval},\mathrm{alt}\}\), take an integer \(k_{\iota_{\mathrm{prior}}}\ge0\), define
\[
\Sigma_{\iota_{\mathrm{prior}}}=\{-1,1\}^{k_{\iota_{\mathrm{prior}}}+1},
\qquad
Q_{\iota_{\mathrm{prior}}}=\frac1{|\Sigma_{\iota_{\mathrm{prior}}}|}
\sum_{\sigma\in\Sigma_{\iota_{\mathrm{prior}}}}P_{\iota_{\mathrm{prior}},\sigma}^{\otimes n},
\]
and suppose that their constituents have common effects
\[
\theta(P_{\mathrm{eval},\sigma})=t_{\mathrm{eval}},
\qquad
\theta(P_{\mathrm{alt},\sigma})=t_{\mathrm{alt}}.
\]
Assume that every evaluation constituent belongs to
\(\mathcal M_r(\alpha,\beta)\), every alternative constituent belongs to
\(\mathcal M(\alpha,\beta)\), and all constituents have uniform covariate law.

Let \(\tau_{\mathrm{lab}}\) be the uniform probability law on the four binary
label pairs, and define the original-record reference and mixture densities by
\[
\tau_{\mathrm{lab}}
=\frac14\sum_{(a_{\mathrm{lab}},y_{\mathrm{lab}})\in\{0,1\}^2}
\delta_{(a_{\mathrm{lab}},y_{\mathrm{lab}})},
\qquad
R_{\mathrm{sample}}=(\lambda\otimes\tau_{\mathrm{lab}})^{\otimes n},
\qquad
q_{\iota_{\mathrm{prior}}}=\frac{dQ_{\iota_{\mathrm{prior}}}}{dR_{\mathrm{sample}}}.
\]
These densities are nonnegative and integrate to one. With total variation defined by
\[
\operatorname{TV}(Q_{\mathrm{eval}},Q_{\mathrm{alt}})
=
\frac12\int
|q_{\mathrm{eval}}-q_{\mathrm{alt}}|\,dR_{\mathrm{sample}},
\]
the squared Hellinger distance satisfies
\[
\mathcal H^2(Q_{\mathrm{eval}},Q_{\mathrm{alt}})
=
\int
(\sqrt{q_{\mathrm{eval}}}-\sqrt{q_{\mathrm{alt}}})^2\,dR_{\mathrm{sample}}
=
2-2\int\sqrt{q_{\mathrm{eval}}q_{\mathrm{alt}}}\,dR_{\mathrm{sample}}.
\]
Cauchy--Schwarz gives
\[
\int\sqrt{q_{\mathrm{eval}}q_{\mathrm{alt}}}\,dR_{\mathrm{sample}}\le1,
\qquad
\int(\sqrt{q_{\mathrm{eval}}}+\sqrt{q_{\mathrm{alt}}})^2
\,dR_{\mathrm{sample}}\le4.
\]
Factoring the density difference and applying Cauchy--Schwarz again yields
\[
\operatorname{TV}(Q_{\mathrm{eval}},Q_{\mathrm{alt}})
\le
\frac12
\sqrt{\mathcal H^2(Q_{\mathrm{eval}},Q_{\mathrm{alt}})}
\left[
\int(\sqrt{q_{\mathrm{eval}}}+\sqrt{q_{\mathrm{alt}}})^2
\,dR_{\mathrm{sample}}
\right]^{1/2}
\le
\sqrt{\mathcal H^2(Q_{\mathrm{eval}},Q_{\mathrm{alt}})}.
\]
The product densities with respect to
\(R_{\mathrm{sample}}\otimes\lambda\) are the same densities, independent of
the seed. Integrating over that seed proves
\[
\operatorname{TV}(Q_{\mathrm{eval}}\otimes\lambda,
Q_{\mathrm{alt}}\otimes\lambda)
=
\operatorname{TV}(Q_{\mathrm{eval}},Q_{\mathrm{alt}}).
\]
Consequently, a squared Hellinger bound of \(1/100\) gives total variation
at most \(1/10\) for the sample-and-seed experiments.

Define the measurable coverage events
\[
\mathcal E_{\mathrm{eval}}
=\{(\mathcal O,U):t_{\mathrm{eval}}\in I(\mathcal O,U)\},
\qquad
\mathcal E_{\mathrm{alt}}
=\{(\mathcal O,U):t_{\mathrm{alt}}\in I(\mathcal O,U)\}.
\]
Global honesty and finite averaging imply
\[
(Q_{\mathrm{eval}}\otimes\lambda)(\mathcal E_{\mathrm{eval}})\ge\frac9{10},
\qquad
(Q_{\mathrm{alt}}\otimes\lambda)(\mathcal E_{\mathrm{alt}})\ge\frac9{10}.
\]
Transferring the second event by total variation, then using that a union has
probability at most one, gives
\[
(Q_{\mathrm{eval}}\otimes\lambda)
(\mathcal E_{\mathrm{eval}}\cap\mathcal E_{\mathrm{alt}})
\ge\frac9{10}+\left(\frac9{10}-\frac1{10}\right)-1
=\frac7{10}.
\]
Connectedness of the returned interval gives the pointwise inequality
\[
|I(\mathcal O,U)|
\ge
|t_{\mathrm{alt}}-t_{\mathrm{eval}}|\,
\mathbf1_{\mathcal E_{\mathrm{eval}}\cap\mathcal E_{\mathrm{alt}}}
(\mathcal O,U).
\]
Integration and finite averaging therefore show
\[
\begin{aligned}
\frac7{10}|t_{\mathrm{alt}}-t_{\mathrm{eval}}|
&\le
\mathbb E_{Q_{\mathrm{eval}}\otimes\lambda}|I(\mathcal O,U)|\\
&=
\frac1{|\Sigma_{\mathrm{eval}}|}
\sum_{\sigma\in\Sigma_{\mathrm{eval}}}
\mathbb E_{P_{\mathrm{eval},\sigma}^{\otimes n}\otimes\lambda}
|I(\mathcal O,U)|\\
&\le\mathcal W_n(\alpha,\beta;r,I).
\end{aligned}
\]
The procedure's measurability and bounded interval length justify these
integrals. This argument includes independently randomized procedures.
% lean: CausalSmith.Stat.LogoddsLowsmoothFrontier.finitePrior_hellinger_length_lower

For the mixed priors, \cref{obj:lem:calibrated-support} supplies model
membership and the effects
\[
\theta(P_{0,\sigma})=0,\qquad
\theta(P_{1,\sigma})=32\eta\zeta.
\]
The null constituents belong to every evaluation slice because \(r\ge0\).
Applying the preceding reduction with \(Q_0\) as evaluation prior gives
\[
\mathcal W_n(\alpha,\beta;r,I)
\ge\frac7{10}\,32\eta\zeta
\ge
\frac7{10}\,32\varepsilon^2d_{\mathrm{ceil}}\,
n^{-2(\alpha+\beta)/(2\alpha+2\beta+1)}.
\]
Since the denominator is positive,
\[
\frac{2(\alpha+\beta)}{2\alpha+2\beta+1}\le\frac12
\quad\Longleftrightarrow\quad
\alpha+\beta\le\frac12.
\]
Thus, defining
\[
c_{\mathrm{mix}}=\frac7{10}\,32\varepsilon^2d_{\mathrm{ceil}}>0,
\]
we have proved
\[
c_{\mathrm{mix}}n^{-\mathsf a(\alpha,\beta)}
\le\mathcal W_n(\alpha,\beta;r,I)
\qquad(\alpha+\beta\le1/2).
\]
% lean: CausalSmith.Stat.LogoddsLowsmoothFrontier.mixed_prior_length_lower
% lean: CausalSmith.Stat.LogoddsLowsmoothFrontier.frontier_mixture_lower_rates

\item For the radius-dependent lower bound, suppose \(r>0\) and define
\[
k_{\mathrm{fair}}=k_{\mathrm{low}}(2\beta,n),
\qquad
\delta=\varepsilon k_{\mathrm{fair}}^{-\beta},
\qquad
t_{\mathrm{fair}}=\frac r2.
\]
Here \(0<2\beta<1/2\), \(0<\delta\le\varepsilon\le1/100\), and
\(t_{\mathrm{fair}}\in(0,1/4]\). The rank estimates imply
\[
\frac{n}{k_{\mathrm{fair}}}\le\kappa,
\qquad
C_{\mathrm{mix}}\frac{n^2}{k_{\mathrm{fair}}}\delta^4
=
C_{\mathrm{mix}}\varepsilon^4
n^2k_{\mathrm{fair}}^{-(4\beta+1)}
\le\frac1{100}.
\]
By \cref{obj:lem:original-record-mixtures},
\[
\mathcal H^2\!\left(
Q_{t_{\mathrm{fair}}},
P_{*,t_{\mathrm{fair}},\delta}^{\otimes n}
\right)\le\frac1{100}.
\]
The fair support assertion in \cref{obj:lem:calibrated-support} supplies
model membership and
\[
\theta(P_{t_{\mathrm{fair}},\sigma})=t_{\mathrm{fair}},
\qquad
\theta(P_{*,t_{\mathrm{fair}},\delta})
=\mathsf T(t_{\mathrm{fair}},\delta).
\]
The random constituents belong to the evaluation slice because
\(0<t_{\mathrm{fair}}\le r\). The comparator can be regarded as a constant
finite prior: averaging identical copies leaves its sample law unchanged.
By \cref{obj:lem:exact-calibrations},
\[
\frac{t_{\mathrm{fair}}}{2}
\le\mathsf T(t_{\mathrm{fair}},\delta)\le t_{\mathrm{fair}},
\qquad
t_{\mathrm{fair}}-\mathsf T(t_{\mathrm{fair}},\delta)
\ge3t_{\mathrm{fair}}\delta^2.
\]
Using the random fair prior as evaluation prior in the testing reduction,
and then the rank separation estimate at \(\gamma_{\mathrm{rank}}=2\beta\),
we obtain
\[
\begin{aligned}
\mathcal W_n(\alpha,\beta;r,I)
&\ge
\frac7{10}
\bigl[t_{\mathrm{fair}}-\mathsf T(t_{\mathrm{fair}},\delta)\bigr]\\
&\ge
\frac7{10}\,\frac32 r\varepsilon^2k_{\mathrm{fair}}^{-2\beta}\\
&\ge
\frac7{10}\,\frac32\varepsilon^2d_{\mathrm{ceil}}\,
r n^{-4\beta/(4\beta+1)}.
\end{aligned}
\]
Define
\[
c_{\mathrm{radius}}
=\frac7{10}\,\frac32\varepsilon^2d_{\mathrm{ceil}}>0.
\]
We have therefore proved
\[
c_{\mathrm{radius}}r n^{-\mathsf b(\beta)}
\le\mathcal W_n(\alpha,\beta;r,I)
\qquad(0<r\le1/2).
\]
% lean: CausalSmith.Stat.LogoddsLowsmoothFrontier.fair_prior_length_lower
% lean: CausalSmith.Stat.LogoddsLowsmoothFrontier.frontier_mixture_lower_rates

\item Choose the positive finite absolute upper constant supplied by
\cref{obj:thm:finite-honest-upper}, and set
\[
C=K_0.
\]
For every admissible engine and naming policy, that result supplies
\[
I^{\mathrm{up}}_{n,\mathsf E,\mathsf N}(\alpha,\beta)
\in\mathcal A_n(\alpha,\beta),
\]
together with
\[
\mathbb E_{P^{\otimes n}\otimes\lambda}
\left|I^{\mathrm{up}}_{n,\mathsf E,\mathsf N}(\alpha,\beta)(\mathcal O,U)\right|
\le C\rho_n(\alpha,\beta;r)
\qquad(P\in\mathcal M_r(\alpha,\beta)).
\]
Here the identification of the diagnostic envelope with \(\rho_n\) is
\cref{obj:def:frontier-handle}. Taking the supremum over the slice and using
the infimum definition in \cref{obj:def:length-objective} gives
\[
J_n(\alpha,\beta;r)
\le
\mathcal W_n\!\left(
\alpha,\beta;r,I^{\mathrm{up}}_{n,\mathsf E,\mathsf N}(\alpha,\beta)
\right)
\le C\rho_n(\alpha,\beta;r).
\]
The slice is nonempty: the law
\[
P_{\mathrm{null}}=\lambda\otimes\tau_{\mathrm{lab}}
\]
has
\[
e(P_{\mathrm{null}},x)=\mu_0(P_{\mathrm{null}},x)
=\mu_1(P_{\mathrm{null}},x)=\frac12,
\qquad
g(P_{\mathrm{null}},x)=\nu(P_{\mathrm{null}},x)
=\theta(P_{\mathrm{null}})=0,
\]
and hence satisfies the model conditions and belongs to every slice with
\(r\ge0\).

The upper result also supplies Borel measurability, global \(90\%\) coverage,
and the ordered rational closed output in the effect envelope. The
construction in \cref{obj:def:upper-handle} depends on the fixed engine,
public inputs, selected names, and original records, and therefore satisfies
\[
I^{\mathrm{up}}_{n,\mathsf E,\mathsf N}(\alpha,\beta)(\mathcal O,U)
=
I^{\mathrm{up}}_{n,\mathsf E,\mathsf N}(\alpha,\beta)(\mathcal O).
\]
Its input list makes it radius independent. For \(n\ge2\),
\cref{obj:lem:fixed-budget-realization} supplies successful termination
within the two prescribed budgets and finite rational evaluation. For
\(n=1\), the construction returns \([-1/2,1/2]\) directly.
% lean: CausalSmith.Stat.LogoddsLowsmoothFrontier.frontier_upper_comparisons

\item Define the common lower constants by
\[
c_{\mathrm{numerator}}
=\min\{c_{\mathrm{mix}},3/16\},
\qquad
c=\frac12\min\{c_{\mathrm{numerator}},c_{\mathrm{radius}}\}.
\]
Both constants are strictly positive. If \(\alpha+\beta\le1/2\), the mixed
bound gives
\[
c_{\mathrm{numerator}}n^{-\mathsf a(\alpha,\beta)}
\le
c_{\mathrm{mix}}n^{-\mathsf a(\alpha,\beta)}
\le\mathcal W_n(\alpha,\beta;r,I).
\]
If \(\alpha+\beta>1/2\), positivity of \(2\alpha+2\beta+1\) gives
\[
\frac12\le\frac{2(\alpha+\beta)}{2\alpha+2\beta+1},
\qquad
\mathsf a(\alpha,\beta)=\frac12.
\]
By \cref{obj:lem:parametric-null-floor} and the infimum definition of \(J_n\),
\[
c_{\mathrm{numerator}}n^{-\mathsf a(\alpha,\beta)}
\le
\frac3{16}n^{-1/2}
\le J_n(\alpha,\beta;r)
\le\mathcal W_n(\alpha,\beta;r,I).
\]
This proves the numerator bound throughout \(\mathcal D\), with equality
\(\alpha+\beta=1/2\) included in the first case.
% lean: CausalSmith.Stat.LogoddsLowsmoothFrontier.frontier_numerator_lower_of_rough

For the radius bound, \(r>0\) is covered above. At \(r=0\), nonnegativity of
interval length gives
\[
c_{\mathrm{radius}}\,0\,n^{-\mathsf b(\beta)}
=0\le\mathcal W_n(\alpha,\beta;0,I).
\]
Thus
\[
c_{\mathrm{radius}}r n^{-\mathsf b(\beta)}
\le\mathcal W_n(\alpha,\beta;r,I)
\qquad(0\le r\le1/2).
\]
% lean: CausalSmith.Stat.LogoddsLowsmoothFrontier.frontier_radius_lower_of_positive

Each summand of \(\rho_n\) is nonnegative. Reducing both coefficients to
their minimum and adding the two lower inequalities yields
\[
\begin{aligned}
c\rho_n(\alpha,\beta;r)
&=
\frac{\min\{c_{\mathrm{numerator}},c_{\mathrm{radius}}\}}2
\left[n^{-\mathsf a(\alpha,\beta)}
+r n^{-\mathsf b(\beta)}\right]\\
&\le
\frac12\left[
\mathcal W_n(\alpha,\beta;r,I)
+\mathcal W_n(\alpha,\beta;r,I)
\right]\\
&=\mathcal W_n(\alpha,\beta;r,I).
\end{aligned}
\]
This is the asserted converse for every honest procedure, including
independently randomized procedures.
% lean: CausalSmith.Stat.LogoddsLowsmoothFrontier.frontier_rate_lower_of_two

\item The honest class is nonempty: the constant procedure returning
\([-1/2,1/2]\) is Borel and covers every ambient-model effect by
\cref{obj:ass:effect-envelope}. Its length is one, while all procedure
objectives are nonnegative. Taking the infimum of the preceding
all-procedure lower bound is therefore legitimate and gives
\[
c\rho_n(\alpha,\beta;r)\le J_n(\alpha,\beta;r).
\]
Combining this with the upper comparison proves
\[
c\rho_n(\alpha,\beta;r)
\le J_n(\alpha,\beta;r)
\le
\mathcal W_n\!\left(
\alpha,\beta;r,I^{\mathrm{up}}_{n,\mathsf E,\mathsf N}(\alpha,\beta)
\right)
\le C\rho_n(\alpha,\beta;r).
\]
For the upper procedure, seed independence implies, for every law \(P\),
\[
\mathbb E_{P^{\otimes n}\otimes\lambda}
\left|I^{\mathrm{up}}_{n,\mathsf E,\mathsf N}(\alpha,\beta)(\mathcal O,U)\right|
=
\mathbb E_{P^{\otimes n}}
\left|I^{\mathrm{up}}_{n,\mathsf E,\mathsf N}(\alpha,\beta)(\mathcal O)\right|,
\]
which gives the displayed comparison in the statement. The constants were
chosen independently of the engine and naming policy.

By \cref{obj:lem:concrete-arithmetic-engine},
\(\mathsf E_0\) is admissible, and
\cref{obj:lem:fixed-budget-realization} supplies admissibility of
\(\mathsf N_0\). Specializing the universal comparison and operational
properties to this pair gives
\[
I_n^\star=I^{\mathrm{up}}_{n,\mathsf E_0,\mathsf N_0},
\qquad c_\star=c,\qquad C_\star=C,\qquad n_0=1.
\]
% lean: CausalSmith.Stat.LogoddsLowsmoothFrontier.universalFrontier_of_procedure_lower

\item Let \(\mathcal K\subseteq\mathcal D\) be nonempty and compact. The constant
functions with values \(c\) and \(1\) have singleton images on \(\mathcal K\), so
\[
\inf_{(\alpha,\beta)\in\mathcal K}c=c>0,
\qquad
\sup_{(\alpha,\beta)\in\mathcal K}1=1.
\]
The constant upper bound gives
\[
\sup_{(\alpha,\beta)\in\mathcal K}C\le C<\infty.
\]
These are the required compact-uniform conditions for the same constants
and threshold.
% lean: CausalSmith.Stat.LogoddsLowsmoothFrontier.frontier_compact_uniform

\item It remains to establish the exponent identities. The domain conditions give
\[
\alpha>0,\qquad
2\alpha+2\beta+1>0,\qquad
4\beta+1>0.
\]
As already verified by comparison of the two entries of the minimum,
\[
\mathsf a(\alpha,\beta)
=
\frac{2(\alpha+\beta)}{2\alpha+2\beta+1}
\quad\text{if }\alpha+\beta\le\frac12,
\qquad
\mathsf a(\alpha,\beta)=\frac12
\quad\text{if }\alpha+\beta\ge\frac12.
\]
In the first range, subtraction over the positive common denominator gives
\[
\begin{aligned}
\mathsf a(\alpha,\beta)-\mathsf b(\beta)
&=
\frac{
2(\alpha+\beta)(4\beta+1)
-4\beta(2\alpha+2\beta+1)}
{(2\alpha+2\beta+1)(4\beta+1)}\\
&=
\frac{2(\alpha-\beta)}
{(2\alpha+2\beta+1)(4\beta+1)}.
\end{aligned}
\]
In the second range,
\[
\mathsf a(\alpha,\beta)-\mathsf b(\beta)
=
\frac12-\frac{4\beta}{4\beta+1}
=
\frac{1-4\beta}{2(4\beta+1)}.
\]
If \(\alpha+\beta\le1/2\), positivity follows from \(\alpha>\beta\) and
the positive denominators. If \(\alpha+\beta>1/2\), it follows from
\(\beta<1/4\) and the positive denominator. Thus
\[
\mathsf a(\alpha,\beta)-\mathsf b(\beta)>0
\qquad((\alpha,\beta)\in\mathcal D).
\]
At \(\alpha+\beta=1/2\), both branch identities apply to the same difference,
and consequently
\[
\frac{2(\alpha-\beta)}
{(2\alpha+2\beta+1)(4\beta+1)}
=
\mathsf a(\alpha,\beta)-\mathsf b(\beta)
=
\frac{1-4\beta}{2(4\beta+1)}.
\]
This establishes agreement at the boundary and completes the proof.
% lean: CausalSmith.Stat.LogoddsLowsmoothFrontier.frontier_elbow_identities
\end{enumerate}
\end{proof}

\begin{proof}[Proof of \cref{obj:thm:full-honest-frontier-solution}]
\begin{enumerate}
\item Choose the absolute constants \(c,C\) supplied by
\cref{obj:thm:sharp-honest-frontier}, so that \(c>0\) and \(C>0\).
For every admissible engine \(\mathsf E\), admissible naming policy
\(\mathsf N\), exponent pair \((\alpha,\beta)\in\mathcal D\), and integer
\(n\ge1\), that result applies to the literal construction
\[
I^\star_{n,\mathsf E,\mathsf N}(\alpha,\beta)
=
I^{\mathrm{up}}_{n,\mathsf E,\mathsf N}(\alpha,\beta).
\]
It supplies membership in \(\mathcal A_n(\alpha,\beta)\), Borel
measurability, and termination on every valid named input within the
prescribed budgets. The inputs are the fixed engine, the public exponents
and their selected names, and the original observed records and their
selected covariate names, as specified in
\cref{obj:def:resolutions,obj:def:upper-handle}. For every sample and seed,
the returned interval has the form
\[
[l,u],\qquad l,u\in\mathbb Q,\qquad
-\frac12\le l\le u\le\frac12.
\]
For every \(r\in[0,1/2]\), the same result gives
\[
c\rho_n(\alpha,\beta;r)
\le J_n(\alpha,\beta;r)
\le
\sup_{P\in\mathcal M_r(\alpha,\beta)}
E_{P^{\otimes n}\otimes\lambda}
\left[
\left|
I^{\mathrm{up}}_{n,\mathsf E,\mathsf N}(\alpha,\beta)
(\mathcal O,U)
\right|
\right]
\le C\rho_n(\alpha,\beta;r),
\]
where \(\lambda\) is the uniform law of the independent seed \(U\).
The construction uses the same inputs for every evaluation radius.

Retain the choices
\[
c_\star(\alpha,\beta)=c,\qquad
C_\star(\alpha,\beta)=C,\qquad
n_0(\alpha,\beta)=1.
\]
The compact-uniform conclusion of
\cref{obj:thm:sharp-honest-frontier} states that, for every nonempty compact
\(\mathcal K\subset\mathcal D\),
\[
0<\inf_{(\alpha,\beta)\in\mathcal K}c_\star(\alpha,\beta),
\qquad
\sup_{(\alpha,\beta)\in\mathcal K}C_\star(\alpha,\beta)<\infty,
\qquad
\sup_{(\alpha,\beta)\in\mathcal K}n_0(\alpha,\beta)=1.
\]
Its exponent conclusions give, throughout \(\mathcal D\),
\[
\mathsf a(\alpha,\beta)-\mathsf b(\beta)>0,
\]
\[
\mathsf a(\alpha,\beta)-\mathsf b(\beta)=
\begin{cases}
\displaystyle
\frac{2(\alpha-\beta)}
{(2\alpha+2\beta+1)(4\beta+1)},
&\alpha+\beta\le1/2,\\[6pt]
\displaystyle
\frac{1-4\beta}{2(4\beta+1)},
&\alpha+\beta\ge1/2,
\end{cases}
\]
and, when \(\alpha+\beta=1/2\),
\[
\frac{2(\alpha-\beta)}
{(2\alpha+2\beta+1)(4\beta+1)}
=
\frac{1-4\beta}{2(4\beta+1)}.
\]
Finally, its procedure-wise converse supplies, for every
\((\alpha,\beta)\in\mathcal D\), \(r\in[0,1/2]\), \(n\ge1\), and
\(I\in\mathcal A_n(\alpha,\beta)\),
\[
c\rho_n(\alpha,\beta;r)
\le
\sup_{P\in\mathcal M_r(\alpha,\beta)}
E_{P^{\otimes n}\otimes\lambda}
\bigl[|I(\mathcal O,U)|\bigr],
\]
including procedures using independent randomization.
% lean: CausalSmith.Stat.LogoddsLowsmoothFrontier.sharp_honest_frontier

\item By \cref{obj:lem:concrete-arithmetic-engine}, the concrete engine
\(\mathsf E_0\) is admissible. By the first assertion of
\cref{obj:lem:fixed-budget-realization}, the dyadic-floor naming policy
\(\mathsf N_0\) is admissible. Specialize the universal conclusion above
to this pair. The resulting procedure is exactly
\[
I_n^\star(\alpha,\beta)
=
I^{\mathrm{up}}_{n,\mathsf E_0,\mathsf N_0}(\alpha,\beta),
\]
as defined in \cref{obj:def:frontier-handle}. Hence, for every
\((\alpha,\beta)\in\mathcal D\) and \(n\ge1\),
\[
I_n^\star(\alpha,\beta)\in\mathcal A_n(\alpha,\beta),
\]
and every sample and seed produces a closed interval with ordered
rational endpoints in \([-1/2,1/2]\). For every \(r\in[0,1/2]\),
\[
c\rho_n(\alpha,\beta;r)
\le J_n(\alpha,\beta;r)
\le
\sup_{P\in\mathcal M_r(\alpha,\beta)}
E_{P^{\otimes n}\otimes\lambda}
\bigl[|I_n^\star(\alpha,\beta)(\mathcal O,U)|\bigr]
\le C\rho_n(\alpha,\beta;r).
\]
This specialization preserves the constants chosen above and completes
the concrete attainment assertion together with the universal assertions.
% lean: CausalSmith.Stat.LogoddsLowsmoothFrontier.full_honest_frontier_solution
\end{enumerate}
\end{proof}

\bibliographystyle{plainnat}

\bibliography{references}

\end{document}
